% 本文件由 LaTeX 排版 Agent 根据有效 translation.json 自动生成。
% author=John Chrysostom; work_id=2402; title=讲道集：希伯来书

\chapter{\Emph{讲道集：希伯来书}}

% structure: p0001 source=h1 target=omitted reason=page-title-already-rendered-by-parent

本文系他在安息后出版，根据安提约基雅的司铎君士坦丁的笔记整理。\par

\Emph{序言} 论题与摘要\newline{} \Emph{讲道一} \BibleRef{希 1:1-4}\newline{} \Emph{讲道二} \BibleRef{希 1:3-5}\newline{} \Emph{讲道三} \BibleRef{希 1:6-2:4}\newline{} \Emph{讲道四} \BibleRef{希 2:5-15}\newline{} \Emph{讲道五} \BibleRef{希 2:16-3:6}\newline{} \Emph{讲道六} \BibleRef{希 3:7-4:10}\newline{} \Emph{讲道七} \BibleRef{希 4:11-16}\newline{} \Emph{讲道八} \BibleRef{希 5:1-14}\newline{} \Emph{讲道九} \BibleRef{希 6:1-6}\newline{} \Emph{讲道十} \BibleRef{希 6:7-12}\newline{} \Emph{讲道十一} \BibleRef{希 6:13-20}\newline{} \Emph{讲道十二} \BibleRef{希 7:1-10}\newline{} \Emph{讲道十三} \BibleRef{希 7:11-28}\newline{} \Emph{讲道十四} \BibleRef{希 8:1-13}\newline{} \Emph{讲道十五} \BibleRef{希 9:1-14}\newline{} \Emph{讲道十六} \BibleRef{希 9:15-23}\newline{} \Emph{讲道十七} \BibleRef{希 9:24-10:9}\newline{} \Emph{讲道十八} \BibleRef{希 10:8-18}\newline{} \Emph{讲道十九} \BibleRef{希 10:19-25}\newline{} \Emph{讲道二十} \BibleRef{希 10:26-31}\newline{} \Emph{讲道二十一} \BibleRef{希 10:32-11:2}\newline{} \Emph{讲道二十二} \BibleRef{希 11:3-6}\newline{} \Emph{讲道二十三} \BibleRef{希 11:7-12}\newline{} \Emph{讲道二十四} \BibleRef{希 11:13-16}\newline{} \Emph{讲道二十五} \BibleRef{希 11:17-19}\newline{} \Emph{讲道二十六} \BibleRef{希 11:20-27}\newline{} \Emph{讲道二十七} \BibleRef{希 11:28-36}\newline{} \Emph{讲道二十八} \BibleRef{希 11:37-12:3}\newline{} \Emph{讲道二十九} \BibleRef{希 12:4-10}\newline{} \Emph{讲道三十} \BibleRef{希 12:11-15}\newline{} \Emph{讲道三十一} \BibleRef{希 12:14-17}\newline{} \Emph{讲道三十二} \BibleRef{希 12:18-29}\newline{} \Emph{讲道三十三} \BibleRef{希 12:28-13:16}\newline{} \Emph{讲道三十四} \BibleRef{希 13:17-25}\par

\SourceRecord{Translated by Frederic Gardiner.；Nicene and Post-Nicene Fathers, First Series；Vol. 14.；Edited by Philip Schaff.；Buffalo, NY: Christian Literature Publishing Co.,；1889.；<http://www.newadvent.org/fathers/2402.htm>.}

% structure: p0001 source=h1 target=section reason=page-title

\section{希伯来书讲道集（论题与摘要）}

\SourceRecord{Translated by Frederic Gardiner.；Nicene and Post-Nicene Fathers, First Series；Vol. 14.；Edited by Philip Schaff.；Buffalo, NY: Christian Literature Publishing Co.,；1889.；<http://www.newadvent.org/fathers/240200.htm>.}

% structure: p0001 source=h1 target=section reason=page-title

\section{\WorkTitle{讲道一：论希伯来人书}}

% structure: p0002 source=h2 target=subsection reason=relative-heading-rank; base=h2

\subsection{希伯来人书 1:1-2}

\BibleQuote{天主在古时，曾多次并以多种方式，藉着先知对我们的祖先说过话；但在这末期内，他藉着自己的儿子对我们说了话。天主立了他为万有的承继者，并藉着他创造了宇宙。}\par

1. 的确，\BibleQuote{罪恶在那里越多，恩宠在那里也越格外丰富}，\BibleRef{罗 5:20} 圣保禄在写给希伯来人的书信的开头，就暗示了这一点。因为由于他们受着苦难，被恶事折磨，并以此判断事物，很可能认为自己的处境比所有其他人更糟——他表明，他们反而成了更大、甚至极丰盛恩宠的分享者；他在讲道的一开始就激动听众。因此他说：\BibleQuote{天主在古时，曾多次并以多种方式，藉着先知对我们的祖先说过话；但在这末期内，他藉着自己的儿子对我们说了话。}\par

为什么他不把\Quote{自己}与\Quote{先知}对立起来呢？当然，他比先知们大得多，因为更大的托付交给了他。但他没有这样做。为什么？首先，为避免谈论自己的大事。其次，因为他的听众尚未完美。第三，因为他更希望提升他们，并表明他们的优越性是大的。就如他所说：先知对我们的祖先来说算得了什么大事？因为对我们，祂差遣了自己的独生子。\par

他这样开始很好：\Quote{多次并以多种方式}，因为这指出连先知自己也没有见过天主；然而，圣子见过祂。因为\Quote{多次并以多种方式}与\Quote{用不同的方式}相同。\BibleQuote{因为我（祂说）增加了异象，并借着先知的职务使用了比喻。}\BibleRef{欧 12:10} 所以，优越之处不仅在于，对他们（祖先）派遣了先知，而对我们派遣了圣子；而且在於，没有一位先知见过天主，只有独生子见过祂。他并没有立刻肯定这一点，但通过随后所说的话，当他谈到祂的人性时，他确立了这一点：\BibleQuote{因为天主曾向哪一位天使说过：你是我的儿子}\BibleRef{希 1:5}，以及\BibleQuote{你坐在我右边}\BibleRef{希 1:13}？\par

请看他的大智慧。首先，他从先知们显示了优越性。然后，当他确认了这一点后，他宣告：对他们，祂通过先知说话；但对我们，通过独生子。然后，祂通过天使向他们说话，这一点他也合理地确立了（因为天使也与犹太人交往）：然而，在这方面我们也占优越，因为主人对我们说话，而对他们是仆人，先知是同仆。\par

2. 他也说得好：\Quote{在这末期内}，因为借此他既激动他们，又鼓励对未来失望的他们。就像他在另一处也说：\BibleQuote{主已经近了，什么也不要挂虑}\BibleRef{斐 4:5}，以及\BibleQuote{因为现在我们的救恩比当初信的时候更近了}\BibleRef{罗 13:11}。这里也是如此。那么他说的是什么呢？那在斗争中耗尽的人，当他听到斗争的结束，就稍微恢复呼吸，知道这确实是劳苦的结束，安息的开始。\par

\Quote{在这末期内通过祂的儿子对我们说了话。}看，他又用了\Quote{在儿子中}的说法，为\Quote{通过儿子}，这是针对那些坚持认为这个短语专属於圣神的人。你看见吗？\Quote{在}就是\Quote{通过}？\par

\Quote{古时}这个说法和\Quote{在末期内}这个说法暗示了另一个意思：——当漫长的时期过去，当我们濒临惩罚的边缘，当恩赐已经失效，当没有得救的指望，当我们以为会比任何人得着更少的时候——那时，我们反而得着了更多。\par

看看他多么谨慎地说了这话。因为他没有说\Quote{基督说了}（尽管确实是祂说的），但由於他们的灵魂软弱，还不能承受有关基督的事，他说\Quote{天主通过祂说了话}。你是什么意思？难道天主通过圣子说话吗？是的。那么，难道你就是这样显示优越性吗？因为你这里只指出新约和旧约同属一位，这个优越性并不大。因此，他从此继续这个论点，说：\Quote{祂通过圣子对我们说了话}。\par

（注意：保禄如何与门徒同心同德，把自己放在与他们同等的地位上，说祂对\Quote{我们}说了话；尽管祂并没有直接对保禄说话，而是对宗徒们，并通过他们向众人说话。但他提升他们（希伯来人），并宣告祂也对他们说了话。到目前为止，他完全没有责备犹太人。因为几乎所有先知所说话的对象，都是一类邪恶污秽的人。但至今还没有论到这些；而是论到从天主而来的恩赐。）\par

3. \Quote{天主立了他}，他说，\Quote{为万有的承继者}。什么是\Quote{立了他为万有的承继者}？他在这里谈到肉身（人性）。正如他在第二篇圣咏中也说：\BibleQuote{你向我请求，我必将万民赐你作产业}\BibleRef{咏 2:8}。因为不再是\Quote{雅各伯是上主所保留的一份}，也不是\Quote{以色列是他自己所有的产业}\BibleRef{申 32:9}，而是所有的人：也就是说，祂使他成为万有之主；这也是伯多禄在宗徒大事录中所说的：\BibleQuote{天主已把他立为主，立为基督了}\BibleRef{宗 2:36}。但他用了\Quote{承继者}这个名字，宣告了两件事：祂真正的儿子身份和祂不可剥夺的主权。\Quote{万有的承继者}，即全世界的承继者。\par

然后，他又把话题带回到原先的观点。\BibleQuote{藉着他，他也创造了诸世界（诸世代）。}那些说\Quote{曾有一时，祂不存在}的人在哪里？\par

然后，他用上升的层级，说出了比这一切大得多的话：\par

% structure: p0016 source=h2 target=subsection reason=relative-heading-rank; base=h2

\subsection{希伯来人书 1:3-4}

\BibleQuote{祂是天主荣耀的光辉，是祂本体的真像，以祂大能的话维系万有；当祂藉着自己洗净了我们的罪，就坐在高天至大者的右边。祂所承受的名字，比天使的名字更尊贵，正如同祂比天使更崇高一样。}\par

啊！宗徒的智慧！或者更确切地说，不是保禄的智慧，而是圣神的恩宠才是值得惊叹的。因为他确实不是凭自己的心思说出这些话，也不是以那种方式找到他的智慧。（从何而来呢？从刀、皮子，或工场吗？）而是来自天主的作为。因为他自己的理智并没有生出这些思想，那时他的理智是如此卑微和浅薄，丝毫不比下等人高明；（又怎能呢？因为它完全消耗在买卖和皮子上吗？）但圣神的恩宠却藉着它所愿意的人彰显其力量。\par

正如一个人，想把一个小孩子领到一个高处，甚至高到天堂之顶，他温和地、逐步地进行，藉着下面的阶梯领他向上——然后当他把他放到高处，吩咐他向下观看，见他头晕眼花、迷糊眩晕时，便抓住他，领他下到较低的台阶，让他喘口气；等他恢复后，又领他上去，再带他下来——同样，蒙福的保禄也正是如此行事，无论对希伯来人还是对所有人，这是他从他老师那里学来的。因为连老师也这样做：有时祂领听众到高处，有时又将他们领下，不让他们停留太久。\par

那么，请看这里——他是藉着多少步骤领他们上升，将他们置于信仰的最高峰，然后在他们尚未头晕眼花、被眩晕抓住之前，又如何再次将他们领下，让他们喘口气，说：\BibleQuote{祂藉着圣子对我们说了话，} \BibleQuote{祂立了祂为万有的承继者。} 因为圣子这个名称至此是普通的。如果这里指的是真正的圣子，祂就高于一切；但尽管如此，目前他证明了祂来自高处。\par

且看他如何说：\BibleQuote{祂立了祂，}他说，\BibleQuote{为万有的承继者。} \BibleQuote{立为承继者}这句话是谦卑的。然后他领他们上到更高的台阶，加上：\BibleQuote{藉着祂也造了世界。}然后再上一个更高的，之后就没有别的了：\BibleQuote{祂是天主荣耀的光辉，是祂本体的真像。}他确实将他们领到了不可接近的光明中，到那光辉本身。在他们尚未被闪耀得眼花之前，看他如何又温和地领他们下来，说：\BibleQuote{祂以祂大能的话维系万有；当祂藉着自己洗净了我们的罪，就坐在高天至大者的右边。}他不简单地说\BibleQuote{祂坐下了，}而是说\BibleQuote{在洗净之后，祂坐下了，}因为他触及了降生成人的安排，他的话又变得谦卑了。\par

然后，在顺道稍讲了一点之后（因为他说\BibleQuote{坐在至高者的右边}），[他转向]谦卑的事：\BibleQuote{祂所承受的名字，比天使的名字更尊贵，正如同祂比天使更崇高一样。}从此之后，他讨论的是按肉身而言的方面，因为\BibleQuote{更崇高}并不表达祂在圣神方面的本质（因为那不是\BibleQuote{被造的}，而是\BibleQuote{受生的}），而是按肉身而言：因为这肉身是\BibleQuote{被造的}。然而，这里所讨论的不是关于受造存在。但正如若望所说：\BibleQuote{那在我以后来的，成了在我以前的}（\BibleRef{若 1:15}），即是在尊荣和评价上更高；同样，这里\BibleQuote{祂比天使更崇高}——就是在评价上更高、更好、更荣耀——\BibleQuote{正如同祂所承受的名字比天使的名字更尊贵一样。}你看到他是在谈论那按肉身而言的吗？因为天主圣言一直拥有这个名字；祂不是后来才\BibleQuote{承受}的；也不是在\BibleQuote{洗净了我们的罪}之后才变得\BibleQuote{比天使更崇高}；而是祂一直是\BibleQuote{更崇高}的，而且崇高得无可比拟。因为这是按肉身而论祂说的。\par

确实，这也是我们谈论人时的方式：有时说崇高之事，有时说卑微。比如，当我们说\Quote{人算什么，} \Quote{人是尘土，} \Quote{人是灰烬，}我们是以较差的部分称呼整体。但当我们说\Quote{人是不死的动物，} \Quote{人是理性的，与天上者同族，}我们又是以较好的部分称呼整体。同样，在基督的情形，保禄有时从较差的方面讲论，有时从较好的方面；既想确认救恩计划，又想教导关于不朽的本性。\par

4. 既然\BibleQuote{祂洗净了我们的罪，}让我们保持纯洁；愿我们不沾染任何污点，而是保护祂所植入我们内的美丽和祂的端庄无玷纯洁，\BibleQuote{没有斑点、皱纹或任何这类东西}（\BibleRef{弗 5:27}）。即使是小罪也是\Quote{斑点与皱纹}，我指的是诸如责骂、侮辱、谎言之类的事。\par

不但如此，这些甚至不是小事，相反却是非常大的事；其大足以使人失去天国。如何及以何种方式？\BibleQuote{凡向弟兄说\InnerQuote{傻子}的，}（祂说）\BibleQuote{必受地狱的火}（\BibleRef{玛 5:22}）。如果连称人为\Quote{傻子}，这在一切事中似乎是最轻微的，更像是孩童的戏言，都如此严重；那么，称人为恶毒、狡猾、嫉妒，并抛出千万种其他辱骂的人，将受何种惩罚？还有什么比这更可怕的呢？\par

如今，我恳求你们容忍这劝勉的话语。因为谁若\Quote{做}了什么给\Quote{最小中的一个}，就是向祂做了\BibleRef{玛 25:40}；谁若不向\Quote{最小中的一个}做，就是不向祂做\BibleRef{玛 25:45}。那么，在言语的善恶上岂不也是如此？辱骂弟兄的，就是辱骂天主；尊敬弟兄的，就是尊敬天主。因此，让我们操练舌头，说出善言。因为经上说：\Quote{要制止你的舌头不出恶言}\BibleRef{咏 33:13}。因为天主赐予舌头，并非要我们说恶言、辱骂、彼此诽谤，而是要用来歌颂天主，说出那些\Quote{赐恩宠给听者}\BibleRef{弗 4:29}的话，即建立人的、有益的话。\par

你是否说过别人的坏话？你得到了什么好处，反倒使自己与他一同陷入祸害？因为你获得了诽谤者的名声。的确，没有哪一种恶事，不，没有任何恶事，只停留在受害者身上，而是连作恶者自己也卷入其中。例如，嫉妒的人似乎是在谋害别人，但自己却先收割了自己罪过的果实，消耗自己、磨损自己，并为众人所憎恨。骗子剥夺别人的钱财，是的，也剥夺了自己在人们心中的善意：使自己被众人说坏话。名声远胜于钱财，因为名声不易洗去，而钱财却易获得。更确切地说，缺少名声不会伤害那缺少它的人；但缺少钱财却使你受人责备和嘲笑，成为众人敌视和攻击的对象。\par

易怒的人，首先惩罚和撕裂自己，然后才惩罚与他生气的人。\par

同样，说坏话的人首先侮辱自己，然后才侮辱被说坏话的人；或者，可能连这也不由他控制，他带着污秽可憎之人的名声离去，反而使对方更受人爱戴。因为当一个人听到别人给他恶名，他不以同样的方式回报，反而称赞和敬佩，他称赞的不是别人，而是自己。因为我先前曾说过，正如针对邻人的诽谤首先触及那些制造祸害的人，同样，向邻人行的善事也首先使行善者自己喜悦。善与恶的根源，都公正地先收获自己的果实。正如水，无论是咸是甜，会充满前来取水之人的器皿，却不会减少涌出的源泉；同样，无论邪恶还是美德，无论出自何人，都将成为他的喜乐或毁灭。\par

以上是关于今世的事；但有什么言语能述说来世的事，无论是福还是祸？没有。因为福乐超越一切思想，不仅是言语；而它们的对立面确实用我们熟悉的词语来表达。因为经上说，那里有火、黑暗、锁链和永不死的虫子。但这不仅代表所说的那些事，还有更难以忍受的。为了说服你们，请首先考虑这一点：如果是火，怎么又是黑暗？你们看到那火比这火更难以忍受吗？因为它没有光。如果是火，怎么又永远燃烧？你们看到比这更难以忍受的事发生了吗？因为它不会熄灭。是的，因此它被称为不灭之火。那么，让我们想想，永远燃烧、处于黑暗、发出无数呻吟、咬牙切齿、甚至不被听见，是多么大的痛苦。因为如果在这里，一个娇生惯养的人被投入监狱，他会说仅仅是恶臭、被置于黑暗、与杀人犯捆绑在一起，就比任何死亡更难以忍受；那么，当我们与全世界的杀人犯一同燃烧，既看不见别人也不被别人看见，在如此众多的人群中却以为自己孤独一人时，那会怎样？因为黑暗和幽暗不允许我们分辨身边的人，每个人都将独自燃烧，仿佛只有他一个人在受苦。此外，如果黑暗本身就能折磨和惊吓我们的灵魂，那么当黑暗伴随着同样巨大的痛苦和燃烧时，又会怎样呢？\par

因此，我恳求你们常常自己反复思考这些事，并忍受话语的痛苦，以免我们遭受实事的惩罚。因为确实，这一切必将发生，而那些行为应受这些酷刑室的人，无人能拯救，没有父亲、母亲或兄弟。因为经上说：\Quote{兄弟不能赎回}，祂说，\Quote{人岂能赎回？}\BibleRef{咏 48:7}，尽管他拥有很大的信心，尽管他在天主面前有极大的能力。因为正是祂自己按每个人的行为报应各人，我们的救恩或惩罚取决于这些行为。\par

那么，让我们为自己结交\Quote{不义的钱财的朋友}\BibleRef{路 16:9}，也就是说：让我们施舍；让我们把财产用在她们身上，好使我们能耗尽那火：熄灭它，使我们在那里有胆量。因为在那里，接待我们的不是他们，而是我们自己的行为；因为并非仅仅他们是我们的朋友就能拯救我们，这从后续的话可以得知。因为为什么祂不说：\Quote{你们要为自己结交朋友，使他们接你们到永远的帐幕}，而是还加上了方式？因为祂说\Quote{不义的钱财}，指出我们必须藉着财产来结交她们，表明仅仅友谊不能保护我们，除非我们有善行，除非我们正当地使用不义聚敛的财富。\par

而且，我们这次的讲论——我指的是\OriginalTerm{哀矜}{Almsgiving}——不仅适合富人，也适合穷人。是的，即使有以乞讨为生的人，这话也是对他说的。因为没有一个人，无论多么贫穷，穷到连\Quote{两个小钱}也拿不出来。\BibleRef{路 21:2}因此，一个人从微薄中奉献少量，可以胜过那些拥有大量财产而奉献更多者，正如那位寡妇所行的。因为哀矜的尺度不是按所施与的数量，而是按施与者的能力和意愿来衡量的。在任何情况下都需要意愿，都需要正直的心意，都需要对天主的爱。如果我们以此行事，虽然我们只有一点点而奉献一点点，天主必不转面不顾，反而会视为伟大而可赞叹的：因为祂看顾的是意愿，而不是礼物；如果祂见意愿伟大，祂就按此判定并赏报，使他们分享祂永恒的恩惠。\par

愿天主恩赐我们众人，藉我们的主耶稣基督的恩宠和爱，偕同祂，与圣父及圣神，永受光荣、权能、尊威，从今世直到永远，世世无尽。阿们。\par

\SourceRecord{Translated by Frederic Gardiner.；Nicene and Post-Nicene Fathers, First Series；Vol. 14.；Edited by Philip Schaff.；Buffalo, NY: Christian Literature Publishing Co.,；1889.；<http://www.newadvent.org/fathers/240201.htm>.}

% structure: p0001 source=h1 target=section reason=page-title

\section{希伯来书讲道集第二篇}

% structure: p0002 source=h2 target=subsection reason=relative-heading-rank; base=h2

\subsection{希伯来书 1:3}

\Quote{祂是祂荣耀的光辉，是祂本体的真像，以祂大能的话托住万有；当祂独自洗净了我们的罪。}\par

的确，处处都需要敬畏的心，尤其当我们谈论或聆听有关天主的事时；因为口舌不能说出，心思不能听见任何与我们的天主相称的事。我为何提及口舌或心思呢？因为甚至连远胜于这些的理智，当我们想要说出有关天主的事时，也无法准确地领悟任何事。因为如果\Quote{天主的平安超越一切理解}（见：\BibleRef{斐 4:7}），并且\Quote{天主为爱祂的人所预备的，是眼所未见，耳所未闻，人心所未想到的}（见：\BibleRef{格前 2:9}）；更何况祂自己，平安的天主，万物的创造者，远远超越我们的理性。因此我们应当以信德和敬畏接受一切，当我们的言词因软弱而失败，不能准确地陈述所说的事时，那时更要光荣天主，因为我们有这样的天主，超越我们的思想和概念。因为我们对天主的许多概念无法表达，同样，我们表达的许多事，却无力去理解它们。例如：我们知道天主无所不在，但如何，我们不再明了。我们知道有一个非物质的权能，是我们一切好事的因由；但它是如何，或它是什么，我们不知道。看！我们说话，却不理解。我说，祂无所不在，但我不理解。我说，祂是无始的，但我不理解。我说，祂从自己生了，我又不知如何理解。有些事我们甚至不能说——例如，思想构思却无法说出。\par

并且为向你显示，连\Person{保禄}也是软弱的，未能精确地运用他的比喻；为使你战栗并避免过度探究，请听他所说的话：他称祂为圣子，也称祂为创造者，\Quote{祂是祂荣耀的光辉，是祂本体的真像。}\par

我们必须以敬畏之心接受这一点，并摆脱一切矛盾。\Quote{祂是祂荣耀的光辉}，他说。但注意他在什么意义上理解这一点，你也当这样接受：——祂是出于祂，毫无苦难；祂既不更大，也不更小；因为有些人从这比喻中得出某些奇怪的东西。因为，他们说，\Quote{光辉}不是自立体，而是存在于另一个之中。但你，人哪，不要这样接受，也不要患上\Person{玛策禄}和\Person{佛提诺}的病症。因为他手头有药给你，使你不陷入那种想象，也不让他把你拖入那致命的疾病。他说什么？\Quote{祂本体的真像}[或\OriginalTerm{自立}{subsistence}]：就是说，正如祂[父]是自立地存在，一无所缺，同样圣子也是如此。因为他在这里说的，是显示那原型的不偏不倚的相似和独特的像，即祂[圣子]是自立的。\par

因为上文说\Quote{藉着祂祂创造了万物}的人，在这里将绝对的权柄归于祂。因为他添加了什么？\Quote{以祂大能的话托住万有}；使我们由此推断，祂不仅是祂本体的真像，也是以绝对权柄治理万有。\par

那么你看，他如何将属于父的事应用于圣子。为此他没有简单地说\Quote{托住万有}，也没有说\Quote{以祂的大能}，而是说\Quote{以祂大能的话}。正如刚才我们看见他逐渐上升和下降；现在也是如此，如同踏着阶梯，他升上高处，然后又下降，并说：\Quote{藉着祂祂也创造了世界}。\par

看哪，他在这里也走在两条路上，一条领我们远离\Person{撒伯里乌斯}，另一条远离\Person{亚略}，是的还有另一条，即祂[基督]不应被视为无始，这也是他贯穿始终的做法，也不应被视为与天主疏远。因为如果即使说了这么多之后，仍有些人断言祂是异己的，并给祂另立一位父亲，还说祂与父不和——倘若\Person{保禄}没有宣告这些事，他们会说出什么来呢？\par

那么他如何做？当他被迫医治时，他也就被迫说出卑微的事：例如，\Quote{祂立了祂}（他说）\Quote{为万有的继承人}，并且\Quote{藉着祂祂创造了世界}（见：\BibleRef{希 1:2}）。但为使祂不以另一种方式受辱，他又将祂提升到绝对权柄，并宣告祂与父同尊，是的，如此同等，以致许多人以为祂就是父。\par

你当观察他伟大的智慧。首先他确立前一点并使之准确无误。当这一点得着证明，即祂是天主子，并非与祂疏远，此后他便安全地道出所有高深的言论，不拘多少。因为任何关于祂的高深言论，曾使许多人陷入刚才提到的观念，他首先写下卑微之处，然后安全地上升到他所愿的高度。在说了\Quote{祂立了祂为万有的继承人}和\Quote{藉着祂祂创造了世界}之后，他又加上\Quote{以祂大能的话托住万有}。因为单凭话语治理万有的那位，在产生万有时不能需要任何人。\par

2. 为证明这一点，请留意他如何再次前进，放下\Quote{藉着祂}，而将绝对权能归于祂。因为他用这句话达成心愿后，便从此搁置它，说什么？\BibleQuote{主啊，你在起初奠定了大地的根基，诸天也是你手所造的。}\BibleRef{希 1:10} 哪里也没有说\Quote{藉着祂}，或说\Quote{藉着祂也创造了世界}。那么怎样？难道不是藉着祂创造的吗？是的，但并非如你所说或所想像的\Quote{如同一个工具}；也不是好像除非圣父向祂伸出一只手，否则祂就不会创造。因为正如祂\BibleQuote{不审判任何人}\BibleRef{若 5:22}，却被说成藉着子审判，因为祂生了祂作为审判者；同样，藉着祂创造，因为祂生了祂作为创造者。如果圣父是祂的原因之本原，因为祂是圣父，那么藉着祂所创造的事物更是如此。因此，当他想要表明祂出自于祂时，他必然提到卑微的事。但当他要讲论崇高的事时，\Person{马塞鲁}和\Person{撒伯流}便抓住把柄；但他避免两者的极端，持守中道。因为他既不沉湎于贬抑，免得\Person{撒摩萨塔的保禄}获得立足之地，也不永远停留在崇高言论中；而是相反，显示祂丰富的亲近，免得\Person{撒伯流}冲向他。他称祂为\Quote{子}，\Person{撒摩萨塔的保禄}立刻来到他面前，说祂是一个儿子，如同众人那样。但他给了致命一击，称祂为\Quote{继承人}。然而，与\Person{亚略}一起，他是无耻的。因为\Quote{祂立他为继承人}这话，两者都坚持：前者说它出于软弱；后者仍然提出异议，试图用随后的子句支持自己。因为藉着说\Quote{藉着祂也创造了世界}，他向后打击了无耻的撒摩萨塔人；而\Person{亚略}似乎仍然强大。然而，看他如何也打击他，再次说：\Quote{祂是祂光荣的辉耀。}但看哪！\Person{撒伯流}又跳出来，与\Person{马塞鲁}和\Person{福提诺}一起：但对他们所有人，他也给予一击，说：\Quote{他是祂本体的真像，并以祂权能的话支撑万有。}这里他也伤害了\Person{马西翁}；不太严重，但确实伤害了他。因为整封书信他都在与他们争战。\par

但他所说的\Quote{光荣的辉耀}，也请听基督亲自说：\BibleQuote{我是世界的光。}\BibleRef{若 8:12} 因此他（宗徒）使用\Quote{辉耀}一词，表明这是按照\Quote{光明从光明}的意义说的。这不仅表明这一点，而且也表明祂照亮了我们的灵魂；祂亲自显明了圣父，藉着\Quote{辉耀}，他指示了（圣父与圣子）存有之亲近。注意他措辞的精细。他取了一个\OriginalTerm{本质}{essence}和一个\OriginalTerm{自立体}{subsistence}，以指明两个自立体。他在关于圣神的知识方面也这样做；因为正如他所说，对圣父的知识与对圣神的知识是同一的，因为它们实在是同一的，与自己毫无差异\BibleRef{格前 2:10}；同样，在这里他也抓住某种确定的东西，以表达两者的自立体。\par

他又加上祂是\Quote{\OriginalTerm{印像}{express image}}。因为\Quote{印像}是不同于其原型的东西；然而并非在各方面都不同，而是在具有实在的自立体方面。因为在这里，\Quote{印像}一词也表明与它所印像之物没有变异：在各方面相似。因此，当他称祂为形式和印像时，他们能说什么？\Quote{是的，}他说，\Quote{人也被称为天主的像。}那么怎样？他是（像祂）如同子那样吗？不（他说），因为\Quote{像}一词并不表明相似。然而，在人被称为像这一点上，它表明相似，如同在人中。因为天主在天上是什么，人在世上就是什么，我指的是在统御方面。正如他掌管地上万物，同样天主掌管天上和地上一切。但另一方面，人不被称为\Quote{印像}，他不被称为形式：这措辞表明实体，或者更确切地说，表明实体和实体中的相似。因此，正如\Quote{奴仆的形状}\BibleRef{斐 2:6}所指的无非是毫无变异的人（相对于人性），同样\Quote{天主的形状}所指的无非是天主。\par

\Quote{祂是}（他说）\Quote{祂光荣的辉耀。}请看\Person{保禄}在做什么。说了\Quote{祂是祂光荣的辉耀}之后，他又加了\Quote{祂坐在尊威的右边}：他用了什么名号，却始终找不到\OriginalTerm{实体}{substance}之名。因为无论\Quote{尊威}还是\Quote{光荣}，都没有说出他所想说的名，但他无法找到一个名。这就是我一开始所说的：我们常常想到某物，却无法表达（它）：因为连\Quote{天主}一词也不是实体的名，而且根本不可能找到那个实体的名。\par

如果对天主如此，那有何可怪？因为即使对一位天使，也找不到一个表达其实体的名。或许对灵魂也是如此。因为在我看来，这名字（灵魂）似乎并非指其实体，而是指呼吸。因为人们可能会看到同一个（东西）被称为灵魂、心和理智：因为他说：\BibleQuote{天主，求你给我再造一颗洁净的心}\BibleRef{咏 50:10}，而且人们常看到它（灵魂）被称为神（灵）。\par

\BibleQuote{藉着祂大能之言托住万有。} 告诉我，\BibleQuote{天主说}（经上记载），\BibleQuote{有光} \BibleRef{创 1:3}：\Quote{有人说，圣父命令，圣子服从}？但是请看，祂也［圣子］藉着言语行事。因为（他说）\BibleQuote{托住万有}——即治理；祂维系那将要瓦解的；因为托住世界与创造世界不相上下，甚至更大（若可作惊人语）。因为前者是从无中引出；后者，当受造物将要归于虚无时，将彼此全然对立的万物维系并紧固在一起：这确实伟大奇妙，是极有权能的明证。\par

随后，显示那轻易，他说：\BibleQuote{托住}（他没有说治理，乃是以那些仅用指头移动任何事物并使之旋转的人为喻。）这里他既显示受造之庞大，又显明这庞大对祂不算什么。然后他又显示无劳苦，说：\BibleQuote{藉着祂大能之言。}他说得好：\BibleQuote{藉着言语。}因为在我们，言语被视为虚空之物，他显示在天主并非如此。但祂如何\BibleQuote{藉着言语托住}，他没有进一步说明：因为也是不可能知道的。然后他论及祂的威严：因为\Person{若望}也是如此：说了\BibleQuote{祂是天主} \BibleRef{若 1:1}，之后便引入创造的工程。因为前者间接表达的，即说\BibleQuote{在起初已有圣言}，和\BibleQuote{万物是藉着祂造的} \BibleRef{若 1:3}，后者也藉着\BibleQuote{圣言}和说\BibleQuote{藉着祂也创造了诸世界}而公开宣告。如此他显明祂既是创造者，又在万世之前。那么，当先知论及圣父说\BibleQuote{从永远直到永远，你是天主} \BibleRef{咏 89:2}，而论及圣子说祂在万世之前，是万有的创造者——他们能说什么呢？不仅如此，当论及圣父的同一件事——\Quote{祂在万世之前}——这岂不是也可见于圣子吗？并且有人所说\BibleQuote{祂是生命} \BibleRef{若 1:4}，指出受造的保存，即祂自己是万有的生命——同样，这位也说\BibleQuote{藉着祂大能之言托住万有}：不像那些希腊人，尽其所能地欺骗祂，剥夺了祂的创造本身和眷顾，将祂的权能局限于只达至月亮。\par

\BibleQuote{藉着祂自己洗净了我们的罪。} 论及那些奇妙而伟大的、远超我们的事之后，他又接着论及祂对人的眷顾。因为先前的表述\BibleQuote{托住万有}也是普遍的；然而这更伟大，因为它也是普遍的：因为就祂而言，所有人都信了。正如\Person{若望}说了\BibleQuote{祂是生命}，并以此指出祂的眷顾，又说\BibleQuote{祂是光}。\par

\BibleQuote{藉着祂自己，洗净了我们的罪，便坐在高天至尊的右边。} 他在这里列出两个关于祂眷顾的极大证据：第一是\BibleQuote{洗净我们的罪}，第二是\BibleQuote{藉着祂自己}完成。你在许多地方看见他极为强调这一点——不仅是我们与天主的和好，而且这和好是通过圣子完成的。因为恩赐真正伟大，甚至因着通过圣子而更加伟大。\par

因为说到\BibleQuote{祂坐在右边}和\BibleQuote{藉着祂自己洗净了我们的罪}——虽然他使我们想起十字架，他很快又加上复活与升天。你看他难以言喻的智慧：他没有说\Quote{祂被命令坐下}，而是\BibleQuote{祂坐下了。} 然后，免得你以为祂站着，他又补充说：\BibleQuote{祂曾对哪一个天使说过：你坐在我右边？}\par

\BibleQuote{祂坐在高天至尊的右边。} 这\BibleQuote{高天}是什么？他将天主局限于地方吗？绝无此念！但正如当他说\BibleQuote{右边}时，他没有描述祂有形像，而是显明祂与圣父同等尊荣；同样，说\BibleQuote{高天}时，他并没有把祂局限于那里，而是表达了祂高于万有、升到万有之上。也就是说，祂到达了圣父的宝座本身：因此正如圣父在高天，祂也在高天。因为\Quote{同坐}无非意味着同等尊荣。但如果他们说，祂说\BibleQuote{坐下}，我们可以问他们：那么怎样？祂对站着的祂说话吗？况且，他没有说祂命令，没有说祂吩咐，而是说\BibleQuote{祂说}：没有别的原因，只是免得你以为祂是无始无终的。因为祂之所以这样说，从祂坐的地方可以明显看出。因为若他意在表明卑下，就不会说\BibleQuote{右边}，而会说左边。\par

% structure: p0023 source=h2 target=subsection reason=relative-heading-rank; base=h2

\subsection{希伯来书 1:4}

\BibleQuote{祂比天使更尊贵，正如祂所承受的名号比他们的更卓越。} 这里的\BibleQuote{成为}是代替\Quote{被显明}，如有人所说。然后他从何大胆推理？从这名号。你看见那\BibleQuote{圣子}的名号常用来宣告真实的关系吗？事实上，如果祂不是真正的圣子（而\Quote{真实}无非是\Quote{出于祂}），他如何能从此大胆推理？因为如果祂仅因恩宠而成为圣子，祂不仅不比天使更尊贵，反而比他们更微小。如何呢？因为义人也称为子；而\BibleQuote{圣子}的名号若不是真实的圣子，就不能显出\Quote{卓越}。当他想要指出受造物与创造者之间有某种区别时，请听他所说：\par

% structure: p0025 source=h2 target=subsection reason=relative-heading-rank; base=h2

\subsection{希 1:5}

\BibleQuote{因为祂曾对哪一个天使说过：\InnerQuote{祢是我的儿子，我今天生了祢。}又说：\InnerQuote{我要作祂的父亲，祂要作我的儿子。}}？这些话确实也指着肉体说的：\BibleQuote{我要作祂的父亲，祂要作我的儿子。}——而\BibleQuote{祢是我的儿子，我今天生了祢。}这句话无非是说\Quote{从天主存在的那一刻起}。因为正如祂从现时开始存在（这比任何其他说法更适合祂），同样，\Quote{今天}这个词似乎也是指着肉体说的。因为一旦祂取得了肉体，祂就从此放胆直言。因为肉体分享崇高的事，正如天主性分享卑微的事。那不以成为人为耻、不拒绝现实的那一位，又怎会拒绝这些表达呢？\par

既然我们知道这些事，就让我们不以任何事为耻，也不存有任何高傲。因为如果祂自身是天主、主、天主之子，却不拒绝取奴仆的形态，那么我们更应该做一切事，即使它们卑微。因为请告诉我，人啊，你从哪里得到高傲？从现世的事物吗？但这些事物尚未显现就已消逝。还是从神性的事物？不，这不正是神性的一个优点——没有高傲。\par

那么你为什么还要心怀高傲呢？因为你行事正直吗？听基督说：\BibleQuote{当你们做完了一切事，要说：\InnerQuote{我们是无用的仆人，因为我们做了我们该做的事。}}\BibleRef{路 17:10}\par

或者因为你的财富而高傲？难道你没有看见你面前的人，他们如何赤裸而孤寂地离去？我们不是赤身进入生命，也将赤身离去吗？拥有别人的东西，谁又会高傲呢？因为那些只为自己享乐而使用它们的人，往往在死亡之前就不情愿地被剥夺，在死亡时更是如此。但（有人说）我们活着时可按己意使用它们。首先，你很难看到任何人按己意使用他所拥有的。其次，即使一个人按己意使用东西，这也不是什么大事：因为现在的时间与无尽的岁月相比是短暂的。人啊，你因为富有而高傲吗？为了什么？为了什么原因？因为这也发生在强盗、小偷、杀人犯、柔弱者、淫乱者和各种恶人身上。那么你为什么要高傲？因为如果你正当地使用了它，你就不应高傲，以免亵渎诫命；但如果你不正当使用，那么你确实成了金钱和财物的奴隶，并被它们支配。因为请告诉我，如果一个人患了热病，喝了大量的水，这水暂时解渴，但随后点燃火焰，他应该高傲吗？如果有人无缘无故地有许多挂虑，他应该因此高傲吗？为什么？因为你有很多主人？因为你有无数挂虑？因为有许多人会奉承你？[当然不是。]因为你甚至是他们的奴隶。为了向你证明这一点，请听清楚。我们内在的其他情感，在某些情况下是有用的。例如，愤怒常常是有用的。因为（他说）\BibleQuote{不义的愤怒不会是无辜的}（\BibleRef{德 1:22}），所以人有可能正义地发怒。再者，\BibleQuote{向弟兄发怒而无故的，难免地狱的火}（\BibleRef{玛 5:22}）。又如，竞争、欲望（是有用的）：当欲望指向生育子女时，当竞争指向卓越之事时。正如保禄也说：\BibleQuote{在好事上常有热心是好的}（\BibleRef{迦 4:18}），以及\BibleQuote{渴慕更好的恩赐}（\BibleRef{格前 12:31}）。因此两者都是有用的；但傲慢的精神绝无益处，总是无益和有害的。\par

然而，如果人必须骄傲，[让它]因贫穷，而非因财富。为什么？因为那能以少量为生的人，远大于且优于不能的人。因为请告诉我，假设有人被召到京城，如果其中一些人既不需要牲畜、奴隶、伞、旅舍、凉鞋、器皿，只要有面包、从井中取水就足够；而另一些人却说：\Quote{除非你们提供交通工具、柔软的床，我们不能来；除非我们有许多随从，除非我们被允许不断休息，我们不能来；除非我们使用牲畜，除非我们每天只走一小段路，我们也不能来；我们还需要许多其他东西。}我们应该钦佩谁？这些人还是那些人？显然，是那些什么都不需要的人。同样，在这里：有些人在今生旅程中需要许多东西；另一些人，什么都不需要。因此，如果骄傲合适的话，更应当因贫穷而骄傲。\par

\Quote{但穷人，}他们说\Quote{是可鄙的。}不是他，而是那些鄙视他的人。因为我为什么不鄙视那些不知道如何钦佩当钦佩之事的人呢？因为，如果一个人是画家，只要他们没有受过教导，他将会嘲笑所有嘲笑他的人；他不理会他们所说的话，而是满足于自己的见证。我们能依赖多数人的意见吗？因此，当人们因我们的贫穷而鄙视我们时，我们值得被鄙视，而我们却不鄙视他们，也不称他们为可怜的人。\par

我不说财富产生了多少罪，贫穷产生了多少善。但不如说，财富与贫穷本身并无可称道之处，而在于使用它们的人。基督徒在贫穷中比在财富中更显光辉。如何？他将更少傲慢，更清醒，更庄重，更公正，更体谅；但富有之人，却有许多障碍妨碍这些美德。那么让我们看看富人所做的，更确切地说，滥用财富之人所做的。这样的人行掠夺、欺诈、暴力。不端的爱情、不圣洁的结合、巫术、毒害，以及所有其他恐怖之事——难道你不发现它们都由财富所生吗？你看见了吗，在贫穷中而非在财富中，追求德行更为轻松？因为，我恳求你们，不要以为富人不受今世的惩罚，他们就不犯罪。倘若富人易于受罚，你们定会发现监狱里挤满了他们。但在财富的其他邪恶之外，它还有这一点：拥有它的人，作恶而不受惩罚，将永不止息，反而会受无法医治的创伤，无人能勒住他。\par

此外，若有人选择，他会发现贫穷甚至为我们提供更多享乐的资源。如何？因为它免于忧虑、仇恨、争斗、争论、纷争，以及数不清的邪恶。\par

因此，我们不要追逐财富，也不要永远嫉妒那些拥有很多的人。但让我们这些拥有财富的，善用它们；那些没有的，不要为此忧伤，反而要为一切感谢天主，因为祂使我们能以较少的辛劳获得与富人同等的赏报，甚至（若我们愿意）更大的赏报；从微薄的资本我们将获得巨大的收益。因为那领了两个塔冷通的人，与那领了五个的人，受到同等的钦佩和尊荣。为何？因为他虽然只受托两个塔冷通，却竭尽全力，使所托付的加倍。那么，当我们能以少许收获同样甚至更大的果实，为什么还渴望受托更多呢？劳动确实更少，而赏报却更多？因为穷人的施舍比拥有许多大财产的富人更容易。什么？你们岂不知，人拥有的越多，他越会倾注其爱吗？因此，免得这事临到我们，我们不要追求财富，也不要对贫穷不耐烦，更不要急于致富；让我们这些拥有（财富）的，就按\Person{保禄}所命令的那样去使用它们。\BibleRef{格前 7:29}好使我们获得所应许的福乐。愿我们众人靠着我们的主耶稣基督的恩宠和爱，都能获得它们；愿光荣、权能、尊崇归于祂，与圣父及圣神，从今世直到永远，世世无尽。阿们。\par

\SourceRecord{Translated by Frederic Gardiner.；Nicene and Post-Nicene Fathers, First Series；Vol. 14.；Edited by Philip Schaff.；Buffalo, NY: Christian Literature Publishing Co.,；1889.；<http://www.newadvent.org/fathers/240202.htm>.}

% structure: p0001 source=h1 target=section reason=page-title

\section{希伯来书讲道词 第三篇}

% structure: p0002 source=h2 target=subsection reason=relative-heading-rank; base=h2

\subsection{希伯来书 1:6-8}

\BibleQuote{再者，当天主引领首生者进入世界时，祂说：\InnerQuote{天主的众天使都要朝拜祂。}论到天使，祂说：\InnerQuote{祂使祂的天使成为风，使祂的仆役成为火焰。}但论到子，祂说：\InnerQuote{天主，祢的宝座永永远远。}}\par

1. 我们的主耶稣基督称祂在肉身中的来临为一次\Quote{出去}：正如祂说：\BibleQuote{撒种的人出去撒种。}\BibleRef{玛 13:3}又说：\BibleQuote{我从父那里出来，来到了世上。}\BibleRef{若 16:28}在众多地方都可以看到这一点。但保禄却称之为\Quote{进入}，说：\BibleQuote{再者，当天主引领首生者进入世界时，}此处\Quote{引领进入}是指祂取了肉身。\par

那么，他为何这样使用这个表达呢？所表明的事是明显的，以及它在哪方面是这样说的。因为基督确实称它为\Quote{出去}，恰如其分；因为我们原是远离天主的。正如在皇家庭院中，囚犯和得罪君王的人站在外面，而那位愿意使他们和解的人，并不带他们进去，而是自己出去与他们交谈，直到使他们适合朝见君王，才带他们进去——基督也这样做了。祂出去到我们这里，即取了肉身，并与我们谈论君王的事，然后祂洁净了罪过，并成就了和解，从而带领我们进入。因此，祂称它为\Quote{出去}。\par

但保禄却从那些继承产业、领受任何一份或一份产业的人的比喻中称之为\Quote{进入}。因为这话：\BibleQuote{再者，当天主引领首生者进入世界时，}意思是\Quote{当祂把世界交在祂手中时。}因为当祂被显明时，祂也就拥有了世界的全部。他说的这些话并不是关于天主圣言，而是关于那按肉身而言的。因为若按若望所说：\BibleQuote{祂在世界，世界也是借着祂造的}\BibleRef{若 1:10}，祂又怎能被\Quote{引领进入}呢？不是藉着肉身吗？\par

他又说：\BibleQuote{天主的众天使都要朝拜祂。}既然他即将说出一件伟大而崇高的事，他便预先预备，使它能为听众所接受，因为他把圣父说成是\Quote{引领进入}圣子的。他上面已经说过，天主\Quote{对我们说话，不是藉着先知，而是藉着祂的圣子}；圣子超越天使，并且他从[儿子]这名号上确立这一点。接着，在以下的内容中，又从另一事实确立。那么，这事实是什么呢？就是来自朝拜。他指出祂的伟大程度，如同主人大于奴仆一样；正如任何一个人把另一个人引进家中，立刻命令家中的人向他致敬一样；论到肉身，他说：\BibleQuote{天主的众天使都要朝拜祂。}\par

那么，只是天使吗？不；听下面的话：\BibleQuote{论到天使，祂说：\InnerQuote{祂使祂的天使成为风，使祂的仆役成为火焰。}但论到子，祂说：\InnerQuote{天主，祢的宝座永永远远。}}看，最大的差别！那些是受造的，但祂是非受造的。论到天使，祂说\Quote{使}；那么为什么论到子，祂没有说\Quote{使}呢？虽然他本可以用以下方式表达区别：\BibleQuote{论到天使，祂说：\InnerQuote{祂使祂的天使成为风}，但论到子，\InnerQuote{上主创造了我}；\InnerQuote{天主已立祂为主为基督。}}\BibleRef{箴 8:22}但前一句不是论到子说的，后一句也不是论到天主圣言，而是论到肉身。因为当他想要表达真正的区别时，他不再只包括天使，而是包括上面所有的服役之力。你看见他如何区分，以及以多大的清晰度区分别受造物与造物主、仆役与主人、继承者与真儿子以及奴仆吗？\par

2. \BibleQuote{但论到子，祂说：\InnerQuote{天主，祢的宝座永永远远。}}看，这是王权的象征。\BibleQuote{公义的权杖是祢国度的权杖。}看，这是另一个王权的象征。\par

% structure: p0010 source=h2 target=subsection reason=relative-heading-rank; base=h2

\subsection{希伯来书 1:9}

然后，论到肉身：\BibleRef{希 1:9}\BibleQuote{祢喜爱正义，憎恨不法，因此天主，祢的天主，用喜乐之油傅了祢。}\par

\Quote{祢的天主}是什么意思？为什么在他说出这样一句伟大的话之后，他又加以限定？在这里，他同时击中了犹太人和撒莫萨塔的保禄的追随者、亚略派、马赛鲁、萨伯略以及马西昂。如何？对于犹太人，他通过指明两位——天主与人——来打击他们；对于其他犹太人，即撒莫萨塔的保禄的追随者，他通过这样谈论祂的永恒存在和非受造的本质来打击他们：因为为了与\Quote{祂使}这个词相对，他放上了\Quote{天主，祢的宝座永永远远}。对于亚略派，既有这一点，又加上祂不是奴仆；但若是一个受造物，祂就是奴仆。对于马赛鲁和其他的，这两位在存有上是有分别的。对于马西昂派，天主性没有被傅油，而是人性被傅油。\par

接着他说：\BibleQuote{超过祢的同伴。}但祂的\Quote{同伴}是谁呢？不是人吗？这就是说，基督领受了\BibleQuote{圣神，没有限量地}\BibleRef{若 3:34}。你看见吗？在关于祂非受造本性的教导中，他常常也结合关于\OriginalTerm{救恩计划}{Economy}的教导。有什么比这更清楚的呢？你看见受造的和生发的不是同一回事吗？否则，他就不会做出区分，也不会为了与\Quote{祂使}这个词相对，加上\Quote{但论到子，祂说：\InnerQuote{天主，祢的宝座永永远远。}}他也不会把\Quote{儿子}这个名字称为一个\Quote{更优越的名号}，如果它是同一回事的标志的话。因为什么是优越呢？因为如果受造的和生发的是同一回事，而他们[天使]是被造的，那么祂有什么\Quote{更优越}呢？看！又有\OriginalTerm{神}{\GreekText{Θεὸς}}，带有冠词。\par

% structure: p0014 source=h2 target=subsection reason=relative-heading-rank; base=h2

\subsection{希伯来书 1:10-12}

3. 祂又说：\BibleRef{希 1:10}\newline{}\BibleQuote{主啊，在起初祢奠定了大地，诸天是祢双手的化工。它们将要灭亡，而祢却常存；它们要像衣裳一样陈旧，祢要把它们卷起，像更换外衣一样，它们就被更换了。但祢却是一样的，祢的岁月不会穷尽。}\par

免得你们听到\BibleQuote{当祂引领首生子进入世界时}这句话，就以为这是后来加给祂的恩赐；祂在上文已预先纠正了这一点，如今又进一步修正说\BibleQuote{在起初}：不是现在，而是从起初。请看，祂又致命地打击了撒摩沙塔的保禄和亚略，将有关圣父的话应用于圣子。同时，祂还顺便暗示了另一件更重大的事。因为祂无意中也指出了世界的改变，说：\BibleQuote{它们要像衣裳一样陈旧，祢要把它们卷起，它们就被更换了。}祂在致罗马人书里也说了同样的话，即祂要改变世界（\BibleRef{罗 8:21}）。为表明其容易，祂接着说，如同人卷起一件衣裳，祂也要卷起并改变它。如果祂如此轻易地完成改变和造化，使之变得更好更完善，难道祂还需要另一位来成就低级的造化吗？你们的无耻要到何时？同时，这也是极大的安慰，知道事物不会永恒不变，而是都要接受改变，都要更新；但祂自己却永远存在，永无止境地生活着：祂说：\BibleQuote{祢的岁月不会穷尽}。\par

% structure: p0017 source=h2 target=subsection reason=relative-heading-rank; base=h2

\subsection{希伯来书 1:13}

4. \Quote{然而，祂曾对哪一个天使说过：\BibleQuote{你坐在我的右边，等我使你的仇敌作你的脚凳}？}\par

这又属于主权、同等尊严、尊荣而非软弱，因为圣父为圣子所受的待遇而发怒。这属于圣父对圣子极大的爱和尊荣，如同父亲对儿子一般。因为为祂发怒的人，岂能与祂疏远？祂在第二篇圣咏中也说：\BibleQuote{那住在天上的，必嘲笑他们，上主必愚弄他们。那时祂必在怒中对他们讲话，在烈怒中使他们惊慌。}（\BibleRef{咏 2:4}）祂自己又说：\BibleQuote{至于那些不要我作他们君王的人，带到这里来，在我面前杀掉他们。}（\BibleRef{路 19:27}）因为这是祂自己的话，也听听祂在另一处所说的：\BibleQuote{我多少次愿意聚集你们的孩子，你们却不愿意！看，你们的房屋将留给你们荒芜。}（\BibleRef{路 13:34}）又说：\BibleQuote{天国必从你们手中夺去，赐给那能结果实的民族。}（\BibleRef{玛 21:43}）又说：\BibleQuote{凡跌在这石头上的，必被摔碎；这石头落在谁身上，就把他压得粉碎。}（\BibleRef{玛 21:44}）此外，祂是在那世界中审判他们的，所以祂在此世也亲自报应了他们。因此，\BibleQuote{等我使你的仇敌作你的脚凳}这句话，只是表达对圣子的尊荣。\par

% structure: p0020 source=h2 target=subsection reason=relative-heading-rank; base=h2

\subsection{希伯来书 1:14}

\BibleQuote{众天使岂不都是服役的、奉派为那些将要承受救恩的人效劳吗？}祂说，如果连他们都为我们的救恩效劳，那么他们为圣子效劳又有什么奇怪呢？请看祂如何提高他们的心志，并显示天主对我们的大爱，因为祂把高于我们的天使派来为我们服务。仿佛有人说，祂使用他们正是为此；这就是天使的职务：为我们的救恩而事奉天主。因此，为弟兄的得救而做一切，是天使的工作；其实更是基督自己的工作，因为基督以主的身分拯救，而天使以仆人的身分拯救。而我们虽然是仆人，却是天使的同工。祂说，你们为什么如此热切地注视天使呢？他们是天主圣子的仆人，为我们的缘故被派遣多方奔走，服务于我们的救恩。所以他们与我们是同工的。\par

看祂如何不给予受造物的种类以巨大差别。虽然天使与人之间的差距很大，祂却把他们拉近我们，几乎是在说：他们为我们劳苦，为我们奔波，可以说是服侍我们。这就是他们的职务，为我们而到处被派遣。\par

这些例子在旧约和新约中比比皆是。当天使向牧羊人、或向玛利亚、或向若瑟报喜时；当他们坐在坟墓旁；当他们被派去对门徒说：\BibleQuote{加里肋亚人，你们为什么站着望天呢？}（\BibleRef{宗 1:11}），当他们释放伯多禄出狱，当他们与斐理伯交谈时，请想想这多么大的尊荣：当天主派遣祂的天使像对待朋友一样服务时；当向科尔乃略显现时；当将众宗徒从监狱中领出并说：\BibleQuote{你们去，站在圣殿里，把一切有关这生命的话讲给百姓听。}（\BibleRef{宗 5:20}）；连保禄自己也有天使显现。你看见他们奉天主的命为我们服务，而且在最重要的事上为我们服务吗？因此保禄说：\BibleQuote{一切都是你们的，无论是生，是死，是现世的，是将来的。}（\BibleRef{格前 3:22}）\par

那么，圣子也被派遣，但不是作为仆人，也不是作为仆人，而是作为子、独生子，并愿与圣父相同。更确切地说，祂不是\Quote{被派遣}：因为祂没有从一处到另一处，而是取了肉身；而这些天使改变地方，离开原来所在之处，来到他们原先不在的地方。\par

借此，祂又顺便鼓励他们说：你们怕什么？天使正为我们服务。\par

% structure: p0026 source=h2 target=subsection reason=relative-heading-rank; base=h2

\subsection{希伯来书 2:1}

5. 在论及圣子时，既讲述了关于\OriginalTerm{救恩工程}{Economy}的事，也讲述了关于创造的事，以及关于祂的主权的事；并且表明了祂同等尊荣，作为绝对的主宰，祂不仅治理人类，也治理天上的权能。之后，当他完成他的论证后，便劝勉他们：我们应当更留心所听的事。\BibleRef{希 2:1} \Quote{因此，我们必须更专心留意所听见的事}（他说）\Quote{，免得随流失去。} 为什么说\Quote{更}？这里他意指\Quote{更}甚于法律；但他在表述上隐去了这一点，然而在推理的过程中，他以劝告或劝勉的方式使之明确。因为这样更好。\par

% structure: p0028 source=h2 target=subsection reason=relative-heading-rank; base=h2

\subsection{\BibleRef{希 2:2-3}}

\Quote{因为藉天使所传的话}（他说）\Quote{既是确定的，凡干犯悖逆的，都受了该受的报应；那么，我们若忽略了这么大的救恩，怎能逃罪呢？这救恩起初是由主亲自宣讲的，后来由听见的人给我们证实了。}\par

为什么我们应当\Quote{更专心留意所听见的事}？难道从前的事不也是出于天主，正如这些事一样吗？那么，他要么是意指\Quote{更}甚于法律，要么是意指\Quote{非常专心}，而非作比较——绝无此意。因为由于旧约历时长久，他们极为看重它，而这些事因为新近而被轻视，他便证明（超过他论证所需）我们更应该留意这些事。如何证明？他实质上是在说：这些和那些都出于天主，但并非同等方式。这一点他后来向我们显明；但目前他处理得较为表面，但后来更清楚地说：\Quote{如果那前约没有瑕疵}\BibleRef{希 8:7}，以及许多其他类似的话：\Quote{因为那渐渐衰残和变旧的，就快要消失了。}\BibleRef{希 8:13} 然而，在他演讲的开头，他尚不敢说任何这类话，直到他先用更充分的论证占据了听众的心。\par

那么，为什么我们应当\Quote{更专心留意}？\Quote{免得随流失去}，他说——也就是说，免得我们灭亡，免得我们坠落。他在这里表明了这坠落的严重性，即坠落后再返回是困难的，因为这是出于故意的疏忽。他借用了箴言的说法。因为他说：\Quote{我儿，〔要谨守〕，免得你疏失}\BibleRef{箴 3:21}，既表明了坠落之易，也表明了毁灭之重。也就是说，我们的悖逆并非没有危险。而通过他的推理方式，他表明惩罚更大，但再次以疑问的形式留下，而非下结论。因为这正是使言论不冒犯人的方式——不总是自己下判断，而是让听众自己来判决；这会使他们更易被说服。先知纳堂在旧约中也是这样做的，而基督在玛窦福音中也说：\Quote{对那园户，他将会怎样做？}\BibleRef{玛 21:40} 祂强迫他们自己下判决：这就是最大的胜利。\par

接着，当他说了\Quote{因为藉天使所传的话既是确定的}——他没有加上\Quote{那么更藉基督所传的}，而是略过这一点，说了较轻微的：\Quote{我们若忽略了这么大的救恩，怎能逃罪呢？} 请看他是如何作比较的。\Quote{因为藉天使所传的话}，他说。那里是\Quote{藉天使}，这里是\Quote{藉主}——那里是\Quote{话}，这里却是\Quote{救恩}。\par

接着，恐怕有人说：保禄啊，你的话是基督的话吗？他证明这些话语的可靠性，既因为他是从主亲自听见的，又因为现在是由天主所说的；因为不仅仅是传下来的声音，如梅瑟的情形，而是有神迹发生，有事实作证。\par

6. 但\Quote{因为藉天使所传的话既是确定的}是什么意思？因为他在致迦拉达人书中也这样说：\Quote{由天使藉中保之手设立的。}\BibleRef{迦 3:19} 又说：\Quote{你们领受了天使所传的法律，却没有遵守。}\BibleRef{宗 7:53} 他在各处都说法律是由天使赐下的。有些人确实认为这指的是梅瑟；但这是没有理由的。因为这里他说的是复数\Quote{天使}；并且他这里所说的天使是那些天上的天使。那么是什么意思呢？要么他只指十诫（因为在那里梅瑟说话，天主回答他\BibleRef{出 19:19}）——要么是天使在场，天主布置他们——或者他这样说是就旧约中所有说过和做过的事而言，仿佛天使在其中都有份。但为什么在另一处说\Quote{法律是由梅瑟颁布的}\BibleRef{若 1:17}，而这里说\Quote{由天使}呢？因为经上说：\Quote{天主在密云中降临。}\BibleRef{出 19:16}\par

\Quote{因为藉天使所传的话既是确定的} 什么是\Quote{既是确定的}？可以说是真实、在其适当季节中信实的；一切所说的话都应验了。或者是这个意思，或者是说它们有效，恐吓即将成就。或者\Quote{话}指诫命。因为除了法律之外，天主差遣的天使也命令了许多事：例如在波津、在\WorkTitle{民长纪}中、在〔三松的故事〕中。\BibleRef{民 2:1} 这就是为什么他不说\Quote{法律}而说\Quote{话}的原因。我似乎觉得他更可能是指这一点，即那些委托给天使管理的事。那么我们要说什么呢？当时奉命管理那民族的使者在场，他们自己吹号，并行了其他事，如火焰、密云等。\BibleRef{出 19:16}\par

\Quote{每样过犯和悖逆}他说。不是这个或那个，而是\Quote{每样}一个。他说，没有一样不得到报复，反而\Quote{得了应得的报偿}，而不是说惩罚。他为何这样说？这是\Person{保禄}的风格，不注重措辞，而是随意使用带有恶意的词语，即使在善意的事情上也是如此。正如在另一处他说：\BibleQuote{将一切思想掳去，使之服从基督}。\BibleRef{格后 10:5} 他又用\Quote{报偿}来表示惩罚，正如这里他将惩罚称为\Quote{奖赏}。\BibleQuote{如果这是公义的事}他说，\BibleQuote{天主将患难报应那加患难给你们的人，而将安宁赐给你们这些受患难的人}。\BibleRef{得后 1:6} 也就是说，公义没有被违反，而是天主出去反对他们，使惩罚临到罪人身上，尽管并非他们所有的罪都被显明，而只是明显违背诫命的地方。\par

\Quote{我们怎能逃脱，}他说，\Quote{如果我们忽略如此伟大的救恩？}由此他表明，先前的救恩并非伟大之事。他也很好地加上了\Quote{如此伟大}。因为（他说）祂现在不会从战争中拯救我们，也不会赐给我们大地和地上的美物，而将是死亡的解体，魔鬼的毁灭，天国，永生。所有这些他都简要表述了，说：\Quote{如果我们忽略如此伟大的救恩}。\par

7. 然后他补充了使这值得相信的事。\Quote{这最初是由主开始传讲的}：也就是说，它的源头本身就是源泉。不是一个人将它带到地上，也不是任何受造之力，而是独生子自己。\par

\Quote{并由听到祂的人向我们证实了。}什么是\Quote{证实}？是指被相信了，或是成就了。因为（他说）我们有了定金；也就是说，它没有熄灭，没有停止，反而是坚强有力的。原因是神圣的能力在其中运行。意思是，那些从主听到的人，他们亲自向我们证实了。这是一件伟大而可信的事：对此\Person{路加}也在他的\WorkTitle{福音}开头说：\BibleQuote{正如那些从起初亲眼见过并为圣言服务的人所传给我们的。}\BibleRef{路 1:2}\par

% structure: p0040 source=h2 target=subsection reason=relative-heading-rank; base=h2

\subsection{希伯来书 2:4}

那么，它是如何被证实的？如果那些听到的人是造假者呢？有人可能会说。然后他推翻了这一异议，并表明这恩宠并非人为。如果他们误入歧途，天主就不会为他们作证；因为他补充说\BibleRef{希 2:4}，\BibleQuote{天主也同他们一起作证。}他们固然作证，天主也作证。祂如何作证？不是用言语或声音（尽管这也值得相信）：而是如何？\Quote{藉着征兆、奇事和各种异能。}（他说得好，\Quote{各种异能，}表明恩赐的丰盛：这在以前的律法中并非如此，也没有如此伟大和多样的征兆。）也就是说，我们不是简单地相信他们，而是通过征兆和奇事：因此我们不是相信他们，而是相信天主自己。\par

\Quote{并藉着圣神的恩赐，按祂自己的旨意。}\par

那么，如果巫士也行征兆，而犹太人曾说祂\Quote{藉着贝耳则步驱魔}呢？\BibleRef{路 11:15} 但他们不行这样的征兆：因此他说\Quote{各种异能}：因为那些不是异能（或能力），而是软弱和幻想，完全虚妄的事。因此他说：\Quote{藉着圣神的恩赐，按祂自己的旨意。}\par

8. 这里他似乎在向我暗示更进一步的事。因为那里有许多人拥有恩赐是不大可能的，但这些恩赐随着他们变得更为懒散而消失了。那么，为了甚至在这件事上安慰他们，不让他们堕落，他将一切归因于天主的旨意。他知道（他说）什么是有益的，为了谁，并相应地分施祂的恩宠。这也正是他（\Person{保禄}）在\WorkTitle{格林多前书}中所做的，说：\BibleQuote{天主将我们每一个人安置在教会中，如祂所喜悦的。}\BibleRef{格前 12:18} 又说：\BibleQuote{圣神的显示是为使众人得益。}\BibleRef{格前 12:7}\par

\Quote{按祂的旨意。}他表明恩赐是按照圣父的旨意。但往往由于他们不洁和懒散的生活，许多人没有领受恩赐，有时那些生活良好纯洁的人也没有领受。为什么？我请问你？免得他们变得傲慢，免得他们自高自大，免得他们变得更疏忽，免得他们更兴奋。因为即使没有恩赐，仅仅是纯洁生活的自觉就足以使人自高，更何况恩宠加添时呢。因此，恩赐较多赐予谦卑的人、纯朴的人，尤其是纯朴的人：因为经上说：\BibleQuote{以纯朴和喜乐的心。}\BibleRef{宗 2:46} 是的，他藉此更催促他们，如果他们变得懒散，就给他们一个刺激。因为谦卑的人，以及不把自己看得过高的人，在领受恩赐后变得更加热切，因为他得到了超出自己应得的东西，并认为自己不配。但那些自认为做得不错的人，将恩赐视为自己的功劳，就会自高自大。因此，天主有益地分施这些：人们也可以在教会中看到：一人有教导的话语，另一人没有能力开口。不应让这个人（他说）因此忧伤。因为\BibleQuote{圣神的显示是为使众人得益。}\BibleRef{格前 12:7} 因为如果一个家主知道该把什么东西托付给谁，更何况天主，祂明白人的心思，\Quote{祂在万物存在之前就知道一切。}只有一件事值得忧伤：罪，此外别无他物。\par

不要说：\Quote{我为什么没有财富？}或\Quote{如果我有了财富，我就会施舍给穷人。}你不知道，如果你有了财富，恐怕你反而会变得贪婪。因为现在你确实说这些话，但一旦受到考验，你就会是另一个样子。正如当我们饱足时，我们以为自己能禁食；但当我们稍有空缺时，别的念头就进入我们里面。同样，当我们不接近烈酒时，我们以为自己能控制食欲；但一旦被它抓住，就不再如此了。\par

不要说：\Quote{我为什么没有教导的恩赐？}或\Quote{如果我有了它，我本该教化无数的灵魂。}你不知道，如果你有了它，它会不会成为你的定案——是否嫉妒、是否懒惰会使你埋藏你的塔冷通。现在，你确实从这一切中自由，尽管你没有给出\BibleQuote{食粮的分量}\BibleRef{路 12:42}，你也不被追究；但那时，你就要为许多人负责。\par

9. 此外，你现在也并非没有恩赐。在微小的事上显出，如果你有了别的恩赐，你会怎样。\BibleQuote{因为}（他说）\BibleQuote{你们在小事上不忠信，谁会把大事给你们呢？}\BibleRef{路 16:11} 要像那位寡妇那样证明自己：她有两个小钱，她把所有的一切都投进去了。\par

你寻求财富吗？证明你轻视这些小事，好让我也在大事上信任你。但如果你连这些小事都不轻视，那么你更不会轻视那些大事了。\par

同样，在言语上，证明你能恰当地运用劝勉和劝告。你没有外在的口才吗？你没有丰富的思想吗？但无论如何，你知道这些平常的事。你有一个孩子，你有一个邻人，你有一个朋友，你有一个兄弟，你有亲属。虽然你不能在教会公开地讲长篇道理，但你可以在私下劝勉这些人。这里不需要修辞，也不需要精心的演说：在这些事上证明，如果你有口才，你本不会忽略它。但如果你在小事上不认真，我怎能在大事上信任你呢？\par

因为，每个人都能做到这一点，请听保禄所说的，他如何嘱咐甚至在俗的人：\BibleQuote{你们要彼此劝勉}（他说）\BibleQuote{正如你们现在所做的。}\BibleRef{得前 5:11} 以及\BibleQuote{你们要用这些话彼此安慰。}\BibleRef{得前 4:18} 天主知道如何分施给每个人。你比梅瑟更好吗？请听他是如何退缩以回避艰难的。\BibleQuote{我怎么能独自承担他们呢？}（他说）\BibleQuote{因为你对我说：你抚养他们，如同乳养的父亲哺育婴儿一样。}\BibleRef{户 11:12} 那么天主做了什么？祂取了他的神，分给了其他人，表明即使他承担他们，那恩赐也不是他自己的，而是圣神的。如果你有了恩赐，你也许会骄傲，也许会走岔路。你不知道自己，正如天主知道你一样。让我们不要说：\Quote{那有什么用？}\Quote{那为什么？}当天主分配时，我们不要向祂要求解释：因为这是极端不虔敬和愚昧的事。我们是仆人，是远离我们主人的仆人，连我们面前的事也不知道。\par

10. 那么，我们不要忙于探究天主的计划，而是凡祂所赐予的，我们就要守护，即使它微小，即使它是最低下的，我们就会完全被悦纳。或者更确切地说，天主的恩赐没有一样是小的：你因为缺乏教导的恩赐而悲伤吗？那么告诉我，你认为哪个更大：是教导的恩赐，还是驱除疾病的恩赐？毫无疑问是后者。但且告诉我：使盲人复明，不也比驱除疾病更大吗？但且告诉我：使死人复活，不也比使盲人复明更大吗？再告诉我：借着影子和手帕做到这一点，不也比用言语做到更大吗？那么告诉我，你愿意选择哪一样？是借着影子和手帕使死人复活，还是拥有教导的恩赐？你肯定会说前者，即借着影子和手帕使死人复活。那么，如果我向你证明，还有另一种恩赐远大于此，而且你可以在自己有能力获得它时却得不到它，那么你被剥夺那些其他的恩赐岂不是公正的吗？而且这恩赐不是一两个人，而是所有人都可以拥有的。我知道你们张大了嘴，感到惊奇，因为你们听说你们有能力获得比使死人复活、使盲人复明更大的恩赐，能做在宗徒时代所做的同样的事。而这对你们来说似乎是难以置信的。\par

那么，这恩赐是什么呢？爱德。不，请相信我；因为这话不是我的，而是基督藉保禄所说的。因为他说了什么？\BibleQuote{你们要渴求那更大的恩赐；但我还要指示你们一条更优越的道路。}\BibleRef{格前 12:31} 这\BibleQuote{更优越的}是什么意思？他意思是如此。格林多人因自己的恩赐而骄傲，那些有舌音恩赐的人，那最小的恩赐，却自高自大，轻视其余的人。因此他说：你们固然渴望恩赐吗？我指示你们一条道路，不是仅仅胜过，而是远远超越。然后他说：\BibleQuote{我若能说天使的语言，却没有爱德，我什么也不算。我若有全备的信心，甚至能移山，却没有爱德，我什么也不算。}\BibleRef{格前 13:1}\par

你们看到了这恩赐吗？要热切渴求这恩赐。这比复活死人更大，这远超其余一切。确实如此，听基督亲自对门徒说：\BibleQuote{众人因此就认出你们是我的门徒。}\BibleRef{若 13:35} 然后祂指出如何认出，祂没有提奇迹，而是什么？\BibleQuote{如果你们彼此相爱。}祂又对圣父说：\BibleQuote{他们因合一，就认出是祢派遣了我。}\BibleRef{若 17:21} 祂对门徒说：\BibleQuote{我给你们一条新命令：你们要彼此相爱。}\BibleRef{若 13:34} 这样的人因此比复活死人者更可敬、更光荣，理当如此。因为那（复活死人）完全出于天主的恩宠，但这（爱德）也出于你自己的热忱。这属于真正的基督徒：这显示基督、被钉者的门徒，与尘世无分的人。没有爱德，连殉道也无益。\par

作为明证，请看这一点。圣保禄举出两种最高的德行，或更准确地说，三种：即奇迹、知识、生活。若没有爱德，其他的，他说，都算不了什么。我要说它们如何算不了什么。他说：\BibleQuote{我若把我所有的财物施舍给人，却没有爱德，我便算不了什么。}\BibleRef{格前 13:3} 因为即使人施舍穷人、耗尽钱财，也可能没有爱德。\par

11. 事实上，我们在论爱德之处已充分说明了这些，请读者参考那里。同时，如我所说，让我们热切渴求这恩赐，让我们彼此相爱；这样，为圆满获得德行，我们便不需别的，一切都将轻松无劳苦，且我们以极大的勤勉完善地做到一切。\par

但是请看，如今人说，我们彼此相爱。因为一人有两个朋友，另一人有三个。但这并非为天主而爱，而是为被爱。为天主而爱不以这个为爱的原则；这样的人会以对待兄弟的态度对待众人；爱那些有同一信仰者如真弟兄；对异端者、外邦人和犹太人，他们固然是本性上的弟兄，却是卑劣无益的——怜恤他们，耗尽心力，为他们哀哭。在此我们若爱众人，甚至仇敌，就相似天主；而不是若行奇迹。因为我们看到天主行奇事时便赞叹，但更赞叹的是祂对人显示爱，祂的忍耐。如果连在天主内，这也值得极大赞叹，那么在人间，更明显的是这使我们成为可赞叹的。\par

让我们热心追求这事：我们将毫不逊色于圣保禄、圣伯多禄和那些复活了无数死者的人，尽管我们可能无法驱退热症。但没有爱德，即使我们行了比宗徒们更大的奇迹，即使我们为信仰冒了无数危险，对我们也没有任何益处。这些话不是我说的，而是那位培养爱德者所知的。那么让我们顺从祂；这样我们才能获得所应许的美善，愿我们都能有分于此，靠我们的主耶稣基督的恩宠，愿光荣归于祂及圣父和圣神，从现在直到永远，万世无疆。阿们。\par

\SourceRecord{Translated by Frederic Gardiner.；Nicene and Post-Nicene Fathers, First Series；Vol. 14.；Edited by Philip Schaff.；Buffalo, NY: Christian Literature Publishing Co.,；1889.；<http://www.newadvent.org/fathers/240203.htm>.}

% structure: p0001 source=h1 target=section reason=page-title

\section{\BibleBook{希伯来人}{书}讲道 第四篇}

% structure: p0002 source=h2 target=subsection reason=relative-heading-rank; base=h2

\subsection{\BibleBook{希伯来人}{书} 2:5-7}

\BibleQuote{祂没有将我们将来的世界——即我们所谈论的世界——置于天使权下；但有人在某处作证说：\InnerQuote{人算什么，祢竟顾念他？人子算什么，祢竟眷顾他？祢使他稍微逊于天使。}}\par

1. 我本来希望确知，是否有人以应有的热诚聆听所讲的话，我们是否没有将种子撒在路旁：因为若是那样，我会更愉快地讲道。因为我们必须宣讲，即使无人听，这是救主加在我们身上的恐惧。因为祂说：\Quote{你向这百姓作证；即使他们不听，你自己也无罪。} \BibleRef{则 3:19} 然而，若我确信你们的热情，我不仅因恐惧而发言，也会因喜乐而为之。因为现在，即使无人听，即使我的工作——只要我尽了自己的本分——不带来危险，劳作也并不全然愉快。有什么益处呢？当我不被责备，却无人得益时？但若有人注意，我们所得的好处与其说是逃避刑罚，不如说是你们的进步。\par

那么我如何知道呢？我注意到你们中有一些人不太专心，我将在私下遇见他们时拷问他们。若我发现他们记得所讲过的任何事（我不是说全部，因为这对你们不容易），即使只是许多中的几件，显然我将不再怀疑其余的人。的确，我们本应不事先通知，趁你们不备时攻击你们。然而，即使以这种方式，若能达成目的，也足够了。不仅如此，即使现在，我也可以趁你们不备攻击你们。因为我已预先警告你们，我\Emph{要}拷问你们；但\Emph{何时}拷问，我尚未表明。因为可能是今天，可能是明天，可能是二十天或三十天后，可能是更少，可能是更多。同样，天主也使我们的死期不确定，没有让我们清楚知道它是今天临到，还是明天，或是一年后，或是多年后；这样，因着期待的不确定性，我们可以在所有时间里坚守德行。我们确实会离开，祂已说过——但何时，祂尚未说。同样，我也说过我将拷问你们，但未加何时，希望你们始终警觉。\par

不可有人说：\Quote{我四五周前或更早听到这些事，不能记住。}因为我愿听众如此记住，使回想持久而不易消退，也不应否认所讲的话。因为我愿你们记住，不是要告诉我，而是为使你们得益；这对我至关重要。所以，没有人这样说。\par

2. 然而，我现在必须从书信的后续部分开始。那么，今天我们面前要讲的是什么呢？\par

\BibleQuote{因为不是交给天使，}他说，\BibleQuote{祂将我们将来的世界——即我们所谈论的——交予管辖。}难道祂在谈论另一个世界吗？不，而是这个。因此祂加上\BibleQuote{我们所谈论的}，以免人的心思游荡寻找另一个。那么，祂为何称之为\BibleQuote{将来的世界}？正如祂在别处所说：\BibleQuote{他是那要求者之预像}\BibleRef{罗 5:14}，当祂在致罗马人书中论及亚当和基督时；按照肉身称基督为\BibleQuote{那要求者}，就亚当的时代而言（因为那时祂要求到来）。所以现在，既然祂说过：\BibleQuote{但当祂再带领首生者进入世界时}，免得你们以为祂在说另一个世界，这从许多考虑以及祂说\BibleQuote{将来}得到确证。因为世界是将来，但天主子始终存有。所以这个即将到来的世界，祂没有交予天使，而是交予基督。因为显然（祂说）这是指着子说的：因为肯定没有人会主张另一种可能，即指着天使说的。\par

然后他又引出另一个证言并说：\BibleQuote{但有某处作证说。}他为何没有提及先知的名字，反而隐藏呢？是的，在其他证言中他也这样做：例如他说：\BibleQuote{但当祂再带领首生者进入世界时，祂说：\InnerQuote{愿天主所有的天使都崇拜祂。}又说：\InnerQuote{我要作祂的父亲。}至于天使，祂说：\InnerQuote{祂使祂的使者成为风。}又说：\InnerQuote{主啊，起初祢奠定了大地的根基。}}\EditorialNote{参见：一6、5、7、10}——所以此处他也说：\BibleQuote{但有某处作证说。}而这一点（我认为）是隐藏自我的人的表现，也表明他们精通圣经；他不列出说话者，而是将其作为熟悉和显然的内容引出。\par

% structure: p0010 source=h2 target=subsection reason=relative-heading-rank; base=h2

\subsection{\BibleBook{希伯来人}{书} 2:8}

\BibleQuote{人算什么，祢竟顾念他？人子算什么，祢竟眷顾他？祢使他稍微逊于天使，以荣耀和尊贵作他的冠冕。} \BibleRef{希 2:8} \BibleQuote{祢使万物都屈服在他的脚下。}\par

虽然这些话是泛指人性，但更恰当地适用于按照肉身的基督。因为这句话：\BibleQuote{祢使万物都屈服在他的脚下，}是属于祂，而非属于我们。天主子在我们一无所有时眷顾了我们；祂在摄取我们的本性，并使之与祂结合后，就变得比一切都高。\par

\Quote{因为，}他说，\Quote{祂使万物都服从在祂脚下，没有留下一件不服从祂的；但现今我们还看不见万物都服从祂。}他的意思是：——既然他曾说过\BibleQuote{直到我把祢的仇敌放在祢的脚下}\BibleRef{希 1:13}——而他们很可能仍会忧伤——于是他在此之后插入了几句，然后补充了这个见证以证实前者。为使他们不致说：祂怎么把祂的仇敌放在祂脚下，当我们遭受了这么多苦难？他在前一处已经充分暗示了（因为\Quote{直到}一词表明，不是立即发生，而是经过一段时间），但这里他继续跟进。因为你不要以为（他说）他们尚未被制服，就不会被制服；因为必须被制服是明显的；为此，这预言才被宣告。\Quote{因为，}他说，\Quote{祂使万物都服从在祂脚下，没有留下一件不服从祂的。}那么怎么万物还没有服从祂？因为它们日后要服从祂。\par

若万物都必须服从祂，但尚未服从，你不要忧伤，也不要烦恼。若确实到了结局，万物都服从了，而你还受苦，你才应当怨尤；\Quote{但现今我们还看不见万物都服从祂。}君王尚未明显得胜。那么你为何在受苦时烦恼？福音的宣讲尚未得胜一切；还不是万物全然服从的时候。\par

% structure: p0015 source=h2 target=subsection reason=relative-heading-rank; base=h2

\subsection{希伯来书 2:9}

3. 然后又有另一个安慰：若那日后要叫万物服从祂的，自己也曾死去，并忍受了无数的苦难。\BibleRef{希 2:9}\Quote{但，}他说，\Quote{我们看见那位暂时比天使低微的耶稣，因为受死的苦，}——然后又是好事——\Quote{得了光荣与尊荣。}你看见这一切如何适用于祂吗？因为\Quote{暂时}一词更适合祂，祂只在阴府停留了三天，而不适合我们这些长久处于腐败中的人。同样，\Quote{光荣与尊荣}一词也更适用于祂，远胜于我们。\par

他又提醒他们十字架，以此达成两件事：既显示祂的眷顾，又劝勉他们勇敢承受一切，仰望主。（他会说）若那受天使朝拜的，为你们的缘故，忍受了比天使低微，那么你这比天使卑微的，更当为祂忍受一切。然后他表明十字架是\Quote{光荣与尊荣}，正如祂自己常说的：\Quote{使人子得光荣}\BibleRef{若 11:5}；以及\Quote{人子已得了光荣}\BibleRef{若 12:23}。若祂把为仆人受苦称为\Quote{光荣}，那么你更该为主受苦。\par

你看见十字架的果实有多么巨大吗？不要惧怕此事；因为它在你看来固然是凄惨的，却带来无数的好处。从这些思考中，他显示出试炼的益处。然后他说：\Quote{好叫祂因天主的恩宠，为众人尝到死味。}\par

\Quote{好叫因天主的恩宠，}他说。祂因天主对我们恩宠的缘故，受了这些苦。\Quote{祂没有怜惜自己的儿子，}他说，\Quote{反而为我们众人把祂交出来}\BibleRef{罗 8:32}。为什么？祂并不欠我们这恩情，而是出于恩宠做了这事。在致罗马人书里他又说：\Quote{天主的恩宠和那因耶稣基督一人而来的恩赐，更加丰富地倾注在众人身上}\BibleRef{罗 5:15}。\par

\Quote{好叫祂因天主的恩宠，为众人尝到死味，}不仅为信者，也为整个世界；因为祂确是为众人死了；但若众人不信呢？祂已完成自己的部分。\par

他正确地说\Quote{尝到死味}，而没有说\Quote{死}。因为祂仿佛真的品尝了一下，在那里度过短暂时间后，立刻复活了。\par

那么他说\Quote{因受死的苦}，表明是真正的死亡；又说\Quote{超越天使}，表明复活。如医师虽然不需要品尝为病人预备的食物，但因关心病人，自己先尝，以劝病人有勇气放心吃那食物。同样，由于人人都害怕死亡，祂为了劝勉众人勇敢面对死亡，自己也尝了死味，尽管祂不需要。\Quote{因为，}祂说，\Quote{这世界的元首来，在祂身上毫无权柄}\BibleRef{若 14:30}。所以\Quote{因恩宠}和\Quote{为众人尝到死味}这两个说法都证实了这一点。\par

% structure: p0023 source=h2 target=subsection reason=relative-heading-rank; base=h2

\subsection{希伯来书 2:10}

4. \Quote{因为那为万物所属、为万物所本的，在引领许多儿子进入光荣的时候，使救恩的元首藉着苦难得以成全，原是合宜的。}他在这里说的是圣父。你看见他如何又把\Quote{为万物所本}也归于祂吗？若这是卑微的表示，并只适用于子，他就不会这样做了。他所说的是：祂行了与祂对人类之爱相称的事，使祂的首生子比一切更为光荣，并将祂立为其他人的榜样，如同一位超群出众的高贵角力者。\par

\Quote{救恩的元首，}即他们救恩的原因。你看见差距有多么大吗？祂是子，我们也是子；但祂拯救，我们被拯救。你看见祂怎样既把我们联合，又将我们分开；他说\Quote{引领许多儿子进入光荣}，这里他联合我们——\Quote{救恩的元首}，又将我们分开。\par

\Quote{以苦难来成全。}那么苦难是一种成全和救恩的原因。你看见了吗？遭受苦难并非被完全遗弃之人的分，因为天主正是藉此首先荣耀了祂的圣子，带领祂经过苦难。确实，祂降生成人，承受祂所受的苦难，比创造世界、从无中引出万有，是更大的事。创造固然是祂仁爱的标记，但前者更甚。宗徒本人也指出这一点，说：\Quote{为在将来的世代，显示祂恩宠的丰盛，祂在基督耶稣内使我们一同复活，一同坐在天上。} \BibleRef{弗 2:7}\par

\Quote{因为那为万物所属、为万物所本的天主，在带领许多儿子进入光荣时，藉苦难来成全那救恩的元首，原是合宜的。}（他说）因为那温柔眷顾、创造万物的那一位，理应为其余的人交出圣子，以一代多。然而他没有这样表达，而是说\Quote{藉苦难来成全}，表明为别人受苦不仅有益于\Quote{他}，而且他自己也变得更为荣耀、更为成全。这也是针对信众说的，以此安慰他们：因为基督正是在受苦时得荣耀。但当我说祂得荣耀时，不要以为祂增添了荣耀：祂本性所固有的，祂始终拥有，并未增添什么。\par

% structure: p0028 source=h2 target=subsection reason=relative-heading-rank; base=h2

\subsection{希伯来书 2:11-12}

5. 他说：\Quote{因为那祝圣人的和被祝圣的，都是出于一源。为此，祂不以称他们为弟兄为耻。}请看，他又如何将他们联系在一起，尊荣并安慰他们，使他们成为基督的弟兄，就此而言，他们\Quote{出于一源}。然后他又加以防范，表明他所指的是按肉身而言，于是引入\Quote{那祝圣人的}，即基督，\Quote{和被祝圣的}，即我们。你看见这差别是何等巨大吗？祂祝圣，我们被祝圣。上面他说\Quote{救恩的元首}。\Quote{因为天主是独一的，万物都出于祂。} \BibleRef{格前 8:6}\par

\Quote{为此，祂不以称他们为弟兄为耻。}你看见他如何再次显示出超越性吗？因为他说\Quote{不以……为耻}，表明这全然不是出于事物的本性，而是出于那\Quote{不以}任何事物\Quote{为耻}者的深情厚爱，出于祂极大的谦卑。虽然我们是\Quote{出于一源}，但祂祝圣，我们被祝圣，差别是巨大的。而且\Quote{祂}出自圣父，作为真子，即出自祂的实体；我们作为受造物，即从无中被引出，所以差别是巨大的。因此他说：\Quote{祂不以称他们为弟兄为耻} \BibleRef{希 2:12}，\Quote{说：\InnerQuote{我要向我的弟兄们宣扬你的名。}} \BibleRef{咏 21:22} 因为当祂披上肉身时，也披上了弟兄之情，同时进入了弟兄之情。\par

% structure: p0031 source=h2 target=subsection reason=relative-heading-rank; base=h2

\subsection{希伯来书 2:13}

他自然引出这一点。但这里的\Quote{我要依靠祂} \BibleRef{撒下 22:3}，是什么意思？因为下文也是自然引入的。\Quote{看，我和天主所赐给我的孩子们。} \EditorialNote{385 8:18} 因为正如这里祂显示自己是父亲，前面则显示是弟兄。\Quote{我要向我的弟兄们宣扬你的名}，祂说。\par

他又通过下文 \BibleRef{希 2:14} 指出这超越性和巨大间隔：\Quote{既然孩子们都有血肉之躯，}（你看见他指的是相像之处？关于肉身），\Quote{祂也照样亲自成了血肉。} 愿所有异端者羞愧，愿那些说祂只是表象而非实体降临的人掩面。因为他不仅说了\Quote{祂成了血肉}，然后就此打住；尽管这样说已经足够，但他还加了更多，说\Quote{照样}，不是指表象或影像（因为那样\Quote{照样}就不成立了），而是指实在，显示弟兄之情。\par

% structure: p0034 source=h2 target=subsection reason=relative-heading-rank; base=h2

\subsection{希伯来书 2:14}

6. 接着他也列出了救恩计划的原因。他说：\Quote{为能藉死亡，毁灭那握有死亡权势的，即魔鬼。}\par

这里他指出奇事：魔鬼所凭借的，竟被用来击败它；那曾是它对抗世界的强力武器，即死亡，基督却用它击打了魔鬼。由此他展示了征服者的大能。你看见死亡带来了何等大的善吗？\par

% structure: p0037 source=h2 target=subsection reason=relative-heading-rank; base=h2

\subsection{希伯来书 2:15}

\Quote{并要解救那些因怕死而一生处于奴役中的人。}（他说）你们为何战栗？为何害怕那已被消灭的？它不再可怖，已被践踏，已被完全鄙视；它卑鄙且无足轻重。 \BibleRef{弟后 1:10}\par

但\Quote{因怕死而一生处于奴役之下}是什么意思呢？他或者是指：畏惧死亡的人就是奴隶，他宁可忍受一切也不愿死；或者是指：所有人都曾是死亡的奴隶，受其控制，因为死亡尚未被消除；或者是指：人活在持续的恐惧中，时刻预料自己将死，并且害怕死亡，所以当这种恐惧伴随他们时，他们无法感受快乐。因为他暗示了这一点，说：\Quote{一生之久。}他在这里表明：遭受痛苦、折磨、迫害、失去家园、财产及其他一切的人，他们的生活比古代那些享受奢侈、不受此类磨难、持续顺遂的人更为甜美和自由——如果后者确实\Quote{一生之久}处于这种恐惧和奴役之中；而前者已被释放，嘲笑他们所畏惧的事。因为现在就如同：当一个人被带往导致死亡的囚禁中，并一直预期死亡时，有人却用丰盛的美味喂养他（古代的死亡大致如此）；但现在，有人除去那恐惧连同美味，应许一场竞赛，提出一场不再导致死亡、而是导向王国的战斗。你愿意成为哪一类人？是在监狱中被喂养、每天等候判决的人，还是那些竭力奋斗、甘愿劳苦、以便为自己戴上王国冠冕的人？你看到他如何振作他们的灵魂，使他们振奋了吗？他也表明：不仅死亡被消灭了，而且那位永远对我们发动不休战争的——我指的是魔鬼——也因此被消灭了；因为不畏惧死亡的人就不在魔鬼的暴政之下。因为若\Quote{人以皮换皮，为了性命，情愿舍弃一切}（\BibleRef{约 2:4}），那么当一个人决心连性命也不顾时，他还会成为谁的奴隶呢？他无所畏惧，无所恐慌，超乎众人之上，比众人更自由。因为连自己性命都不顾的人，更不顾其他一切。当魔鬼找到这样的灵魂时，他无法在其中做成任何工作。因为什么？告诉我，他会以损失财产、贬黜、流放来威胁吗？但这对那\Quote{不以性命为念}（\BibleRef{宗 20:24}）的人来说，都是小事，依照蒙福的保禄的话。你看到吗？在推翻死亡的暴政时，他也推翻了魔鬼的力量。因为那学会思考无数关于复活真理的人，怎么会怕死呢？他怎么会再战栗呢？\par

7. 所以，你们不要忧愁，说：为什么我们遭受这样那样的事？因为这样胜利就更加光荣。若不是借着死亡祂摧毁了死亡，胜利就不算光荣；最奇妙的是，祂恰恰借着敌人所依仗的力量战胜了他，处处显示祂丰富的手段和巧妙的安排。那么，我们不要辜负所赐予的恩惠。因为他说：\Quote{我们所领受的，不是恐惧的神，而是能力、爱德和刚健之心的神。}（\BibleRef{罗 8:15}）让我们高昂地站立，嘲笑死亡。\par

但我停一下，因为我不禁悲伤地叹息[想到]基督将我们提升到何处，以及我们将自己贬抑到何处。因为我看到公共场中的哀哭，为临终者的悲叹，嚎啕大哭，以及其他不当的行为，相信我，我在那些外教人、犹太人和异端者面前感到羞愧，他们在看到这一切，以及所有因此公开嘲笑我们的人面前都感到羞愧。因为以后无论我说什么，当我论及复活时，都将无人领会。为什么？因为外教人不听从我所说的话，却留意你们所做的事。因为他们会立刻说：\Quote{这些家伙中哪一个能蔑视死亡？连看别人死都做不到。}\par

保禄讲论了美妙的事，美妙且配得上天堂及天主对人的爱。因为他怎么说？\Quote{祂要释放那些因怕死而一生处于奴役之下的人。}但你们的行为却与这些话相悖，仿佛在攻击它们，使人们无法相信。事实上，天主为对抗这一点而建造了堡垒，为摧毁这恶习，已安排了许多事。因为告诉我：明亮的火炬代表什么？我们难道不是把它们当作运动员的先导？颂歌又代表什么？我们难道不是光荣天主，感谢祂终于为逝者加冕，使他脱离劳苦，除去不定，将他带至自己身边？颂歌不正是为此吗？圣咏不正是为此吗？这些都是欢庆之人的行动。经上说：\Quote{有谁欢乐？让他咏唱圣咏。}（\BibleRef{雅 5:13}）但外教人不理会这些事。因为有人会说：不要对我讲那人在痛苦之外表现得有哲学智慧，因为这不是什么伟大或惊人的事——请给我看一个正是痛苦之中仍有哲学智慧的人，那时我才相信复活。\par

的确，从事世俗事务的妇女这样做并不奇怪。然而，这甚至也是可怕的，因为她们也被要求同样的哲学智慧。因此保禄也说：\Quote{关于沉睡的人，我不愿意你们不知道，免得你们悲伤，像那些没有希望的人一样。}（\BibleRef{得前 4:13}）他写这话不是给独修士，也不是给永贞童女，而是给世俗中的男女。但这还不算太可怕。但当任何一个声称与世界同钉十字架的男人或女人扯头发、尖声哭喊时——还有什么比这更不像样的呢？相信我，如果事情按应有的方式进行，这样的人应长期被排除在教会的门槛之外。因为真正值得为她们悲伤的，是那些仍然惧怕死亡、对复活没有信心的人。\par

\Quote{但是我不怀疑复活}（有人说）\Quote{但我渴望与他交往。} 那么请告诉我，当他离家外出，且长期不归时，你难道不也是这样吗？\Quote{是的，但那时我也哭泣}（她说）\Quote{并因思念他而哀伤。} 但那真正渴望同伴者的行为是这样的，而这是那绝望于他归来者的表现。\par

想一想，你在那种场合歌唱的是什么：\BibleQuote{返回你的安息，我的灵魂，因为上主厚待了你。}\BibleRef{咏 115:7} 以及：\BibleQuote{我不怕凶险，因为你与我同在。}\BibleRef{咏 22:4} 以及：\BibleQuote{你是我逃避困厄的避难所。}\BibleRef{咏 31:7} 想想这些圣咏的意思。但你却不留心，反而因悲伤而沉醉。\par

仔细思想别人在葬礼上的哀哭，好让你在自己的情况中得到药方。\Quote{返回你的安息，我的灵魂，因为上主厚待了你。} 告诉我，你说上主厚待了你，却又哭泣？这岂非只是演戏，岂非伪善？因为如果你真相信你所说的话，你的悲伤就是多余的；但如果你是在戏弄、做戏，把这些当作神话，那你为什么要唱圣咏？为什么要忍受那些送葬者？你为什么不赶走那些歌唱者？但那是疯子才会做的事。然而，另一件事更甚。\par

目前，我劝告你们；但随着时间推移，我将更严肃地处理此事：因为我非常害怕这种习俗可能给教会带来某种严重的疾病。我们稍后将纠正哭泣的情况。同时，我告诫并见证，无论对富人还是穷人，对妇女还是男人。\par

愿天主确实恩赐你们所有人无哭无泪地离开生命，并按照合宜的规矩，年迈的父亲由儿子、母亲由女儿、孙辈、曾孙辈送到坟墓，在健康的老年中，而不虞早逝。愿如此成就，我这样祈祷，并劝勉主教们和你们全体互相恳求天主，并共同作此祈祷。但是，如果（天主不许，但愿永不发生）发生任何痛苦的死亡——我说痛苦，不是就其本质而言（因为从此不再有痛苦的死亡，它与睡眠无异），而是就你们的心绪而言——如果发生有人雇用这些哭丧的妇女，请相信我（我不是随便说说，而是已经决定，谁若生气，让他生气），我们将把那人长时间逐出教会，如同对待偶像崇拜者一样。因为如果保禄称\Quote{贪婪的人就是偶像崇拜者}\BibleRef{弗 5:5}，那么对于那将偶像崇拜者的做法引入信徒中的人，更是如此。\par

因为，请告诉我，你为何邀请司铎和歌唱者？难道不是为了提供安慰？不是为了尊敬亡者？那么你为什么侮辱他？你为什么把他当作公开的展览？你为什么像在舞台上一样戏弄他？我们前来，谈论有关复活的事，教导所有的人，甚至那些尚未遭受打击的人，通过对他的尊敬，如果发生类似的事，要勇敢承担。而你却带来那些尽力推翻我们[教导]的人？有什么比这嘲弄和讽刺更糟？有什么比这矛盾更严重？\par

8. 你们应羞愧并敬畏；但如果你们不愿，我们不能容忍将这些如此有害的做法引入教会。因为经上说：\Quote{犯罪的人，要在众人面前加以责备。}\BibleRef{弟前 5:20} 至于那些可怜而悲惨的妇女，我们通过你们禁止她们再进入信徒的葬礼，免得我们迫使她们认真地为自己的邪恶哀哭，并教导她们不要在别人的不幸中做这些事，而要为她们自己的不幸哭泣。因为一位慈爱的父亲，当他有一个不听话的儿子时，不仅劝告他不要亲近恶人，还使他们害怕。看，我劝告你们，并通过你们劝告那些妇女：不要邀请这类人，她们也不要参加。愿天主使我的话产生一些效果，使我的威胁生效。但如果（天主不许）我们被忽视，那么我们别无选择，只能将威胁付诸实施，用教会的法律惩罚你们，并按她们应得的方式惩罚那些妇女。\par

若有人固执而藐视，让他听到基督现在所说的话：\Quote{如果你的兄弟得罪了你，去，要在你和他独处的时候，指出他的错来。} 但如果他不听劝，\Quote{另外带一两个人去。} 如果即便如此他仍反对，\Quote{告诉教会，如果他连教会也不听从，你就把他看作外教人或税吏。}\BibleRef{玛 18:15} 如果当一个人得罪了我，不肯听从时，[主]命令我这样远离他，请你们判断，对于那得罪自己、得罪天主的人，我该怎样看待他？因为当我们如此温和地对待你们时，你们自己不也谴责我们吗？\par

但如果有人藐视我们所施行的束缚，让基督再次教导他，说：\Quote{凡你们在地上所束缚的，在天上也要被束缚；凡你们在地上所释放的，在天上也要被释放。}\BibleRef{玛 18:18} 因为即使我们自己可怜、无用、理应受藐视，正如我们确实是的，但我们并非为自己报仇，也不是抵挡愤怒，而是关心你们的得救。\par

请受敬畏感动，我恳求你们，并尊重。因为如果有人容忍一位朋友用超过必要的激烈攻击他，如果他知道他的目的，知道他出于善意而非傲慢，那么对于一位责备他的教师，且这位教师自己并非以权威，也不是以统治者的地位，而是以慈祥监护人的地位说这些话，就更应容忍。因为我们说这些话，并非希望显示我们的权威，（因为我们怎么可能呢？我们祈祷永不必试验它们？）而是为你们忧伤、哀哭。\par

请你们宽恕我，但愿没有人轻视教会的束缚。因为束缚人的并不是人，而是基督，祂把这权柄赐给了我们，并使人成为如此伟大尊荣的主人。事实上，我们确实愿意使用这权柄来解开束缚；或者更确切地说，我们甚至希望不需要那样做，因为我们希望在我们中间没有任何人被束缚——我们并非如此可怜和不幸，即使我们中有一些极端无用的人。但如果我们被迫如此行事，请宽恕我们。因为我们不是出于自愿，也不是希望如此，而是出于对你们这些被束缚者的忧伤，才把锁链放在你们身上。如果有人轻视这些锁链，审判的时刻将会到来，那时他将明白。至于之后的事，我不愿多谈，免得伤害你们的心。首先，我们确实不愿意被带入这种必要；但如果被带入，我们便尽我们的本分，把锁链套上。如果有人挣脱它们，我已经尽了我的本分，从此免于责难，而你们将向那位命令我束缚人的主交账。\par

因为当一位君王公开坐朝时，如果侍立在他身旁的侍卫之一奉命捆绑某个侍从，并给他套上锁链，而这人不仅推开那侍从，还把锁链打成碎片，那么受侮辱的难道是侍卫，而不是那发令的君王吗？因为如果祂将那些对信徒所做的事视为自己的事，那么当那位被任命教导的人受侮辱时，祂更会感到是自己受了侮辱。\par

但愿天主使这教会中的任何首领都不必被迫施加这些锁链。因为不犯罪固然是好事，忍受责斥也是有益的。让我们接受责斥，并努力不犯罪；如果我们犯了罪，就让我们承受责斥。因为正如不被受伤是极好的，但若受伤发生，就要治疗伤口，同样，在此事上也是如此。\par

但愿没有人需要这样的治疗。\BibleQuote{但我们深信你们有更好的事，且是关乎救恩的事，虽然我们这样说。}\BibleRef{希 6:9}但我们讲了更激烈的话，是为了更大的保障。因为宁可我被你们看作是一个严厉、苛刻、固执己见的人，也胜过你们做出天主不喜悦的事。但我们信赖天主，这次斥责对你们并非无益，相反你们将如此改变，以致这些讲论可成为对你们的赞美和颂扬：使我们众人堪当获得那些美好事物，就是天主为爱祂的人在基督耶稣我们的主内所预许的。愿光荣、权能、尊崇归于祂及圣父和圣神，从现在直到永远，永世无尽。阿们。\par

\SourceRecord{Translated by Frederic Gardiner.；Nicene and Post-Nicene Fathers, First Series；Vol. 14.；Edited by Philip Schaff.；Buffalo, NY: Christian Literature Publishing Co.,；1889.；<http://www.newadvent.org/fathers/240204.htm>.}

% structure: p0001 source=h1 target=section reason=page-title

\section{致希伯来人书讲道集 第五篇}

% structure: p0002 source=h2 target=subsection reason=relative-heading-rank; base=h2

\subsection{希伯来书 2:16-17}

\Quote{\Emph{因为祂确实没有握住天使，而是握住了亚巴郎的后裔。} \Emph{因此，祂必须在一切事上成为与祂的弟兄相似。}}\par

1. 保禄为了显示天主对人类的大恩，以及祂对人类的爱，先说道：\BibleQuote{因为孩子们既同有血肉，祂也照样取了同样的血肉} \BibleRef{希 2:14}，接着论述这个主题。因为不要轻视所说的，也不要以为这只是小事，即祂取了我们的血肉。祂没有将此恩赐给天使；\Quote{因为祂确实没有握住天使，而是握住了亚巴郎的后裔。} 他是什么意思？祂没有取天使的本性，而是取了人的本性。但所谓\Quote{祂握住}是什么意思？他是说，祂没有抓住那属于天使的本性，而是抓住了我们。但为什么他不说\Quote{祂取了}，而用了\Quote{祂握住}这个说法？这源于一个比喻：有人追赶那些远离他们的人，尽一切努力在他们逃跑时赶上他们，在他们跳开时抓住他们。因为当人的本性从祂面前逃跑，跑得很远时（\BibleRef{弗 2:13}），祂追赶并抓住了我们。祂表明祂这样做完全是出于仁慈、爱和温柔的关怀。正如他所说：\BibleQuote{他们不都是服役的使者，被派遣为那些将要承受救恩的人服务吗？} \BibleRef{希 1:14}——他显示了祂对人性极大的关心，以及天主非常重视它。同样，在此处他通过比较更进一步说明了这一点，因为他说道：\Quote{祂没有握住天使。} 事实上，这确实是一件伟大而奇妙的事，充满了惊叹：我们的血肉竟能高居天上，受到天使、总领天使、革鲁宾和色辣芬的朝拜。因为我多次思考这事，便感到惊奇，对人类产生了伟大的想法。因为我看到开始是伟大而辉煌的，天主对我们的本性怀有巨大的热忱。\par

此外，他没有说\Quote{祂握住了人}，而是为了抬高他们（希伯来人）并显示他们的种族是伟大而可敬的，他说：\Quote{而是握住了亚巴郎的后裔。}\par

\Quote{因此，祂必须在一切事上成为与祂的弟兄相似。} 这是怎么回事，\Quote{在一切事上}？他是说：祂被生，被养育，成长，经历了所有必要的事，最后死去。这就是\Quote{在一切事上成为与祂的弟兄相似}。因为在他详细论述了祂的尊威和天上的荣耀之后，他开始论述救恩计划。请看他以多么大的能力这样做。他如何将祂描绘成具有巨大的热忱\Quote{要与我们一起}：这表明了极大的关怀。因为他在上面说过：\Quote{孩子们既同有血肉，祂也同样取了相同的血肉}；在此处他也说：\Quote{在一切事上成为与祂的弟兄相似}。这几乎是说：那如此伟大的一位，那\Quote{是祂荣耀的光辉}，那\Quote{是祂本体的真像}，那\Quote{创造了诸世界}，那\Quote{坐在圣父的右边}，祂愿意并热切成为我们的一切弟兄，并为此离开了天使和其他权能，降到我们中间，抓住了我们，并成就了无数美善。祂摧毁了死亡，从暴政中驱逐了魔鬼，把我们从捆绑中释放出来：不仅通过兄弟关系尊荣我们，而且还通过其他无数方式。因为祂也愿意成为我们在圣父面前的大司祭：因为他接着说道：\par

2. \Quote{为要成为仁慈忠信的大司祭，以事奉天主。} 为此原因（他是说）祂取了我们的血肉，只是出于对人的爱，好能怜悯我们。因为救恩计划的别的原因没有，只有这一个。因为祂看见我们被抛弃在地上，灭亡，被死亡所暴虐，就怜悯了我们。他说：\Quote{为百姓赎罪。为要成为仁慈忠信的大司祭。}\par

\Quote{忠信}是什么意思？真实、有能力。因为圣子是忠信的大司祭，能够把祂作为大司祭所代祷的人从罪恶中拯救出来。为要献上能净化我们的祭献，为此原因祂成了人。\par

因此他补充道：\Quote{以事奉天主}——即为了与天主有关的事。我们完全变成了天主的仇敌（他要说），已被定罪、降级，没有人能为我们献祭。祂看到我们的处境，就怜悯了我们，不是为我们指定一位大司祭，而是亲自成为大司祭。\Quote{忠信}是什么意思？他补充说：\Quote{为百姓赎罪。}\par

% structure: p0010 source=h2 target=subsection reason=relative-heading-rank; base=h2

\subsection{希伯来书 2:18}

\Quote{因为祂亲自受试探而受苦，就能扶助那些受试探的人。} 这完全是低下的、卑微的，与天主不相称。\Quote{因为祂亲自受了苦}，他说。这里他谈论的是那成了血肉的一位，这话是为了听众的完全确信，并因他们的软弱而说。他是说：祂亲身经历了我们所受的一切；\Quote{现在}祂并非不知道我们的苦难；不仅作为天主祂知道，而且作为人也知道，借着祂所受的试炼；祂受了许多苦，祂知道如何同情。然而天主是不能受苦的；但这里他描述的是降生成人的事，好像他说道：就连基督的血肉也受了许多可怕的事。祂知道什么是患难；祂知道什么是试探，不亚于我们这些受苦的人，因为祂自己也受了苦。\par

（那么这\Quote{祂能扶助那些受试探的人}是什么？就如有人会说：祂要极其热切地伸出援手，祂会同情。）\par

3. 既然他们渴望得到某种伟大之物，并比外邦人的归化者更有优势，他便表明，在此处他们确有优势，同时又丝毫未损害那些外邦人。这优势在何处？因为救恩出自他们，因为祂首先抓住了他们，因为祂从那个种族取了肉身。他说：\Quote{因为，} \BibleQuote{祂并不援助天使，而是援助亚巴郎的后裔。} 由此，他既尊敬了圣祖，也指出了什么是\Quote{亚巴郎的后裔}。他提醒他们，曾向圣祖所作的恩许，说：\BibleQuote{我要将这地赐给你和你的后裔} \BibleRef{创 13:15}；以极小之事，表明他们彼此亲近，因为他们是\Quote{出于一体}。但这亲近并不算大；于是他返回这一点，从此讲述了那按肉身的经纶，并说：仅仅是愿意成为人，就是极大关怀与爱的证明；但现在，并非仅此而已，还有那藉着祂赐予我们的不朽恩惠，因为他说：\BibleQuote{为百姓的罪过作赎罪。}\par

为什么祂不说\Quote{世界}，反而说\Quote{百姓}？因为祂背负了所有人的罪。因为至此祂的论述是关于他们（希伯来人）。当日天使也向若瑟说：\BibleQuote{你要给祂起名叫耶稣，因为祂要将自己的百姓从罪恶中救出来。} \BibleRef{玛 1:21} 因为这也应当首先发生，祂为此而来，先拯救他们，然后藉着他们拯救其余的人，虽然结果恰恰相反。宗徒们起初也说：\BibleQuote{天主既兴起自己的儿子，就先差祂到你们这里来，赐福给你们} \BibleRef{宗 3:26}；又说：\BibleQuote{这救恩的道是传给你们的。} \BibleRef{宗 13:26} 此处，祂说\BibleQuote{为百姓的罪过作赎罪}，以此显明犹太人的高贵出身。暂时祂如此说话。因为正是祂赦免众人的罪，祂在瘫子的事上表明了，说：\BibleQuote{你的罪赦了} \BibleRef{谷 2:5}；也在圣洗中：祂对门徒说：\BibleQuote{你们要去使万民成为门徒，因父、及子、及圣神之名给他们授洗} \BibleRef{玛 28:19}。\par

4. 但保禄一旦谈论起肉身，便毫无畏惧地述说一切微贱之事：请看下文：\par

% structure: p0016 source=h2 target=subsection reason=relative-heading-rank; base=h2

\subsection{希伯来书 3:1-2}

\Quote{因此，圣洁的弟兄，同蒙天上召叫的人，应当思考我们所宣认的宗徒和大司祭——基督耶稣，祂对委派（或：使）祂的那一位是忠信的，正如摩西在祂的全家中一样。}\par

他既将要拿祂与摩西比较，便将论述引向大司祭职的法律；因为所有人都对摩西怀有崇高的敬意；此外，他早已预先播下优越性的种子。因此，他从肉身开始，上达天主性，在那里再无任何比较。他从肉身（祂的人性）开始，暂时采用平等，说：\Quote{正如摩西在祂的全家中一样}：他起初并未显明祂的优越性，以免听者惊退，立刻掩耳。虽然他们是信徒，然而他们仍对摩西抱有强烈的敬意。他说：\Quote{祂对使祂的那一位是忠信的}——使祂成为什么？\Quote{宗徒和大司祭。} 他这里毫不涉及祂的本质或天主性；而只是论及人的尊荣。\par

\Quote{正如摩西在祂的全家中一样}，即或在百姓中，或在圣殿中。但他此处使用\Quote{在祂家中}这一表达，正如人谈及家中的人一样；正如某位家中的监护人和管家，摩西对百姓也是如此。因为藉着\Quote{家}他意指百姓，他接着说：\BibleQuote{这家就是我们} \BibleRef{希 3:6}；即我们是在祂的创造中。然后便是优越性。\par

% structure: p0020 source=h2 target=subsection reason=relative-heading-rank; base=h2

\subsection{希伯来书 3:3-4}

（再次，他指的是肉身），\Quote{因为这人被算作比摩西更配得荣耀，} \Quote{正如建造房屋的比房屋更尊荣}；（他意指）摩西自己也属于这家。（再者，他没有说：一个是仆人，另一个是主人，而是暗指了这一点。）如果百姓是房屋，而他属于百姓，那么他当然属于这家。因为我们通常也这样说：某人是某人家里的。因为他这里说的是房屋，不是圣殿，因为圣殿不是由天主建造，而是由人建造。但使祂如此的，是天主。他指的是摩西。请看他是如何暗中显示优越性的。他说：\Quote{忠信} \Quote{在祂的全家中}，他自己也属于这家，即属于百姓。建造者比房屋更尊荣，但他并没有说\Quote{工匠比他的作品更尊荣}，而是说\Quote{建造房屋的比房屋更尊荣} \BibleRef{希 3:4}。\Quote{但建造万有的是天主。} 你看，他说的不是圣殿，而是全体百姓。\par

% structure: p0022 source=h2 target=subsection reason=relative-heading-rank; base=h2

\subsection{希伯来书 3:5}

\Quote{摩西在祂的全家中固然忠信，如同仆人，为给那些要传讲的事作见证。} 再看另一个优越之处，即源自圣子与仆人的。你又看到，藉着\Quote{圣子}的称呼，他暗示了真正的亲属关系。\BibleRef{希 3:6} \Quote{但基督作为圣子，治理自己的家。} 你岂不明白，他如何将受造物与造物主、仆人与儿子区分开来？而且，祂确实以主人的身份进入祂父的产业，而另一个人则如同仆人。\par

\Quote{谁的} [即] 天主的 \Quote{房屋就是我们，如果我们坚持信心和望德的欢乐直到结束。} 这里他又鼓励他们勇往直前，不要跌倒：因为我们将是天主的 \Quote{房屋}（他说），正如\Person{梅瑟}一样，\Quote{如果我们坚持我们的信心和欢乐直到结束。} 然而，那位在考验中痛苦并跌倒的人（他要说），并不光荣；那羞愧、隐藏自己的人没有信心；那困惑的人也不光荣。\par

接着他又称赞他们说：\Quote{如果我们坚持信心和望德的欢乐直到结束，} 暗示他们已经开始了；但需要坚持到底，并且不仅要站立，还要使他们的望德坚定 \Quote{在信德的确信之中，} 不被考验所动摇。\par

5. 不要惊讶，\Quote{他亲自受试探} \BibleRef{希 2:18} 这话是按人的方式说的。因为如果圣经论到那并未降生成人的圣父说：\Quote{主从天上垂看，察看一切世人} \BibleRef{咏 13:2}，意思是准确地知晓一切；又说：\Quote{我要下去，察看他们是否完全像那达到我耳中的呼声一样} \BibleRef{创 18:21}；又说：\Quote{天主不能容忍人的恶行} \BibleRef{创 6:5}——神圣的圣经彰显了他愤怒的伟大；更何况那在肉身内受苦的，这些事都是论基督说的。因为既然许多人认为经验是最可靠的知识途径，他愿意表明：那受苦者知道人的本性所受的是什么。\par

\Quote{因此，圣洁的弟兄们}（他说\Quote{因此}代替\Quote{为此}），\Quote{同蒙天召的}——（在此不寻求任何事物，如果你已被召到那里——那里是报酬，那里是赏报。那么怎样？）\Quote{应当思想我们所承认为宗徒和大司祭的耶稣基督，他对那委派他的忠信，如同\Person{梅瑟}在\Quote{他的全家}一样。}（\Quote{他对那委派他的忠信}是什么意思？即是，善加管理，保护属于他的一切，不让他们轻易被带走，\Quote{如同\Person{梅瑟}在他的全家}）也就是说，要知道你们的大司祭是谁，他是怎样的，你们就不再需要别的安慰或鼓励。现在他称他为\Quote{宗徒}，因为他是被\Quote{派遣}的，并称他为\Quote{我们所承认为大司祭的}，即信德的大司祭。这一位也被托付了一群人，如同另一位被托付了人民的领导，但这是更大的、更崇高的。\par

\Quote{作为以后将要谈论之事的证据。} 你是什么意思？难道天主接受人的见证吗？是的，当然。因为如果他呼天唤地作证（如先知所说：\Quote{天哪，请听！地啊，请侧耳！因为上主发言了} \BibleRef{依 1:2}——以及\Quote{山谷啊，地的根基啊，请听！因为上主与他的百姓有争辩} \BibleRef{米 6:2}），何况人呢？也就是说，当那些人[犹太人]无耻时，他们可以作见证。\par

% structure: p0029 source=h2 target=subsection reason=relative-heading-rank; base=h2

\subsection{希伯来书 3:6}

\Quote{但基督身为儿子。} 一位照管他人的财产，而这一位照管自己的。\Quote{和望德的欢乐。} 他说\Quote{望德}说得好。因为既然一切美善都在望德中，而我们还应当\Quote{坚持它}，甚至如同已经实现的荣耀一样；因此他说：\Quote{望德的欢乐。}\par

他又加上：\Quote{让我们坚持到底。} \BibleRef{罗 8:24} 因为\Quote{我们得救是在乎望德}；因此，如果\Quote{我们得救是在乎望德}，并且\Quote{耐心等待} \BibleRef{罗 8:25}，就不要为眼前的事忧伤，也不要现在就寻求以后应许的事；\Quote{因为}（他说）\Quote{所见的望德不是望德。} 既然美善是伟大的，我们无法在此短暂的生命中领受它们。那么他向我们将来的事预先告诉我们，却并不在此赐下，有何目的呢？为的是藉着应许抚慰我们的灵魂，藉着承诺坚固我们的热忱，为我们的竞赛傅油，激发我们的心智。因此，这一切都成就了。\par

6. 因此，我们不要烦恼，谁也不要因看见恶人亨通而烦恼。报应不在这里，无论对恶行还是德行；若有时对恶行或德行有报应，那也不是按功过，而只是审判的预尝，使那些不信复活的人甚至能藉着这里发生的事醒悟。所以，当我们看见恶人富有，不要沮丧；看见善人受苦，不要烦恼。因为冠冕在那里，惩罚在那里。\par

是的，从另一个角度看，一个坏人不可能是全然坏的，他或许也有一些善；一个好人也可能是全然好的，他或许也有一些罪。因此，当恶人兴旺时，这对他自己是灾祸，为使他在这里接受了那少许善事的赏报，日后在那边完全受罚；正因此他在今生得到报应。最幸福的是在这里受罚的人，为使他除去一切罪恶，得以离去时被认可、纯洁、无需被追究。保禄教导我们这一点，他说：\Quote{因此你们中有许多软弱多病，长眠的也不少。} \BibleRef{格前 11:30} 又说：\Quote{我把这样的人交给了撒殚。} \BibleRef{格前 5:5} 先知说：\Quote{因为她从主的手中领受了双倍的罪。} \BibleRef{依 40:2} 达味也说：\Quote{看，我的仇敌比我的头发还多，他们都用无理的仇恨恨我}：\Quote{赦免我一切的罪。} \BibleRef{咏 25:19} 又有另一位说：\Quote{上主，我们的天主，赐给我们和平，因为你把一切再次赐给了我们。} \BibleRef{依 26:12}\par

然而这些是一个人的话，表明好人在这里接受他们罪的惩罚。但是哪里有提到恶人在这里接受好事，而在那边完全受罚呢？请听亚巴郎对那个富人说：\Quote{你领受了好事，} 和 \Quote{拉匝禄领受了坏事。} \BibleRef{路 16:25} 什么好事？因为在这里通过说\Quote{你领受}，而不是\Quote{你曾经领受}，他表明各人按照自己所应得的被对待，一人享福，一人受苦。他说：\Quote{所以他在这里得安慰}（因为你看到他无罪）\Quote{而你受折磨。} 因此，当我们看到罪人在此亨通时，不要困惑；但当我们自己受苦时，让我们欢喜。因为这正是偿付罪的惩罚。\par

7. 因此，我们不要寻求安逸，因为基督应许给祂的门徒苦难，保禄说：\Quote{凡愿意在基督耶稣内虔敬生活的，都会遭受迫害。} \BibleRef{弟后 3:12} 没有一个高尚的角斗士在赛场上寻求沐浴、满是食物和酒的桌子。这不是角斗士所为，而是懒汉的行为。因为角斗士与尘土、油、烈日的光线、大量汗水、压力和束缚争战。这是比赛和战斗的时刻，因此也是受伤、流血和痛苦的时刻。请听有福的保禄所说：\Quote{我战斗正是这样，不像打空气的。} \BibleRef{格前 9:26} 让我们想，我们整个生命都在战斗中，这样我们就永远不会寻求休息，当我们受苦时，也永不会感到奇怪：就像拳击手在战斗时不感到奇怪一样。有另一个休息的季节。我们必须通过苦难得以成全。\par

即使没有迫害或苦难，仍有其他痛苦天天临到我们。如果我们不能承受这些，我们几乎不能忍受那些。\Quote{你们所遇到的诱惑，} \Quote{无非是人所常见的。} \BibleRef{格前 10:13} 因此，我们确实要向天主祈祷，让我们不陷于诱惑；但如果我们陷于诱惑，让我们高尚地承受它。因为谨慎的人的本分是不自投危险；而高尚的人和真正的哲人却这样行。因此，我们不要轻率地投入危险，因为那是鲁莽；但如果我们被带入危险，被环境召唤，也不要屈服，因为那是懦弱。但如果福音确实召唤我们，我们不要拒绝；但在简单情况下，当没有理由、需要或必要性因敬畏天主而召唤我们时，我们不要冲进去。因为这纯粹是炫耀和无用的野心。但如果任何损害宗教的事情发生，那么即使必须忍受一万次死亡，我们也什么都不要拒绝。当你发现虔诚之事如你所愿兴旺时，不要挑战考验。为什么引来不必要的、无益的危险呢？\par

我说这些话，因为我希望你们遵守基督的诫命，祂命令我们\Quote{祈祷不要进入诱惑} \BibleRef{玛 26:41}，并命令我们\Quote{背起十字架跟随}祂。\BibleRef{玛 16:24} 因为这些事并不矛盾，不，它们反而极其和谐。你要像勇敢的士兵一样做好准备，常穿着盔甲，清醒、警醒，时刻寻找敌人：但不要发动战争，因为这不是士兵的行为，而是煽动叛乱者的行为。但另一方面，如果虔诚的号角召唤你，立即出发，不顾自己的生命，以极大的热忱进入比赛，打破敌人的方阵，打伤魔鬼的脸，竖起你的胜利纪念碑。但如果虔诚丝毫未受伤害，没有人破坏我们的教义（我指的是与灵魂相关的那些），也没有强迫我们做任何不悦乐天主的事，就不要多事。\par

基督徒的生活必须充满流血牺牲；我不是说流别人的血，而是准备流自己的血。因此，让我们为了基督的缘故，如同倒水一样（因为身体里流动的血是水）随时准备倾倒自己的血，并且以如同脱衣服一样的好性情脱下我们的肉体。我们这样做，如果我们不被财富束缚，不被房屋束缚，不被情感束缚，如果我们超脱于万物。因为如果那些过世俗士兵生活的人告别一切，无论战争召唤他们到哪里，他们就出现在那里，旅行，以准备好心态忍受一切；那么，我们作为基督的士兵，更应该如此预备自己，并坚定地对抗情欲的战争。\par

8. 如今没有迫害，愿天主永不再有：但有一场别的战争，即对金钱的贪欲、嫉妒、情欲的战争。保禄描述这场战争时说：\Quote{我们不是对抗血肉之人。}\BibleRef{弗 6:12}这场战争随时在旁。因此他愿我们永远武装。因为他愿我们永远武装，他说：\Quote{要站住，束上腰。}\BibleRef{弗 6:14}这句话本身也属于现在时，表达我们应当常备不懈。因为舌头的战争很大，眼睛的战争很大；因此我们必须压制它——情欲的战争也很大。\par

因此，他在那一点上开始装备基督的士兵：因为他说：\Quote{要站住，束上你们的腰，}又加上了\Quote{用真理}\BibleRef{弗 6:14}。为什么\Quote{用真理}？因为情欲是嘲弄和谎言：因此先知说：\Quote{我的腰充满了讥讽。}\BibleRef{咏 37:7}那不是快乐，而是快乐的影子。他说：\Quote{束上你们的腰，用真理束上}；也就是说，用真正的快乐，用节制，用有序的行为。他给出这个劝告，是因为他知道罪的不合理性，并愿意我们所有的肢体都被围护起来；因为经上说：\Quote{不义的忿怒，不能免罚。}\BibleRef{德 1:22}\par

此外，他愿我们周围有胸甲和盾牌。因为欲望是一只轻易冲出的野兽，我们需要无数墙壁和栅栏来克服和抑制它。为此，天主特别用骨头，如同用石头一样建造了身体的这一部分，在其周围设了支撑，使欲望不至于在任何时候冲破或割穿，轻易伤害整个人。因为它是一团火（经上说）和一场大风暴，身体的其他部分无法承受这种暴力。医生们也说，为此肺被铺在心脏下面，使心脏本身置于柔软嫩滑之物中，如同击打在一块海绵上，得以不断休息，而不至于因冲击抵抗的、坚硬的胸骨，在剧烈搏动中受伤。因此我们需要一个坚固的胸甲，以便使这头野兽常常安静。\par

我们也需要头盔；因为推理官能在那里，从中我们可能得救，当我们做正确的事时，也可能被毁灭——因此他说：\Quote{救恩的头盔}\BibleRef{弗 6:17}。因为大脑本性柔软，因此上面覆盖着头骨，如同壳一样。它对我们是一切善恶的原因，知道什么是适宜的，什么不是。是的，我们的脚和手也需要盔甲，不是这些手脚，而是如前所述灵魂的手脚——前者用于做正确的事，后者使它们走在当走之处。因此，让我们彻底武装自己，我们就能战胜敌人，并在我们的主基督耶稣那里为自己戴上冠冕，愿光荣、权能、尊崇归于他及圣父和圣神，自今至永远，世世无穷。阿们。\par

\SourceRecord{Translated by Frederic Gardiner.；Nicene and Post-Nicene Fathers, First Series；Vol. 14.；Edited by Philip Schaff.；Buffalo, NY: Christian Literature Publishing Co.,；1889.；<http://www.newadvent.org/fathers/240205.htm>.}

% structure: p0001 source=h1 target=section reason=page-title

\section{希伯来书讲道集 第六篇}

% structure: p0002 source=h2 target=subsection reason=relative-heading-rank; base=h2

\subsection{希伯来书 3:7-11}

\Quote{ \Emph{因此，正如圣神所说：\BibleQuote{今天如果你们听从祂的声音，不要再心硬，像在激怒之时，像在旷野中试探之日，你们的祖先在那里试探了我，考验了我，并看见了我的作为四十年。因此，我厌恶那一世代，说：他们心中常常迷误，不认识我的道路。}所以} \Emph{我在忿怒中起誓：\BibleQuote{他们绝不得进入我的安息。}} }\par

1. \Person{保禄}论及望德，并说过\BibleQuote{我们就是祂的家，如果我们持守直到末了的信德和因望德而来的喜乐}之后（见：\BibleRef{希 3:6}），接着便显示我们应该坚定地期盼，并以此从圣经加以证明。但请留意，因为他表达这一点的方式有些困难，不易理解。因此我们必须先陈述自己的观点，在简要解释整个论证后，再阐明书信的言词。因为如果你们明白了宗徒的宗旨，就再也不需要我们了。\par

他的谈论涉及望德，以及我们必须盼望未来的事，并且凡在此劳苦的人，必定会有某种赏报、果实和安息。于是他从先知那里证明这一点；先知说了什么呢？\Quote{因此，正如圣神所说：\BibleQuote{今天如果你们听从祂的声音，不要再心硬，像在激怒之时，像在旷野中试探之日：你们的祖先在那里试探了我，考验了我，并看见了我的作为四十年。因此，我厌恶那一世代，说：他们心中常常迷误，不认识我的道路。所以我在忿怒中起誓：他们绝不得进入我的安息。}}\par

他说有\Quote{三}种安息：一种是安息日的安息，天主在其中歇了祂的工程；第二种是巴勒斯坦的安息，犹太人进入那里后，便从艰辛劳苦中得安息；第三种是真正的安息，即天国；获得它的人，确实从劳苦和烦恼中进入安息。因此他在这里提到这三种安息。\par

为何在只讨论一种安息时，他却提到三种？为要表明先知所说的是关于这一种。因为先知说的（他说）不是第一种。怎么可能呢？因为那早已发生了。也不是第二种，即巴勒斯坦的安息。怎么可能呢？因为他说：\BibleQuote{他们绝不得进入我的安息。}所以剩下的就是这第三种。\par

2. 但也有必要讲述历史，使论证更加清晰。当他们离开埃及，走完漫长的路程，并领受了无数次天主大能的证明——无论是在埃及、在红海（见：\BibleRef{宗 7:36}），还是在旷野——他们决定派遣探子去探查那地的性质。这些人去了，回来时虽然赞美那地，说它盛产上好的果实，但那里的人强壮而不可战胜。忘恩负义、愚昧无知的犹太人本应追忆天主先前的恩惠：当他们在众多埃及军队中被围困时，祂如何拯救他们脱离危险，使他们成为敌人战利品的主人；在旷野中，祂又裂开磐石，赐给他们丰沛的水，赐给他们玛纳，以及祂所行的其他奇事；〔我说，本应记住这些〕并信赖天主，但他们毫不顾及，反而惊恐万分，好像什么都没有发生过一样，说：\Quote{我们要回埃及去，\InnerQuote{天主把我们带出来}（据说）\InnerQuote{是要杀我们，连我们的妻子儿女一起。}}（见：\BibleRef{户 14:3}）因此，天主恼怒他们如此迅速地忘掉所行的事，就起誓说，说了这话的那一代人不得进入安息；他们都死在旷野。于是他说：当\Person{达味}后来、在这些事之后、在那代人之后说话时，说：\Quote{今天，如果你们听从祂的声音，不要再心硬，}免得你们遭受祖先所受的、并被剥夺安息；他显然〔说的是〕某种〔将来的〕安息。因为（他说）如果他们已获得了安息，为何祂又对他们说：\Quote{今天如果你们听从祂的声音，不要再心硬，}像你们的祖先一样？那么，除了安息日是其预象和预表的天国之外，还有什么别的安息呢？\par

3. 接着，他引证了全部的见证（这就是：\BibleQuote{今天如果你们听从祂的声音，不要再心硬，像在激怒之时，像在旷野中试探之日，你们的祖先在那里试探了我，考验了我，并看见了我的作为四十年。因此，我厌恶那一世代，说：他们心中常常迷误，不认识我的道路。所以我在忿怒中起誓：他们绝不得进入我的安息。}），随后补充道：\par

% structure: p0010 source=h2 target=subsection reason=relative-heading-rank; base=h2

\subsection{希伯来书 3:12}

\BibleQuote{弟兄们，你们要小心，免得你们中间有人存着邪恶不信的心，离弃了永生的天主。}因为不信出自心硬：就像身体上长茧变硬的部分不听从医生的手一样，变硬的灵魂也不听从天主的话。此外，很可能有些人甚至不信，好像所行的事不真实似的。\par

因此他说：\BibleQuote{你们要小心，免得你们中间有人存着邪恶不信的心，离弃了永生的天主。}因为来自将来的论证不如来自过去的论证有说服力，所以他要他们回想历史上他们缺乏信德的事。因为（他说）如果你们的祖先因为不像所应许的那样盼望而遭受这些，你们更会如此。因为对他们也说了这话：因为\BibleQuote{今天}（他说）\Quote{永远存在，}只要世界存留。\par

% structure: p0013 source=h2 target=subsection reason=relative-heading-rank; base=h2

\subsection{希伯来书 3:13}

4. 为此，\Quote{你们要天天彼此劝勉，趁今天还叫\Quote{今天}的时候。} 这就是说，要互相建立，振作自己，免得同样的事发生在你们身上。\Quote{免得你们中间有人被罪的诡计所硬化。} 你看见罪产生不信吗？因为不信生出邪恶生活，同样，灵魂\Quote{陷入邪恶深处时，便变得轻蔑} \BibleRef{箴 18:3}，变得轻蔑后，就甚至不再相信，以此免于恐惧。因为\Quote{他们说}（有人说），\Quote{主看不见，雅各伯的天主也不留意。} \BibleRef{咏 93:7} 又说，\Quote{我们的嘴唇属于我们：谁是我们主上？} \BibleRef{咏 11:4}；又说\Quote{恶人为什么激怒天主？} \BibleRef{咏 9:13}；又说\Quote{愚妄人心里说：没有天主；他们都腐败，行为可憎。} \BibleRef{咏 13:1} \Quote{天主不在他眼中，因为他在他面前行诡诈，要发现他的罪孽而憎恨。} \BibleRef{咏 36:1} 是的，基督也说了同样的话：\Quote{凡作恶的，恨光，不来就光。} \BibleRef{若 3:20}\par

% structure: p0015 source=h2 target=subsection reason=relative-heading-rank; base=h2

\subsection{希伯来书 3:14}

然后他引用 \BibleRef{希 3:14}，\Quote{因为我们已经成了基督的分享者。} 这是什么意思，\Quote{我们已经成了基督的分享者}？我们分享祂（他意思是），我们与祂成为一体——祂是元首，我们是身体，\Quote{同为继承者，同一身体；我们是一个身体，属于祂的肉和祂的骨。} \BibleRef{弗 3:6}\par

\Quote{如果我们把起初的信心（或：我们存在的\OriginalTerm{根基}{subsistence}）坚持到底。} 什么是\Quote{我们存在的根基}？就是使我们站立、得以存在、赖以存立的信德，如同人们所说的。\par

% structure: p0018 source=h2 target=subsection reason=relative-heading-rank; base=h2

\subsection{希伯来书 3:16-4:2}

5. 然后他引用 \BibleRef{希 3:15}，\Quote{当经上说：今天如果你们听见祂的声音，不要心硬，像在悖逆时那样。} 这是一个\OriginalTerm{倒置}{transposition}，\Quote{当经上说：今天如果你们听见祂的声音，不要心硬。} [必须这样读：]\par

\BibleRef{希 4:1} \Quote{我们该存畏惧，免得尽管进入祂安息的诺言尚存，你们中有人似乎达不到；因为福音同样传给了我们，像传给了他们一样，当经上说：今天如果你们听见祂的声音}（因为\Quote{今天}是\Quote{每一个时候}）。\par

然后[他补充说]\Quote{但所听的话对他们无益，因为他们没有用信德与听见的人结合。} 如何无益？然后要警告他们，他通过下面的话表明同样的事：\par

\BibleRef{希 3:16} \Quote{有些人虽然听见，却悖逆了；然而不是所有从埃及经梅瑟出来的人都是这样。四十年之久，祂向谁发怒呢？岂不是向那些犯罪、尸骨倒在旷野的人吗？祂向谁起誓，他们不得进入祂的安息呢？岂不是向那些不信的人吗？这样看来，他们不能进入是因为不信。} 在再次重复证言之后，他还加上问题，使论证更清晰。因为他说（他重复），\Quote{今天如果你们听见祂的声音，不要心硬，像在悖逆时那样。} 他说[是]论到谁硬了心？论到谁不信？岂不是犹太人吗？\par

他所说大意如此：他们也听见了，像我们听见一样；但对他们没有益处。那么，不要以为因\Quote{听}宣讲就会得益；看他们也听了，却因不信而无所得。\par

那时，加肋布和若苏厄，因为他们不与那些不信的人同意，逃脱了临到他们的灾祸。并且你看他说得好巧妙，不是说\Quote{他们不同意}，而是说\Quote{他们没有\OriginalTerm{混合}{mixed}}——也就是说，他们分开，但不是分裂，当时所有其他人都持同一心意。这里，在我看来也暗示了分裂。\par

% structure: p0025 source=h2 target=subsection reason=relative-heading-rank; base=h2

\subsection{希伯来书 4:3-5}

6. \BibleRef{希 4:3} 因为\Quote{我们相信了的人，} 他说，\Quote{进入安息。} 从何显然，他补充说：\Quote{正如祂说：我在忿怒中起誓：他们绝不能进入我的安息，虽然工程从创世以来已经完成。} 这确实不是证明我们将进入，而是证明他们未能进入。那么呢？到此为止，他的目的是要说明，正如那个安息不排斥另一个安息的提及，同样，这个[安息]也不排斥天堂的[安息]。到此为止，他希望表明他们[以色列人]没有到达安息。因为他意味着这一点，他说 \BibleRef{希 4:4}，\Quote{因为祂在某一处论到第七日这样说：天主在第七日安息了，歇了祂一切的工作。而在此处又说：他们绝不能进入我的安息。} 你看这如何并不妨碍这是一个安息？\par

% structure: p0027 source=h2 target=subsection reason=relative-heading-rank; base=h2

\subsection{希伯来书 4:6-8}

\Quote{由此可见，}（他说）\Quote{必须有人进入，而原先听到福音的，因不信未能进入；于是重新限定一个日子，藉达味说：今天，过了这么久；正如前所说。} 但他是什么意思？\Quote{既然}（他意思是）\Quote{必须有人}必定\Quote{进入}，而\Quote{他们没有进入。} 并且宣告有一个进入，且\Quote{必须有人进入}，让我们听凭何而清楚。因为过了许多年（他说），达味又说：\Quote{今天如果你们听见祂的声音，不要心硬} \BibleRef{希 4:8}，\Quote{因为若苏厄若给了他们安息，以后就不会再提另一个日子。} 显然，他说这些，针对那些要得到某种报酬的人。\par

% structure: p0029 source=h2 target=subsection reason=relative-heading-rank; base=h2

\subsection{希伯来书 4:9}

7. \BibleQuote{因此，为天主的子民留下了一个安息。} 这从哪里显明呢？从劝勉的话：\BibleQuote{不要硬着你们的心}：因为如果没有安息，就不会有这样的劝勉。他们也不会被劝诫不要做同样的事 [像犹太人那样]，免得遭受同样的事，除非他们将要遭受同样的事。但是，那些拥有巴勒斯坦的人怎么会遭受同样的事 [即被排除在安息之外] 呢？除非有其他安息？\par

他巧妙地把论证归结。因为他不是说\Quote{安息}，而是说\Quote{安息日之遵守}；用恰当的称谓称王国为\Quote{安息日之遵守}，那是他们所欢喜并被吸引的。因为，在安息日，他命令戒绝一切恶事，只应做那些与敬拜天主有关的事——这些事是司铎们惯常完成的，以及任何有益灵魂的事，除此之外别无他事——同样地，那时也将如此。\par

% structure: p0032 source=h2 target=subsection reason=relative-heading-rank; base=h2

\subsection{\BibleRef{希 4:10}}

然而，说这话的不是他，而是什么呢？\BibleRef{希 4:10}，\BibleQuote{因为那进入他的安息的，也歇了自己的工，正如天主歇了他的工一样。} 他说，正如天主歇了他的工，那进入他的安息的也歇了工。因为他向他们讲论的是关于安息的事，他们渴望听到这将在何时发生，他便以此结束论证。\par

8. 他又说\BibleQuote{今日}，好叫他们永不失去希望。他说，\BibleQuote{天天彼此劝勉，}[\BibleQuote{趁著还有今日，}]也就是说，即使人犯了罪，只要还在\BibleQuote{今日}，他就有希望；因此，只要人活着，谁也不要绝望。首先，他说，\BibleQuote{不要有不相信的恶心。} \BibleRef{希 3:12} 但即使有，也不要绝望，而要自新；因为只要我们还在这个世界，\BibleQuote{今日}的时期就存在。但在这里，他不仅指不信，也指怨言：他说，\BibleQuote{他们的尸体倒在旷野里。}\par

然后，免得有人认为他们只会被剥夺安息，他又加上了惩罚，说\BibleRef{希 4:12}，\BibleQuote{因为天主的圣言是生活的，是有效的，比一切双刃的剑更锐利，甚至能刺入灵魂和精神的分离，关节与骨髓的分离，且能辨明心中的思想和意向。} 这里他是在谈论地狱和惩罚。\BibleQuote{它刺入}（他说）我们心的隐秘处，并斩断灵魂。这里不是尸体的倒下，也不是像那里被剥夺国土，而是被剥夺天国，并被交付给永恒的地狱和不死的惩罚与报应。\par

\BibleRef{希 3:13} \BibleQuote{但彼此劝勉。} 请看（这话的）温柔和温和：他没有说\BibleQuote{责备}，而是说\BibleQuote{劝勉}。我们应当这样对待那些因患难而困苦的人。他写信给得撒洛尼人时也说：\BibleQuote{警戒那不守规矩的人} \BibleRef{得前 5:14}，但在说到软弱的人时，却不是这样，而是什么呢？\BibleQuote{抚慰那心性怯懦的人，扶持软弱的人，对众人要忍耐}；也就是说，不要停止希望；不要绝望。因为那不鼓励因患难而困苦的人，会使他更加硬心。\par

9. \BibleQuote{免得你们中间有人，} 他说，\BibleQuote{被罪的欺骗所硬化。} 他意指或指魔鬼的欺骗（因为确实是一种欺骗：不仰望将来的事，以为我们无责任，不会为今世的行为受罚，也没有复活）；或另一种意义：麻木（或）绝望是欺骗。因为说\Quote{还有什么剩下呢？我已一次犯了罪，没有希望自新了}，这是欺骗。\par

然后他向他们提出希望，说\BibleRef{希 3:14}，\BibleQuote{我们已经成了有分于基督的人}；几乎是在说：那如此爱我们、使我们配得如此大事、以致使我们成为他身体的那一位，必不容我们丧亡。让我们想想（他说）我们被视为什么：我们和基督是一体；那么，不要不信他。他又暗示了另一处所说的话：\BibleQuote{如果我们与他同受苦难，也必与他一同为王。} \BibleRef{弟后 2:12} 因为这就是\BibleQuote{我们成了有分于他的人}所隐含的，我们分享基督所分享的同样事物。\par

他从美善的事劝勉他们；他说，\BibleQuote{因为我们是}他说，\BibleQuote{有分于基督的人。} 然后，又从阴暗的事：\BibleRef{希 4:1} \BibleQuote{所以，我们应存戒惧，免得在进入他安息的应许还存留的时候，你们中间有人落空了。} 因为那是明显且已被承认的。\par

\BibleRef{希 3:9} \BibleQuote{他们试探了我，} 他说，\BibleQuote{且看见了我在四十年间的作为。} 你看，不应该向天主问理由，而是无论他是否辩护 [我们的案件]，都信靠他。因为现在他控告的是那些人 [古人] \Quote{他们试探了天主}。因为那要求对他的能力或对他的眷顾或对他的慈爱得到证明的人，还不相信他是大能的或对人仁慈的。他在写信给这些 [希伯来人] 时也暗示了这一点，他们很可能在患难中已经希望获得经验和确凿的证据，证明他的能力和他对他们的眷顾。你看，在所有情况下，挑衅和惹怒都源于不信。\par

那麼他說什麼？\BibleRef{希 4:9}：\BibleQuote{為天主的子民，仍保留了一個安息。}請看他如何總結了整個論證。他說，\Quote{他發誓，}對先前那些人，\Quote{他們不得進入}那\Quote{安息，}而他們沒有進入。然後在他們之後很久，對猶太人講話時，他說：\Quote{不要硬著心，}像你們的祖先一樣，顯示還有另一個安息。因為我們不必談巴勒斯坦，他們已經佔有了它。他也不是在說安息日；肯定他不是在講很久以前發生的事。因此，他暗示了另一個，那才是真正的安息。\par

10. 因為那才是真正的安息，在那裡\BibleQuote{痛苦、憂傷和嘆息都已逃逸} \BibleRef{依 35:10}：沒有憂慮、辛勞、鬥爭，也沒有驚嚇和動搖靈魂的恐懼；只有那充滿喜樂的敬畏天主。沒有\BibleQuote{你必汗流滿面才得餬口}，也沒有\BibleQuote{荊棘和蒺藜} \BibleRef{創 3:19}；不再有\BibleQuote{你要在痛苦中生產兒女，你要戀慕你丈夫，他要管轄你} \BibleRef{創 3:16}。一切是和平、喜樂、歡欣、愉快、良善、溫柔。沒有嫉妒、沒有怨恨、沒有疾病、沒有死亡，無論是身體的或靈魂的死亡。沒有黑暗，沒有黑夜；全是白晝，全是光明，萬物明亮。不會疲倦，不會厭煩：我們將永遠保持對美善事物的渴望。\par

你願意我也給你一個那情景的圖像嗎？這是不可能的。但無論如何，在可能範圍內，我將試著給你一個圖像。讓我們仰望天空，當沒有任何雲彩遮擋，它展現出它的星冠時。然後，當我們長久地凝視其外觀的美麗，讓我們思想，我們也將有一個地面，當然不是這樣的地面，而是比金子之於泥土更為美麗，並且（讓我們思想）更高的屋頂，那又在之上；然後是天使、總領天使、無數的非物質力量，天主自己的宮殿，聖父的寶座。\par

但語言（如我所說）太軟弱，無法完全描述。經驗是必要的，以及由經驗而來的認識。告訴我，你認為亞當在樂園中是怎樣的？這生命之路遠比那更好，如同天比地更好。\par

11. 但無論如何，讓我們再尋找另一個圖像。假設現在統治者是全世界的君王，他既不被戰爭也不被憂慮所困擾，只受尊崇，生活豐裕；有大量貢物，四面八方黃金湧向他，他被人仰望，你認為他會有什麼感受，如果他看到世界各地所有的戰爭都停止了？將會是這樣的事。但更確切地說，我甚至還沒有達到那個我所尋求的圖像；因此，我也必須再尋找另一個。\par

那麼，我請你思考：因為如同一個王室的孩子，只要他在母腹中，就對一切毫無知覺，但如果他忽然從那裡出來，登上王座，不是漸進地，而是一下子得到萬物；現在和未來的光景也是如此。或者，如果一個俘虜，受了無數的苦難，忽然被提到王座上。\par

但即使如此，我也沒有準確地達到那個圖像。因為在這裡，確實，一個人可能得到任何好東西，即使你說是王國本身，在第一天，他的慾望是充沛的，第二天、第三天也是如此，但隨著時間的推移，他雖然繼續享受快樂，但不再那麼巨大。因為無論是什麼，因熟悉而總是停止。但在那邊，不僅不減弱，反而增加。因為請思考，這是多麼偉大的事：一個靈魂在離開那裡之後，不再期待那些好事的終結，也不會改變，而是增加，且生命沒有終結，生命從一切危險、沮喪和憂慮中解脫，充滿歡樂和無數的祝福。\par

因為如果我們出去到平原，看到士兵的帳篷掛著帷幔，長矛、頭盔、盾牌的中心閃爍發光，我們驚奇不已；但如果我們也碰巧看到國王本人在中間奔跑，甚至騎著馬，穿著金色盔甲，我們就覺得擁有一切；你認為當你看到聖人們永恆的帳幕支搭在天上時，會是怎樣？（因為經上說：\BibleQuote{他們要接你們到永恆的帳幕裡去。} \BibleRef{路 16:9}）當你看到他們每一個都發光，比太陽的光芒更強烈，不是來自黃銅或鋼鐵，而是來自人的眼睛無法直視的榮耀？這當然是關於人。但是，如果有人要說成千上萬的天使、總領天使、革魯賓、色辣芬、寶座、宰制、率領、掌權者，他們的美麗無可比擬，超越一切理解？\par

但我要追趕那無法追趕的事，要到何時呢？\BibleQuote{眼睛未曾看見，}經上說，\BibleQuote{耳朵未曾聽見，人心也未曾想到的，就是天主為愛祂的人所預備的。} \BibleRef{格前 2:9} 因此，沒有比那些錯過的人更可憐，也沒有比那些得到的人更有福。那麼，讓我們成為有福的人，好能得到在我們主耶穌基督內的永恆福樂；願光榮、權能、尊崇歸於祂，連同聖父及聖神，從現在直到永遠，世世無窮。阿們。\par

\SourceRecord{Translated by Frederic Gardiner.；Nicene and Post-Nicene Fathers, First Series；Vol. 14.；Edited by Philip Schaff.；Buffalo, NY: Christian Literature Publishing Co.,；1889.；<http://www.newadvent.org/fathers/240206.htm>.}

% structure: p0001 source=h1 target=section reason=page-title

\section{《希伯来书》第七篇讲道}

% structure: p0002 source=h2 target=subsection reason=relative-heading-rank; base=h2

\subsection{希 4:11-13}

\BibleQuote{所以，我们应努力进入那安息，免得有人随同那不信的榜样而跌倒。因为天主的圣言是生活的、有效的，比双刃的剑更锐利，直穿透灵魂和神魂、关节与骨髓，且能辨别心中的思想和意向。}\par

1. 信德诚然伟大，带来救恩，没有它，人永远不能得救。然而，仅靠信德本身并不足成事，还需要有相应的品行。为此，保禄也劝勉那些已蒙受奥秘的人说：\BibleQuote{我们应努力进入那安息。} \BibleQuote{我们应努力}（他说），信德不足，还须加上生活，并须极其恳切；因为要升入天堂，确实需要极大的努力。如果那些在旷野中遭遇如此大苦难的人，尚且不配进入（应许的）土地，未能得着那地，是因为他们抱怨并犯了奸淫；那么，我们若漫不经心、懒散度日，又怎能配得天堂呢？因此，我们需要极大的努力。\par

请注意，惩罚不止于不能进入（他说：\BibleQuote{我们应努力进入那安息}，免得我们失去如此大的福分），但他还加上了最能警醒人的话。那是什么呢？\BibleQuote{免得有人随同那不信的榜样而跌倒。} 这是什么意思？意思是：我们应将心思、希望、期待放在那边，免得我们失败。因为（不然）我们将会失败，榜样已经表明：\BibleQuote{免得〔等等〕随同那同样的}，他说。\par

2. 其次，免得你听到\BibleQuote{随同那同样的〔榜样〕}，就以为惩罚是相同的，请听他接着说了什么：\BibleQuote{因为天主的圣言是生活的、有效的，比双刃的剑更锐利，直穿透灵魂和神魂、关节与骨髓，且能辨别心中的思想和意向。} 这些话表明，那圣言——天主——也成就了先前的事，且活着，并未熄灭。\par

所以，当你听到圣言时，不可轻看。因为\BibleQuote{祂比剑更锐利}，他说。请注意祂的俯就，并由此思考为何先知们也需要提及刀剑、弓和剑。经上说：\BibleQuote{你若不肯回头，祂必磨快祂的剑，祂已拉紧弓、预备妥当了。} \BibleRef{咏 7:12} 因为如今，经过这么长的时间，在他们得以成全之后，祂尚且不能单凭圣言的名号击倒人，而需要使用这些表达来显示（福音与法律）比较所得的优越性：更何况（古时）呢？\par

\BibleQuote{直穿透}，他说，\BibleQuote{灵魂和神魂之间}。这是什么？他暗示了更可怕的事。要么是指祂将神魂与灵魂分开，要么是指祂甚至能穿透无形体的它们，而不仅是像剑穿透躯体。这里他表明，灵魂也要受罚，并且祂彻底搜查最内在的东西，完全穿透整个人。\par

\BibleQuote{且能辨别心中的思想和意向，没有一个受造物在祂面前不是明显的。} 这些话最令他们恐惧。因为（他说），不要因为你们仍坚守信德就自信，而毫无确证。祂审判内心，因为祂穿过那里，既惩罚又搜查。\par

何必说人呢？他说。因为即使你谈论天使、总领天使、革鲁宾、色辣芬，甚至任何\BibleQuote{受造物}：一切都向那眼目敞开，一切都清楚明了；没有什么能逃过它；\BibleQuote{万物在那与我们有关系者的眼前，都是赤露敞开的。}\par

但\BibleQuote{敞开}是什么意思？〔这是〕一个比喻，取自从牺牲身上剥下的皮。因为如同在那情况下，人宰杀它们，将皮从肉上剥下，便显露出所有内在的部分，使之显明在我们眼前；同样，万有在天主面前也是赤露的。请你留意，他是何等不断地需要身体的比喻；这源于听者的软弱。因为他们软弱，他在说他们是\BibleQuote{迟钝}、\BibleQuote{仍需吃奶，不能吃干粮}时已表明。\BibleQuote{万物都是赤露的}，他说，\BibleQuote{在那与我们有关系者的眼前。} \EditorialNote{参 希 4:11-12}\par

3. 但\BibleQuote{随同那不信的榜样}是什么意思？就好比说，为什么古时的人没有看见那地？他们已经领受了天主能力的凭据，本应相信，却过于畏惧，对天主没有伟大的想象，并且灰心丧气——于是他们灭亡了。还有更多可说：在完成了旅程的大部分之后，正当他们在门口、在港湾时，却沉入了海中。我担心（他说）你们也会如此。这就是\BibleQuote{随同那不信的榜样}的意思。\par

因为连这些人（他所写信的对象）也受了许多苦，他后来作证说：\Quote{请你们回想先前的时日，那时你们才蒙光照，就忍受了痛苦的大考验。}\BibleRef{希 10:32} 所以谁也不要灰心，也不要因疲倦而接近终点时跌倒。因为确实有人，在开始时以全部热忱投入战斗，但稍后不肯坚持到底，就失去了一切。你们的祖先（他说）足以教导你们不要陷入相同的罪行，不要遭受他们所遭受的同样的事。这就是\Quote{同样背信的榜样}。我们不可灰心，他意思是（这也是他在书信末尾所说的）：\Quote{你们要举起下垂的手，和疲乏的膝盖}：\Quote{免得有人}，他说，\Quote{有人同样背信的榜样而跌倒}。\BibleRef{希 12:12} 因为这才是真正的跌倒。\par

然后，免得你们听见\Quote{有人同样背信的榜样而跌倒}时，以为他们也是遭受了同样的死亡，请看他说什么：\Quote{因为天主的圣言是生活的，是有效的，比各种双刃剑更锐利。}因为圣言落在这些人的灵魂上，比任何剑更严厉，造成重伤，并给予致命的一击。这些事无须他提供证明，也不必用论证确立，因为有如此可怕的历史。因为（他会说）什么样的战争毁灭了他们？什么样的剑？难道他们不是自己跌倒的吗？我们不可因没有遭受同样的事而疏忽。当\Quote{被称为今天}时，我们还有能力恢复自己。\par

因为恐怕我们在听到有关灵魂的事时变得疏忽，他也加上有关身体的事。因为那时，如同一位君王，当他的官员犯了重大过失时，先剥夺他们的指挥权（比方说），然后除去他们的腰带、官职和传令官，最后才惩罚他们：同样，圣神的剑也这样工作。\par

4. 接着他谈论圣子：\Quote{我们必要向祂交账}他说。什么是\Quote{我们必要向祂交账}？意思是我们要为我们所做的事向祂交账。那么，我们该如何行动，才能不跌倒、不灰心呢？\par

这些事（他会说）足以教导我们。但我们还有\Quote{一位伟大的大司祭，已进入了诸天，就是耶稣基督，天主子。}因为他补充了这一点，所以他接着说：\Quote{因为我们所有的，不是一位不能同情我们软弱的大司祭。}因此他上面说：\Quote{祂既然亲自受过试探，就能扶助受试探的人。}请看此处他也这样做。他所说的话大意是：祂走了我们现在所走的路，甚至更崎岖的路。因为祂经历了一切人性（的痛苦）。\par

% structure: p0018 source=h2 target=subsection reason=relative-heading-rank; base=h2

\subsection{希 4:14}

他上面说：\Quote{没有一个受造物在祂面前不是明显的}，暗示祂的天主性；然后，既然谈到了肉身，他再次以更迁就我们的方式讲论，说 \BibleRef{希 4:14}，\Quote{我们既然有一位伟大的大司祭，已进入了诸天}：并显示祂的关怀更大，保护他们如同自己的子民，不愿他们跌倒。因为梅瑟（他说）没有进入安息，而祂（基督）却进入了。奇妙的是，他从未明确陈述这一点，免得他们似乎找到借口；但他暗示了，但为了不显得控告那人，他没有公开说。因为如果，当这些事都未说过时，他们尚且提出这些指控，说\Quote{这人反对梅瑟和法律}\BibleRef{宗 21:21}；何况他若说\Quote{不是巴勒斯坦，而是天堂}，他们会说出比这些更厉害的话。\par

5. 但他没有把一切都归于大司祭，也要求我们自身的努力，即我们的宣认。因为他说：\Quote{我们既然有一位伟大的大司祭，已进入了诸天，就是耶稣基督，天主子，就要坚持我们的宣认。}他指的是哪种宣认？即：有复活，有报应；有无数的好东西；基督是天主，信仰是正确的。让我们宣认这些，让我们坚持这些。因为它们的真实性从大司祭在内的事实得到证明。我们没有失望，让我们宣认；虽然现实尚未显现，但让我们宣认；如果已经显现，那就是谎言。所以这也是真实的：我们的（好事）被延迟了。因为我们也的确有一位伟大的大司祭。\par

% structure: p0021 source=h2 target=subsection reason=relative-heading-rank; base=h2

\subsection{希 4:15}

\Quote{因为我们所有的，不是一位不能同情我们软弱的大司祭。}他不是（他意思是）对我们的处境无知，不像许多大司祭，他们不认识在苦难中的人，也不知道任何时候有苦难。因为在人的情况中，一个未经历实际感受的人不可能知道受苦者的痛苦。我们的大司祭经历了所有的事。因此祂先受苦，然后上升，以便能同情我们。\par

\Quote{在各方面都与我们相似受过试探，只是没有罪。}请注意他上面用了\Quote{同样}一词，此处用了\Quote{相似}。\BibleRef{希 2:14} 这就是说，祂受过迫害，被吐唾沫，被控告，被嘲笑，被诬告，被驱逐，最后被钉在十字架上。\par

\Quote{相似我们，没有罪。}这些话也暗示另一件事：即使在苦难中，人也能没有罪地度过。所以当他说\Quote{肉身的相似}\BibleRef{罗 8:3} 时，意思不是祂只取了\Quote{肉身的相似}，而是\Quote{肉身}。那么他为什么说\Quote{相似}？因为他是在谈论\Quote{罪恶的肉身}：因为它相似我们的肉身，因为其本质与我们相同，但在罪上不再相同。\par

% structure: p0025 source=h2 target=subsection reason=relative-heading-rank; base=h2

\subsection{希 4:16}

6. \Quote{那么，让我们怀着信心，走近祂恩宠的宝座，以获得仁慈，并找到恩宠，作及时的扶助。}\par

祂所说的 \Quote{恩宠的宝座} 是指什么？就是那座王室的宝座，经上论及它说：\Quote{主对我主说：你坐在我的右边。} \BibleRef{咏 109:1}\par

什么是 \Quote{让我们怀着信心前来}？因为有一位无罪的\OriginalTerm{大司祭}{High Priest} 与世界抗争。因为祂说：\Quote{放心，我已战胜了世界} \BibleRef{若 16:33}；因为这就是经受一切，却仍纯洁无罪的。虽然我们（他意指）处于罪下，但祂却是无罪的。\par

我们应当如何 \Quote{大胆地前来}？因为 \Emph{现在} 是恩宠的宝座，不是审判的宝座。因此大胆地，\Quote{使我们获得仁慈}，正是我们所寻求的。因为这事是慷慨的，是王室的恩赐。\par

\Quote{并找到恩宠，作及时的扶助。} 祂说得对：\Quote{作及时的扶助。} 如果你现在前来（他意指），你将获得恩宠和仁慈，因为你是在 \Quote{适当的时机} 前来；但如果你 \Emph{那时} 前来，便不再〔得到它〕。因为 \Emph{那时} 前来是不合时宜的，因为那时不再是 \Quote{恩宠的宝座}。直到那时，祂坐着施予宽恕，但当终结〔来临〕时，祂便起来审判。因为经上说：\Quote{天主，求你起来审判大地。} \BibleRef{咏 81:8}（\Quote{让我们怀着信心前来}，或者他又说，没有 \Quote{邪恶的良心}，即不怀疑，因为这样的人不能 \Quote{怀着信心前来}。）为此，经上说：\Quote{在悦纳的时候，我俯允了你；在救恩的日子，我援助了你。} \BibleRef{格后 6:2} 因为甚至 \Emph{现在}，对于那些在圣洗后犯罪的人，能找到悔改，也是出于恩宠。\par

但为了不让你听到\OriginalTerm{大司祭}{High Priest}时，以为祂是站着的，他立即引到宝座。然而司祭不是坐着，而是站着。你看到吗？祂被立为\OriginalTerm{大司祭}{High Priest}，并非出于本性，而是出于恩宠、屈尊和谦卑？\par

现在我们也该说：\Quote{让我们走近}，并 \Quote{怀着信心} 祈求：只要我们带来信德，祂就赐予一切。现在是恩赐的时刻；愿无人对自己绝望。那时将是绝望的时刻——当婚宴厅关闭，当君王进来观看宾客，当那些被算作相称的人，将得到祖宗的怀抱作为他们的份：但现在还不是这样。因为观众仍在聚集，比赛仍在进行，奖品仍在悬念中。\par

7. 那么，让我们努力。因为连保禄也说：\Quote{我奔跑，不像无定向的。} \BibleRef{格前 9:26} 需要奔跑，且要奋力奔跑。那奔跑的，看不见迎面而来的人；无论他经过草地，还是干旱之地：奔跑的人不看观众，只看奖品。无论他们是富是贫，无论有人讥笑他或称赞他，无论有人侮辱、向他扔石头、或掠夺他的房屋，无论他看见儿女、妻子或任何事物。他只专注于一件事：奔跑，赢得奖品。奔跑的人从不站住，因为即使稍一松懈，就失去了一切。奔跑的人不仅在终点前毫不松懈，反而那时更加奋力加快速度。\par

我说这话，是针对那些说：在我们年轻时我们操练，在我们年轻时我们禁食，\Emph{现在}我们老了。\Emph{现在}你更应该使你的谨慎更加严格。不要向我数算从前做得特别好的事：现在要变得年轻而充满活力。因为那奔跑这有形的赛跑的人，当白发临到他时，很可能不能像以前那样跑：因为整个竞赛依赖身体；但你呢——为什么你减慢了速度？因为在这赛跑中，需要灵魂，一个完全醒悟的灵魂：而灵魂在老年时反而更坚强；那时它处于全盛时期，那时它处于巅峰。\par

因为正如身体，当被热病和接二连三的疾病所压迫时，即使强壮，也会疲惫；但当它摆脱了这场攻击，就恢复它本有的力量；同样，灵魂在青年时期是狂热的，主要被荣耀的爱、奢华的生活、感官的欲望和许多其他幻想所占据；但老年到来时，驱散了所有这些激情，有些是由于满足，有些是由于智慧。因为老年放松了身体的力量，并且不让灵魂甚至想使用时也使用它们，而是抑制它们如同各种敌人，将它安置在无忧之境，并产生极大的平静，并带来更大的敬畏。\par

因为即使没有别的，据说，那些年长的人知道他们正走向终点，并且他们确实临近死亡。因此，当此生的欲望在消退，而对审判座的期待临近时，软化灵魂的顽固，它难道不会变得更专心，如果人愿意的话？\par

8. 那么，你说，当我们看见老人比年轻人更难对付时？你告诉我这是邪恶的过度。因为在疯子身上，我们也看见他们跌落悬崖，无人推他们。因此，当一个老人患有年轻人的疾病时，这是邪恶的过度；此外，即使在青年时，这样的人也不会有借口：因为他不能说：\Quote{不要记念我青年时的罪过和我的无知。} \BibleRef{咏 24:7} 因为那些在老年时保持不变的人显明，即使在青年时，他之所以如此不是由于无知，也不是由于缺乏经验，也不是由于年龄阶段，而是由于懈怠。因为那个能说\Quote{不要记念我青年时的罪过和我的无知}的人，是那些行为如老人的人，那些在老年中改变的人。但如果即使在老年他继续同样的不当行为，这样的人怎能配得上老人的名号？他连年龄本身也不尊重。因为那个说\Quote{不要记念我青年时的罪过和我的无知}的人，说这话时是作为在老年中行为正直的人。那么，不要借着老年时的行为，也剥夺自己青年时罪过的宽恕。\par

因为所做的事怎能不是不合理且不可原谅的呢？一个老人坐在酒馆里。一个老人赶去看赛马——一个老人走进剧院，像孩子一样与人群奔跑。真是羞耻与嘲弄，外表用白发装饰，内心却有孩童的心思。\par

如果年轻人侮辱他，他立刻提出他的白发。你自己首先要尊重它们；然而，如果你不尊重你自己的白发，即使在老年，你又怎能要求年轻人尊重它们呢？你不尊重白发，反而使它蒙羞。天主用白发尊荣了你：祂赐予你崇高的尊严。你为何背叛这尊荣？年轻人怎能尊重你，当你比他更放荡？因为白发之所以可敬，是当它行事与白发相称时；但当它扮演青年时，它会比年轻人更可笑。那么，你们这些老人，当你们被自己的混乱沉醉时，又怎能向年轻人提出这些劝诫呢？\par

9. 我说这些事，不是指责老人，而是指责年轻人。因为在我看来，那些行事如此的人，即使活到百岁，也是年轻的；正如年轻人，如果还是小孩子，但若谨慎，也比老人更好。这教训不是我的，而是圣经也承认同样的区别。因为经上说：\Quote{可敬的年龄并不是长寿，无瑕的生命才是老年。} \BibleRef{智 4:8}\par

因为我们尊敬白发，并非我们认为白色比黑色更珍贵，而是因为它是美德生活的证据；当我们看见它们时，我们由此推断内在的成熟。但如果人们继续做出与白发不相称的事，他们反而会更加可笑。因为我们也尊敬皇帝，以及紫袍和冠冕，因为它们是职位的象征。但如果看见他穿着紫袍，被唾弃，被卫兵踩在脚下，被掐住喉咙，被投入监狱，被撕碎，我们还会尊敬紫袍或冠冕吗？我们岂不反而为这排场本身哀哭？那么，不要因你的白发要求被尊敬，当你自己委屈了它。因为它本当从你那里得到满足，因为你使如此高贵和可敬的形像蒙羞。\par

我们不是说这些事反对所有老人，我们的论说也不是简单地反对老年（我没有那么疯狂），而是反对使老年蒙羞的年轻精神。也不是关于那些已经年老的人我们悲伤地说这些事，而是关于那些使白发蒙羞的人。\par

因为老人是一位君王，如果你愿意，他比穿着紫袍的人更王者，如果他控制自己的激情，并使它们服从，如同卫兵的等级。但如果他被拖拽，被从宝座上推下，成为贪财、虚荣、装饰、奢侈、醉酒、愤怒、感官享乐的奴隶，用油梳妆头发，并以生活方式显示受辱的年龄，这样的人应得什么惩罚呢？\par

10. 但希望你们不是那样，年轻人！因为对你们也没有犯罪的借口。为什么？因为在青年时可以变老：正如在老年时也有青年，反之亦然。因为如同白发不能拯救一个人，同样黑发也不是障碍。因为如果老人做我所说的这些事是可耻的，那对年轻人来说就更可耻，但年轻人仍未免于指责。因为年轻人只有一个借口，即他被召去管理事务，当他仍无经验，需要时间和实践时；但当他需要表现节制和勇气时，或需要保护自己的财物时，就不再有了。\par

因为有时年轻人比老人更被责备。因为老人需要许多照顾，年老使他衰弱；但年轻人，如果愿意，能够自己照顾自己，当他比老人更掠夺，记得伤害，轻蔑，不挺身保护他人，说出许多不合时宜的话，傲慢，谩骂，醉酒时，他应该得到什么借口呢？\par

如果他在贞洁方面认为自己不能被指控，那么要想想他在这里也有许多帮助，如果他愿意。尽管欲望比他搅动老人更猛烈，但无论如何，他能做的比老人多得多，从而驯服那野兽。这些是什么？劳动、阅读、彻夜不眠、禁食。\par

11. 那么，这些事与我们有何相干？有人说：我们又不是隐修士。你对我说这话？你对保禄说吧，当他说：\Quote{以各种祈求和哀恳，时时醒寤祈祷}\BibleRef{弗 6:18}，当他说：\Quote{不要为肉体作准备，以纵情恣欲}\BibleRef{罗 13:14}。诚然，他写这些事不是只写给隐修士，而是写给所有住在城市里的人。因为生活在世俗中的人，除了与妻子同住这一点外，岂应在其他事上胜过隐修士吗？在这件事上他获得许可，但在其余事上并无不同，反而他有义务与隐修士一样去行一切事。\par

再者，基督所宣讲的\OriginalTerm{真福八端}{Beatitudes}并非只针对隐修士；因为若果如此，整个世界早已丧亡，我们就要归咎于天主的残忍了。如果这些真福只对隐修士讲述，而世俗人无法实行，但天主又允许婚姻，那么祂就毁灭了所有人。因为如果结婚不可能履行隐修士的本分，那么一切都已丧亡毁灭，美德的功能就被局限在狭窄的范围内。\par

而且，婚姻如何能是尊贵的，如果它如此阻碍我们？那么怎么样呢？即使我们有妻子，仍可能，非常可能，追求美德，只要我们愿意。如何做？如果我们拥有\Quote{妻子}，却\Quote{像没有一样}；如果我们不为我们的\Quote{财产}而欢喜；如果我们\Quote{享用世界，却非滥用它}\BibleRef{格前 7:29}。\par

如果有人因婚姻状态受到阻碍，让他们知道：障碍不是婚姻，而是那滥用婚姻的意向。因为醉酒并非由酒造成，而是由邪恶的意向和过量饮用所致。节制地使用婚姻，你将在天国为首，并享受一切美善。愿我们众人因我们的主耶稣基督的恩宠和慈爱，偕同圣神，归光荣于圣父，于无穷世。阿们。\par

\SourceRecord{Translated by Frederic Gardiner.；Nicene and Post-Nicene Fathers, First Series；Vol. 14.；Edited by Philip Schaff.；Buffalo, NY: Christian Literature Publishing Co.,；1889.；<http://www.newadvent.org/fathers/240207.htm>.}

% structure: p0001 source=h1 target=section reason=page-title

\section{论希伯来人书第八篇讲道}

% structure: p0002 source=h2 target=subsection reason=relative-heading-rank; base=h2

\subsection{希伯来书5:1-3}

\Quote{ \Emph{凡从人间选取的大司祭，是奉派为人行关于天主的事，为奉献供物和牺牲，以赎罪过，} \Emph{他能同情无知和迷失的人，因为他自己也为软弱所困；因此，他怎样为人民奉献赎罪祭，也当怎样为自己奉献。}}\par

1. 有福的保禄希望在接下来的论证中表明，这盟约远比旧约更优越。他先提出一些间接的考量来做到这一点。因为既然没有任何有形或显耀的事物——没有圣殿，也没有至圣所，没有穿着如此华服的大司祭，没有法律上的规条——反而一切都是更高超、更成全的，且没有属体的事物，一切都是属神的，而属神的事物并不像属体的事物那样吸引软弱的人；于是他彻底审察了整个问题。\par

请注意他的智慧：他首先从司祭开始，并不断称祂为大司祭，从这第一点就显出（两个时期的）区别。为此他首先界定了司祭是什么，并表明祂是否具有司祭应有的特质，以及是否有司祭职的标记。然而有一个反对意见摆在他面前：祂（基督）甚至并非出身高贵，也不属于司祭支派，且在地上不是司祭。那么，祂如何是司祭呢？有人可能会说。\par

正如在\WorkTitle{罗马书}中，他提出一个他们不容易被说服的论点——即信德成就了法律辛劳和日常生活的汗水所不能成就的事——他便转向了圣祖，并将整个问题回溯到那个时代。同样，现在在这里，他开启了司祭职的另一条路径，通过先前发生的事来显示其优越性。如同在惩罚的问题上，他不仅将地狱摆在他们面前，也摆在他们祖先身上发生的事；同样，现在在这里，他也首先从当前的事物来确立这个立场。因为固然属天的事物应由属地的事物来证明，但当听众软弱时，则采取相反的做法。\par

2. 他先设定了那些共同（于基督和他们的司祭）的事物，然后表明祂是更优越的。因为比较性的卓越正是这样产生的：在某些方面有共同之处，在另一些方面有优越之处；否则就不再是比较了。\par

\Quote{凡从人间选取的大司祭，}这是与基督共同之处；\Quote{是奉派为人行关于天主的事，}这也是共同的；\Quote{为奉献供物和牺牲为人民赎罪，}这也是共同的，却不完全；接下来的就不再是这样了：\Quote{他能同情无知和迷失的人，}从这点开始是优越之处，\Quote{因为他自己也为软弱所困；因此，他怎样为人民奉献赎罪祭，也当怎样为自己奉献。}\par

% structure: p0009 source=h2 target=subsection reason=relative-heading-rank; base=h2

\subsection{希伯来书5:4-5}

接着还有另一些（论点）：祂是由另一位所立（他说），并非擅自闯入（这职分）。这也是共同的\BibleRef{希 5:4}：\Quote{谁也不能自己取得这尊荣，而应蒙天主召选，如同亚郎一样。}\par

在这里，他又在另一点上安抚他们：因为祂是由天主派遣的；这是基督常对犹太人说的。\Quote{派遣我的父比我大}，以及\Quote{我不是凭自己而来的。}\BibleRef{若 12:49}\par

在我看来，这些话也暗示了犹太人的司祭，他们不再是司祭，而是擅自闯入司祭职并败坏它的人；\BibleRef{希 5:5}\Quote{照样，基督也没有自取作大司祭的光荣。}\par

那么祂是如何被任命的呢（有人说）？因为亚郎多次被任命，譬如藉着棍杖，又当有火降下焚毁那些想要擅闯司祭职的人。但在此情形中，他们（犹太司祭）不仅未受任何伤害，反而享有崇高的声望。那么，祂的任命从何而来？他从预言中指明。他没有任何感官可察觉、可见的事物。为此，他藉着预言，藉着未来的事来肯定这一点：\Quote{但是那向祂说的：\InnerQuote{祢是我的儿子，我今日生了祢}。}这与圣子有什么关系？是的（他说），这是为祂由天主任命作预备。\par

% structure: p0014 source=h2 target=subsection reason=relative-heading-rank; base=h2

\subsection{希伯来书5:6}

\Quote{又如祂在另一处说：\InnerQuote{祢照默基瑟德的品位永为司祭。}}这话如今是对谁说的呢？\par

谁是\Quote{照默基瑟德的品位}？没有别的（祂）。因为他们都处于法律之下，他们都守安息日，他们都受割损；不能指出任何别的（人）。\par

% structure: p0017 source=h2 target=subsection reason=relative-heading-rank; base=h2

\subsection{希伯来书5:7-8}

3. \Quote{祂在血肉之身时，以大呼声和眼泪，向那能救祂脱离死亡者，献上祈祷和恳求，就因祂的虔敬而获得了俯听；祂虽然是儿子，却由所受的苦难学习了服从。}你看见了吗？他不过是在陈述祂的照顾和祂爱的伟大。\Quote{以大呼声}是什么意思？福音从未如此说，也没有说祂祈祷时哭泣，也没有说祂发出呼声。你看见了吗？这是迁就。因为他不能仅仅说祂祈祷，还加上\Quote{以大呼声}。\par

% structure: p0019 source=h2 target=subsection reason=relative-heading-rank; base=h2

\subsection{希伯来书5:9-10}

\Quote{祂虽是儿子，却由所受的苦难学习了服从。} \Quote{祂既成为成全的，便为一切服从祂的人成了永远救恩的根源，被天主称为照默基瑟德品位的大司祭。}\par

既是\Quote{呼声}，为何又是\Quote{大呼声和眼泪}？\par

\BibleQuote{奉献了祭品}（他说），\BibleQuote{又因祂的敬畏而获得了俯听}。你说什么？让异端者羞愧。天主子\BibleQuote{因祂的敬畏而获得了俯听}。至于先知们，人还能说什么呢？又有什么关联呢，说\BibleQuote{祂因敬畏而获听，虽然是圣子，却由所受的苦难学习了服从}？谁会将这些话用于天主？哪有如此疯狂的人？即使神志不清，谁会说出这些话？\BibleQuote{因祂的敬畏而获得了俯听}（他说），\BibleQuote{祂由所受的苦难学习了服从}。什么服从？在此之前祂已顺从至死，如同圣子对圣父，后来祂怎能学习呢？你看，这话是针对降生成人说的吗？\par

现在告诉我，祂是否祈求圣父救祂脱离死亡？正是为此祂\BibleQuote{极度忧伤，说：\InnerQuote{若是可能，就让这杯离开我}}\BibleRef{玛 26:38}？然而祂从未为复活祈求圣父，反而公开宣称：\BibleQuote{拆毁这座圣殿，三天内我要把它重建起来}\BibleRef{若 2:19}。以及\BibleQuote{我有权捨掉我的生命，也有权再取回它；没有人夺去我的生命，是我自己捨掉的}\BibleRef{若 10:18}。那么这是怎么回事？祂为何祈祷？（祂又说：\BibleQuote{看，我们上耶路撒冷去，人子要被交给司祭长和经师，他们要定祂死罪，并把祂交给外邦人戏弄、鞭打、钉死，但第三天祂要复活}\BibleRef{玛 20:18}，却没有说\BibleQuote{我的圣父要使我复活}。）那么祂为此祈祷了吗？但祂为谁祈祷？为那些信祂的人。\par

他的意思是：\Quote{祂很容易被垂听。}因为当时他们对祂还没有正确的看法，所以他说祂被垂听了。正如祂自己在安慰门徒时所说的：\BibleQuote{如果你们爱我，就该喜欢，因为我往圣父那里去}\BibleRef{若 14:28}，以及\BibleQuote{圣父比我大}。然而，祂\BibleQuote{使自己空虚}\BibleRef{斐 2:7}、自愿交付的那一位，又怎能不自显荣耀呢？因为经上说：祂\BibleQuote{为我们捨了自己}\BibleRef{迦 1:4}。又说：\BibleQuote{祂捨了自己作众人的赎价}\BibleRef{弟前 2:6}。那么这是怎么回事？你看出，关于祂的卑微之事都是指肉身而言。同样，这里说\BibleQuote{祂虽然是圣子，却因敬畏而获听}。祂想表明，成功是出于祂自己，而非天主的恩宠。祂的敬畏如此之大（他说），以至于天主因此看重祂。\par

\BibleQuote{祂学习了}，他说，服从天主。这里他又展示了受苦的巨大益处。\BibleQuote{祂既被成全}，他说，\BibleQuote{便成了所有服从祂的人永远救恩的根源}。（参见前文，第384、391页。）但祂既是圣子，尚且从苦难中学习了服从，何况我们呢？你看见他讲论了多少关于服从的事，为劝服他们服从？因为在我看来，他们不肯受约束。\BibleQuote{从所受的苦难}，他说，祂\BibleQuote{学会了}不断服从天主，并且\BibleQuote{因苦难被成全}。这就是成全，我们必须借此达到成全。因为祂不仅自己得救，还为他人成为救恩的丰富供应。因为\BibleQuote{祂既被成全，便成了所有服从祂的人永远救恩的根源}。\par

% structure: p0026 source=h2 target=subsection reason=relative-heading-rank; base=h2

\subsection{希 5:11-12}

4. \BibleQuote{被称为天主按照\Person{默基瑟德}品位的大司祭}\BibleRef{希 5:11}；\BibleQuote{关于这事我们有许多话要说，且难以解释。}当他准备论述司祭职的区别时，他先责备他们，指出这样的迁就就如\BibleQuote{奶}，并且因为他们还是婴孩，他更长久地停留在卑微、关乎肉身的主题上，像谈论任何义人一样谈论祂。注意，他既没有完全沉默不言，也没有完全说出教义；一方面为提升他们的思想，劝勉他们成为成全的人，使他们不至于失去伟大的教义；另一方面，又不至于压垮他们的心智。\par

\BibleQuote{关于这事}，他说，\BibleQuote{我们有许多话要说，且难以解释，因为你们听觉迟钝。}因为他们不听从，教义便\BibleQuote{难以解释}。因为当人面对那些不与你同行、不留意所讲之事的人时，他很难向他们解释清楚。\par

但也许你们中间有人困惑，认为由于希伯来人，他自己被妨碍听到更成全的教义。反之，我认为或许在这里，除了少数人，也有许多像他们一样的人，所以这话也可以针对你们自己；但为了那少数人，我将说话。\par

那么他是否完全沉默，还是后来继续了这个主题，就像在致罗马人书中所做的一样？因为在那里，他先堵住了反驳者的口，说道：\BibleQuote{然而，人哪，你是谁，竟敢向天主辩论？}\BibleRef{罗 9:20}，然后才附上了解答。在我看来，他既没有完全沉默，也没有完全说出来，为引导听众渴望知识。因为他提到了这个主题，并说教义中储存着某些伟大的事物，请看他是如何将责备与赞颂结合在一起。\par

因为\Person{保禄}的智慧总有一部分是混合痛苦与仁慈。他在\WorkTitle{迦拉达书}中也这样做，说：\BibleQuote{你们跑得好，谁阻碍了你们？} \BibleRef{迦 5:7}又说：\BibleQuote{你们受了这许多苦，难道都是徒然的吗？如果真是徒然的。} \BibleRef{迦 3:4}以及\BibleQuote{我在主内对你们有信心。} \BibleRef{迦 5:10}他对这些（希伯来人）也说：\BibleQuote{但我们深信你们有更好的事，那些与救恩相伴的事。} \BibleRef{希 6:9}因为他达到两个效果：既不过分紧张他们，也不让他们跌倒；因为如果别人的例子足以激发听者并引导他效法，那么当一个人以自己为榜样并被要求效法自己时，可能性也随之而来。因此他也表明这一点，不让他们像完全被定罪的人那样跌倒，也不像总是作恶的人，而是（说）他们曾经甚至是良善的；\BibleRef{希 5:12}因为他说：\Quote{按时间来说，你们本应做导师。}这里他表明他们曾经是信者有很长时间，也表明他们本应教导别人。\par

5. 总之，看他如何不断努力引入关于大司祭的论述，却又一再推迟。因为请听他如何开始：\BibleQuote{既然我们有一位伟大的大司祭，已经穿越了诸天} \BibleRef{希 4:14}；没有说他如何伟大，他又说：\BibleQuote{每一位从人间选取的大司祭，是为人在关乎天主的事上被设立的。} \BibleRef{希 5:1}又说：\BibleQuote{同样，基督也没有自取荣耀成为大司祭。} \BibleRef{希 5:5}在说了\BibleQuote{你永远按照默基瑟德的品位作司祭} \BibleRef{希 5:6}之后，他又推迟（主题），说：\BibleQuote{祂在血肉的日子，献上祈祷和恳求。} \BibleRef{希 5:7}因此，当他多次被拒绝后，他仿佛为自己辩解般说：错在你们。唉！多么大的差别！当他们应该教导别人时，他们甚至不是简单的学习者，而是学习者中的最末等。\BibleRef{希 5:12}\BibleQuote{因为按时间来说，你们本应做导师，却还需要有人再把天主圣言初阶教导你们。}这里他指的是（基督的）人性。因为在世俗文学中，必须先学习字母，同样在这里，他们首先被教导关于人性的事。\par

你看他说卑微之事的原因是什么。\Person{保禄}也对\Place{雅典}人这样做，讲论说：\BibleQuote{天主曾容忍了这无知的时代，如今却命令各处的人都要悔改，因为祂已定下日子，要借祂所设立的人按正义审判天下，并且叫祂从死人中复活，给众人一个确据。} \BibleRef{宗 17:30}因此，如果他提到任何崇高之事，就简短表达，而卑微的陈述则散布在书信的许多部分。这样他也表明了崇高；因为所说的卑微本身就禁止人们怀疑这些事涉及天主性。所以这里也保持了稳妥的立场。\par

但这种迟钝是由什么引起的呢？他在\WorkTitle{格林多前书}中特别指出了这一点，说：\BibleQuote{因为在你们中间有嫉妒、纷争和分裂，你们岂不是属血肉的吗？} \BibleRef{格前 3:3}但我恳请你们注意他极大的智慧，看他如何总是按照眼前的病症行事。因为那边的软弱更多来自无知，或者说来自罪恶；但这里不仅来自罪恶，也来自持续的苦难。因此他也使用能表明区别的表达，不是说\Quote{你们成了属血肉的}，而是说\Quote{迟钝}；那时是\Quote{属血肉的}，而这里痛苦更大。因为那些人（格林多人）确实不能忍受（他的责备），因为他们是属血肉的；但这些（希伯来人）却能忍受。因为他说\BibleQuote{你们成了听觉迟钝} \BibleRef{希 5:11}，这表明他们从前是健全的、强壮的、热忱炽热，这也是他后来对他们所作的见证。\par

6. \Quote{并且成了需要奶，而不是固体食物的人。}他常把卑微的道理称为\Quote{奶}，无论在这里还是别处。他说：\Quote{按时间来说，你们本应做导师}：正因为那时间，你们尤其应当强壮，却偏偏因此退步。现在他称它为\Quote{奶}，因为它适合较单纯的人。但对于较成全的人，它是有害的，在这些事上停留是有损的。所以现在不应引入法律的事，也不应从它们做出比较，（例如）祂是一位大司祭，献了祭，需要呼喊和恳求。因此，看这些事对我们如何不健康；但那时它们滋养了他们，对他们绝非不健康。\par

所以，天主的圣言是真正的滋养。因为他说：\BibleQuote{我必不降下缺粮的饥荒，也不缺水的口渴，而是缺了听天主言语的饥荒。} \BibleRef{亚 8:11}\par

\BibleQuote{我给你们喝奶，没有给固体食物} \BibleRef{格前 3:2}；他没有说\Quote{我喂养了你们}，表明这样的（滋养）不是食物，而是（情形）像那些不能用面包喂养的小孩子。因为这样的人不是被给饮料，而是他们的食物代替了饮料。\par

再者，他没有说\Quote{你们需要}，而是说\Quote{你们成了需要奶而不是固体食物的人}。也就是说，你们愿意了；你们使自己陷入这种境地，这种需要。\par

% structure: p0039 source=h2 target=subsection reason=relative-heading-rank; base=h2

\subsection{\BibleRef{希 5:13}}

\Quote{凡只能吃奶的，对正义的言语都是生疏的，因为他还是个婴孩。} 什么是\Quote{正义的言语［教义］}？在我看来，他似乎也在此暗示品行。这也是基督所说的：\Quote{除非你们的正义超过经师和法利塞人的正义\BibleRef{玛 5:20}}，他同样说：\Quote{对正义的言语生疏}，即，对超越的哲学生疏的人，无法接受完美而精确的生活。或者，他在这里用\Quote{正义}指基督以及关于祂的高深教义。\par

他说他们\Quote{变得迟钝}，但他没有说明原因，让他们自己知道，也不愿使自己的言论难以忍受。但在迦拉达书的事上，他既\Quote{惊奇}\BibleRef{迦 1:6}，又\Quote{疑惑}\BibleRef{迦 4:20}，这更激励人，如同一个从未预料到此事会发生的人所说的话。因为疑惑所表明的就是这个。\par

你看，另有一种婴孩，另有一种成年。让我们成为这意义上的\Quote{成年}吧：即使是儿童和少年，也能达到那\Quote{成年}：因为这不是本性的，而是德行的。\par

% structure: p0043 source=h2 target=subsection reason=relative-heading-rank; base=h2

\subsection{\BibleRef{希 5:14}}

7. \Quote{唯独长大成人的才能吃硬食，他们的官能因习用而操练出来，能分辨善恶。} 那些人没有\Quote{操练他们的官能}，也不\Quote{知道善恶}。他在这里说\Quote{分辨善恶}，并非指生活［品行］，因为这是人人都可能且容易知道的，而是指健康崇高与败坏低劣的教义。婴孩不懂得辨别好坏的食物。至少，他常常把脏东西放进嘴里，取用有害之物；他做一切事都不加判断；但成年人则不然。这样的人就是那些轻易听从一切、不加选择地交出耳朵的人：在我看来，这也责备了这些［希伯来人］，因为他们轻易\Quote{被飘来飘去}，有时跟随这些，有时跟随那些。他在［这封书信］结尾也暗示了这一点，说：\Quote{不要被各种异端道理所迷惑。\BibleRef{希 13:9}} 这就是\Quote{分辨善恶}的意思。\Quote{口舌尝食物，灵魂辨言语。\BibleRef{约 34:3}}\par

8. 让我们学习这教训。当你听说一个人不是外邦人也不是犹太人时，不要立刻相信他是基督徒；还要检查所有其他方面；因为甚至\OriginalTerm{摩尼派}{Manichæans}和所有异端都戴上了这个面具，以便欺骗更单纯的人。但如果我们灵魂的\Quote{官能}已\Quote{操练出来，能分辨善恶}，我们就能辨别这样的［教师］。\par

但我们的\Quote{官能}如何变得\Quote{操练}？通过持续聆听；通过圣经的经验。因为当我们陈述那些［异端者］的错误，你今日、明日聆听，并证明那是不对的，你就学会了全部，就知道了全部：即使你今天不明白，明天也会明白。\par

他说：\Quote{他们的官能操练出来。} 你看，我们必须通过神圣的研究操练我们的听觉，使它们不感到陌生。\Quote{操练出来}，他说，\Quote{为了分辨}，即，为了熟练。\par

有人说没有复活；另一个人不期待任何未来的事；另有人说有一位不同的神；另有人说祂始于玛利亚。你看他们如何因缺乏节制而全部跌倒，有些过度，有些不足。例如，最早的异端是\OriginalTerm{马西昂}{Marcion}的异端；它引入了另一位不同的神，并不存在。看那过度。之后是\OriginalTerm{撒伯流}{Sabellius}的异端，说圣子、圣神和圣父是一体。然后是\OriginalTerm{马策禄}{Marcellus}和\OriginalTerm{福提诺}{Photinus}的异端，陈说相同的事。此外还有\OriginalTerm{撒莫萨塔的保禄}{Paul of Samosata}的异端，说祂始于玛利亚。之后是\OriginalTerm{摩尼派}{Manichæans}的异端；因为这是所有之中最近的。在这些之后是\OriginalTerm{阿黎乌}{Arius}的异端。还有其他异端。\par

正因如此，我们领受了信德，使我们不必被迫攻击无数的异端并处理它们，而是无论任何人试图增加或减少什么，我们都视之为虚假。因为正如那些给予标准的人不迫使［人们］忙于无数的度量，而是吩咐他们持守所给予的；教义方面也是如此。\par

9. 但没有人愿意留意圣经。因为如果我们留意，不仅我们自己不会被欺骗所缠住，还能帮助那些被欺骗的人获得自由，并将他们从危险中拉出来。因为强壮的士兵不仅能自救，还能帮助同伴，并使同伴脱离敌人的恶意。但事实上，有些人甚至不知道有圣经。然而，圣神做了如此多的明智安排，以便它们能被安全地保存。\par

从起初看，好让你学到天主不可言说的爱。祂感动了有福的\Person{梅瑟}；祂刻了石版；祂使他在山上停留了四十天；又同样多的天数来颁赐法律。之后祂派遣了先知，他们遭受了无数的苦难。战争来了；他们全都杀了他们，把他们砍成碎片，书卷被烧了。祂又感动了另一位可钦佩的人，即是\Person{厄斯德拉}，出版它们，并使它们从残存中汇编起来。之后祂安排它们由\OriginalTerm{七十贤士}{the seventy}翻译。他们翻译了它们。基督来了，祂接受了它们；宗徒们将它们分给人间。基督行了征兆和奇事。\par

那么，在如此巨大的辛劳之后呢？宗徒们也写了，正如保禄所说的：\BibleQuote{他们写下来，是为了劝诫我们这些活在末世的人。}\BibleRef{格前 10:11} 基督又说：\BibleQuote{你们错了，因为你们不明白圣经。}\BibleRef{玛 22:29} 保禄又说：\BibleQuote{藉着圣经所生的忍耐和安慰，我们可以怀着望德。}\BibleRef{罗 15:4} 又说：\BibleQuote{凡圣经都是天主所默示的，于教导是有益的。}\BibleRef{弟后 3:16} 又说：\BibleQuote{让基督的话丰富地住在你们内。}\BibleRef{哥 3:16} 先知说：\BibleQuote{他昼夜默思天主的法律。}\BibleRef{咏 1:2} 又在另一处说：\BibleQuote{愿你们的一切谈论，都在至高者的法律中。}\BibleRef{德 9:15} 又说：\BibleQuote{你的言语在我的上颚多么甘甜！}（他没有说在我的耳中，而是说在我的上颚）；\BibleQuote{比蜂蜜和蜂巢更甜。}\BibleRef{咏 118:103} 梅瑟说：\BibleQuote{你们要不断默想它们，无论起来、坐下、躺卧。}\BibleRef{申 6:7} \BibleQuote{你要专务这些事。}\BibleRef{弟前 4:15} 关于它们，可以说数不清的事。然而，尽管如此，有些人甚至根本不知道有圣经存在。为此，相信我，从我们这里生不出任何健全、有益的东西。\par

10. 然而，如果有人想学习军事，他必须研习军法。如果有人想学航海、木工或其他技艺，他必须学习那门技艺的原则。但在此事上，他们却不肯这样做，虽然这是一门需要高度警觉的科学。因为这也是一门需要教导的艺术，请听先知说：\BibleQuote{孩子们，前来听我，我必教你们敬畏上主。}\BibleRef{咏 33:11} 显然，敬畏上主需要教导。然后他说：\BibleQuote{谁愿得享长寿？}\BibleRef{咏 33:12} 他是指那永生；又说：\BibleQuote{谨守口舌，不说恶言，嘴唇不说欺诈的话；离恶行善，寻求和平，并追随不懈。}\BibleRef{咏 34:13}\par

你们知道是谁说了这些话吗？是一位先知、一位史家、一位宗徒，还是一位圣史？就我而言，我认为你们不知道，只有少数人知道。而且这些人自己，如果我们从别处引用一段经文，他们也会和你们其余的人一样。请看，我用另一种说法重复同样的陈述。\BibleQuote{你们要洗濯，自洁，从我眼前除掉你们的恶行，学习行善，寻求正义。谨守口舌，不说恶言，学习行善。}\BibleRef{依 1:16} 你们看见德行需要教导吗？因为一位说：\BibleQuote{我要教导你们敬畏上主}，另一位说：\BibleQuote{学习行善}。\par

那么，你们知道这些话出自哪里吗？就我而言，我认为你们不知道，只有少数人知道。然而，每个星期，这些经文都读给你们听两次乃至三次；当读经者走上读经台时，他首先说这是哪本书，是哪位先知的，然后才读出内容，为使你们更能理解，不仅知道书的内容，还知道写作的原因，以及谁说了这些话。但一切都是徒然，毫无益处。因为你们的热情都放在今生的事物上，对属神的事却毫不在意。因此，连那些俗务也不能如意，反而困难重重。因为基督说：\BibleQuote{你们先寻求天主的国，这一切自会加给你们。}\BibleRef{玛 6:33} 祂说这些是附加的；但我们却颠倒了次序，寻求尘世和地上的美物，好像那些天上的事物是加给我们的。因此，我们两者都没有得到。让我们终于醒悟，成为将来之事的渴望者；因为这样，这些也会随之而来。因为寻求天主之事的人，不可能不获得人的幸福。这是真理本身所说的。让我们不要反其道而行，而要坚持基督的训诲，免得我们失去一切。愿天主能赐予你们悔改的心，使你们变得更好，在我主耶稣基督内，祂与圣父及圣神，是光荣、权能、尊威，现今及永远，至于无穷之世。阿们。\par

\SourceRecord{Translated by Frederic Gardiner.；Nicene and Post-Nicene Fathers, First Series；Vol. 14.；Edited by Philip Schaff.；Buffalo, NY: Christian Literature Publishing Co.,；1889.；<http://www.newadvent.org/fathers/240208.htm>.}

% structure: p0001 source=h1 target=section reason=page-title

\section{论《希伯来书》第九篇讲道}

% structure: p0002 source=h2 target=subsection reason=relative-heading-rank; base=h2

\subsection{《希伯来书》六章1至3节}

\Quote{\Emph{因此，我们应离开基督初步的道理，} \Emph{努力向成全的境界前进，不必再打下悔改死行、信德于天主的根基；关于各种洗礼、覆手、死人复活和永远审判的道理。如果天主准许，我们就这样做。}}\par

1. 你们已经听见保禄多么责备希伯来人，因为他们总想学习同样的事。理由充分：\Quote{按时间说，你们本该作教师了，却还需要有人教导你们天主旨意的初步基础。}\BibleRef{希 5:12}\par

我担心这话也可恰当地对你们说：\Quote{按时间说，你们本该作教师了，}却仍停留在初学者的地步，总是听同样的事，在同样的问题上，你们与没听人讲论过并无分别。如果有人考问你们，除了极少可数的人外，无人能够回答。\par

但这并非小损失。因为教师常想前进，谈论更高深更奥秘的主题，但受教者的不专心阻碍了他。\par

正如文法教师，如果一个孩子虽然不断听讲初阶，但未掌握，教师就必须不断地向他重复同样的内容，直到他准确学会为止；因为若没有把基础打好就领他学别的东西，是大愚昧。同样，在教会里，如果我们常说同样的话而你们再无长进，我们也不会停止说同样的话。\par

因为如果我们的宣讲是为了显示和野心，那么从一个题目跳到另一个题目，不断变换，不顾你们，只求你们的掌声，才是合理的。但我们献身并非为此，我们的辛劳全为你们的益处，所以我们不会停止向你们讲论同样的主题，直到你们成功学会。因为我本可多说外邦人的迷信，论及\OriginalTerm{摩尼教徒}{Manichæans}和\OriginalTerm{马西昂派}{Marcionists}，靠着天主的恩宠给他们沉重打击，但这类言论不合时宜。因为对那些尚不准确知道自己事务的人，那些尚未学会贪婪是邪恶的人，谁会说那些话，在时机未到之前就带他们涉及其他主题呢？\par

因此，我们不会停止说同样的话，无论你们是否被说服。然而我们担心，因为不断说同样的话，如果你们不听，可能使悖逆者的罪罚更重。\par

但我必须不针对你们全体说这话；因为我知道许多人因来此而受益，他们可公正地谴责那些别人，因为他们的无知和不专心暗中损害别人。但即便如此，他们也不会受害。因为不断听同样的事，甚至对知道的人也有益处，因为通过多次听我们所知的，我们会更深受感动。例如，我们知道谦逊是极好的事，基督也多次谈论它；但当我们听到这些话语本身及对它们的反省时，我们更加感动，即使听上万次。\par

2. 现在也是我们合适向你们说这话的时候：\Quote{所以，我们应离开基督初步的道理，努力向成全的境界前进。}\par

什么是\Quote{初步的道理}？他接着说：\Quote{不必再打下}（他说的是）\Quote{悔改死行、信德于天主的根基；关于各种洗礼、覆手、死人复活和永远审判的道理。}\par

但如果这是\Quote{初步}，那么我们的道理除了悔改\Quote{死行}，并通过圣神领受\Quote{信德}，相信\Quote{死人复活和永远审判}之外，还有什么呢？但什么是\Quote{初步}？他说，\Quote{初步}不是别的，就是当没有严谨生活的时候。因为正如必须教导学文法的人首先学字母，同样基督徒也必须准确知道这些事，并毫不怀疑。如果他再需要教导，他就还没有基础。因为一个根基稳固的人应当固定而站立稳定，不被移动。但如果一个受过要理并受洗的人十年后又来听论及\Quote{信德}，说我们应当\Quote{相信}\Quote{死人复活}，他就还没有基础，他又在寻求基督宗教的初步。因为信德是基础，其余是建筑。请听他［宗徒］说：\Quote{我奠立了根基，别人在上面建筑。}\BibleRef{格前 3:10}\Quote{若有人在这根基上建造金、银、宝石、木、草、禾秸。}\BibleRef{格前 3:12}\par

\Quote{不必再打下}（他说）\Quote{悔改死行的根基。}\par

3. 但什么是\Quote{努力向成全的境界前进}？从此以后（他意思是）甚至上到屋顶，即让我们有至善的生活。因为在字母上，\OriginalTerm{阿尔法}{Alpha}包含了全部；在根基上，包含了整座建筑；同样，对信德的完全确信蕴含了生活的纯洁。没有这个，人不能成为基督徒，正如没有基础就没有建筑；没有字母就没有文理。但如果一个人总在字母周围打转，或总在根基而非建筑上，他将一无所获。\par

但不要以为信德被称为初步就是被贬低了：因为它确实是全部力量。因为当他说：\BibleQuote{凡只能吃奶的，在正义的话上还是婴孩} \BibleRef{希 5:13}，他称为\Quote{奶}的并非此意。但对这些事仍存疑惑，乃是心智软弱、需要许多讲论的[标记]。因为这些是健全的道理。我们称那有信德又生活正直的人为\Quote{成全的人}[即\Quote{成年}]；但若有人有信德，却行为邪恶，且对信德本身存疑，使道理蒙羞，我们理当称他为\Quote{婴孩}，因为他已退回起点。因此，即使我们在信德中已有万年之久，若在其中不坚定，仍是婴孩；当我们展示不符合信德的生活时；当我们仍在奠基时。\par

4. 但除了生活的方式，他还对这些人[希伯来人]提出另一控诉，说他们摇摆不定，需要\Quote{奠立悔改的根基，远离死亡的行为}。因为那从一转向另一、放弃这个而选择那个的人，理应先谴责这个，并从那体系分离，然后转向另一个。但如果他打算再抓住第一个，他如何能触及第二个呢？\par

那么，关于法律（他说）呢？我们谴责了它，却又跑回它那里。这不是摇摆不定，因为在这里[在福音之下]我们也有法律。\Quote{那么，我们}（他说）\BibleQuote{因信德就废了法律吗？断乎不是！我们更是坚固法律。} \BibleRef{罗 3:31} 我原是在谈论恶行。因为那追求德性的人，理应先谴责邪恶，然后追求它。因为悔改不能证明他们洁净。因此他们立刻领受圣洗，使得他们靠自己不能成就的，能藉基督的恩宠成就。因此，悔改不足以净化，人必须先领受圣洗。无论如何，必须来到圣洗前，如此谴责了罪恶并定了他们的罪。\par

但什么是\Quote{多种圣洗的道理}？并不是说有许多圣洗，而是只有一个。那么他为什么用复数表达？因为他曾说过：\Quote{不重新奠立悔改的根基。}因为如果他再次给他们领受圣洗、重新教理讲授，而他们在起初已领受圣洗、又再次被教导哪些事该做、哪些不该做，他们将永远无法纠正。\par

\Quote{以及按手。}因为这样他们领受了圣神，\BibleQuote{当\Person{保禄}按手在他们身上时} \BibleRef{宗 19:6}，经上这样说。\par

\Quote{以及死人的复活。}因为这既在圣洗中实现，也在宣认中被肯定。\par

\Quote{以及永罚。}但他为什么说这个？因为很可能他们已经相信了，却会动摇[失去信德]，或过邪恶怠惰的生活，他说：\Quote{要警醒。}\par

他们不能这样说：如果我们生活懈怠，我们将再次领受圣洗，我们将再次接受教理讲授，我们将再次领受圣神；即使现在我们失去信德，我们也能再次藉着圣洗洗去罪恶，达到与从前相同的境况。你们（他说）这样想是受骗了。\par

% structure: p0024 source=h2 target=subsection reason=relative-heading-rank; base=h2

\subsection{\BibleRef{希 6:4-5}}

5. \BibleQuote{因为那些曾一次蒙光照，尝过天上的恩赐，与圣神有分，并尝过天主美好的言语和来世的能力，而竟背离了的人，要使他们重新更新到悔改，是不可能的，因为他们自己又把天主子钉在十字架上，公然羞辱他。}\par

请看，他是如何羞辱他们并以警告的口吻开始的。\Quote{不可能。}他（说）不要再指望那不可能的事；（因为他不是说不适宜，或不合宜，或不合律法，而是说\Quote{不可能}，为使[他们]陷入绝望），如果你们曾完全蒙光照。\par

% structure: p0027 source=h2 target=subsection reason=relative-heading-rank; base=h2

\subsection{\BibleRef{希 6:6}}

然后他加上：\BibleQuote{并尝过天上的恩赐。}他（说）\BibleQuote{若你们尝过天上的恩赐}，即罪的赦免。\BibleQuote{并与圣神有分，尝过天主美好的言语}（他这里是在谈论道理）\BibleQuote{和来世的能力}（他说的是什么能力？或是行奇迹，或是\BibleQuote{圣神的凭质} \BibleRef{格后 1:22}）\BibleQuote{而竟背离了，要使他们重新更新到悔改，是不可能的，因为他们自己又把天主子钉在十字架上，公然羞辱他。}他（说）\Quote{更新他们} \Quote{到悔改}，即藉着悔改，因为\Quote{到悔改}就是藉着悔改。那么，悔改被排除了吗？不是悔改，绝非如此！而是藉着圣洗再次更新。因为他并没有只说\Quote{不可能}更新\Quote{到悔改}就停止，而是加上了如何\Quote{不可能，[因]重新钉十字架。}\par

\Quote{更新}，即成为新的，因为使人成为新的是唯独圣洗的工作：因为（经上说）\BibleQuote{你的青春如鹰得以更新。} \BibleRef{咏 103:5} 但使那些已成为新的人后来因罪变旧，从这旧态中得到释放、变得坚强，则是悔改的工作。然而，使他们恢复到先前的光辉却是不可能的，因为那完全是恩宠。\par

6. \Quote{亲自重钉于己}，他说，\Quote{把圣子重新钉在十字架上，并公开羞辱祂。}他的意思是这样：圣洗是一个十字架，\Quote{我们的旧人已与祂同钉十字架}\BibleRef{罗 6:6}，因为我们\Quote{被接合于祂死亡的相似形}\BibleRef{罗 6:5}，又说：\Quote{我们藉着圣洗与祂同葬于死}\BibleRef{罗 6:4}。因此，正如基督不可能第二次被钉十字架，因为那是\Quote{公开羞辱祂}。因为\Quote{死不再管辖祂}\BibleRef{罗 6:9}，若祂已复活，藉着复活超越了死亡；若祂藉着死亡战胜了死亡，然后又被钉，这一切便成了虚构和嘲弄。所以，第二次施洗的人，是将祂重钉。\par

但什么是\Quote{重钉}？就是重新钉死。因为正如基督死于十字架上，我们在圣洗中也是如此，不是肉身上的死，而是死于罪恶。看，有两种死亡。祂死于肉身；在我们，旧人被埋葬，新人兴起，被接合于祂死亡的相似形。因此，若必须再受洗，就必须这同一基督再死。因为圣洗无非是使受洗者被置于死地并使他复活。\par

他说的很对：\Quote{重钉于己}。因为这样做的人，好像忘记了先前的恩宠，并随意安排自己的生活，在各方面都好像有另一个圣洗。因此，我们应当留意，使自己安全。\par

7. 什么是\Quote{尝过天恩的滋味}？就是\Quote{罪过的赦免}：因为这本是唯独天主才能赐予的，恩宠是一次性的恩宠。\Quote{那么，我们该怎么说？我们可留在罪中，好叫恩宠丰盈吗？绝不可！}\BibleRef{罗 6:1}但若我们总是要藉恩宠得救，我们就永远不会善。因为只有一次恩宠，我们尚且如此怠惰，若我们知道罪过可能再被洗去，我们还会停止犯罪吗？依我看来，不会。\par

他在这里表明恩赐很多：为了解释，他说：你已蒙受如此大的赦免；因为那坐在黑暗中、那为仇敌、那公开敌对、那疏远、那为天主所憎恨、那丧亡的，他忽然被光照，蒙受了圣神、天恩、义子名分、天国，以及其他美物、不可言传的奥迹；而这人竟不因此而变得更好，反而虽当受丧亡，却获得了救恩和荣耀，好像成就了大事；他怎能再受洗呢？\par

于是，他基于两个理由说这事不可能，并把更强的放在后面：第一，因为那被认为配得如此祝福的人，却背叛了所赐予的一切，他不配再被更新；基督也不可能再被重钉：因为这是\Quote{公开羞辱祂}。\par

因此，没有第二次圣洗：确实没有。若有，就有第三次、第四次；因为前者不断被后者废去，后者又被另一个废去，如此无穷无尽。\par

他说：\Quote{尝过天主美善的话}；他没有展开说明；\Quote{和来世的能力}，因为像天使一样生活，无需属地之物，知道这是引领我们进入来世福乐的方法；这些我们可以藉着圣神学习，并进入那些神圣的深处。\par

什么是\Quote{来世的能力}？永生、天使般的交往。这些我们已经藉着信德从圣神领受了凭据。那么告诉我，若你被引入一座宫殿，并受托管理其中一切，然后你背叛了一切，还会再托付给你吗？\par

8. 那么（你说）？没有悔改吗？有悔改，但没有第二次圣洗：然而有悔改，且大有力量，能释放那受洗后犯许多罪者的罪担，只要他愿意，甚至使处于危险中的人得安全，即使他已陷入极深的邪恶。这从许多地方可以显明。因为有人说：\Quote{那跌倒的，岂不再起来吗？那转离的，岂不转回向天主吗？}\BibleRef{耶 8:4}若我们愿意，基督可以在我们内再次成形：因为听保禄说：\Quote{我的孩子们，我为你们再受生产之苦，直到基督在你们内成形}\BibleRef{迦 4:19}。我们只要抓住悔改。\par

请看天主对人的爱！我们在各方面都当受最初的惩罚；因为我们接受了自然律，享受了无数祝福，却不承认我们的主，并过着不洁的生活。然而祂不仅没有惩罚我们，反而使我们有分于无数的祝福，好像我们成就了大事。我们又堕落了，即便如此祂也没有惩罚我们，反而赐予悔改的良药，足以除去并抹去我们所有的罪；只要我们认识这药的性质，以及我们该如何应用它。\par

那么，悔改的良药是什么，又如何配制？第一，谴责自己的罪；因为（经上说）\Quote{我没有隐藏我的罪}\BibleRef{咏 31:5}；又说：\Quote{我要向主承认我的过犯，祢就赦免了我心中的罪。}以及\Quote{你首先要陈明你的罪，好使你成义}\BibleRef{依 43:26}。又说：\Quote{义人首先开口控告自己}\BibleRef{箴 18:17}。\par

第二，极大的谦卑：因为就如一条金链；若有人抓住开头，其余都会跟随。因为如果你按应有的方式承认你的罪，灵魂就谦卑了。因为良心转向自身，使它被制服。\par

此外，还必须向谦卑的心加上其他事物，如果这是真福达味所认识的那种谦卑，当他说：\Quote{天主决不轻视一颗破碎和痛悔的心。}\BibleRef{咏 50:17} 因为破碎的东西不会升起，不会攻击，而是准备好受苦，本身并不升起。痛悔之心也是如此：即使受侮辱，即使被恶待，它仍保持安静，不急于报复。\par

在谦卑之后，还需要恳切的祈祷、许多眼泪，白日的眼泪和夜晚的眼泪；因为他说：\Quote{每夜我要洗涤我的床铺，以眼泪浸湿我的褥子；我因呻吟而疲惫。}\BibleRef{咏 6:6} 又说：\Quote{我以灰尘当食，以眼泪调饮。}\BibleRef{咏 102:9}\par

在如此恳切的祈祷之后，还需要大量的施舍；因为尤其这一点给悔改的药物增添力量。正如在医生的帮助中，有一种药物包含多种草药，但其中一种是必不可少的；同样，在悔改中，这种草药也是必不可少的，甚至就是一切。因为请听圣经的话：\Quote{施舍吧，一切便都洁净了。}\BibleRef{路 11:41} 又说：\Quote{藉施舍和信实，罪恶得以消除。}\BibleRef{箴 16:6} 以及：\Quote{水能熄灭燃烧的火，施舍能消除大罪。}\BibleRef{德 3:30}\par

其次，不对任何人发怒，不怀恨在心；宽恕所有人的过犯。因为经上说：\Quote{人对人怀怒，却向主寻求医治。}\BibleRef{德 28:3} \Quote{宽恕，你们便得宽恕。}\BibleRef{谷 11:25}\par

还有，使我们的弟兄从迷途中悔改。因为经上说：\Quote{去，使你们的弟兄悔改，好使你们的罪得赦免。} 以及，与司铎保持密切关系，因为经上说：\Quote{若……犯了罪，也必得赦免。}\BibleRef{雅 5:15} 挺身维护受冤屈的人。不怀怒气。以温和承受一切。\par

9. 如今，在你们得知藉悔改洗去罪恶是可能的之前，你们岂不因为再没有第二次洗礼而痛苦，且对自身绝望吗？但现在我们既然得知悔改和赦免如何得以成功实现，并且若我们愿意善用悔改，便能全然逃脱，那么，如果我们甚至不思考自己的罪，还能得到什么宽恕呢？因为如果这被做到了，一切都将成就。\par

因为正如进入门内的人就在里面；同样，清点自己罪恶的人也必然会得到医治。但如果他说：我是罪人，却不具体清点它们，不说：我犯了这罪和那罪；他将永远不会停止，确实不断告明，却从不真心关心改正。因为一旦他有了一个开始，其余的一切无疑都会随之而来，只要他在某一点上有了开始；因为在一切事情上，开始和初步都是困难的。让我们以此为根基，一切便都会顺利容易。\par

因此，我恳求你们，让我们开始：一人以恳切的祈祷，另一人以不断的哭泣，再一人以垂头丧气。因为即使如此微小的事也不是无益的：因为\Quote{我看见（经上说）他忧伤垂头，我便医治了他的路途。}\BibleRef{依 57:17}\par

但让我们所有人都藉着施舍和宽恕邻人的过犯，不记仇，不报复，来谦卑自己的灵魂。如果我们不断反思自己的罪恶，任何外在环境都不能使我们得意：无论是财富、权力、权威还是荣誉；甚至即使我们坐在皇帝的车上，我们也会深深叹息。因为就连真福达味也是一位君王，然而他说：\Quote{每夜我要洗涤我的床铺，}\BibleRef{咏 6:6} 他丝毫没有因紫袍和王冠而受伤害；他没有自高自大；因为他知道自己是一个人，而且因他的心已破碎，他过着哀伤的生活。\par

10. 因为一切人间之物是什么？灰烬和尘土，有如风前的飞沫；烟云和阴影，随风飘荡的叶子；花朵、梦、故事、寓言、虚妄膨胀而消散的风和空气；没有根基的羽毛、流动的溪流，或者还有比这些更虚空的东西。\par

因为告诉我，你以什么为大？你认为什么尊位为大？是执政官吗？因为许多人认为没有比这更大的尊位了。不是执政官的人丝毫不亚于那个如此显赫、如此受人钦佩的人。两者地位相同；两者不久之后都不复存在。\par

他在何时成为〔执政官〕？多久？告诉我：两天？但这种事也会在梦中发生。但你说：那是〔只是〕梦。而这是什么？因为（告诉我）白天的事物，就〔因此〕不是梦吗？为何我们不把这些事物称为梦？因为正如梦在白天来临时被证明是虚无，同样这些事物在夜晚来临时也被证明是虚无。因为白昼和黑夜各自获得了相等的时间，平分了所有时间。因此，正如一个人白天不会因夜里发生的事而欢乐，同样，在夜里他也不可能收获白天所做之事的果实。你做了执政官？我在夜里也做了；只是我在夜里，你在白天。这又怎样？即使如此，你也没有比我有什么优势，除了或许有人说：某某是执政官，以及从言语中产生的快乐给你优势。\par

我的意思是某种类似的东西，因为我要更清楚地表达：如果我说\Quote{某某是执政官}，把这个名字给他，它是否一说出就消失了？事物本身也是如此：执政官一出现，就不再存在。但让我们假设〔他执政〕一年、两年、三年或四年。那些十次担任执政官的人在哪里？无处可寻。\par

但\Person{保禄}并非如此。因为过去和现在他都持续活着；他不是只活了一天、两天、十天、二十天、三十天，也不是十年、二十年或三十年——然后死去。甚至四百年已经过去，他至今仍然显赫，甚至比他活着时更加显赫。这些事确实在地上；但天上圣人的光荣，有什么言语能表达呢？\par

因此我恳求你们，让我们寻求这光荣；让我们追求它，好能获得它。因为这才是真光荣。让我们从此远离今世的事物，好能在我们的主\Person{耶稣基督}内获得恩宠和仁慈：愿光荣、权能、尊威和崇拜，借着祂（我们的主\Person{耶稣基督}），归于圣父并与圣神一起，从现在直到永远，永无穷尽。阿们。\par

\SourceRecord{Translated by Frederic Gardiner.；Nicene and Post-Nicene Fathers, First Series；Vol. 14.；Edited by Philip Schaff.；Buffalo, NY: Christian Literature Publishing Co.,；1889.；<http://www.newadvent.org/fathers/240209.htm>.}

% structure: p0001 source=h1 target=section reason=page-title

\section{论希伯来人书 第十篇讲道}

% structure: p0002 source=h2 target=subsection reason=relative-heading-rank; base=h2

\subsection{\BibleRef{希 6:7-8}}

\BibleQuote{\Emph{因为土地吮吸时常降在它上面的雨水，为耕种它的人生出合宜的菜蔬，就从天主领受祝福。但它若长出荆棘和蒺藜，就被废弃，几近咒骂，它的结局就是焚烧。}}\par

1. 让我们以敬畏、以敬畏和极大的战兢来聆听天主的圣言。因为（经上说）：\BibleQuote{应当以敬畏事奉主，且以战兢向祂欢欣。} \BibleRef{咏2:11} 但若连我们的喜乐和我们的欢跃都应当\BibleQuote{带着战兢}，那么，当我们不以惊恐来聆听所说的话，而所说的话，像现在这样，本身就是可畏的，我们该当什么惩罚呢？\par

因为他说过\Quote{为那些背弃了的人，是不可能}再受第二次洗礼，并藉着洗礼池获得赦免，又指出了情况的可怕之后，他接着说：\BibleQuote{因为土地吮吸时常降在它上面的雨水，为耕种它的人生出合宜的菜蔬，就从天主领受祝福。但它若长出荆棘和蒺藜，就被废弃，几近咒骂，它的结局就是焚烧。}\par

那么，让我们恐惧吧，亲爱的！这警告不是保禄的，这些话不是人的；而是圣神的，是基督在他内说话的。那么，难道有谁没有这些荆棘吗？即使我们清除了，也不该因此自信，而要恐惧战兢，免得荆棘随时在我们中间长出来。但当我们从头到脚都是\BibleQuote{荆棘和蒺藜}时，（告诉我）从哪里我们自信呢？而且正在变得麻木不仁？是什么使我们惰怠？如果\BibleQuote{那自以为站得稳的}该恐惧\BibleQuote{免得跌倒}；因为（他说）\BibleQuote{那自以为站得稳的，要小心免得跌倒} \BibleRef{格前10:12}；那么跌倒的人，该多么急切要重新起来！如果保禄恐惧\BibleQuote{免得我传福音给别人，自己反被弃绝} \BibleRef{格前9:27}；而这曾如此被认可的人害怕自己成为被弃绝的：那么我们已经是被弃绝的，若我们毫无恐惧，只当作习惯和形式来履行我们的基督徒生活，我们将得到什么宽恕呢？因此，让我们恐惧吧，亲爱的！\BibleQuote{因为天主的愤怒从天上显明出来} \BibleRef{罗1:18}。让我们恐惧，因为它不仅\BibleQuote{显明}\BibleQuote{对不虔敬}，而是\BibleQuote{对一切不义}。什么是\Quote{对一切不义}？就是对一切，无论大小。\par

2. 在此他暗示天主对人的慈爱；他把（福音的）教导称为\Quote{雨水}；而他上面所说的\BibleQuote{按时期来说，你们本应做导师} \BibleRef{希5:12}，在这里他也说了。的确，圣经在许多地方将教导称为\Quote{雨水}。因为（经上说）\BibleQuote{我必命云不再降雨在它上面} \BibleRef{依5:6}，是说到\Quote{葡萄园}。这同一件事在另一处称为\BibleQuote{饥荒，无水之渴} \BibleRef{亚8:11}。又说：\BibleQuote{天主的河满了水} \BibleRef{咏64:9}。\par

他说：\BibleQuote{因为土地吮吸时常降在它上面的雨水。}在此他表明他们领受了并饮用了圣言，是的，且多次享受了这雨水，但即使如此，他们仍无得益。因为（他意指）若你们未曾被耕耘，若你们没有享受过雨水，邪恶就不会这么大。因为（经上说）\BibleQuote{我若没有来向他们说话，他们就没有罪} \BibleRef{若15:22}。但若你们常常饮用并得（滋养），为何你们结出别的而不是果实呢？因为（经上说）\BibleQuote{我指望它结葡萄，它却结了荆棘} \BibleRef{依5:2}。\par

你们看，处处圣经将罪称为\Quote{荆棘}。因为达味也说：\BibleQuote{我转为哀悼，当一根荆棘扎入我中} \BibleRef{咏31:4}。对于荆棘不是简单地临到我们，而是扎入；即使只有一点点留在内，即使我们不完全拔出来，那一点点本身就照样引起疼痛，如同荆棘的情形。我为什么说\Quote{那一点点本身}？即使在拔出来之后，它仍在里面留下伤口的疼痛很长时间。需要大量的照顾和治疗，我们才能完全摆脱它。因为仅仅除去罪是不够的，还必须医治受伤之处。\par

但我担心所说的话适用于我们更甚于别人。因为他说：\BibleQuote{土地吮吸时常降在它上面的雨水。}我们常常饮用，常常聆听，但\BibleQuote{太阳一升起} \BibleRef{玛13:6}，我们立刻就失去了水分，因此结出荆棘来。那么荆棘是什么呢？让我们听基督说：\BibleQuote{世俗的焦虑和财富的欺骗窒息了圣言，它就结不出果实} \BibleRef{玛13:22}。\par

3. \BibleQuote{因为土地吮吸时常降在它上面的雨水，并生出合宜的菜蔬。}因为没有什么比生活的纯洁更合宜，没有什么比至善的生活更适宜，没有什么比德行更合宜。\par

\BibleQuote{并生出为耕种它的人合宜的菜蔬，就从天主领受祝福。}在此他说天主是万物的原因，给了外邦人一击，他们将果实的出产归功于土地的能力。因为（他说）不是农夫的手促使土地结果实，而是来自天主的命令。因此他说：\Quote{从天主领受祝福。}\par

并且请看，在说到荆棘时，他没有说\Quote{结出荆棘}，也没有使用这个表达有用的词；而是说什么？\Quote{长出}（字面意思是\Quote{生出}）\Quote{荆棘}，如同有人说：\Quote{逼迫而出}，\Quote{抛掷而出}。\par

\Quote{被弃}（他说）\Quote{近于诅咒}。哦！这句话中的安慰何其大！因为他说的是\Quote{近于诅咒}，而不是\Quote{诅咒}。那人尚未陷入诅咒，却已临近[它]，也能远离[它]。\par

他不仅以此激励他们，也以下文鼓励他们。因为他没有说\Quote{被弃，近于诅咒}，\Quote{必被焚烧}，而是说\Quote{它的结局就是焚烧}，如果那人一直如此（他意指）直到终了。因此，如果我们割除并焚烧荆棘，我们就能享有那无数美善，成为蒙悦纳的，并分享祝福。\par

他恰当地称罪为\Quote{蒺藜}，说\Quote{那长出荆棘和蒺藜的}；因为无论你从哪一面抓住它，它都会刺伤你，甚至看上去也不悦目。\par

4. 在充分责备、警告和刺伤他们之后，他转而医治他们，免得过分打击他们，使他们萎靡不振。因为打击一个\Quote{迟钝}的人，会使他更加迟钝。因此，他既不完全奉承他们，免得他们自满，也不彻底伤害他们，而是插入少许使他们受伤的话，然后用大量医治的话语来恢复。\par

因为他怎么说？我们说这些，并非定你们的罪，也不是认为你们长满了荆棘，而是恐怕这事发生。因为用言语恐吓你们，免得你们在现实中受苦，这样更好。这尤其体现了保禄的智慧。\par

% structure: p0019 source=h2 target=subsection reason=relative-heading-rank; base=h2

\subsection{\BibleBook{希伯来}{书} 6:9}

而且，他没有说\Quote{我们认为}或\Quote{我们猜想}或\Quote{我们期望}或\Quote{我们希望}，而是说什么？\BibleRef{希 6:9} \BibleQuote{可是，亲爱的，我们虽这样说，却深信你们身上有更好的事，以及关乎救恩的事。}这句话他也用在写给迦拉达人的书信中：\BibleQuote{我在主内深信你们不会另有心思。}\BibleRef{迦 5:10} 因为在那个例子中，因为他们很该受责备，他不能从现状称赞他们，于是从未来称赞（他说\Quote{你们不会另有心思}）；他没有说\Quote{你们是}，而是说\Quote{你们不会另有心思}。但在这里，他从现状鼓励他们：\BibleQuote{我们曾深信你们有更好的事，就是关乎救恩的事，尽管我们这样说。}既然他不能从现状说太多，他便从过去的事来坚固他的安慰，说：\par

% structure: p0021 source=h2 target=subsection reason=relative-heading-rank; base=h2

\subsection{\BibleBook{希伯来}{书} 6:10}

\BibleQuote{因为天主并非不公义，以致忘记你们的工作和你们为祂的名所显示的爱德，就是你们过去和现在对圣徒的服事。}哦，他多么恢复了他们的精神，使他们重新振作，借着提醒他们从前的事，促使他们不能以为天主已经忘记了。（因为人若不完全确信他的希望，说天主不公义，他就不能不犯罪。因此他迫使他们无论如何要期待那将来的事。因为一个对现在绝望、放弃努力的人，可能因[将来的事]而恢复。）正如他写信给迦拉达人时所说：\BibleQuote{你们跑得好}\BibleRef{迦 5:7}；又说：\BibleQuote{你们受了这么多苦，难道都是白费的吗？如果真是白费的话。}\BibleRef{迦 3:4}\par

并且，正如在这里他将赞美与责备放在一起，说：\BibleQuote{按时间说，你们本该做师傅}\BibleRef{希 5:12}，同样在那里他也说：\BibleQuote{我惊奇你们这么快就离弃了。}\BibleRef{迦 1:6} 责备中包含着赞美。因为对于伟大的事，当它们失败时，我们感到惊奇。你看，赞美隐藏在指责和责备之下。他不仅对自己说这话，也是针对众人。因为他没有说\Quote{我深信}，而是说\Quote{我们深信你们有更好的事}，即好事（他意指）。他这样说或是指行为方面，或是指报偿方面。接着，在上面说了\Quote{被弃，近于诅咒}，且\Quote{必被焚烧}之后，他说：我们绝不是说你们。\BibleQuote{因为天主并非不公义，以致忘记你们的工作和爱德。}\BibleRef{希 6:10}\par

% structure: p0024 source=h2 target=subsection reason=relative-heading-rank; base=h2

\subsection{\BibleBook{希伯来}{书} 6:11-12}

5. 那么，我们为什么说了这些呢？\BibleQuote{但是，我们切愿你们每人显出同样的勤勉，以达到望德的圆满，直到终了；使你们不怠惰，而是效法那些因信德和忍耐而继承恩许的人。}\par

\Quote{我们切愿，}他说，而不是因此仅仅劳苦，或仅在言词上希望。而是什么？\Quote{我们切愿}你们持守德行，不是谴责你们过去的行为（他意指），而是为将来担心。他没有说\Quote{不是谴责你们过去的行为，而是现在的；因为你们已经灰心，你们变得太懈怠了}，但请看他是如何温柔地指出，而没有使他们受伤。\par

因为他说了什么？\BibleQuote{但是，我们切愿你们每人都显出同样的勤勉，直到终了。}这正是保禄智慧的杰出之处：他没有明确表示他们\Quote{已经}放弃，已经变得疏忽。因为当他说\Quote{我们切愿你们每人}——就像一个人说：我希望你们始终热忱；如同你们以前那样，现在和将来也是如此。因为这使他的责备更温和，更容易接受。\par

他没有说\Quote{我愿意}，这是表达教师的权威，而是说\Quote{我们切愿}，这是表达父亲的亲情，且比\Quote{愿意}更强烈。无异于说：请原谅我们，即使我们说了不中听的话。\par

\BibleQuote{我们切愿你们每人显出同样的勤勉，以达到你们望德的圆满，直到终了。}望德（他意指）带领我们前行，它再次恢复我们。不要疲倦，不要绝望，免得你们的望德落空。因为行善的人也期盼善事，绝不绝望于自己。\par

\Quote{使你们不致懈怠。} 仍是\Quote{成为}；但他在上面说：\BibleQuote{你们已成听觉迟钝的。}\BibleRef{希 5:11} 请注意他却如何将迟钝仅限于听觉。这里他也暗示同样的事；不说\Quote{使你们不坚持其中}，而说这话。但他又引领他们走向那未来还未负责的时候；其意是说\Quote{使你们不致过于懈怠}：因为那尚未到来的，我们无法负责。因为那在当前受劝勉要认真的，既是松懈，或许会变得更懈怠；但那以未来为劝勉的，则不然。\par

\Quote{我们愿}（他说）\Quote{你们各人。} 他对他们的爱是大的：对大小一视同仁；而且他知道一切，不忽略任何人，对每人表现同样的慈爱，对众人同等珍视；由此他也更能说服他们接受他言语中令人不快之处。\par

\Quote{使你们不懒惰}，他说。因为懒惰伤害身体，同样对善事的懒惰使灵魂更迟钝软弱。\par

6. \Quote{却效法}（他说）\Quote{那些凭信德和耐心继承恩许的人。} 他们是谁，他以后说明。他先前说：\Quote{效法你们从前的善行。} 然后，免得他们说：什么？他引导他们回到先祖：从他们自己的历史中带来善行的榜样，但从先祖带来被遗弃之感的榜样；使他们不以为自己被轻视、被遗弃、毫无价值，而知道最尊贵的人是在试炼中行走人生旅程；天主对伟大可钦的人这样行事。\par

现在我们当（他说）以耐心承受一切：因为这也是相信；如果他给予而你立刻领受，你怎算相信了呢？因为这样就不再是你的信德，而是我、给予者的了。但如果我许诺给，经过一百年才给，而你不失望；那么你便认为我值得相信，你对我有正确的看法。你看，不信常不仅来自缺乏希望，也来自气馁与缺乏耐心，而非来自对许诺者的谴责。\par

\Quote{因为天主}（他说）\Quote{并非不义，以致忘记你们的爱德}以及\Quote{你们为祂的圣名所表现的殷勤，就是你们服事了圣徒，并且仍在服事。} 他为她们作证大事，不仅是行为；而且是欣然行之的行为，正如他在另一处所说：\BibleQuote{并且不但如此，他们也把自己献给主，也献给我们。}\BibleRef{格后 8:5}\par

\Quote{这}（他说）\Quote{你们为祂的圣名所表现的，就是服事了圣徒，并且仍在服事。} 瞧他又如何安抚他们，加上\Quote{并且仍在服事。} 他说：即使现在，你们仍在服事，并且他因表明他们所作的不是给圣徒，而是给天主而振作他们。\Quote{你们所表现的}（他说）；他没说\Quote{给圣徒}，而是\Quote{向着天主}，因为这是\Quote{为祂的圣名。} 意思是：你们为祂的名而行一切。因此，那位从你们得享如此大热心和爱的主，绝不会轻视你们或忘记你们。\par

7. 听见这些，我恳求你们，\Quote{服事圣徒。} 因为每个信徒因其是信徒就是圣人。即使他是在俗的人，也是圣人。\Quote{因为}（他说）\BibleQuote{不信的丈夫因妻子而成为圣的，不信的妻子因丈夫而成为圣的。}\BibleRef{格前 7:14} 你看信德如何造成圣德。如果我们看见甚至世俗人在患难中，让我们伸出手。不要只热衷于那些住在山上的人；他们固然在生活方式和信德上是圣人；但这些人因信德是圣人，其中许多人也因生活是圣人。如果我们看见隐修士被下监，便进去；但若是在俗的人，便拒绝进去，不要这样。他也是圣人和弟兄。\par

那么（你说）如果他不洁污秽呢？听基督说：\BibleQuote{你们不要判断，免得受判断。}\BibleRef{玛 7:1} 你为天主而行。不，我说什么？即使看见外教人在患难中，我们也当向他显示仁慈，并向一切患难中的人，更向那在信德中的人。听保禄说：\BibleQuote{向众人行善，尤其向信德之家的人。}\BibleRef{迦 6:10}\par

但我不知这（观念）从何传入，或这习俗如何盛行。因为那只追求独修者、只愿向他们行善、而对别人反而好奇追究，并说：\Quote{除非他值得，除非他正义，除非他行奇迹，我不伸手}；这样的人已夺去了爱德的大部分，甚至日久会摧毁这事本身。然而爱德正是向罪人、向有罪者显示的。因为这是爱德：不是怜悯行善的，而是怜悯行恶的。\par

8. 为使你明白这事，请听这个比喻：\BibleQuote{有一个人}（经上记载）\BibleQuote{从\Place{耶路撒冷}下到\Place{耶里哥}，落在强盗手中} \BibleRef{路 10:30}；他们打了他，把他丢在路旁，把他打得遍体鳞伤。有一个肋未人经过，看见他，就绕过去了；有一个司祭经过，看见他，也匆匆过去了；有一个撒玛黎雅人经过，对他施以极大的照顾。因为他\BibleQuote{包扎了他的伤} \BibleRef{路 10:34}，倒上油，扶他骑上自己的驴，\BibleQuote{带他到客店，对店主说：照顾他} \BibleRef{路 10:35}；而且（请注意他的慷慨大方），\BibleQuote{我，}他说，\BibleQuote{凡你所花费的，都要还给你。} \BibleRef{路 10:35} 那么谁是他的近人呢？\BibleQuote{那怜悯他的。你去，}祂说，\BibleQuote{也照样做吧。} \BibleRef{路 10:37} 请看祂讲的比喻多么精彩。祂没有说一个犹太人对撒玛黎雅人做了某事，而是说一个撒玛黎雅人显示了那样的大方。既然听了这些，我们不要只关心\BibleQuote{信仰家里的人} \BibleRef{迦 6:10}，而忽略了别人。所以你也一样，如果你看见任何人受苦，不要好奇地进一步查问。他的苦难本身就对你的援助提出了合理的要求。因为如果你看见一头驴窒息，你把它扶起来，并不好奇地问这是谁的驴，更何况对一个人，更不应过分好奇地问他是谁。他是天主的人，无论是外邦人还是犹太人；因为即使他不信，他仍然需要帮助。因为如果真的让你去查问和判断，你这样说也许有道理，但事实上，他的不幸不容许你追究这些。因为即使对健康的人，也不应过分好奇，不要管别人的闲事，何况对那些受苦的人。\par

9. 从另一方面看，我们该说什么呢？你看见他兴旺、受尊重，你就说他邪恶、无价值吗？但如果你看见他受苦，不要说他是邪恶的。因为当一个人很有声望时，我们这样说还算公平；但当他遭遇灾难、需要帮助时，说他是邪恶的就不对了。因为这是残忍、不人道和傲慢。告诉我，有什么比犹太人更不义的？然而，虽然天主惩罚了他们，而且是公义的，是的，非常公义，但祂仍赞许那些怜悯他们的人，而惩罚那些对他们幸灾乐祸的人。\BibleRef{亚 6:6} 因为经上记载：\BibleQuote{他们不为\Person{若瑟}的苦难而悲伤。}\par

此外，经上又说：\BibleQuote{赎回（赎价）那将被杀的人，不可吝惜。} \BibleRef{箴 24:11} （祂没有说，好奇地查问，并了解他是谁；然而，大多数被带去处决的人都是邪恶的，）因为这就是爱德。因为对朋友行善，并不完全是为了天主；但对陌生人行善，这个人才是纯粹为了天主而行善。\BibleQuote{不可吝惜}你的钱财，即使必须花光所有，也要给予。\par

但是，当我们看见人处于极度痛苦中，哀哭自己，受比死一万次还更痛苦的事，而且常常是不公正的，我们却吝惜钱财，而不顾惜我们的弟兄；我们关心无生命的东西，却忽视活的灵魂。然而\Person{保禄}说：\BibleQuote{以温和开导那些反对的人，或许天主会赐给他们悔改，得以认识真理，使他们从魔鬼的网罗中解脱出来，因为他们已被魔鬼俘虏，去随从他的旨意。} \BibleRef{弟后 2:25} \BibleQuote{或许，}他说；你看到这话充满了多么大的忍耐。\par

我们也该效法祂，不要对任何人失望。因为渔夫们，当他们多次撒网（设想一下）没有成功时，但后来再撒一次，就得了全部。同样，我们也期望你们立刻向我们显示成熟的果实。因为农夫，在他播种之后，等一天或两天，长久地期待；然后一下子他看到果实从四面八方涌现。我们也期望在你们身上也发生这样的事，借着我们的主\Person{耶稣基督}的恩宠和慈爱，愿荣耀、权能、尊荣归于祂，以及圣父和圣神，从今世直到永远，世世无穷。阿们。\par

\SourceRecord{Translated by Frederic Gardiner.；Nicene and Post-Nicene Fathers, First Series；Vol. 14.；Edited by Philip Schaff.；Buffalo, NY: Christian Literature Publishing Co.,；1889.；<http://www.newadvent.org/fathers/240210.htm>.}

% structure: p0001 source=h1 target=section reason=page-title

\section{希伯来书第十一篇讲道}

% structure: p0002 source=h2 target=subsection reason=relative-heading-rank; base=h2

\subsection{希 6:13-16}

\Quote{天主既然向亚巴郎应许，因为没有比自己更大的可以指着发誓，就指着自己发誓说：\InnerQuote{我必多多祝福你，使你昌盛繁茂。} 这样，他因坚忍而获得了应许。人确实指着比自己大的发誓，并且发誓作保证，是了结一切争论的。}\par

1. 在勇敢地反思了希伯来人的缺点并充分警告他们之后，他先以赞美安慰他们，其次——也是更有力的理由——借着[他们必定会达成希望的对象]的想法安慰他们。此外，他并非从未来的事，而是再次从过去的事汲取安慰，这反而更能说服他们。因为在惩罚的事上，他更多地以那些（即未来的事）警告他们，同样地，在奖赏方面，他也以这些（即过去的事）鼓励他们，就此显明天主的行事方式。那就是，不是立即实现对所应许的，而是在漫长的时间之后。祂这样做，既是展示祂权能的最大证据，也是为了引导我们进入信德，使那些在患难中生活却未领受应许或赏报的人，不至于在困苦中灰心。\par

他略过所有[其他]——尽管他本可提及许多人——而举出亚巴郎，既因其人格的尊贵，也因这事在他身上以特殊方式发生。\par

然而，在书信的[结尾]他说：\BibleQuote{这一切人，虽看见应许从远处而来，并拥抱了它们，却没有领受，为使我们若没有他们，就不能达到完美。}\BibleRef{希 11:13} \BibleQuote{天主既然向亚巴郎应许}（他说）\BibleQuote{因为没有比自己更大的可以指着发誓，就指着自己发誓说：\InnerQuote{我必多多祝福你，使你昌盛繁茂。} 这样，他因坚忍而获得了应许。}\EditorialNote{见：希 11:39-40} 那么，他如何在[结尾]说\Quote{他没有领受应许}，而在这里说\Quote{他因坚忍而获得了应许}呢？他是怎样没有领受？又是怎样获得了呢？他在这里与那里所说的不是同一件事，而是使安慰具有双重性。天主向亚巴郎作了应许，并在漫长的时间之后，赐下了此处所提及的事物，但其他的那些事物尚未赐下。\par

\Quote{这样，他因坚忍而获得了应许。}你看出，单是应许并不成就一切，还需要坚忍的等待吗？在这里，他警告他们，表明应许时常因灰心而落空。他确实曾借着[犹太人]百姓的[例子]表明这点：因为他们灰心，所以没有获得应许。但现在他借着亚巴郎的例子表明相反的情况。此后，在书信的[末尾]，他还证明了更多：[即]即使他们坚忍了，也没有获得；然而即便如此，他们也不忧伤。\par

2. \Quote{人确实指着比自己大的发誓，并且发誓作保证，是了结一切争论的。但天主因为没有比自己更大的可以指着发誓，就指着自己发誓。}那么，是谁向亚巴郎发誓呢？不是圣子吗？有人说，不。的确是他；不过，我不争论此事。所以，当圣子发誓时说\Quote{我实实在在告诉你们}，难道不明显是因为他不能指着更大的发誓吗？因为正如圣父指着自己发誓，同样圣子也指着自己发誓说：\Quote{我实实在在告诉你们。}在这里，他也提醒他们基督常发的誓言。\BibleQuote{我实实在在告诉你们：信从我的人，永远不死。}\BibleRef{若 11:26}\par

什么是\Quote{并且发誓作保证，是了结一切争论的}？即是，\Quote{借此每一个可疑问题都得到解决}：不是这个或那个，而是每一个。\par

% structure: p0010 source=h2 target=subsection reason=relative-heading-rank; base=h2

\subsection{希 6:17}

然而，天主本来即使没有誓言也应受到相信：\BibleRef{希 6:17} \Quote{故此}（他说）天主愿意更丰盛地向应许的承继者表明祂旨意的不变性，就借着誓言加以保证（\OriginalTerm{中保}{mediated}）。在这些话中，他也包括了信徒，因此提到这个\Quote{应许}，就是与我们共同（与他们）所立的。\Quote{祂借着誓言作了中保}，他说。这里他又说圣子是人与天主之间的中保。\par

% structure: p0012 source=h2 target=subsection reason=relative-heading-rank; base=h2

\subsection{希 6:18}

\BibleQuote{好叫我们借着这两件不可更改的事——天主绝不能在这两件事上说谎——} 这两件是什么？就是说话和应许，以及在应许上加添誓言。因为在人当中，借着誓言所确认的事被认为更可靠，为此祂也加上了誓言。\par

你看出祂不看重自己的尊严，而是如何说服人，并忍受有关自己的不体面之言吗？也就是说，祂愿意赐予充分的保证。在亚巴郎的事上，[宗徒]指出一切都出于天主，并非出于他的坚忍，因为祂甚至愿意加上誓言，因为人指着祂发誓，天主也\Quote{指着自己发誓}，即\Quote{指着自己}。他们确实是指着比自己大的，但祂不是指着比自己大的。然而祂仍然做了。因为人指着自己发誓与天主指着自己发誓并不相同，因为人没有权柄掌管自己。那么，你看出这话不只对亚巴郎，也是对我们说的：\Quote{使我们}（他说）\Quote{这些逃往避难所、紧握摆在面前的希望的人，能得到强有力的安慰。} 这里再次提到，\Quote{他因坚忍而获得了应许。}\par

\Quote{如今}他解释，他并没有说\Quote{当他发誓的时候}。但他借着谈及指着更大的发誓，表明了誓言是什么。但因人类生性难以相信，祂屈尊俯就，与我们[同列]。正如祂为了我们发誓，尽管祂不被相信是有损祂的尊严，同样[宗徒]也说了另一句话：\BibleQuote{祂从所受的苦难学会了服从}\BibleRef{希 5:8}，因为人认为经过经验更可靠。\par

\Quote{摆在我们面前的希望}是什么？他说：从这些过去的事我们推测未来。若这些事在如此长久之后实现了，那么其他事也必会实现。如此，关于亚巴郎所发生的事，也给予我们关于未来之事的信心。\par

% structure: p0017 source=h2 target=subsection reason=relative-heading-rank; base=h2

\subsection{希伯来书 6:19-20}

3.\newline{}\Quote{我们把这希望当作灵魂的锚，既稳妥又坚固，且进入幔内的那地方；耶稣已作前驱为我们先进去了，照默基瑟德的品位永为大司祭。}他表明，当我们还在世上、尚未离开此世时，我们已经处于应许之中。因为借着希望，我们已在天堂。他说：\Quote{等待吧，因为必定实现。}随后给他们充分的保证，他说：\Quote{不，是借着希望。}他没有说\Quote{我们在里面}，而是说\Quote{它已进入里面}，这更真实、更有说服力。因为如同锚从船只抛下，即使万般风浪摇动，也不让船只飘荡，反而因着依赖使之稳固，希望也是如此。\par

请看，他发现了多么恰当的比喻：他没有说\Quote{基础}，这并不适合，而是说\Quote{锚}。因为那在汹涌的海上、似乎不很稳固的东西，却立在水上如同在陆地，虽被摇动却不动摇。关于那些非常坚定、有智慧的人，基督很有理由说了那句话：\BibleQuote{凡把房子建在磐石上的}。\BibleRef{玛 7:24}但对于那些动摇、需要靠希望带领的人，保禄贴切地设定了这个比喻。因为巨浪和大风暴使船摇摆，但希望不让它到处漂流，尽管无数风浪摇动它。所以，若无此希望，我们早就沉没了。不仅在属灵之事上，就是在现世事务中，也能发现希望的力量何等巨大。无论经商、农耕、从军，若人不把希望摆在面前，甚至不会开始工作。但他没有简单地说\Quote{锚}，而是说\Quote{稳妥坚固}，即不动摇的。\Quote{进入幔内的那地方}，意即\Quote{它直达天堂}。\par

4.\newline{}此后，他进而引入信德，为叫希望不仅存在，而且非常真实。因为起誓之后，他又加上另一件事，即事实的证明，因为\Quote{耶稣已作前驱为我们进去了。}但前驱是为某人作前驱，正如若翰是基督的前驱。\par

他没有简单地说\Quote{他进去了}，而是说\Quote{他为我们作前驱进去了}，好像我们也应当达到一样。因为前驱与跟随者之间没有很大间隔；否则就不是前驱了；因为前驱与跟随者应当走同一条路，并先后到达。\par

\Quote{照默基瑟德的品位永为大司祭}，他说。这里还有另一个安慰：如果我们的大司祭在高天，远比犹太人中的大司祭更卓越，不只在于司祭的品类，也在于地方、帐幕、盟约和人物。这是按血肉所说的。\par

5.\newline{}那么，祂作大司祭的人，应当远远超越。亚郎与基督之间有多大的差别，我们与犹太人也当有同样大的差别。请看，我们有我们的牺牲在高天、司祭在高天、祭献在高天；让我们献上能在那祭台上供奉的祭品，不再是绵羊和牛犊，不再是血和脂肪。这一切都已废除；取而代之的是\BibleQuote{合理的敬礼}。\BibleRef{罗 12:1}但什么是\Quote{合理的敬礼}？是通过灵魂所作的祭献，是通过精神所作的祭献。（经上记载：\BibleQuote{天主是神，朝拜祂的人，应当以心神和真理朝拜祂}。\BibleRef{若 4:24}）这些祭品不需要身体，不需要器具，也不需要特殊地方，每人自己就是司祭，例如：节制、克己、慈悲、忍受虐待、恒忍、谦卑。\par

这些祭品在旧约中也预先描绘出来。经上说：\BibleQuote{向天主献上正义的祭品}\BibleRef{咏 4:5}；\BibleQuote{献上赞颂的祭品}\BibleRef{咏 49:14}；\BibleQuote{赞颂的祭品要光荣我}\BibleRef{咏 49:23}；\BibleQuote{天主所要的祭品，就是痛悔的精神}\BibleRef{咏 50:17}；\BibleQuote{上主要求你什么？无非是听从祂？}\BibleRef{米 6:8}；\BibleQuote{燔祭和赎罪祭，非你所喜悦，于是我来了，天主，我来承行你的旨意}\BibleRef{咏 40:6}；又说：\BibleQuote{你们从舍巴带来乳香有何用？}\BibleRef{耶 6:20}；\BibleQuote{你们歌唱的声音远离我，我决不听从你们琴瑟的乐声}\BibleRef{亚 5:23}；但代替这些的是\BibleQuote{我喜欢仁慈胜过祭献}\BibleRef{欧 6:6}。你看，\BibleQuote{这样的祭献是天主所喜悦的}\BibleRef{希 13:16}。你也看到，从一开始，前一类祭品已经让位，这些替代了它们。\par

因此让我们献上这些，因为前一类确是富足之人和有财产之人的祭品，但这类是德行的祭品；前者来自外部，后者来自内心；前者任何人都可献上，后者只有少数人可行。一个人比一只羊优越多少，这祭品就比那祭品优越多少；因为在这里你将自己的灵魂作为牺牲献上。\par

6.\newline{}还有其他一些祭品，确是全燔祭，即殉道者的身体；那里灵魂和身体都献上。这些有馨香之气。你也能，若你愿意，献上这样的祭品。\par

因为即使你不在火中焚烧自己的身体，那又怎样？然而你可以在另一种火中焚烧，例如，在自愿的贫穷中，在苦难中。因为有能力在奢侈和挥霍中度日，却选择劳苦和痛苦的生活，并克制身体，这不正是全燔祭吗？克制你的身体，把它钉在十字架上，你也将获得这殉道的荣冠。因为在另一种情况下由剑完成的事，在此应由自愿的心意实现。不要让对财富的贪爱焚烧或占据你，而是让这不合理的欲望本身被圣神之火吞噬和熄灭；让它被圣神之剑斩断。\par

这是一种卓越的祭献，不需要其他司铎，只要献祭者本人。这是一种卓越的祭献，虽然在地上举行，却立即被提升到高处。我们岂不惊叹古时火从天降，烧尽一切吗？如今也可能有比那更奇妙的火降下，烧尽所有献上的供物：不，不是烧尽，而是将它们提升到天上。因为它不将它们化为灰烬，而是将它们作为礼物献给天主。\par

7. 科尔乃略的奉献就是如此。因为（经上说）\Quote{你的祈祷和你的施舍已上升到天主面前，成为纪念。}\BibleRef{宗 10:4} 你看到一种最卓越的结合。当我们自己也听从来到我们身边的穷人时，我们就被垂听。经上说：\Quote{谁掩耳不听穷人的哀呼，他呼求时，天主也不俯听。}\BibleRef{箴 21:13} \Quote{眷顾贫困和可怜人的人是有福的：在上主之日，他必得救。}\BibleRef{咏 39:1} 但哪一天是恶的，除了那对罪人是恶的一天呢？\par

\Quote{眷顾}是什么意思？就是理解什么是穷人，彻底了解他的痛苦。因为了解他痛苦的人，必定立即怜悯他。当你看见一个穷人时，不要匆匆走过，而要立即反思，如果你是他，你会怎样。你难道不希望所有人都为你这样做吗？他说：\Quote{眷顾的人}。你要想，他和你是同样自由的人，享有同样的高贵出身，与你拥有一切共同之物；然而，他常常连你的狗都不如。相反，狗吃饱了，他却常常饿着肚子躺着睡觉，自由人变得比你的奴隶更低贱。\par

但他们为你做必要的服务。是什么？他们好好服侍你吗？那么让我指出，这个穷人也为你做远超过他们的必要服务。因为在审判之日，他将站在你身边，把你从火中救出。你的所有奴隶能做什么像这样的事？当塔彼达死时，是谁使她复活？是站在周围的奴隶，还是穷人？但你甚至不愿让自由人与你的奴隶平等。霜冻酷烈，穷人衣衫褴褛，几乎冻死，牙齿打颤，容貌和气息都足以感动你；而你却浑身暖和、酒足饭饱地走过，你又怎能期望天主在你遭遇不幸时拯救你呢？\par

而且你常常这样说：\Quote{如果是我自己，我遇到一个做了许多错事的人，我会原谅他；难道天主不原谅吗？}不要这样说。那个没有得罪你、你能够解救的人，你忽视了他。那么，得罪天主的你，祂怎能原谅你呢？这难道不配下地狱吗？\par

多么令人惊讶！你常常用无数色彩斑斓、金线绣制的衣服装饰一具尸体，它毫无知觉，不再感受这份尊荣；而那个因饥饿和寒冷而痛苦、哀叹、受折磨、被折磨的人，你却忽视他，更多地献给虚荣，而不是敬畏天主。\par

8. 但愿只到此为止；但立刻就对求助者提出指控。因为为什么不工作（你说）？为何他要无所事事被供养？但（告诉我）你所拥有的是靠工作得来的吗？难道不是你从你父亲那里继承来的吗？即使你工作，这怎能成为你指责别人的理由？你没有听到保禄说什么吗？因为他说过：\Quote{谁若不愿意工作，就不应当吃饭}\BibleRef{得后 3:10}，接着又说：\Quote{你们行善不要厌倦。}\BibleRef{得后 3:13}\par

但他们说什么？他是个骗子。你说什么，人哪？为了一块面包或一件衣服，你就称他为骗子？但（你说）他会立刻卖掉它。难道你所有的事情都处理得好吗？但什么？难道所有穷人都是因为懒惰吗？难道没有人因为船难而贫穷吗？没有人因为诉讼？没有人因为被抢劫？没有人因为危险？没有人因为疾病？没有人因为其他困难？然而，如果我们听见有人哀叹这样的灾祸，大声呼喊，赤裸地望天，长发破衣，我们立刻称他为：骗子！欺诈者！骗子！你不感到羞愧吗？你称谁为骗子？什么也不给，也不要去指控那人。\par

但（你说）他有办法，是装假。这是你对自己的指控，不是对他的。他知道他要面对残酷的人，面对野兽而非人，而且，即使他讲一个可怜的故事，也引不起任何人的注意：因此他不得不也装扮成更可怜的样子，以融化你的心。如果我们看见一个人穿着体面的衣服来乞讨，你会说：\Quote{这是骗子，他这样来是为了让人以为他出身好。}如果我们看见一个人穿着相反的样子，我们也指责他。那么他们该怎么办呢？哦，残忍！哦，不人道！\par

而你（说）为什么他们要暴露伤残的肢体？是因为你。如果我们有同情心，他们就不需要这些诡计；如果他们第一次求助就说服我们，他们就不会想出这些办法。有谁如此不幸，愿意如此大声喊叫，愿意做出不雅的行为，愿意公开哀哭，带着衣衫褴褛的妻子和孩子，往自己身上撒灰？这些事比贫穷更糟糕多少？然而，因为他们做了这些，他们不仅得不到怜悯，反而被我们指责。\par

9. 那么，我们还会因向天主祈祷时未蒙垂听而愤慨吗？我们还会因恳求时未得应允而烦恼吗？我亲爱的，难道我们不因恐惧而战栗吗？\par

但是（你说），我常常施舍。难道你不是常常吃饭吗？难道你常把向你乞讨的子女赶走吗？啊，无耻！你称穷人为无耻？而你自己在掠夺时却不无耻，他乞求面包时却无耻！难道你不思想肚腹的需求何等巨大？你岂不为此而行一切？你岂不为此而忽略灵性之事？不是天堂和天国摆在你面前吗？而你因惧怕那（食欲）的暴虐，忍受一切，却轻视那（天国）。这\Emph{才}是无耻。\par

你难道没看见残废的老人吗？但啊，何等琐碎！\InnerQuote{某人}（你说）\InnerQuote{出借了那么多金币，某人那么多，却仍在乞讨。}你在重述孩童的故事和闲话；因为他们也总是从保姆那里听到这类故事。我不信。我不相信。绝无此事。一个人会出借金钱，却在富足时乞讨？为了什么目的，告诉我？还有比乞讨更可耻的吗？宁愿死也不乞讨。我们的残忍何时休？那么怎样？人人都放贷吗？人人都是骗子吗？难道没有真正贫穷的人？\Quote{是的}（你说）\Quote{很多。}那么，既然你严格查问他们的生活，为何不帮助那些人呢？这是一个借口和托词。\par

\BibleQuote{凡求你的，就给他；向你借贷的，不要推辞。} \BibleRef{玛 5:42} 伸出你的手，不要握紧。我们没有被任命为人们生活的审查者，因为那样我们就不会怜悯任何人。当你呼求天主时，为什么说：不要记念我的罪？那么，即使那人是个大罪人，对他也同样宽容，不要记念他的罪。这是仁爱的季节，不是严格查问的季节；是慈悲的季节，不是算账的季节。他希望能得到供养：如果你愿意，就给予；如果不愿意，就打发他走，不要引起怀疑。你为什么悲惨痛苦？为什么你不但不怜悯他，反而还赶走那些愿意施舍的人？因为当这样的人从你口中听到：这个（家伙）是个骗子，那个是伪君子，另一个放贷，他就既不施舍给这个，也不给那个；因为他怀疑所有人都是那样。因为你晓得：我们容易怀疑恶，但不容易怀疑善。\par

10. 让我们\BibleQuote{要慈悲，不仅仅是如此，而是如同我们的天父一样。} \BibleRef{路 6:36} 祂甚至喂养奸夫、淫妇、巫术者，我还要说什么呢？那些拥有各种邪恶的人。因为在如此大的世界里，必然有很多这样的人。但尽管如此，祂喂养所有人；祂给所有人衣服。没有人曾因饥饿而灭亡，除非是出于自己的选择。所以让我们慈悲。如果有人贫穷困乏，就帮助他。\par

但如今我们竟到了如此不合理的地步，不仅对在巷间行走的穷人如此行，甚至对居住在（宗教）独居中的人也如此。你说，这样的人是个骗子。我岂不一开始就说过：如果我们不加区别地给所有人，我们就总能慈悲；但如果我们开始过分好奇地查问，我们就永远不会慈悲？你是什么意思？一个人会为了得到一块面包而成为骗子吗？如果他索要金银塔冷通，或昂贵衣服，或奴隶，或其它这类东西，人们有理由称他为骗子。但如果他什么都不求，只求食物和住所，这些适合哲学家的生活，告诉我，这难道是骗子的所为吗？让我们停止这种不合时宜的好管闲事，这是撒殚般的，是具有毁灭性的。\par

因为确实，如果有人说他在圣职名单上，或自称司铎，那么你就去（查问），大费周章：因为在这种情况下，不经查问就与之共融并非没有危险。因为危险关乎重要事项，因为你并非给予，而是接受。但如果他需要食物，就不要查问。\par

如果你愿意，就查问亚巴郎是如何款待所有来到他那里的人。如果他过分好奇地查问那些逃到他那里避难的人，他就不会\BibleQuote{款待了天使}。 \BibleRef{希 13:2} 因为也许他没有想到他们是天使，就会把他们和其余人一起赶走。但因为他惯常接待所有人，他甚至接待了天使。\par

什么？难道天主是根据那些接受（你的施舍）者的生活来赏报你的吗？不，（赏报是）出于你自己的心意，出于你丰厚的慷慨；出于你的慈爱；出于你的良善。让这（心意）常存，你必获得一切美善；愿我们众人也获得这一切，藉着我们的主耶稣基督的恩宠和慈爱，愿光荣、权能、尊崇归于祂，并与圣父及圣神，自今至永远，世世无穷。阿们。\par

\SourceRecord{Translated by Frederic Gardiner.；Nicene and Post-Nicene Fathers, First Series；Vol. 14.；Edited by Philip Schaff.；Buffalo, NY: Christian Literature Publishing Co.,；1889.；<http://www.newadvent.org/fathers/240211.htm>.}

% structure: p0001 source=h1 target=section reason=page-title

\section{《希伯来人书》第十二篇讲道}

% structure: p0002 source=h2 target=subsection reason=relative-heading-rank; base=h2

\subsection{希 7:1-3}

\Quote{因为这默基瑟德，撒冷王，至高天主的司祭，当亚巴郎击败众王回来时，前去迎接他，并祝福了他；亚巴郎也将一切所得的十分之一分给了他。默基瑟德，按名字的解说，首先是\InnerQuote{正义之王}，其次也是撒冷王，即\InnerQuote{和平之王}。他无父、无母、无族谱，无生日之始，亦无生命之终，却相似天主子，常任司祭。}\par

1. 保禄为显明新约与旧约之间的区别，处处散布此事；他从远处射箭，广为宣扬，并预先准备。因为从一开始，在引言中，他就立下了这话：\Quote{对他们，天主藉先知说话，但对我们，却藉圣子说话}\EditorialNote{见一1-2}；对他们，是\Quote{多次并以多种方式}，但对我们，却是藉圣子。随后，他论及圣子是谁、做了何事，并劝勉我们要听从祂，免得遭受与犹太人相同的遭遇；他又说祂\Quote{是按默基瑟德的品位，成为大司祭}\BibleRef{希 6:20}；他多次想进入这区别的讨论，并做了许多预备工作；他责斥他们软弱，又安慰他们、恢复他们的信心；最后，他才将关于两约区别的讨论引入到精神饱满的听众耳中。因为精神沮丧的人不会是个乐意听者。为让你们明白这一点，请听圣经所言：\Quote{他们因心苦而不能听梅瑟的话。}\BibleRef{出 6:9}因此，保禄先藉许多考虑——有些令人恐惧，有些较为温和——除去了他们的沮丧，然后才从这一点进入两约区别的讨论。\par

2. 他说了什么？\Quote{因为这默基瑟德，撒冷王，至高天主的司祭。}尤其值得注意的是，他藉类型本身显示区别之大。如我所说，因着听众的软弱，他不断从事物、从过去之事来确认真理。\Quote{因为这默基瑟德，撒冷王，至高天主的司祭，当亚巴郎击败众王回来时，前去迎接他，并祝福了他；亚巴郎也将一切所得的十分之一分给了他。}简洁地叙述了整个故事后，他神秘地注视了它。\par

首先从名字入手。\Quote{按名字的解说，首先是\InnerQuote{正义之王}}因为\OriginalTerm{正义}{Sedec}意为\InnerQuote{正义}；\OriginalTerm{王}{Melchi}意为\InnerQuote{王}；默基瑟德即\InnerQuote{正义之王}。你们看见他在名字上的精确吗？但谁是\InnerQuote{正义之王}，除非是我们的主耶稣基督？\InnerQuote{正义之王}。\Quote{其次也是撒冷王}，从他的城得名，\Quote{即\InnerQuote{和平之王}}，这又属于基督的品性。因为祂使我们成义，并\Quote{缔造了和平}，使\Quote{天上的和地上的万物}\BibleRef{哥 1:20}和好。哪个人是\InnerQuote{正义与和平之王}？没有，唯有我们的主耶稣基督。\par

3. 他又加上了另一特点：\Quote{无父、无母、无族谱，无生日之始，亦无生命之终，却相似天主子，常任司祭。}既然在经文\Quote{你永为司祭，按照默基瑟德的品位}中有一个困难摆在他面前——默基瑟德已经死了，并非\InnerQuote{永为司祭}——请看他是如何神秘地解释了这一点。\par

\InnerQuote{谁能对人说这话呢？}他说：\InnerQuote{我并非在事实上断言这一点；意思是，我们不知道他何时、有何父亲，有何母亲，也不知道他的开始或死亡。}那么，有人会说：\InnerQuote{因我们不知道，就因此他没有死、没有父母吗？}你说得对：他既死了，也有父母。那他又如何是\InnerQuote{无父、无母}？如何\InnerQuote{无生日之始、无生命之终}？何以如此？因为经上没有记载。这又怎样？正如此人因未记其家谱而如此，基督则因事实本身的实况而如此。\par

请看\InnerQuote{无始}，请看\InnerQuote{无终}。如同此人之例，我们不知道他的\InnerQuote{生日之始}或\InnerQuote{生命之终}，因未曾记载；同样，在耶稣身上，我们也不知道，不是因为未曾记载，而是因为不存在。因为前者是类型，故我们说\InnerQuote{因未曾记载}；后者是实体，故我们说\InnerQuote{因不存在}。正如在名字上（\InnerQuote{正义之王}和\InnerQuote{和平之王}是称谓，但这里是实体），同样，那些称谓在那里，实体在这里。那么，祂如何有开始？你们看见：圣子\InnerQuote{无始}，并非指祂没有原因（因为这是不可能的：祂有圣父，否则如何是子？），而是指祂\InnerQuote{无生日之始与生命之终}。\par

\Quote{相似天主子。}相似点在哪里？在于我们不知道两者的终或始。对那人，因未记载；对基督，因不存在。这就是相似。若相似在各方面都存在，就不再是类型与实体，而两者都是类型了。此处正如在绘画中，有相似之处，也有不相似之处。在线条上，特征相似，但加上颜色时，就清楚显出差异，既有相似也有不相似。\par

% structure: p0011 source=h2 target=subsection reason=relative-heading-rank; base=h2

\subsection{希 7:4}

4. \BibleQuote{Now consider}（他说）\BibleQuote{这位何等伟大，连圣祖亚巴郎也将战利品的十分之一给了他。}至此，他一直在运用预像；从此处起，他大胆地显示他【默基瑟德】比犹太的实在更荣光。但若那承载基督预像的人远比司铎们甚至司铎们的祖先本人更优越，那么对这实在本身又该说什么呢？你看他是如何极力表明这优越性。\par

\BibleQuote{现在请想一想}（他说）\BibleQuote{这位何等伟大，连圣祖亚巴郎也将精选之物中的十分之一给了他。}战场上夺取的战利品被称为\Quote{精选之物。}不能说他给他这些是因为他在战争中有份，因为（他说）他遇见他\BibleQuote{从击杀众王回来}，因为他（意思是）留在家中，然而【亚巴郎】将自己劳苦的初果给了他。\par

% structure: p0014 source=h2 target=subsection reason=relative-heading-rank; base=h2

\subsection{\BibleRef{希 7:5-6}}

\BibleQuote{凡由肋未子孙中领受司铎职的，依照法律有向人民，即向他们的弟兄，收取十分之一的诫命，尽管他们都是出自亚巴郎的腰。}\par

（他意谓）司铎职的优越性如此之大，以致那些从祖先起就有同样尊位、并有同一祖先的人，仍远比其他人为优。无论如何，他们从他们那里\Quote{收取十分之一。}那么，当发现有一位从这些人本身收取十分之一时，他们岂不确实居于平信徒之列，而他则居于司铎之中吗？\par

不仅如此，而且他与他们并不具有同样的尊位，而是出自另一个民族：这样，除非他的尊位很大，否则他不会向一个陌生人缴纳十分之一。令人惊讶！他成就了什么？他清楚显明了比他在《致罗马人书》中论信德所处理的那些论题更重大的事。因为在那里他确实宣布亚巴郎是我们政体以及犹太政体的祖先。但在这里他大胆地针对他，并表明那未受割损的人远为优越。那么他如何显明了肋未缴纳了十分之一？亚巴郎（他说）缴纳了。\Quote{这与我们何干？}它特别与你们相干：因为你们不会争辩说肋未子孙优于亚巴郎。\BibleRef{希 7:6}\BibleQuote{至于那与肋未子孙没有族谱关系的，却从亚巴郎收取了十分之一。}\par

% structure: p0018 source=h2 target=subsection reason=relative-heading-rank; base=h2

\subsection{\BibleRef{希 7:7}}

而后他没有简单地略过，反而加上说：\BibleQuote{并且祝福了那领受恩许的。}因为既然这事普遍被尊崇，他显明【默基瑟德】比亚巴郎更应受尊崇，这是根据众人的共同判断。\BibleRef{希 7:7}\BibleQuote{毫无疑问地，}他说，\BibleQuote{那较低的由那较高的祝福。}即，依众人之见，是位份低的被位份高的祝福。那么，基督的预像甚至优于\Quote{那领受恩许的}。\par

% structure: p0020 source=h2 target=subsection reason=relative-heading-rank; base=h2

\subsection{\BibleRef{希 7:8-10}}

\BibleRef{希 7:8}\BibleQuote{此处收受十分之一的是有死的人，但那里却是被证明活着的一位。}但为了避免我们说：告诉我们，你为什么追溯到这么远？他说：\BibleRef{希 7:9}\BibleQuote{并且，容我这样说}（他这样缓和了语气）\BibleQuote{就连那收受十分之一的肋未，也因亚巴郎而缴纳了十分之一。}如何？\BibleRef{希 7:10}\BibleQuote{因为默基瑟德遇见他的时候，肋未还在他祖先的腰中。}即，肋未在他内，尽管他尚未出生。他不是说肋未子孙，而是说肋未。\par

你看到了优越性吗？你看到亚巴郎与那承载我们大司铎预像的默基瑟德之间何等大的间隔吗？他还表明这优越性是由权柄而非必然性造成的。因为一方缴纳十分之一，这表示司铎；另一方给予祝福，这表示优越者。这优越性也传到后代。\par

他以奇妙而得胜的方式驱逐了犹太【体制】。为此他说：\BibleQuote{你们成了迟钝的人，}\BibleRef{希 5:12}，因为他愿意奠定这些基础，使他们不致退缩。这就是保禄的智慧，先妥善预备他们，然后引导他们进入他所愿意的。因为人类本性难以说服，需要许多注意，甚至比植物更需要。因为在那情形中，只有物质身体的本性和土壤，它服从农人的手；但在此处却有意志，它容易有许多变化，时而偏爱这个，时而偏爱那个。因为它迅速转向恶。\par

5. 因此我们应当时常\Quote{警醒}自己，免得我们偶然沉睡。因为（经上记载）\BibleQuote{看，那保护以色列的，不瞌睡，也不睡觉。}\BibleRef{咏 120:4}，\BibleQuote{不容你的脚动摇。}\BibleRef{咏 120:3} 他没有说\Quote{不要动摇}，而是说\Quote{不容你}等等。因此动摇与否取决于我们自己，而不取决于其他。因为如果我们愿意\BibleQuote{坚定不移，毫不动摇}\BibleRef{格前 15:58}，我们便不会被动摇。\par

那么如何？难道什么都不依赖天主吗？一切都依赖天主，但并非如此以致我们的自由意志受阻。'如果这样依赖天主，'（有人说）'为什么他责备我们？'为此我说：'并非如此以致我们的自由意志受阻。'因此它既依赖我们，也依赖他。因为我们必须先选择善；然后他引领我们到他那里。他不预先我们的选择，以免我们的自由意志受侵犯。但当我们选择了，那么他带给我们的帮助是大的。\par

那么保禄怎么说：\BibleQuote{不在乎那愿意的}如果也依赖我们自己，\BibleQuote{也不在乎那奔跑的，而是在乎那仁慈的天主。}\BibleRef{罗 9:16}\par

首先，他不是以己见提出，而是从他面前和已提出的【讨论】中推论出来的。因为他说\BibleQuote{经上记载：我要仁慈对待我要仁慈的人，我要怜悯我要怜悯的人}\BibleRef{罗 9:15}之后，他说\BibleQuote{这样看来，不在乎那愿意的，也不在乎那奔跑的，而是在乎那仁慈的天主。}\BibleQuote{那么你要对我说：为什么他还指责呢？}\BibleRef{罗 9:16}\par

其次，也可以给出另一种解释，即他称所有一切都是祂的，因为主要部分属于祂。因为选择与愿望是我们的；但完成与终结是天主的。所以，既然主要部分出于祂，他就说一切都是出于祂，这是按照人的习惯说法。因为我们自己也这样做。例如：我们看到一座房子建造得很好，就说整座房子是建筑师的杰作，但当然并不全是他的，也有工人的，还有提供材料的主人以及其他许多人的，然而由于他贡献了最大的部分，我们就说整座房子是他的。所以在这种情况下也是如此。再者，对于一群人，当人数众多时，我们说\Quote{所有的人}是；人数少时，则说\Quote{没有人}。所以保禄也说：\Quote{不在乎那愿意的，也不在乎那奔跑的，只在乎那施仁慈的天主。}\par

在此他建立了两个伟大的真理：一是我们不应自高自大；即使你在奔跑（他会说），即使你非常热心，也不要认为善行是属于你自己的。因为如果你没有得到来自上天的推动，一切都是徒劳的。然而，你将获得你努力追求的目标是非常明显的；只要你奔跑，只要你愿意。\par

他并非断言我们奔跑是徒然的，而是说，如果我们认为全部都是我们自己的，如果我们不将更大的部分归于天主，我们就白跑了。因为天主也不愿完全是祂的，免得祂显得无缘无故地给我们加冕；也不愿完全是我们的，免得我们陷入骄傲。因为如果当我们拥有较小的一份时，我们就很看重自己，那么如果全部取决于我们，我们又会怎样呢？\par

6. 的确，天主为了削除我们的自夸，已除去了许多事物，但高傲的势力仍然存在。祂用多少苦难包围我们，以削除我们骄傲的精神！祂用多少野兽环绕我们！因为当有些人说：\Quote{为什么这样？}\Quote{这有什么用？}他们是在说违背天主意愿的话。祂将你置于如此巨大的恐惧之中，但即使如此，你也不谦卑；但如果你一旦取得一点成功，你就骄傲地高及上天。\par

因此，迅速的变迁和逆转发生；但即使如此，我们也不受教。因此，持续不断地有不合时宜的死亡，但我们却好像自己是不死的、永远不会死一样。我们掠夺、欺诈，好像永远不必交账。我们建造，好像要永远住在这里。甚至连每日传入我们耳中的天主圣言，以及事件本身都不能教导我们。没有一天、没有一个时辰，我们看不到接连不断的出殡。但一切都是徒然：没有什么能触及我们的铁石心肠；我们甚至不能因他人的灾祸而变得更好；或者不如说，我们不愿如此。只有当我们自己受苦时，我们才被制服，但如果天主松了手，我们再次举手：没有人考虑什么是人的本分，没有人轻视地上的事物；没有人仰望上天。但如同猪低着头，俯向它们的肚子，在泥沼中打滚；同样，人类的大多数也以最不能容忍的污秽玷污自己，却没有意识到这一点。\par

7. 因为被不洁的泥土玷污，比被罪恶玷污更好；因为被泥土玷污的人，片刻就能洗掉，变得像从未掉进泥坑一样；但掉进罪恶深坑的人，沾染的污秽不是水所能洗净的，而是需要长时间、严格的补赎、眼泪和哀哭，以及比我们对最亲爱的人所表现的更热切的悲叹。因为这污秽是从外部附着于我们，所以我们也能迅速除掉它；但另一种是从内部产生的，因此我们难以洗净，并洁净自己。\BibleQuote{因为从心里}（经上说）\BibleQuote{发出恶念、奸淫、邪淫、偷盗、妄证。}\BibleRef{玛 15:19} 因此先知也说：\BibleQuote{天主，求你给我再造一颗纯洁的心。}\BibleRef{咏 50:10} 另一位说：\BibleQuote{耶路撒冷啊，从你心中洗掉邪恶。}\BibleRef{耶 4:14}（你看，这既是我们的工作，也是天主的。）又说：\BibleQuote{心里纯洁的人是有福的，因为他们要看见天主。}\BibleRef{玛 5:8}\par

让我们尽最大努力变得洁净。让我们擦去我们的罪恶。如何擦去，先知教导说：\BibleQuote{你们应洗涤，自洁，从我眼前除去你们的恶行。}\BibleRef{依 1:16} 什么是\Quote{从我眼前}？因为有些人似乎免于邪恶，但只是对人而言，而在天主面前，他们显明是\Quote{粉饰的坟墓}。所以祂说，要照我所看见的除去它们。\BibleQuote{学习行善，寻求公理，为穷人和卑微人伸冤。}\BibleQuote{现在你们来，我们彼此辩论——上主说：你们的罪虽如朱红，我将使它们白如雪；虽红如丹砂，我将使它们白如羊毛。}\BibleRef{依 1:17} 你看，我们必须首先洁净自己，然后天主洁净我们。因为祂先说了\BibleQuote{你们应洗涤，自洁}，然后加上了\BibleQuote{我将使你们变白。}\par

那么，即使是那些达到了极端邪恶的人，也不要对自身绝望。因为（祂说）即使你已经养成了习惯，甚至几乎变成了邪恶的本性，也不要害怕。因此，祂以那些不是表面的、几乎成为材料本质的颜色为例，说祂要把它们变为相反的状态。因为祂不是简单地说祂要\BibleQuote{洗净}我们，而是说祂要使我们\BibleQuote{白如雪，白如羊毛}，以便在我们面前摆出美好的希望。那么，补赎的力量是多么伟大啊，至少它能使我们如雪，使我们如羊毛般洁白，即使罪恶已先占据并染黑了我们的灵魂。\par

那么，让我们殷勤努力，成为洁净的；祂并没有命令任何沉重的事。\Quote{你们当为孤儿伸冤，为寡妇主持正义。}\BibleRef{依 1:17}你们随处都看见，天主何等重视怜悯，以及为受冤屈的人挺身而出。让我们追求这些善行，这样我们也能因天主的恩宠，得以获得将来的祝福。愿我们所有人都配得这些恩赐，在我们的主\Person{基督耶稣}内，愿光荣、权能、尊威归于祂及圣父与圣神，自今至永远，及无穷世。阿们。\par

\SourceRecord{Translated by Frederic Gardiner.；Nicene and Post-Nicene Fathers, First Series；Vol. 14.；Edited by Philip Schaff.；Buffalo, NY: Christian Literature Publishing Co.,；1889.；<http://www.newadvent.org/fathers/240212.htm>.}

% structure: p0001 source=h1 target=section reason=page-title

\section{希伯来书讲道集 第十三篇}

% structure: p0002 source=h2 target=subsection reason=relative-heading-rank; base=h2

\subsection{希伯来书 7:11-14}

\BibleQuote{\Emph{如果因此成全曾借着}\Emph{肋未司祭职；（因为选民就是借着这司祭职承受了法律？）}\Emph{又有什么需要兴起另一位司祭，按照默基瑟德的品级，而不称为按照亚郎的品级呢？因为司祭职一更改，法律也必须变更。}\Emph{因为所论及的那一位，原属于另一支派，从没有人在祭台前服务过。}\Emph{显然我们的主是由犹大支派出身的，关于这一支派，梅瑟从未提到司祭的事。}}\par

1. \Quote{如果因此}（他说）\Quote{成全借着肋未司祭职}。论及\Person{默基瑟德}，并显示他比\Person{亚巴郎}优越多少，以及陈述了他们之间的巨大差异后，从现在起他开始证明盟约本身的巨大差异，以及一个是怎样的不完全，另一个是完全的。然而，他还没有直接进入这些事本身，而是首先在司祭职和会幕的基础上辩论。因为这些事，当证明来自已被认可和相信的事实时，更容易被不信者接受。\par

他曾显示\Person{默基瑟德}远优越于\Person{肋未}和\Person{亚巴郎}，在司祭的等级中高于他们。他又从不同的角度论证。那么这是什么？为什么（他说）他没有说\Quote{按照\Person{亚郎}的品级}？并且，我请求你注意，[他的论证的]极大优越性。因为正是从那个原本排除他的司祭职的境况，即他不是\Quote{按照\Person{亚郎}的品级}，他从那一境况建立了他，并排除了其他。因为这就是我所说的（他宣称）；为什么他\Quote{没有被按照\Person{亚郎}的品级所造}？\par

并且，\Quote{还有什么需要}这句话有很重的强调。因为如果\Person{基督}曾按照肉体\Quote{按照\Person{默基瑟德}的品级}，然后后来法律被引入，以及所有属于\Person{亚郎}的事，人们可以合理地说，后者作为更完全的，废除了前者，既然它是在其后出现的。但是，如果\Person{基督}后来才来，并且采取了不同的类型，如同他的司祭职的类型，那么显然是因为那些是不完全的。因为（他会说）我们为了论证起见，假设一切都已实现，司祭职中没有任何不完全。那么，在这种情况下，\Quote{有什么需要}使他被称为\Quote{按照\Person{默基瑟德}的品级而不按照\Person{亚郎}的品级}？为什么他废弃了\Person{亚郎}，而引入了另一个司祭职，即\Person{默基瑟德}的？\Quote{如果因此成全}，即事情本身、教义、生活的成全，\Quote{曾是借着肋未司祭职。}\par

并注意他是如何在他路径上向前的。他曾说[他是]\Quote{按照\Person{默基瑟德}的品级}，暗示着[司祭职]\Quote{按照\Person{默基瑟德}的品级}是更优越的：因为[他]远为优越。之后他从时间上也表明这一点，即他在\Person{亚郎}之后；显然作为更好的。\par

2. 并且，下面的话是什么意思？\Quote{因为}（他说）在它之下（或\Quote{在它之上}）人民承受了法律（或\Quote{曾为之立法}）。什么是\Quote{在它之下}等？它以自己为秩序，通过它行一切事。你不能说它是赐给别人的，\Quote{在它之下的人民承受了法律}，即，他们使用了它，也确实使用了它。你确实不能说它是完全的，它没有治理人民；\Quote{他们在它之上被立法}，即，他们使用了它。\par

那么，为什么还需要另一个司祭职呢？\Quote{因为司祭职一更改，法律也必须变更。}但如果必须有另一位司祭，或者更确切地说，另一个司祭职，那么也必须另有一部法律。这是针对那些说\Quote{为什么需要一个新盟约？}的人。因为他本也可以从预言援引证据。\Quote{这是我与你们的祖先所立的盟约}等。\BibleRef{希 8:10} 但现在他是在司祭职的基础上争辩。并且，注意他是如何从一开始就说出这话的。他说：\Quote{按照\Person{默基瑟德}的品级。}借此他排除了\Person{亚郎}的品级。因为如果另一个更好，他就不会说\Quote{按照\Person{默基瑟德}的品级。}因此，如果另一个司祭职被引进来，也必须有一个[另一个]盟约；因为既不可能有司祭而没有盟约、法律和规条，也不可能他接受了不同的司祭职却使用前一个[盟约]。\par

其次，关于异议的基础：\Quote{如果他不是\Person{肋未}人，他怎能做司祭？}他既然通过上面所说的话推翻了这一点，甚至认为不值得回答，只是在顺便提一下。他说（他的意思是）司祭职被更改了，因此盟约也被更改了。它不仅在性质或规条上被更改，也在支派上被更改。因为必须[它被更改]也在支派上。如何？因为司祭职被更改（或\Quote{转移}），从支派到支派，从司祭支派到君王支派，使同一人同时是君王又是司祭。\par

并且注意其中的奥迹。首先它是君王的，然后它成为司祭的：因此在\Person{基督}身上也是如此：因为他确实是永远为君王，但从他取了肉体、献了祭献那一刻起，他成为司祭。你看到这变化，而那些本是异议根据的事，他却引入它们，好像事物的自然秩序要求它们一样。\Quote{因为}（他说）\Quote{这些事所论及的那一位属于另一支派。}我自己也这样说，我知道这个支派[犹大支派]与司祭职毫无关系。因为有一个转移。\par

3. 是的，我还要指出另一个区别（他会说）：不仅从支派，也不仅从位格，也不仅从司铎职的性质，也不仅从盟约，而且从预像本身。\BibleRef{希 7:16}祂成为（\Quote{成了}如此），不是按照属肉体的诫命的法律，而是按照无穷生命的能力。祂成了（他说）\BibleQuote{一位司铎，不是按照属肉体的诫命的法律}：因为那法律在许多方面是无效的。\par

什么是\Quote{属肉体的诫命}？它说：给肉体行割损；给肉体傅油；洗涤肉体；洁净肉体；剃肉体；绑在肉体上；滋养肉体；使肉体安息。再者，它的祝福是什么？肉体的长寿；肉体的奶与蜜；肉体的平安；肉体的奢华。从这法律，亚郎领受了司铎职；默基瑟德却并非如此。\par

% structure: p0014 source=h2 target=subsection reason=relative-heading-rank; base=h2

\subsection{希伯来书 7:15-17}

\BibleQuote{且更明显的是，若有一位像默基瑟德一样的另一位司铎兴起。} 什么明显？两种司铎职之间的间隔，差别；祂是多么优越，\BibleQuote{祂成了司铎，不是按照属肉体的诫命的法律。}（谁？默基瑟德？不；而是基督。）\BibleQuote{但是按照无穷生命的能力。因为祂证明说：你永为司铎，按照默基瑟德的品位}；也就是说，不是为一时，也没有任何限制，\BibleQuote{而是按照无穷生命的能力}，即借着能力，借着\Quote{无穷生命}。\par

然而这并非接着\Quote{祂被立为，不是按照属肉体的诫命的法律}：因为接着要说的应是\Quote{而是按照属神的法律}。但借着\Quote{属肉体的}，他暗示了暂时的。正如他在另一处也说，属肉体的规例，直到改革的时候才被强加。\BibleRef{希 9:10}\par

\Quote{按照生命的能力}，这就是说，因为祂凭自己的能力活着。\par

4. 他曾说过，法律也有变更，到目前为止他已证明了这一点；此后他探究原因，这原因最能使人的心智得到完全的确证，[我是指]彻底认识原因；当我们也知道了原因，以及事物发生的原理时，它更能引导我们趋向信德。\par

% structure: p0019 source=h2 target=subsection reason=relative-heading-rank; base=h2

\subsection{希伯来书 7:18}

\BibleQuote{确实有}（他说）\BibleQuote{先前诫命的废弃，因其软弱和无益。}在这里异端者紧逼。但要仔细听。他并没有说\Quote{因其邪恶}，也没有说\Quote{因其败坏}，而是说\Quote{因其软弱和无益}；他在别处也显示这软弱；正如当他说\BibleQuote{在于它因肉体而软弱。}\BibleRef{罗 8:3} 那么法律本身并不软弱，而是我们。\par

% structure: p0021 source=h2 target=subsection reason=relative-heading-rank; base=h2

\subsection{希伯来书 7:19}

\BibleQuote{因为法律未成全任何事物。} 什么是\Quote{未成全任何事物}？它没有使任何人成全，由于被违抗。此外，即使它被听从，也不会使人成全和德善。但此时他在这里还没有说这个，而是说它没有力量：这是有充分理由的。因为那里写下的诫命是：做这个，不要做那个，只是命令，而没有赐予内在的力量。但\Quote{这盼望}却不是如此。\par

什么是\Quote{废弃}？是驱逐。\Quote{废弃}是对有效之物的废弃。所以，他暗示它[曾]有效，但从此以后无效，因为它一事无成。那么法律毫无用处吗？它确实有用，而且有大用：但在使人成全上却毫无用处。因为在这方面他说，\BibleQuote{法律未成全任何事物}。一切都是预像，一切都是影子；割损、祭献、安息日。因此它们不能深入灵魂，故此它们消逝并渐渐退出。\BibleQuote{但一个更好的盼望的引入做到了，借着它我们亲近天主。}\par

% structure: p0024 source=h2 target=subsection reason=relative-heading-rank; base=h2

\subsection{希伯来书 7:20}

5. \BibleRef{希 7:20}\BibleQuote{而且既不是没有起誓的。}你看见起誓的事在这里对他成为必要的。因此他先前为此讨论了很多，即天主如何发了誓；并且为了[我们]更充分的确定而发誓。\par

\BibleQuote{但一个更好的盼望的引入。}因为那体系也有一个盼望，但不像这个。他们盼望，如果中悦天主，他们应拥有土地，不受任何可怕之事。但在这一时代，我们盼望，如果我们中悦天主，我们将拥有不是地，而是天；或者更确切地说（这远胜于此）我们盼望接近天主，来到圣父的宝座前，与天使一同事奉祂。你看他是如何逐步引入这些事的。因为上面他说\BibleQuote{进入帷幕之内的}，\BibleRef{希 6:19}，但这里，\BibleQuote{借此我们亲近天主。}\par

% structure: p0027 source=h2 target=subsection reason=relative-heading-rank; base=h2

\subsection{希伯来书 7:21-24}

\Quote{况且不是没有誓言。}\Quote{况且不是没有誓言}是什么意思？就是说，请看另一个区别。这些事也不只是应许的（他说）。\Quote{因为那些司祭没有誓言就任职的，} \BibleRef{希 7:21} \Quote{但这一位却藉着那对祂说：\Quote{主起了誓，决不后悔，\InnerQuote{你按照默基瑟德的品位，永为司祭。}}的那一位，以誓言任职。既然如此，耶稣就成了更美盟约的保证人。} 他提出了两个区别：它没有终结，不像法律的（盟约）；他从（执行司祭职的）是基督这一点来证明这事；因为他说\Quote{按照不可消灭的生命的能力}。他也从誓言来证明，因为\Quote{祂起了誓}等等，并从事实来证明；因为如果另一个被抛弃，因为它软弱，这个却因它有力而稳固。他也从司祭来证明。如何？因为祂是唯一的。除非祂是不死的，就不会是唯一的。因为正如有许多司祭，因为他们会死，同样（这里）是唯一的一位，因为祂是不死的。\Quote{既然如此，耶稣就成了更美盟约的保证人，}因为祂向祂发誓说祂将永远（为司祭）；如果祂不是永生的，祂不会这样做。\par

% structure: p0029 source=h2 target=subsection reason=relative-heading-rank; base=h2

\subsection{希伯来书 7:25}

6. \BibleRef{希 7:25} \Quote{因此，凡靠祂进到天主面前的，祂都能拯救到底，因为祂永远活着，为他们转求。} 你们看，这话是针对那按肉体而言的。因为当祂（显现）为司祭时，祂也转求。因此当保禄说 \Quote{祂也为我们转求} \BibleRef{罗 8:34}，他暗示同一件事；大司祭转求。因为\Quote{祂随己意叫死人复活，使他们活过来} \BibleRef{若 5:21}，并且\Quote{如同父}（一样），（当需要拯救时，祂）怎么\Quote{转求}呢？\BibleRef{若 5:22} 祂有\Quote{一切审判}，怎么\Quote{转求}呢？祂\Quote{派遣祂的天使} \BibleRef{玛 13:41}，使天使把一些人\Quote{投在} \Quote{火窑}里，而拯救另一些人，祂怎么\Quote{转求}呢？因此（他说）\Quote{祂能拯救}。为此祂拯救，因为祂不死。既然\Quote{祂永远活着}（他意味）祂没有继承人；如果祂没有继承人，祂就能帮助所有的人。因为（在律法下）确实，大司祭尽管在他（为大司祭）的期间值得钦佩（例如撒慕尔，以及其他类似的人），但此后就不再是这样；因为他们死了。但这里不是这样，而是\Quote{祂}拯救\Quote{到底}。\par

\Quote{到底}是什么意思？他暗示某种奥迹。他说，不仅在这里，在那里也拯救那些\Quote{靠祂进到天主面前的人}。祂如何拯救？\Quote{因为祂永远活着，为他们转求}。你们看到屈尊就卑了吗？你们看到人性了吗？因为他说，祂不是藉着一次转求而获得这个，而是持续地，每当需要为他们转求时。\par

\Quote{到底}是什么？不是暂时的，而是在那里，在来世。那么祂总需要祈祷吗？但这怎么合理呢？就连义人也常常藉一次恳求成就一切，而祂总在祈祷吗？那么祂为何与（父）同坐呢？你们看，这是俯就。意思是：不要害怕，不要说：是的，祂确实爱我们，而且祂对父有信赖，但祂不能永远活着。因为祂永远活着。\par

% structure: p0033 source=h2 target=subsection reason=relative-heading-rank; base=h2

\subsection{希伯来书 7:26}

7. \Quote{因为这样的大司祭才适合我们，祂是圣洁的、无恶的、无玷的、从罪人中分别出来的。} 你们看，这一切都是针对人性而言的。（但当我论及\InnerQuote{人性}时，我指的是具有天主性的人性；不是将二者分开，而是让你们推论出合适的意思。）你们注意到大司祭的区别了吗？他总结了前面所说的，\Quote{在一切事上受过试探，像我们一样，只是没有罪。} \BibleRef{希 4:15} \Quote{因为}（他说）\Quote{这样的大司祭才适合我们，祂是圣洁的、无恶的。} \Quote{无恶的}是什么意思？没有恶毒；即另一位先知所说的：\Quote{祂口中没有诡诈} \BibleRef{依 53:9}，也就是说，（祂）不狡猾。有人能这样论及天主吗？难道不羞于说天主不狡猾、不诡诈？然而，关于祂，就肉体而言，这样说可能是合理的。\Quote{圣洁的、无玷的。} 这也有人能这样论及天主吗？因为祂的本性可能被玷污吗？\Quote{从罪人中分别出来的。}\par

% structure: p0035 source=h2 target=subsection reason=relative-heading-rank; base=h2

\subsection{希伯来书 7:27}

8. 难道只有这一点显示区别，还是祭献本身也显示区别？如何？\BibleRef{希 7:27} \Quote{祂不需要像大司祭那样每日先为自己的罪献祭，因为祂一次而永远地奉献了自己，完成了这事。} \Quote{这事}是什么？这里下文是有关超然祭献的伟大和距离的序曲。他曾提及司祭的一点；曾提及信德的一点；曾提及盟约的一点；虽没有完全，但仍提及了。这里下文是有关祭献本身的序曲。所以，听到祂是司祭，不要以为祂常行司祭之职。因为祂一次执行了，然后\Quote{坐下}了。\BibleRef{希 10:12} 免得你以为祂在高处站着，是仆役，他就表明这事是一份（救恩）计划的一部分。因为正如祂成了仆人，同样（祂也成了）司祭和仆役。但正如成为仆人后，祂没有继续做仆人，同样成为仆役后，祂也没有继续做仆役。因为坐不属于仆役，而是站立。\par

他在这里暗示了这一点，也暗示了祭献的伟大，即虽是唯一的，且仅被献上一次，却成就了其他所有祭献无法成就的事。但他尚未处理这些要点。\par

% structure: p0038 source=h2 target=subsection reason=relative-heading-rank; base=h2

\subsection{希伯来书 7:28}

\Quote{因为祂行了这事}他说。\Quote{这事}；什么事？\Quote{因为}（他说）\Quote{这个【人】必须也要有所奉献}\BibleRef{希 8:3}；不是为自己；因为祂如何奉献自己？而是为百姓。你怎么说？祂能这样行吗？是的（他说）。\Quote{因为法律立了那些有软弱的人为大司祭}\BibleRef{希 7:28}难道祂不需要为自己奉献吗？不，他说。因为，免得你以为\Quote{这一事}\Quote{祂一次而竟全功}这些话也是指祂自己，请听他说：\Quote{因为法律立了那些有软弱的人为大司祭}为此，他们既连续不断地奉献，又为自己奉献。然而，祂是强有力的，祂没有罪，为什么祂要为自己、或多次为他人奉献呢？\par

\Quote{但自法律以后起誓的话，立了那永远成为完善的圣子。}\Quote{成为完善}：那是什么？\Person{保禄}没有列出对比的通用术语；因为说了\Quote{有软弱}之后，他没有说\Quote{圣子}是强有力的，而是说\Quote{成为完善}：即，正如有人所说，强有力的。你看，\Quote{圣子}这个名称是用来与\Quote{仆人}对比的。而\Quote{软弱}这个词，他是指罪或死亡。\par

\Quote{永远}是什么意思？不仅现在无罪，而且永远无罪。如果祂是完美的，如果祂从未犯罪，如果祂永远活着，为什么祂要为我们献上许多祭献？但目前他并不极力强调这一点；他所极力强调的是，祂不为自己奉献。\par

既然我们有这样一位大司祭，让我们效法祂；让我们跟随祂的足迹。没有别的祭献：只有一位独自洁净了我们，在此之后，就是火与地狱。的确，正因如此，他一再重复说：\Quote{一位司祭}，\Quote{一次祭献}，免得有人以为有许多祭献而无所畏惧地犯罪。那么，凡被堪当承受那神印的、凡享用了那祭献的、凡分享了那不朽之筵席的，让我们继续守护我们的高贵出身和尊严，因为堕落并非没有危险。\par

而那些尚未被堪当承受这些恩惠的人，也不可因此自恃。因为当一个人继续犯罪，指望临终时领受圣洗，往往无法获得。相信我，我说这话不是为了吓唬你。我自己认识许多人遇到了这种情况，他们因期待那光照而大肆犯罪，在死亡之日空手而去。因为天主赐给我们圣洗，是为了除去我们的罪，而不是增加我们的罪。如果有人把圣洗当作更多犯罪的保障，它就成了疏忽的原因。因为如果没有那洗濯，他们本会更加谨慎地生活，因为没有赦罪的方法。你看，正是我们导致有人说：\BibleQuote{让我们作恶，为得善果}\BibleRef{罗 3:8}\par

因此，我也劝告你们未入门的人要清醒。愿没有人像雇工一样追求德行，没有人像愚昧的人，没有人像追求沉重负担一样。让我们以乐意的心和喜乐去追求它。因为即使没有赏报存留，我们岂不应该行善吗？但至少，为赏报，让我们成为善良。然而这岂不是羞耻和极大的谴责吗？除非你给我赏报，我就不会节制。那么我敢说：你永远不会节制，即使你过着节制的生活，如果你是为了赏报而做。你根本不重视德行，如果你不爱它。但由于我们极大的软弱，天主愿意暂时让人甚至为赏报而实践德行，然而即使如此，我们也不去追求它。\par

但假设，你若愿意，一个人死了，他行了无数恶事，也堪当承受了圣洗（但我认为这并不容易发生），告诉我，他将如何离开那里？固然不再为所作的事被追究，然而仍无信心；这是合理的。因为当他活了一百年，没有善行可以展示，只知他没有犯罪，甚至这也没有，只是靠恩宠得救，当他看见别人戴上冠冕，光彩照人，极受赞许：即使他不堕入地狱，告诉我，他能忍受沮丧吗？\par

为了使事情清楚，举一个例子。假设有两个士兵，其中一个偷窃、伤害、欺诈，而另一个没有做这些事，却尽显勇士本色，善做大事，在战争中树立战利品，右手沾满鲜血；然后当时候来到，假设（与那窃贼相同的等级中）他立刻被引向皇帝的宝座和紫袍；但假设另一个仍留在原处，仅因皇恩不受惩罚，但他却在最后的位置，被安置在王的下面。告诉我，当他看见与自己同等的人上升到最崇高的尊严，变得如此荣耀，成了世界的主人，而他自己仍在下层，甚至没有被光荣地免除刑罚，只是靠王的恩宠和仁慈，他能忍受失望吗？因为即使王赦免他，免除对他的指控，他仍将生活在羞耻中；当然连别人也不会钦佩他：因为在这样的宽恕中，我们钦佩的不是领受礼物的人，而是施予礼物的人。而且礼物越大，领受者就越羞愧，当他们的过犯很大时。\par

那么，这样的人将用怎样的眼目去看那些在君王殿庭中的人呢？当他们展示自己无数的汗水和伤口时，他却一无所献，只靠天主的慈爱而得救恩本身？因为就如有人为一个将被逮捕的凶手、盗贼、奸夫求情，并命令他留在君王宫殿的门廊处，此后即使他已免于刑罚，也再不能正视任何人的脸：这人的情况也正是如此。\par

因为，我恳求你们，不要以为只因被称为宫殿，所有人都达到同样的境界。因为如果在君王的宫廷里，有总督，并所有在君王身边的人，还有那些地位极低、占据所谓\OriginalTerm{十人团}{Decani}之位的人（虽然总督与\OriginalTerm{十人长}{Decanus}之间的间隔如此之大），那么在上的天廷中岂不更是如此吗？\par

我说这话并非出于自己。因为保罗指出了另一个更大的差异。因为他说：太阳、月亮、星辰与最小的星辰之间有多少差异，天国中的人之间也有多少差异。而且太阳与最小的星之间的差异远大于所谓的\OriginalTerm{十人长}{Decanus}与总督之间的差异，这是显而易见的。因为太阳同时照耀整个世界，使其明亮，并遮掩月亮和星辰，而最小的星甚至常常不出现，即使在黑暗中也不出现。有许多星辰我们看不见。那么，当我们看见别人成为太阳，而我们的地位却是最小甚至看不见的星辰时，我们还有什么安慰呢？\par

我恳求你们，我们不要如此怠惰，不要如此麻木，不要为了安逸的生活而丢弃天主的救恩，而要善用它，并增加它。因为即使一个人是\OriginalTerm{望教者}{Catechumen}，他仍然认识基督，仍然明白信仰，仍然聆听天主的圣言，仍然离知识不远；他知道自己主人的旨意。他为何拖延？为何迟延耽搁？没有什么比善生更好了，无论在此世或来世，无论对\OriginalTerm{已受洗者}{Enlightened}还是望教者而言。\par

11. 请告诉我，我们吩咐过什么沉重的诫命呢？\Quote{有妻子，并要贞洁。}这难道困难吗？怎么？许多人，不仅基督徒，就连外教人也无妻而贞洁地生活。外教人为了虚荣所超越的，你竟连为了敬畏天主也不遵守。\par

\Quote{把你所有的给穷人。}祂说。这难道沉重吗？但在这一点上，外教人也谴责我们，他们仅仅为了虚荣就倾尽所有财产。\par

\Quote{不要使用污秽的言语。}这难道困难吗？因为即使没有命令，我们难道不应该在这事上做正确的事，以避免显得卑鄙吗？因为相反的行为是麻烦的，我指的是使用污秽的言语，这从以下事实显而易见：灵魂感到羞耻并脸红，如果它被引导说出任何此类话语，除非醉酒，否则它不会这样做。因为当坐在公共场所时，即使你在家里这样做，为什么你不在那里做？因为那些在场的人。为什么你不在你妻子面前轻易做同样的事？为了不侮辱她。那么你不做，以免侮辱你的妻子；而你不因侮辱天主而脸红吗？因为祂无所不在，且听见一切。\par

\Quote{不可醉酒，}祂说。因为这事本身，难道不是一种惩罚吗？祂没有说，把你的身体放在刑架上，而是说什么？不要放任它以致剥夺心灵的权威；相反，\BibleQuote{不要为肉体安排去满足私欲}\BibleRef{罗 13:14}。\par

\Quote{不要强夺不属于你的；不要欺诈；不要发虚誓。}祂说。这些事需要什么劳苦！什么流汗！\par

\Quote{不要毁谤任何人，也不要诬告，}祂说。相反的行为确实是劳苦。因为当你说了别人的坏话，你立即处于危险、猜疑之中，心想：\InnerQuote{我所谈论的人听见了吗？}无论他是大人物还是小人物。因为如果他是大人物，你立即确实处于危险之中；但如果是小人物，他会同样报复你，或者更糟的是，他会更厉害地毁谤你。如果我们有意志，就没有什么困难、没有什么沉重的诫命。如果我们没有意志，即使最容易的事在我们看来也是沉重的。有什么比吃饭更容易呢？但由于极度的懦弱，许多人甚至对此感到厌恶，我听到许多人说，甚至连吃饭也是累人的。只要你有了意志，这些事没有一样是累人的。因为在来自天上的恩宠之后，一切取决于意志。让我们愿意善事，以便也能获得永恒的善事，在我们的主耶稣基督内，愿光荣、权能、尊崇归于祂——与圣父及圣神——自今至永远，世世无穷。阿们。\par

\SourceRecord{Translated by Frederic Gardiner.；Nicene and Post-Nicene Fathers, First Series；Vol. 14.；Edited by Philip Schaff.；Buffalo, NY: Christian Literature Publishing Co.,；1889.；<http://www.newadvent.org/fathers/240213.htm>.}

% structure: p0001 source=h1 target=section reason=page-title

\section{讲道第十四篇：论希伯来人书}

% structure: p0002 source=h2 target=subsection reason=relative-heading-rank; base=h2

\subsection{希伯来书 8:1-2}

\Quote{我们所讲论的事，其要旨是这样的：我们有一位这样的大司祭，祂已坐在天上威严宝座的右边；是圣所和真帐幕的服务者，这帐幕是主所支搭的，不是人所支搭的。}\par

\Person{保禄}常将卑微之事与崇高之事混合，始终效法他的师傅，使卑微成为通往崇高的路径，并藉着前者被引向后者；当我们身处伟大事物之中时，我们得知这些（卑微事物）是降卑。他在这里也是如此做的。在宣告了\Quote{祂献上了自己}，并表明祂是一位\Quote{大司祭}之后，他说了什么？\Quote{我们所讲论的事，其要旨是这样的：我们有一位这样的大司祭，祂已坐在威荣宝座的右边。}而这并非司铎的职务，乃是司铎所应服务的那一位的职务。\par

\Quote{圣所的服务者}，不单是服务者，而是\Quote{圣所的服务者，又是那真帐幕的服务者，这帐幕是主所支搭的，不是人所支搭的。}你们看这降卑。他岂不是在前面作过区分，说：\Quote{众位岂不都是服役的神灵吗？}\BibleRef{希 1:14}，因此（他说）并没有对他们说：\Quote{你坐在我的右边}\BibleRef{希 1:13}，因为坐着的不是服务者。那么，这里怎么说\Quote{服务者}和\Quote{圣所的服务者}呢？因为他在这里指的是帐幕。\par

看他是如何提升信主的犹太人的思想的。因为他们很容易想象我们并没有像他们那样的帐幕，看这里（他说）是大司祭，伟大的，是的，比那位更伟大，并且献上了更奇妙的祭献。但这岂不都是空谈吗？岂不都是夸口，仅为赢得我们的心而说的吗？为此，他首先从誓言确立了这一点，随后又从\Quote{帐幕}。因为这一区别也是明显的：但宗徒还想到另一个帐幕，\Quote{即}（他说）主所支搭的（或\Quote{所建立的}），不是人所支搭的。那些说天穹转动的人在哪里？那些宣称它是球形的人在哪里？因为这两种观念都被推翻了。\par

\Quote{现在}（他说）\Quote{我们所讲论的事，其要旨是这样的。}论到\Quote{要旨}，总是指最重要的部分。他再次降低他的论述；在说了崇高之事以后，从此他便无所畏惧地讲话。\par

% structure: p0008 source=h2 target=subsection reason=relative-heading-rank; base=h2

\subsection{希伯来书 8:3-5}

2. 其次，为使你明白他用\Quote{服务者}一词是指人性，请看他是如何再次表明的：\Quote{因为}\BibleRef{希 8:3}（他说）\Quote{凡大司祭都是被设立来奉献礼品和祭献的，因此这人也必须有所奉献。}\par

现在，不要因为你听见祂坐着，就以为祂被称为大司祭是空谈。因为前者，即祂的坐姿，属于天主性的尊威；而后者则属于祂极大的慈爱和对我们温柔的关怀。为此，他反复强调这一点，并更多停留于此：因为他担心另一个真理会推翻它。因此他再次将论述降至此点：因为有人探问祂为什么死了。祂是一位司铎。但没有司铎不献祭的。因此，祂也必须有一祭献。\par

另一方面；在说了祂在高处之后，他从各方面肯定并证明祂是一位司铎：从默基瑟德、从誓言、从奉献祭献。他也由此构建了另一个必要的三段论。\Quote{因为如果}（他说）\Quote{祂在地上，祂就不是司铎，因为已有按法律奉献礼品的司铎。}如果祂是司铎（正如祂确实是），我们必须为祂另寻一个地方。\Quote{因为如果}祂真的\Quote{在地上，祂就不应当是司铎。}又怎能呢？祂没有奉献祭献，没有执行司铎之职。而且理由充分，因为已有司铎。此外，他指出，祂在地上成为司铎是不可能的。又怎能呢？他是说，不可反对（被指定的司铎）。\par

3. 在此我们必须专心留意，并思考宗徒的智慧；因为他再次显明了司铎职的差别。\Quote{他们}（他说）\Quote{事奉天上事物的模型和阴影。}\par

他在这里所说的天上事物是什么？是指属灵的事物。因为虽然它们在地上完成，却仍然配得上诸天。因为当我们的主耶稣基督被宰杀（作为祭献）时，当圣神与我们同在时，当那坐在圣父右边的那位在这里时，当藉着圣洗成为儿子时，当他们成为天上居民的同胞时，当我们在那里有家乡、城邦和公民身份时，当我们对这里的事物是异乡人时，这一切怎能不是\Quote{天上事物}呢？然而什么！我们的圣歌难道不是天上的吗？我们这些在地上的人难道不是与他们一同发出那些无形体权能的神圣歌咏团在上所唱的同一赞美吗？祭台难道不是天上的吗？如何呢？它没有任何属肉体的，一切属灵的事物都成为祭品。祭献不会化为灰烬、烟或蒸汽的香气，它使放在那里的物品光明灿烂。再说，我们所举行的礼仪怎能不是天上的呢？因为当祂说：\Quote{你们赦免谁的罪，就给谁赦免；你们留存谁的罪，就给谁留存}\BibleRef{若 20:23}，当他们拥有天堂钥匙时，这一切怎能不是天上的呢？\par

\BibleQuote{他们}（他说）\BibleQuote{所事奉的乃是天上之事的模型和影子，正如梅瑟将要建造帐幕时，曾受天主的训示说：\InnerQuote{你们要小心，要照山上指示你的样式，制造一切。}} 既然我们的听觉不如视觉那样敏捷（因为我们所听到的，不像亲眼所见那样能储存在灵魂里），天主便向他显示了一切。所以，他要么就是指\Quote{模型和影子}，要么就是指圣殿而言。因为他接着说：\BibleQuote{看}（他的话是）\BibleQuote{你要照山上指示你的样式制造一切}。那么，他所见的仅仅是圣殿的器具吗？还是也涉及祭献和所有其他事物呢？不，即使这样说也不算错：因为教会是属天的，不外乎就是天堂。\par

% structure: p0015 source=h2 target=subsection reason=relative-heading-rank; base=h2

\subsection{希伯来书 8:6}

4. \BibleQuote{但现在祂已获得了更卓越的职分，因此祂也就是更美盟约的中保。} 你看（他的意思是），一种职分若是模型和预像，另一种是真理〔现实〕，那么前者比后者差多少？但这并没有使听众受益，也没有鼓励他们。因此他说了尤其鼓励他们的话：\BibleQuote{这盟约原是建立在更美的恩许上。} 他通过谈论地点、司祭和祭献来提升他们，然后也列出了盟约的广泛差异，先前还说过它是\BibleQuote{软弱无益的}。\BibleRef{希 7:18}\par

请注意他在打算指责它时，设置了何等保障。因为在前面，他先说了\BibleQuote{按照不可消逝的生命的能力} \BibleRef{希 7:16}，然后说\BibleQuote{先前那条诫命就被废除了} \BibleRef{希 7:18}，之后他又说出了某件伟大的事：\BibleQuote{借此我们得以亲近天主}。\BibleRef{希 7:19} 在此处，他先带领我们升到天上，并表明我们拥有天堂而非圣殿，且那些事物是我们的预像，借此抬高了〔新盟约的〕职分，然后才恰当地抬高司祭职。\par

% structure: p0018 source=h2 target=subsection reason=relative-heading-rank; base=h2

\subsection{希伯来书 8:7-9}

但是（如我所说），他写下了尤其鼓励他们的话：\BibleQuote{这盟约原是建立在更美的恩许上。} 从何显现呢？因为前者被废弃，后者取而代之：它之所以有效，正是因为更美。正如他说：\BibleQuote{如果藉着它得以成全}它，\BibleQuote{那么何必还需要另一位按照默基瑟德品位兴起的司祭呢？} \BibleRef{希 7:11}；同样在此处，他使用了相同的三段论，说 \BibleRef{希 8:7} \BibleQuote{因为如果第一个盟约没有瑕疵，那么就不会为第二个寻求余地了}；也就是说，如果它使人\BibleQuote{没有瑕疵}。因为他正是为此而说话，他没有说\BibleQuote{指责它}，而是\EditorialNote{8, 9节} \BibleQuote{指责他们，主说：看，时日将到，我要与以色列家和犹大家订立新约，不像我拉着他们祖先的手，领他们出离埃及地的那天，与他们所立的约；因为他们没有持守我的盟约，我也就不顾念他们了——主说。}\par

是的，确实如此。从何处可见〔第一个盟约〕终结了呢？他固然也从司祭职上显示了，但现在他用明确的话语更清楚地表明它已被废弃。\par

但为何是\BibleQuote{建立在更美的恩许上}？因为，请告诉我，如何能拿地土与天堂相比？但请你思考，他在那里〔即第一个盟约〕也提到了恩许，免得你责备它。因为那里他也说\BibleQuote{一个更美的希望，借此我们得以亲近天主}\BibleRef{希 7:19}，显示那里也有一个希望；而在此处是\BibleQuote{更美的恩许}，暗示那里祂也曾做出恩许。\par

但是，由于他们不断提出异议，他说：\BibleQuote{看！时日将到，主说，我要与以色列家和犹大家订立新约。} 他并非谈论某个旧约：为了阻止他们声称这一点，他也确定了时间。因此，他没有简单地说\BibleQuote{像我与他们祖先所立的约}，免得你会说〔是〕与亚巴郎所立的，或是与诺厄所立的：而是说明〔是哪个约〕，\BibleQuote{不像我拉着他们祖先的手，领他们出离埃及地的那天，与他们所立的约；因为他们没有持守我的盟约，我也就不顾念他们了——主说。} 你看，恶事首先从我们自己开始（\Quote{他们}自己先是，他说，没有持守〔盟约〕），疏忽来自我们自己，而善事来自祂；我是指恩惠之举。他在这里如同辩护，显示祂离弃他们的原因。\par

% structure: p0023 source=h2 target=subsection reason=relative-heading-rank; base=h2

\subsection{希伯来书 8:10}

5. \BibleRef{希 8:10} \BibleQuote{因为，}他说，\BibleQuote{这是主说，在那些日子之后，我要与以色列家所立的约：我要将我的律法放在他们的思想里，写在他们心上；我要做他们的天主，他们要做我的子民。} 祂这样说关于新〔约〕，因为祂的话是\BibleQuote{不像我所立的约。}\par

但除此之外还有别的区别吗？若有人说\Quote{区别不在此，而在乎它被放在他们心里；他没有提到法规上的任何区别，而是指出了赐予的方式：因为不再}（他说）\Quote{盟约写在石版上，而是写在心版上}，那么请犹太人指出这事曾实现过；但他不能，因为在从巴比伦归来之后，盟约又第二次写在石版上。但我指出，宗徒们没有领受任何文字，而是藉着圣神在心裡领受了。因此基督也说：\Quote{当祂来到时，祂必将一切提醒你们，并且祂会教训你们。}\BibleRef{若 14:26}\par

% structure: p0026 source=h2 target=subsection reason=relative-heading-rank; base=h2

\subsection{希伯来书 8:11-12}

6. \Quote{他们不必再教导}（他说）\Quote{各人自己的邻人，各人自己的兄弟说：该认识主，因为从最小的到最大的，他们都要认识我。因为我将宽恕他们的不义，不再记念他们的罪恶和过犯。}看，还有另一个标记。\Quote{从最小的到最大的}（他说）\Quote{他们都要认识我，他们也不必说：该认识主。}这事除了现在，何时曾应验？因为我们的信仰是显明的，但他们的（即犹太人的）并不显明，只被封闭在角落里。\par

盟约之所以称为\Quote{新的}，是指它既不同，又显示出比旧的优越。\Quote{诚然}，有人说，\Quote{它也是新的，当它一部分被除去，一部分未被除去。例如，当一座旧房子快要倒塌时，若有人撇开整体，修补了地基，我们立刻说他使之成为新的；因为取走了某些部分，用别的部分替换。甚至天也因此被称为\Quote{新的}，当它不再是\Quote{铜的}，而是降下雨水时；地同样也是新的，当它不再贫瘠，而非改变时；房子也是新的，当它一部分被取走，一部分存留。同样，他说，他恰当称之为\Quote{新盟约}。}\par

那么，若我指出那盟约在这点上已变得\Quote{旧了}，即它不结果实？为使你确知此事，请读读\Person{哈盖}、\Person{匝加利亚}和\Person{玛拉基亚}所说的，那时从充军归回尚未完全实现；以及\Person{厄斯德拉}所责备的。那么他（人民）是如何接待他的？为何无人求问主？因为司祭们自己也违背（法律），甚至他们自己也不知晓。你看，你的解释如何被推翻，而我却维持自己的：即这盟约按本意应称为\Quote{新的}？\par

此外，我不承认\Quote{天将是新的}（\BibleRef{依 65:17}）这句话是针对此说的。因为当他在申命纪中说\Quote{天将如铜}时，为何他没有在对比的经文中写下\Quote{但若你们听从，它将是新的}？\par

并且为此，他说祂将赐予\Quote{另一个盟约，因为他们没有恒心守第一个。}我以他的话证明这点：\Quote{因为法律所不能行的，因它由于肉身而软弱，}（\BibleRef{罗 8:3}）以及\Quote{你们为什么试探天主，要把重轭放在门徒的颈上，连我们的祖先和我们自己都不能负荷？}（\BibleRef{宗 15:10}）但他说\Quote{他们没有恒心守在里面。}\par

此处他指出，他使我们堪当更大的、属神的特恩：因为经上说：\Quote{他们的声音传遍普世，他们的言语传到地极。}（\BibleRef{咏 19:5}）这就是\Quote{他们不必各人对自己的邻人说：该认识主}的意思。又说：\Quote{大地将充满对主的认识，如同海水覆盖海洋。}（\BibleRef{依 11:9}）\par

% structure: p0033 source=h2 target=subsection reason=relative-heading-rank; base=h2

\subsection{希伯来书 8:13}

7. \Quote{称它是新的}（他说）\Quote{他使先前成为旧的；但那衰残和变旧的，已临近消失。}请看那隐藏的，他如何揭示了先知的深意！他尊重法律，不愿明确称它为\Quote{旧的}；但尽管如此，他仍称它为旧的。因为若先前的是新的，他就不会称这后来的也为\Quote{新的}。因此，由于承认了更多和不同的东西，他宣告\Quote{它已变旧了}。故此它已被废止、在消逝中，不再存在。\par

他从先知那里获得胆量，更适宜地攻击它，表明我们的制度现在正兴旺。就是说，他表明那另一制度是旧的；然后取起\Quote{旧的}一词，自己又附加了另一情况，即老年的特征，他补充了别人所遗漏的，并说\Quote{临近消失}。\par

那么，新的并非单纯使旧的停止，而是因为它已衰老，不再有用。为此他说：\Quote{因为它的软弱和无益}（\BibleRef{希 7:18}），\Quote{法律不能使什么成为完美的}（\BibleRef{希 7:19}），以及\Quote{若第一个盟约是无缺陷的，就不必寻求第二个了。}（\BibleRef{希 8:7}）\Quote{无缺陷的}意即有用的；并非说那旧盟约有什么可指责之处，而是说它不足。他使用了一种通俗的措辞。如同有人说：这房子不是无缺陷的，即它有些毛病、朽坏了；这衣服不是无缺陷的，即它正在破裂。因此他此处并非说它是邪恶的，只是说它有些缺陷和不足。\par

8. 因此，我们也是新的，或者更确切地说，我们已被更新，但现在却变旧了；所以我们就\Quote{接近消亡}，趋于毁灭。让我们除去这种老朽。通过圣洗固然不再可能，但在此世通过悔改却是可能的。如果我们身上有任何旧的东西，让我们抛弃它；如果有任何\Quote{皱纹}，任何污点，任何\Quote{瑕疵}，让我们洗掉它，成为纯洁的\BibleRef{弗 5:27}：使\BibleQuote{君王恋慕我们的美丽}\BibleRef{咏 44:11}。\par

即使那陷入极度丑陋的人，也能恢复那样的美丽，关于这美丽达味曾说，君王必恋慕你的美丽。\BibleQuote{女儿，请听，请侧耳细听；忘记你本族和你父家：如此君王必恋慕你的美丽。}\BibleRef{咏 45:10}然而忘记并不产生美丽。是的，美丽是属于灵魂的。这是什么样的忘记？是忘记罪恶。因为他所指的是来自外邦人的教会，劝戒她不要记念她父辈的事，即那些向偶像献祭的人；因为教会正是从这样的人中被聚集起来的。\par

他没有说：\Quote{不要追随他们}，而是更进一步说：不要让他们进入你的思想；他在另一处也说：\BibleQuote{我口中不提起他们的名字。}\BibleRef{咏 15:4}又说：\BibleQuote{愿我的口不再谈论人的事。}\BibleRef{咏 17:3}这还算不上什么大德行；不，其实这德行很大，但还不及此处所论及的。因为他在那里说什么？他不是说：\Quote{不要谈论人的事，也不要讲论你们父辈的事}；而是说：连记念也不要，不要让他们进入你的思想。你们看，他愿意我们与邪恶保持多么远的距离。因为不记念某事的人就不会去想它；不想它的人就不会去说它；不说它的人就不会去做它。你们看，他从多少条路上把我们围堵起来？他用多么大的间隔把我们移开，甚至移到了极远之处？\par

9. 那么，让我们也\Quote{听从并忘记}我们自己的邪恶。我不是说我们的罪，因为祂说：\BibleQuote{你先承认，然后我就不再记念。}\BibleRef{依 43:26}我的意思是，比如，让我们不再记念贪婪，甚至归还先前所掠夺的。这就是忘记邪恶，驱散贪婪的思想，任何时候都不再接纳它，并且也要抹去已经做错的事。\par

对邪恶的遗忘怎样能临到我们？从对善事的记念，从对天主的记念。如果我们不断记念天主，我们就不能同时也记念那些事。因为他说：\BibleQuote{我在床上记念了祢，在清晨的曙光中思念了祢。}\BibleRef{咏 62:6}因此，我们应当常将天主置于记念中，但特别是在心思不受打扰的时候，当人藉着这种记念能自责、能记住事情的时候。因为在白天，即使我们记念，其他的牵挂和烦恼进入，又把那思想驱散了；但在夜里，当灵魂平静安息时，当它处于港弯和晴空下时，就能不断记念。\BibleQuote{你们在床上要心中思想，并要哀伤。}\BibleRef{咏 4:4}确实，白天也应当保持这种记念。但既然你们总是充满牵挂，在今生的事物中分心，那么至少要在床上记念天主；在清晨的曙光中默想祂。\par

如果在清晨我们默想这些事，我们将以极大的安全进入我们的事务。如果我们首先藉着祈祷和恳求使天主仁慈，这样出去，我们就不会有敌人。即使有，你也会嘲笑他，因为天主是仁慈的。市场上充满战争；每天的事务是一场战斗，是一场暴风雨和风暴。因此我们需要武器；祈祷是一件伟大的武器。我们需要顺风；我们需要学习一切，以便在一整天中毫无沉船、毫无创伤地度过。因为每一天都有许多岩石，船常常触礁沉没。因此，我们特别需要在清早和夜晚祈祷。\par

10. 你们中有许多人常观看奥林匹克运动会，不仅观看，而且热心作斗士的拥护者和崇拜者，有人支持这位【斗士】，有人支持那位。你们知道，在比赛的那些日子和那些夜晚，传令官整夜不思考别的，也没有别的忧虑，只担心斗士出场时不要丢脸。因为那些坐在号手旁边的人告诫他不要与任何人说话，免得他浪费气息而被人嘲笑。因此，如果那即将在人前奋斗的人如此事先考虑，那么我们更应当不断深思熟虑、谨慎小心，因为我们整个生命就是一场竞赛。那么，让每个夜晚成为守夜，让我们小心，当我们白天出去时不要使自己成为笑柄。但愿只是成为笑柄。但现在竞赛的审判者坐在圣父的右边，仔细倾听，以免我们发错任何一个音符，发出任何不和谐的声音。因为他不只是行为的审判者，也是言语的审判者。让我们保持警醒，亲爱的；我们也有那些渴望我们成功的人，如果我们愿意的话。我们每人身边都坐着天使；但我们却整夜打鼾。但愿只是这样。但许多人甚至做了许多放荡的事，有些人甚至直接去了妓院，另一些人则把妓女带回家，使自己的房屋成为卖淫之所。是的，千真万确。难道不是吗？他们真关心自己的竞赛啊。另一些人醉酒胡言；另一些人喧闹。另一些人整夜邪恶地守夜，编织着比睡觉的人更坏的欺骗诡计；另一些人计算高利贷；另一些人用忧虑折磨自己，做任何事，就是不做适合竞赛的事。因此，我劝告你们，让我们放下一切【其他】事物，只关注一件事：我们如何获得奖赏，【我们如何】戴上桂冠；让我们做一切使我们能够获得所应许的福分的事。愿我们众人都在我们的主耶稣基督内获得这些，愿光荣归于他、父及圣神，权能、尊威，从现在直到永远，永世无尽。阿们。\par

\SourceRecord{Translated by Frederic Gardiner.；Nicene and Post-Nicene Fathers, First Series；Vol. 14.；Edited by Philip Schaff.；Buffalo, NY: Christian Literature Publishing Co.,；1889.；<http://www.newadvent.org/fathers/240214.htm>.}

% structure: p0001 source=h1 target=section reason=page-title

\section{关于\BibleBook{}{希伯来书}第十五篇讲道}

% structure: p0002 source=h2 target=subsection reason=relative-heading-rank; base=h2

\subsection{希 9:1-5}

\Quote{ \Emph{如此，第一个盟约确实也有敬礼的规条和世俗的圣所。因为制成了一个帐幕：第一间，其中设有灯台、供桌和供饼，称为圣所；在第二层帐幔后面，有一间称为至圣所，设有金香炉和全包金的约柜，柜内有盛玛纳的金罐、亚郎发芽的棍杖和约版，其上有光荣的革鲁宾，以赎罪盖为阴影：关于这些，现在不能一一详述。} }\par

1. 他已从司铎、司铎职和盟约表明，那一个制度当有结局。现在他从帐幕本身的样式表明这一点。如何？他说，这一个是\Quote{圣所}，那一个是\Quote{至圣所}。那么圣所是前一个时期的象征（因为那里一切都是藉祭献完成的）；而至圣所则是现在这个时期的象征。\par

他用以指至圣所的是天堂，以帐幔指天堂，并且基督的血肉\BibleQuote{进入帐幔内}：也就是说，\BibleQuote{藉祂血肉的帐幔}。\BibleRef{希 6:19}\par

最好从开端来讨论这段经文。那么，他说什么？\Quote{如此，第一个也有}——第一个什么？\Quote{盟约}——\Quote{敬礼的规条。}什么是\Quote{规条}？象征或礼仪。然后（他意指）现在已没有。他表明它已让位了，因为他说，那时它\InnerQuote{有}；所以现在，它虽存立，却已不是。\par

\Quote{和世俗的圣所。}他称之为\Quote{世俗的}，因为允许所有人踏足其间，且在同一殿中，司铎、犹太人、归依者、希腊人、纳齐尔人所站之处是明显的。既然连外邦人也允许踏足，他便称之为\Quote{世俗的}。因为犹太人当然不是\Quote{世界}。\par

\Quote{因为}（他说）\Quote{制成了一个帐幕：第一间，称为圣所，其中设有灯台、供桌和供饼。}这些是世界的象征。\par

\Quote{在第二层帐幔后面}——那时不只一层帐幔，外面也有帐幔——\Quote{有一间帐幕，称为至圣所。}注意他处处称它为帐幕，因为天主在那里安营。\par

\Quote{其中设有}（他说）\Quote{金香炉和全包金的约柜，柜内有盛玛纳的金罐、亚郎发芽的棍杖和约版。}所有这些都是犹太人心硬的可敬而显著的纪念物：\Quote{和约版}（因为他们打碎了它们），\Quote{和玛纳}（因为他们抱怨过；为此，祂为后代保留其记忆，命令把它放在金罐里），\Quote{和亚郎发芽的棍杖。其上有光荣的革鲁宾。}什么是\Quote{光荣的革鲁宾}？他要么指\Quote{光荣的}，要么指那些在天主下面的。\Quote{以赎罪盖为阴影。}\par

此外，在另一层意义上，他在讲论中夸耀这些事，以表明后来的事更伟大。\Quote{关于这些}（他说）\Quote{现在不能一一详述。}这些话暗示这些不只是所见之物，而是某种谜语。\Quote{关于这些}（他说）\Quote{现在不能一一详述}，或许因为它们需要长篇论述。\par

% structure: p0012 source=h2 target=subsection reason=relative-heading-rank; base=h2

\subsection{希 9:6}

2. \Quote{这些事既如此安排，司铎们就常进第一间帐幕，尽敬礼之职。}这就是说，这些事固然存在，但犹太人并未享用它们：他们看不见。所以这些事并不比那些为他们预言的更属于他们。\par

% structure: p0014 source=h2 target=subsection reason=relative-heading-rank; base=h2

\subsection{希 9:7}

\BibleRef{希 9:7}\BibleQuote{但第二间帐幕，大司铎独自一年一次进去，没有不带血的，他为自己和百姓的错误献上这血。}你看，预像早已确立了吗？免得他们问：\Quote{怎么只有一次祭献？}他表明从起初就是这样，因为至少更神圣可畏的祭献只有一次。并且大司铎如何一次而永远献祭？他们从起初就是这样做的，因为那时（他说）大司铎也\Quote{一次而永远}献祭。\par

他说\Quote{没有不带血的}很恰当。（当然不是没有血，却非这血，因为事情没有那么重大。）他指明将有一个祭献，并非被火焚烧，而是以血为特征。既然他称十字架为祭献，虽无火，无木柴，也非多次奉献，却一次在血中奉献；他便表明古代的祭献也属此类，一次在血中奉献。\par

\Quote{他为自己和百姓的错误献上。}他未说\Quote{罪}，而说\Quote{错误}，免得他们傲慢。因为即使你没有故意犯罪，却无意中有了错误，没有人能免于此。\par

并且处处他加上\Quote{为自己}，表明基督更伟大。因为若祂与我们的罪分开，祂如何\Quote{为自己奉献}？那么，为什么你说了这些（有人说）？因为这是优越者的标记。\par

% structure: p0019 source=h2 target=subsection reason=relative-heading-rank; base=h2

\subsection{希 9:8}

3. 至此尚无思辨。但从这里他开始思考，说：\BibleRef{希 9:8} \Quote{圣神指明这一点：当首一帐幕尚且存在时，进入至圣所之路尚未显露。}为此，（他说）这些事就这样被\Quote{规定}，好叫我们得知\Quote{至圣所}，即天堂，至今仍不可进入。那么，我们不可因没有进入就以为它们不存在；正如我们也没有进入至圣所一样。\par

% structure: p0021 source=h2 target=subsection reason=relative-heading-rank; base=h2

\subsection{希 9:9}

\Quote{这}（他说）\Quote{是为当时的时代所立的预象。}他说的\Quote{当时的时代}是什么意思？即基督来临之前：因为基督来临之后，它就不再是\Quote{当时的时代}了：因为它既然已经来到并结束，又怎能是当时呢？\par

他还指出另一件事，当他说到：\Quote{这曾是为当时的时代所立的预象}，即成了预像。\Quote{在其中献上礼物和祭品，却不能按良心使行敬拜者完全。}你现在看到了\BibleRef{希 7:19}\Quote{法律不能使人完全}和\BibleRef{希 8:7}\Quote{若那前一个盟约没有瑕疵}是什么意思？怎么说的？\Quote{按良心而言。}因为祭品不能除掉灵魂的污秽，仍只涉及身体：\Quote{按照肉体的条例。} \BibleRef{希 7:16} 因为的确它们不能除掉奸淫、谋杀或亵渎。你明白了吗？你吃了这个，你没吃那个，这些是无关紧要的事。[\Quote{所立之物}]只在于饮食和各样的洗涤。\Quote{你喝了这个，}他说：然而关于喝并未有任何规定，但他这样说，是视之为微不足道的事。\par

% structure: p0024 source=h2 target=subsection reason=relative-heading-rank; base=h2

\subsection{希 9:10}

\Quote{以及种种洗涤和肉体的规条，加在他们身上，直到改革的时候。}因为这是肉体的义。在这里他贬低祭品，表明它们无效，并且它们存在是\Quote{直到改革的时候}，即它们等候那改革万物的时期。\par

% structure: p0026 source=h2 target=subsection reason=relative-heading-rank; base=h2

\subsection{希 9:11}

4. \Quote{但基督已经来到，为已实现的福乐的大司祭，经过那更大更全备的帐幕，不是人手所造的。}这里他指的是肉身。他说的好，\Quote{更大更全备}，因为天主圣言和圣神的一切权能居住其中；\Quote{因为天主不限定地赐予圣神（给祂）。} \BibleRef{若 3:34} 而\Quote{更全备}，因为它既无可指责，又能纠正更大的事。\par

\Quote{就是说，不是出于受造的。}你看它如何是\Quote{更大}。因为若由人建造，就不会是\Quote{出于圣神} \BibleRef{玛 1:20}。它也不是\Quote{出于受造的}；就是说，不是出于这些受造之物，而是出于圣神，出于圣神的。\par

你看见他如何称呼身体为帐幕、幔子和天堂吗？\Quote{经由那更大更全备的帐幕，通过幔子，即祂的肉身。} \BibleRef{希 10:20} 又说，\Quote{进入幔内。} \BibleRef{希 6:19} 又说，\Quote{进入至圣所，来显现在天主面前。} \BibleRef{希 9:24} 那么他为何这样说？乃因所象征的对象不同。我的意思是，天堂是一道幔子，因为它像幔子一样隔开了至圣所；肉身是隐藏天主性的幔子；帐幕也容纳天主性。同样，天堂是帐幕：因为司祭在那里，在里面。\par

\Quote{但基督}（他说）\Quote{已经来到为大司祭}：他没有说\Quote{成了}，而是\Quote{已经来到}，也就是说，为这目的而来，并非继承他人。他不是先来而后成为大司祭，而是来与成为同时。他没有说\Quote{已经来到为大司祭}论及牺牲之物，而是\Quote{为已实现的福乐}，仿佛他的言谈无力将整体摆在我们面前。\par

% structure: p0031 source=h2 target=subsection reason=relative-heading-rank; base=h2

\subsection{希 9:12}

\Quote{不是借着山羊和牛犊的血}，他说，\Quote{而是借着祂自己的血}，（他说）\Quote{祂一次进入圣所。}你看，他这样称呼天堂。\Quote{一次}（他说）\Quote{祂进入圣所，获得了永恒的救赎。}而\Quote{获得了}这个表述，表达了极其困难、超乎意料的事，即一次进入，祂\Quote{获得了永久的救赎}。\par

% structure: p0033 source=h2 target=subsection reason=relative-heading-rank; base=h2

\subsection{希 9:13-14}

5. 接下来是劝服的内容。\par

\Quote{因为倘若公牛和山羊的血，以及母牛的灰烬，洒在不洁的人身上，圣化而洁净肉体；那么基督的血，祂借着圣神将自己无瑕疵地奉献给天主，岂不更能洁净你们的良心，除去死行，以事奉生活的天主？}\par

因为（他说）如果\Quote{公牛的血}能洁净肉体，基督的血岂不更能洗去灵魂的污秽。为使你听到\Quote{圣化}一词时不以为是什么大事，他指明并显示这些洁净之间的区别，一种高，一种低。他说这很有道理，因为那是\Quote{公牛的血}，而这是\Quote{基督的血}。\par

他不但满足于称号，还陈明了奉献的方式。\Quote{祂}（他说）\Quote{借着圣神将自己无瑕疵地奉献给天主。}就是说，祭牲毫无瑕疵，无罪纯洁。因为\Quote{借着圣神}是这个意思，不是借着火，也不是借着其他事物。\par

\Quote{洁净你们的良心}（他说）\Quote{除去死行}。 他说的\Quote{除去死行}很好；若有人触碰死尸，他就被污染；同样，若有人触碰\Quote{死行}，他就因良心而被玷污。\Quote{事奉}（他说）\Quote{永生真实的天主}。在此他声明，当一个人有\Quote{死行事奉永生真实的天主}时，就不可能事奉永生真实的天主，因为它们是既死又假的；他说得很有道理。\par

6. 因此，没有人可以带着\Quote{死行}进入这里。因为若碰触死尸的人不可进入，何况拥有\Quote{死行}的人呢？因为这是最严重的污染。\Quote{死行}是指所有没有生命、散发恶臭的事物。正如死尸对任何感官都无益，反而令靠近的人厌恶，同样，罪也立刻打击理性官能，不让理智本身平静，而是扰乱并折磨它。\par

也有人说，瘟疫一开始就腐蚀活体；罪也是如此。它与瘟疫毫无区别，不是先污染空气然后污染身体，而是直接刺入灵魂。难道你没有看见患瘟疫的人如何发热、如何辗转反侧、如何充满臭气、面容如何扭曲、如何全然不洁吗？犯罪的人也是如此，虽然他们看不见。因为，请告诉我，被贪财或肉欲支配的人，岂不比任何发烧的人更糟糕？当他行事并屈服于一切无耻之事时，岂不比这些都更不洁？\par

7. 有什么比贪爱金钱的人更卑贱呢？妓女或舞台上女子所不肯做的事，他也无不做。相反，很可能她们倒会拒绝做某事，而他却不会。他甚至屈尊做适合奴隶的事，谄媚那些他不该谄媚的人；又在不该倨傲的地方倨傲，在各方面自相矛盾。他会坐在那里谄媚恶人，常常是比自己贫穷和卑贱得多的堕落老人；却对善良和各方面有德的人傲慢无礼。你在两方面都看到卑贱和无耻：他既过分谦卑，又自夸。\par

然而妓女站在自己的房子前，她们的罪名是用身体换钱：但也许有人说，贫穷和饥饿迫使她们这样做（尽管这最多不是充分的借口，因为她们本可以靠工作谋生）。但贪财的人站着，不是站在自己房前，而是站在城中央，将不是自己的身体而是灵魂出卖给魔鬼；以致魔鬼与他为伴，进入他里面，如同进入妓女一般；满足了一切情欲后就离开；全城的人都看见，不只是两三个人。\par

这又是妓女的特性：她们属于给金钱的人。即使他是一个奴隶或角斗士，或任何什么人，只要他付给她们的报酬，她们就接待他。但自由人，即使比所有人都高贵，没有钱她们也不接纳。这些人也做同样的事。当思想不带钱时，他们便拒绝；但为了金钱，他们与可憎的人交往，甚至与斗兽的人交往，无耻地与他们结伴，毁坏灵魂的美。因为正如那些妇女天生面容可憎、黝黑、笨拙、粗野、无形、畸形，各方面令人厌恶，这些人的灵魂也变得如此，无法用外在的脂粉掩饰其畸形。因为当丑陋到了极点，无论他们怎样设计，也无法伪装成功。\par

因为那无耻造成了妓女，请听先知所说：\BibleQuote{你向众人无耻；你有妓女的面容。} \BibleRef{耶 3:3} 这话也可对贪财的人说：\Quote{你向众人无耻，}不是向这些或那些人，而是\Quote{向众人。}如何呢？这样的人不尊敬父亲、儿子、妻子、朋友、兄弟、恩人，或绝对任何人。我为什么说朋友、兄弟和父亲呢？他不尊敬天主自己，而我们所相信的一切在他看来都是寓言；他大笑，被自己的巨大情欲所陶醉，甚至不让任何可能有益于他的话进入耳中。\par

啊！他们的荒谬！然后他们说什么！\Quote{祸哉，你，玛门，以及没有你的人。} 对此我因愤怒而撕心裂肺：因为说这些话的人有祸了，尽管他们说时是玩笑。因为请告诉我，天主不是说了这样的威胁吗？祂说：\BibleQuote{你们不能服事两个主人。} \BibleRef{玛 6:24} 你竟无视这威胁吗？保禄不是说这是崇拜偶像，并称\BibleQuote{贪财的人是崇拜偶像者}吗？\BibleRef{弗 5:5}\par

8. 而你站着发笑，效法世上舞台上的女子那样引人发笑。这已倾覆了一切，打倒了一切。我们的事务，无论是我们的生意还是礼貌，都变成了笑谈；没有什么是稳定的，没有什么是严肃的。我说这些不只是对世俗人说的；但我知道我暗示的是谁。因为教会已充满了笑声。无论谁说了一句巧妙的话，在场的人立刻发笑；奇妙的是，许多人在祈祷时也笑个不停。\par

处处都是魔鬼在领舞，它已侵入一切，主宰一切。基督受辱，被弃置一旁；教会无人重视。你们没有听见保禄说：\Quote{污秽的言语、愚妄的戏笑}要从你们中间除掉吗？\BibleRef{弗 5:4} 他把\Quote{戏笑}与\Quote{污秽}并列，而你们竟发笑？什么是\Quote{愚妄的言语}？就是毫无益处的话。你，一个隐修士，竟发笑并放松面容吗？你这被钉十字架的人？你这哀恸的人？告诉我，你发笑吗？你哪里听说过基督这样做？没有的事：倒是祂常怀忧伤。因为当祂望着耶路撒冷时，祂哭了；当祂想起叛徒时，祂忧愁；当祂将要复活拉匝禄时，祂哭了；而你还发笑？若那不为他人之罪忧伤的人尚该受责，那对自己的罪毫无悲伤反而嘲笑的人，又当如何呢？这是哀恸和苦难、折磨和制服肉身、争战和出汗的季节，而你竟发笑？你没有看见撒辣是怎样受责的吗？你没有听见基督说：\Quote{你们现今欢笑的人是有祸的，因为你们将要哀恸哭泣}？\BibleRef{路 6:25} 你天天咏唱这些，告诉我，你说什么？\Quote{我笑了？}决不是；而是什么？\Quote{我因呻吟而劳苦。}\BibleRef{咏 6:6}\par

也许有些如此放荡和愚蠢的人，甚至在这番斥责之时仍发笑，因为我们就这样谈论笑。因为他们的错乱、他们的疯狂竟如此之大，以至于感受不到这斥责。\par

天主的司铎正站立着，献上众人的祈祷：而你竟发笑，毫无畏惧？当他战战兢兢地为你献上祈祷时，你竟藐视一切？你没有听见圣经说：\Quote{藐视的人，你们有祸了！}\BibleRef{宗 13:41} 你不战栗吗？你不谦卑吗？即使你进入王宫，你也会在衣着、神态、举止等各方面约束自己；而在这里，有真正的王宫，有如天上般的景象，你却发笑？你，我知道，确实看不见，但你应听：天使无处不在，尤其在主的殿中，他们侍立在君王前，整个殿宇充满这些无形的天使。\par

我这番话也是向妇女们说的：她们在丈夫面前往往不敢随意发笑，即使笑，也并非总是，只是在闲暇之时；但在这里，她们却总是笑。告诉我，女人，你蒙着头在圣堂里发笑吗？你到这里来是为了告解罪过，在天主前俯伏哀求，为你的不幸所犯的过犯恳求，而你是带着笑做这一切的吗？那么你怎能平息祂的怒气呢？\par

9. 但有人会说：笑有什么害处呢？笑本身没有害处；害处在于过度，不合时宜。笑被植入我们内，使我们遇到久别的朋友时可以一笑，见到垂头丧气的人时可以用微笑安慰他们；但并非要我们纵声大笑，总是笑。笑被植入我们的灵魂，是为了让灵魂有时得到舒缓，而非完全松懈。因为肉身的欲望也被植入我们内，但绝不是说因为它被植入，我们就应当使用它，或过度使用它；而应加以节制，不能说因为被植入，我们就该使用它。\par

要以眼泪事奉天主，好能洗净你的罪。我知道许多人讥笑我们，说：\Quote{眼泪马上就来了。}因此现在是流泪的时候。我也知道那些说\Quote{我们吃喝吧，因为明天我们要死了}\BibleRef{格前 15:32}的人厌恶这样做。\Quote{虚而又虚，万事皆虚。}\BibleRef{训 1:2} 不是我说的，而是那曾经历一切的人这样说：\Quote{我建造了房屋，栽种了葡萄园，开凿了水池，拥有仆婢。}\BibleRef{训 2:4} 那么这一切之后呢？\Quote{虚而又虚，万事皆虚。}\BibleRef{训 12:8}\par

因此，亲爱的，让我们哀恸，让我们哀恸，好使我们真正欢笑，在纯粹喜乐的时刻真正欢乐。因为这里的喜乐总是掺杂着悲伤：绝不可能找到纯粹的喜乐。而那喜乐是单纯无欺的：它毫无诡诈，毫无掺杂。让我们在那喜乐中欢愉；让我们追求那喜乐。而获得这喜乐别无他法，唯有在此世选择不是安逸而是有益的事，甘愿受一点苦，凡事感恩。这样我们便能获得天国，愿我们众人皆堪当承受，在基督耶稣我们的主内。愿光荣归于祂及圣父，偕同圣神，今世直到永远，及世之世。阿们。\par

\SourceRecord{Translated by Frederic Gardiner.；Nicene and Post-Nicene Fathers, First Series；Vol. 14.；Edited by Philip Schaff.；Buffalo, NY: Christian Literature Publishing Co.,；1889.；<http://www.newadvent.org/fathers/240215.htm>.}

% structure: p0001 source=h1 target=section reason=page-title

\section{《希伯来书》第十六篇讲道}

% structure: p0002 source=h2 target=subsection reason=relative-heading-rank; base=h2

\subsection{希 9:15-18}

\BibleQuote{\Emph{为此，祂作了新盟约的中保，以祂的死亡救赎了在先前的盟约下所犯的过犯，使那些蒙召的人能领受所应许的永恒产业。} \Emph{因为凡有遗嘱，必须立遗嘱者死亡才生效；} \Emph{人活着时，遗嘱无效，} \Emph{所以} \Emph{前一个盟约也不是不用血而成立的。}}\par

1. 很可能许多较软弱的人，因基督已死而特别怀疑祂的许诺。因此保禄从世俗习惯中引出这个比喻，且是出于丰富。他说：\Quote{的确，正因如此，我们应当鼓起勇气。} 为何呢？因为遗嘱是在立遗嘱者死后才确立并生效的。\Quote{为此，} 他说，\Quote{祂是新盟约的中保。} 遗嘱是在最后的日子，即死亡之日制定的。\par

基督说：\BibleQuote{我愿我在哪里，他们也同我在哪里。} \BibleRef{若 17:24} 关于被剥夺继承权的人，请听祂说：\BibleQuote{我不为所有人祈求，只为那些因他们的话而信从我的人祈求。} \BibleRef{若 17:20} 遗嘱又与立遗嘱者和受遗人相关，因此他们有些事要接受，有些事要做。在此亦然：祂在作了无数许诺后，也要求他们做些什么，说：\BibleQuote{我给你们一条新诫命。} \BibleRef{若 13:34} 再次，遗嘱应有见证人。请再听祂说：\BibleQuote{我是为自己作证，那派遣我的父也为我作证。} \BibleRef{若 8:18} 又说：\BibleQuote{祂要为我作证} \BibleRef{若 15:26}，这是指护慰者。祂又派遣十二位宗徒说：\Quote{你们要在天主面前作证。}\par

2. \Quote{为此} （他说）\Quote{祂是新盟约的中保。} 什么是\Quote{中保}？中保并非其所中保之事物的主人，事物属于一人，中保是另一人：例如婚姻的中保不是新郎，而是帮助即将结婚之人的那位。同样，在此圣子成了圣父与我们之间的中保。圣父不愿将这产业留给我们，而向我们发怒，因我们远离祂而不悦；于是圣子成为我们与祂之间的中保，说服了圣父。\par

那么，祂如何成为中保？祂带来了来自圣父的话，也带来给我们，将来自圣父的传给我们，并加上祂自己的死亡。我们犯了罪，本应死亡；祂为我们而死，使我们配得这盟约。这样，盟约便稳固了，因为从今以后它不再为不配者而立。起初，祂像父亲为儿子那样安排；但当我们变得不配后，就不再需要盟约，而要受惩罚。\par

那么（他会说），你们为何看重法律呢？因为它使我们陷入如此大的罪中，以至于若非我们的主为我们而死，我们绝不能得救；法律没有能力，因为它是软弱的。\par

3. 他不再仅从世俗习惯，也从旧约中所发生的事来确立这一点——这尤其影响了他们。那里没有一个人死亡；那么那个盟约如何能稳固呢？同样（他说）。如何？因为那里也有血，正如这里有血。如果那不是基督的血，不必惊讶；因为那是预象。\Quote{所以，} 他说，\BibleQuote{前一个盟约也不是不用血而成立的。}\par

\Quote{成立} 是什么意思？即被确认、被批准。\Quote{所以} 一词意为\Quote{为此原因}。盟约的象征也必须与死亡有关。\par

% structure: p0011 source=h2 target=subsection reason=relative-heading-rank; base=h2

\subsection{希 9:19-20}

为何（请告诉我）盟约的书要被洒血？\Quote{因为}（他说）\BibleQuote{梅瑟按法律向全体百姓宣布了每一条诫命后，就用牛犊的血，同水、朱红羊毛和牛膝草，洒在书卷本身和全体百姓身上，说：\InnerQuote{这是天主命令你们的盟约之血。}} 那么告诉我，为何盟约的书被洒血，百姓也被洒血？岂不是因为那宝贵的血，从起初就预表了？为何用\Quote{牛膝草}？它紧密而能吸附。为何\Quote{水}？也显出水的洁净。为何\Quote{羊毛}？这也是为了保留血。此处血和水显明同一件事，因为圣洗圣事就是祂的苦难。\par

% structure: p0013 source=h2 target=subsection reason=relative-heading-rank; base=h2

\subsection{希 9:21-22}

4. \BibleQuote{他又用血洒了会幕和所有礼仪用的器皿。而且按法律，几乎一切都是用血洁净的，没有流血，就没有赦免。} 为何\Quote{几乎}？他为何加上修饰语？因为那些规条不是完美的净化，也不是完美的赦免，而是半完全且程度很小的。但在此祂说：\BibleQuote{这是我的血，新盟约的血，为你们倾流，以赦免罪过。} \BibleRef{玛 26:28}\par

那么\Quote{书卷}在哪里？祂洁净了他们的心思。他们自己就是新约的书卷。\Quote{礼仪用的器皿}在哪里？他们自己就是。\Quote{会幕}在哪里？又是他们自己；因为\BibleQuote{我要住在他们中间，} 祂说，\BibleQuote{在他们中行走。} \BibleRef{格后 6:16}\par

5. 但他们并没有被\Quote{朱红羊毛}所洒，也没有被\Quote{牛膝草}所洒。这是为什么呢？因为那洁净并非属肉体的，而是属灵性的，那血也是属灵性的。何以见得？它并非流自无理性动物的身体，而是流自圣神所预备的身体。用这血洒我们的不是梅瑟，而是基督，借着所传的话：\Quote{这是新约的血，为赦免罪过。}这句话代替牛膝草，蘸在血中，洒遍众人。在那里，身体外表被洗净，因为那洁净是属肉体的；但在这里，既然洁净是属灵性的，它就进入灵魂，洁净灵魂，不是简单地喷洒，而是在我们灵魂中涌流。已受洗的人明白所说的。在他们的情况中，确实只洒了表面；但被洒的人又把它洗掉了；因为他肯定不是一直带着血迹四处走动。但在灵魂的情况中，却不是这样，而是那血与灵魂的本质融合，使之强健、纯洁，并引领它达到那不可接近的美。\par

6. 然后他表明祂的死不仅是确证的原因，也是洁净的原因。因为既然死亡被认为是一件可憎的事，尤其是十字架的死，他说这血洁净了，甚至是一种珍贵的洁净，而且是针对更伟大的事物。因此，先前的祭献是为了这血。因此，羔羊的一切都是为了这个原因。\par

% structure: p0018 source=h2 target=subsection reason=relative-heading-rank; base=h2

\subsection{希伯来书 9:23}

\Quote{因此，那些天上之物的模型}（他说）\Quote{必须用这些祭物洁净，但天上的事物本身需要用比这些更好的祭献来洁净。}\par

那么，它们如何是\Quote{天上之物的模型}呢？他现在所说的\Quote{天上之物}又是什么意思？是指天堂吗？还是指天使？都不是，而是指我们的事物。由此可知，我们的事物是在天上，天上的事物也是我们的，即使它们在地上完成；因为虽然天使在地上，却被称作属天的。革鲁宾出现在地上，但仍然是属天的。为什么我说\Quote{出现}？相反，它们居住在地上，如同在乐园里一样；但这无关紧要；因为它们是属天的。并且，\Quote{我们的家乡在天上}\BibleRef{斐 3:20}，而我们却生活在这里。\par

\Quote{但这些是天上之物}，也就是说，存在于我们中间的哲学，那些蒙召的人。\par

\Quote{用比这些更好的祭献。}所谓\Quote{更好}是指比某种善的事物更好。因此，\Quote{天上之物的模型}也成为善的；因为模型本身并非邪恶，否则它们所模型的事物也会是邪恶的。\par

7. 既然我们是属天的，并且获得了如此伟大的祭献，让我们敬畏吧。我们不要再停留在地上；因为现在，凡是愿意的人，都有可能不在地上。因为在地上或不在是道德倾向和选择的后果。例如：天主被称为在天上。何以如此？不是因为他被空间所限，绝非如此，也不是因为他使地上没有祂的临在，而是因为祂与天使的关系和亲密。那么，如果我们亲近天主，我们就在天上。当我看见天上的主，当我本身成为天时，我还关心什么天呢？祂说：\Quote{我们要来，我和父，并要在他那里作我们的住所。}\BibleRef{若 14:23}\par

让我们使自己的灵魂成为天。天本质上是明亮的；因为即使在风暴中，它也不会变黑，因为它本身并不改变面貌，而是云层聚集遮盖了它。天有太阳；我们也有正义的太阳。我说可以成为天；我看甚至可以变得比天更好。怎样？当我们拥有了太阳的主。天始终纯洁无瑕，在风暴或黑夜中不变。因此，我们也不要被苦难或\Quote{魔鬼的诡计}\BibleRef{弗 6:11}所影响，而要持续无瑕纯洁。天是高的，远离大地。我们也应当这样（就我们自身而言）；让我们远离大地，升到那高度，远远地离开大地。天高于雨水和风暴，不受任何这些东西的影响。如果我们愿意，我们也能做到。\par

它似乎受到影响，但其实并非如此。因此，即使我们似乎受到影响，我们也不要受影响。因为如同在风暴中，大多数人不知道天的美丽，以为它改变了，但哲人们知道它根本没有受影响；同样，在我们自身遇到苦难时，大多数人以为我们随着苦难改变了，以为苦难触及了我们的内心，但哲人们知道苦难没有触及我们。\par

8. 那么，让我们成为天，让我们登上那高度，这样我们就会看到人与蚂蚁毫无区别。我不是说只有穷人，也不是说众人，即使那里有一位将军，即使那里有皇帝，我们也不会分辨皇帝与平民。我们不会知道什么是金子，什么是银子，什么是丝绸或紫色衣服：我们会看到一切都如同苍蝇，如果我们坐在那高度上。那里没有骚动，没有混乱，没有喧嚷。\par

有人问：行走在地上的人，如何能升到那高度呢？我不只用言语告诉你，我还要用事实向你展示那些已经达到那高度的人。那么，他们是谁呢？\par

我是指像\Person{保禄}这样的人，他们在地上却生活在天堂。但为什么我说\Quote{在天堂}？他们比那重天更高，甚至比另一重天还高，且攀升至天主自己。因为（他说）\BibleQuote{谁能使我们与基督的爱分开呢？是困苦、窘迫、迫害、饥荒、赤身、危险或刀剑吗？}\BibleRef{罗 8:35} 又说：\BibleQuote{我们并不注视那看得见的，而是注视那看不见的。}\BibleRef{格后 4:18} 你看见他连这里的事物也不注视了吗？但为向你显示他比诸天更高，请听他亲自说：\BibleQuote{我深信无论是死亡、是生命、是现在的事、是将来的事、是高处、是深处、或是任何其他受造物，都不能使我们与基督的爱分开。}\BibleRef{罗 8:38}\par

你看见思想如何超越一切事物，使他不只高于这个受造界，也不只高于这些诸天，甚至高于任何其他（如果有其他）？你看见他心灵的崇高吗？你看见这个帐棚匠成了什么吗？因为他有意志，他一生都在市场度过。\par

9. 因为没有任何阻碍——完全没有——能使我们不能超越众人，只要我们有意志。因为如果我们能在常人无法企及的艺术中如此成功，更何况在那不需要如此大劳苦的事上呢？\par

因为，请告诉我，有什么比在绷紧的绳索上行走——如履平地——并在高处穿衣脱衣，如同坐在长椅上更困难的呢？这表演难道不是令我们如此惊恐，以至于我们甚至不愿观看，一看到就恐惧战栗吗？再告诉我，有什么比把一根杆子放在脸上，并在上面放一个孩子，表演无数技巧来取悦观众更困难的呢？有什么比用剑玩球更困难的呢？请告诉我，有什么比彻底探查海底更困难的呢？还可以举出无数其他的技艺。\par

但比所有这些更容易的是德行和登上天堂，只要我们有意志。因为在这里只需要有意志，其余的一切就随之而来。因为我们不可以说\Quote{我不能}，也不可责怪造物主。因为如果祂使我们无能，却又命令，那是对祂自己的控诉。\par

10. 那么（有人说）为什么许多人不能呢？为什么许多人不愿意呢？因为如果愿意，所有人都能。因此\Person{保禄}也说：\BibleQuote{我愿众人都像我一样}\BibleRef{格前 7:7}，因为他知道所有人都能像他一样。因为他不会这样说，如果这是不可能的。你希望成为这样吗？只要抓住开始。\par

现在请告诉我，在任何技艺上，当我们想要掌握它们时，我们满足于愿望吗？还是也亲自参与？例如，有人想成为舵手；他不只是说\Quote{我愿意}就满足了，而是也动手工作。他想成为商人；他不只是说\Quote{我愿意}，而是也动手工作。他又想出国旅行，他也不只是说\Quote{我愿意}，而是动手工作。那么在一切事上，只有愿望是不够的，还必须加上工作；而你希望登上天堂时，难道只是说\Quote{我愿意}吗？\par

那么（他说）你怎么说\Quote{愿意}就够了呢？[我意思是]意愿与行为结合，抓住事情本身，劳苦。因为我们有\Person{天主}与我们一同工作，与我们一同行动。只要我们做出选择，只要我们专注于事情如同工作，只要我们认真思考，只要我们铭记在心，一切就随之而来。但如果我们沉睡，打着鼾指望进入天堂，我们怎能获得天上的产业呢？\par

因此，我劝你们，要愿意，要愿意。为什么我们要进行所有与现世生命有关的交易，而我们明天就要离开？那么让我们选择那种足以使我们度过整个永恒的德行吧：在其中我们将持续不断，并享受永远的福乐；愿我们都能获得，在我们的主基督耶稣内，愿光荣、权能、尊荣归于祂，同圣父和圣神，从今世直到永远，世世无尽。阿们。\par

\SourceRecord{Translated by Frederic Gardiner.；Nicene and Post-Nicene Fathers, First Series；Vol. 14.；Edited by Philip Schaff.；Buffalo, NY: Christian Literature Publishing Co.,；1889.；<http://www.newadvent.org/fathers/240216.htm>.}

% structure: p0001 source=h1 target=section reason=page-title

\section{希伯来书讲道 第十七篇}

% structure: p0002 source=h2 target=subsection reason=relative-heading-rank; base=h2

\subsection{希伯来书 9:24-26}

\Quote{ \Emph{因为基督并非进入人手所造的圣所，那些是模型} \Emph{真圣所的模型，而是进入了天堂本身，如今为我们出现在天主面前。也不是祂应当多次奉献自己，如同大司祭每年带着别人的血进入圣所，否则，祂从创世以来就必须多次受苦了。但是如今，一次} \Emph{在世界的末期，祂出现了，为要除去} \Emph{罪过藉着祂自己的祭献。} }\par

1. 犹太人极其以圣殿和会幕为荣。因此他们说：\Quote{上主的圣殿，上主的圣殿，上主的圣殿。} \BibleRef{耶 7:4} 因为在全地上，没有一座像这样的圣殿，无论是造价、美丽还是其他方面。因为天主命定了它，命令要极其壮丽地建造，因为他们也被物质事物所吸引和推动。因为它的墙壁上有金砖；任何愿意的人可以在列王纪下卷和厄则克耳书中，以及当时花费了多少塔冷通的金子中，学到这一点。\par

但第二座圣殿是一座更荣耀的建筑，无论在其美丽还是其他方面。它受到尊敬不仅因为这个原因，也因为它是一座。因为他们习惯从地极前往那里，无论是从巴比伦还是从埃塞俄比亚。路加在宗徒大事录中显示这一点，他说：\Quote{那里有} \Quote{帕提雅人、玛待人、厄蓝人，以及居住在美索不达米亚、犹太和卡帕多细雅、本都和亚细亚、夫黎基雅和旁非里雅、埃及和靠近基勒乃的利比亚地区的人。} \BibleRef{宗 2:5} 那时，住在世界各地的人都聚集在那里，圣殿的名声很大。\par

那么，保禄做了什么呢？他对待祭献的方式，在这里也做了。因为正如在那里，他以基督的死亡来对照（它们），同样在这里，他也以整个天堂来对照圣殿。\par

2. 他不只以此指出了区别，而且通过补充说，这位司铎更接近天主：因为他说了：\Quote{出现在天主面前。} 因此，他使事情显得庄严，不仅藉着（对）天堂（的考虑），而且藉着（基督的）进入（那里）。因为不仅像这里通过象征，而且在那里祂亲自看见天主。\par

你看到吗，谦卑的事物贯穿始终所表达的屈尊就卑？那么你为什么还惊异祂在那里代祷，在那里祂将自己置于大司祭的地位？\Quote{也不是，祂应当多次奉献自己，如同大司祭。}\par

\Quote{因为基督并非进入人手所造的圣所}（他说）\Quote{那些是真圣所的模型。}（这些是真实的，那些是模型，因为圣殿也被这样安排，如同天上的天堂。）\par

你说什么？那位无所不在、充满万有的，难道除非祂进入天堂，祂才\Quote{出现}吗？你看，所有这些都属于血肉。\par

\Quote{出现}，他说，\Quote{为我们出现在天主面前。} 什么是\Quote{为我们}？祂上去（他的意思是）带着一个能够平息圣父愤怒的祭献。为什么（告诉我）？祂是仇敌吗？天使们是仇敌，祂不是仇敌。因为天使们是仇敌，听他说：\Quote{祂成就了和平，涉及地上的事物和天上的事物。} \BibleRef{哥 1:20} 因此祂也\Quote{进入了天堂，如今为我们出现在天主面前。} 祂\Quote{如今出现}，但\Quote{为我们。}\par

3. \Quote{也不是祂应当多次奉献自己，如同大司祭每年带着别人的血进入圣所。} 你看有多少区别？\Quote{多次}对应\Quote{一次}；\Quote{别人的血}对应\Quote{祂自己的血。} 距离很大。那么祂自己既是牺牲品，又是司铎和祭献。因为如果不是这样，而且必须奉献许多祭献，祂就必须被多次钉死了。\Quote{因为那时，}他说，\Quote{祂从创世以来就必须多次受苦了。}\par

在此处，他也隐晦地提到了某事。\Quote{但是如今，在末期，再一次。} 为什么\Quote{在末期}？在许多罪之后。因此，如果它发生在起初，那么没有人会相信；而且祂不必死第二次，一切都是无用的。但是因为后来有许多过犯，祂有理由那时出现：他在另一处也表达了，\Quote{罪恶越多，恩宠越发丰盛。但是如今，在末期，祂出现了一次，藉着祂自己的祭献来除去罪过。} \BibleRef{罗 5:20}\par

% structure: p0014 source=h2 target=subsection reason=relative-heading-rank; base=h2

\subsection{希伯来书 9:27}

4. \BibleRef{希 9:27} \Quote{就如人注定只死一次，之后就是审判。} 他接着也说为什么祂只死一次：因为祂藉着一次死亡成为了赎价。\Quote{它已被注定}（他说）\Quote{人只死一次。} 那么，这就是\Quote{祂死了一次}的意义，为所有人。（那么怎么样？我们不再死那次死亡吗？我们确实要死，但我们不能停留在死亡中：这根本就不是死亡。因为死亡的专制，真正的死亡，是当一个人死后永远不能再回到生命。但是当死后仍然活着，而且是更好的生命，这不是死亡，而是睡眠。）既然死亡要占有所有人，因此祂死了，为要拯救我们。\par

% structure: p0016 source=h2 target=subsection reason=relative-heading-rank; base=h2

\subsection{希伯来书 9:28}

\Quote{所以基督是一次被奉献的。} 被谁奉献？显然被祂自己。在此处他说，祂不仅是司铎，而且是牺牲品，也是祭品。因此，\Quote{被奉献了}这些话。\Quote{一次被奉献}（他说）\Quote{为承担许多人的罪过。} 为什么\Quote{许多人的}，而不是\Quote{所有人的}？因为并非所有人都相信。因为祂确实为所有人死了，这是祂的部分：因为那死亡是对抗所有人毁灭的平衡。但是祂没有承担所有人的罪过，因为他们不愿意。\par

\Quote{祂背负了罪}是什么意思？正如在祭献中，我们提起自己的罪并说：\Quote{无论我们故意或无意犯了罪，求祢宽恕}，也就是说，我们先提及它们，然后请求宽恕。这里也是如此。基督在哪里做了这事？请听祂自己说：\Quote{我为了他们祝圣我自己。}\BibleRef{若 17:19}看！祂背负了罪。祂从人那里取来，带到父前；并非为了定他们（人类）的罪，而是为宽恕他们。\par

\Quote{祂要第二次显现给那期待祂的人，与罪无关，为救恩}（他说）。\Quote{与罪无关}是什么意思？这就是说，祂不犯罪。因为祂并非像欠了死亡之债而死去，也不是因罪而死。但\Quote{祂要显现}是什么意思？你要说，为惩罚。然而他并没有这样说，而是说了令人欣慰的话：\Quote{祂要第二次显现给那期待祂的人，与罪无关，为救恩}。因此，今后他们不再需要祭献来自救，而是藉着行为去做。\par

% structure: p0020 source=h2 target=subsection reason=relative-heading-rank; base=h2

\subsection{\BibleRef{希 10:1}}

5. \BibleRef{希 10:1}\Quote{因为}（他说）\Quote{法律拥有未来美物的影子，不是事物的真像}；即不是真正的实在。就如绘画，只要一个人只画轮廓，那只是某种\Quote{影子}；但当加上明亮的色彩和填充颜色，就变成了\Quote{像}。法律也类似。\par

\Quote{因为}（他说）\Quote{法律拥有未来美物的影子，不是事物的真像}，即指祭献和赦免：\Quote{绝不能藉着那些每年常献的祭，使前来敬拜的人得以完全。}\par

% structure: p0023 source=h2 target=subsection reason=relative-heading-rank; base=h2

\subsection{\BibleRef{希 10:2-9}}

\Quote{难道它们不会停止被献吗？因为敬拜者一次被洁净，就不再觉得有罪了。但在这些祭物中，每年都叫人想起罪来。因为公牛和山羊的血绝不能除去罪。所以，当基督来到世界时，祂说：\InnerQuote{祭物和供物祢不想要，但祢为我预备了身体。燔祭和赎罪祭是祢不喜欢的。那时我说：看！我来了，书卷上已记载了我，为要遵行祢的旨意，天主。}上面说：\InnerQuote{祭物、供物和燔祭、赎罪祭，祢既不要，也不喜欢}（这些是按法律所献的），然后祂说：\InnerQuote{看！我来了，为要遵行祢的旨意，天主。}祂除去前者，为要立定后者。}\BibleRef{希 10:2}\par

你再次看到他论据的充实？这个祭献（他说）是一次；而其他却是多次：因此它们没有力量，因为它们多。因为，请告诉我，如果一次已足够，何需多次？所以它们多，并且\Quote{不断}被献，证明他们（敬拜者）从未被洁净。就如一种药物，如果它有效、能促进健康、能完全除去疾病，一次施用就达到全部效果；因此，如果一次施用就完成全部，就证明它的力量，因为它不再被施用；而它的作用就在于不再被施用；反之，如果不断施用，这明显证明它没有力量。因为药物的优点是一次施用，而非多次。同样，在此例中也是如此。为什么他们确实用\Quote{同样的祭物}不断被治愈？因为如果他们从所有罪中得释放，祭献就不会每天继续被献。因为它们被定为要在每日早晚不断为全体百姓奉献。所以那是罪的控告，而非罪的释放；是软弱的控告，而非力量的展示。因为第一次没有力量，就献上另一个；而这个无效，又献上一个；因此这是罪的证据。的确，\Quote{献祭}正是罪的证据，\Quote{不断}是软弱的证据。但关于基督，正好相反：祂是\Quote{一次被献}。因此，预像只包含外形，而不包含力量；就如肖像，肖像有人的外形，却没有力量。所以实体和预像彼此有些共同点。因为外形在两者中相同，但力量不同。同样，关于天堂和会幕也是如此，因为外形相同：那里有至圣所，但力量和其他事物并不相同。\par

\Quote{祂藉着献上自己而显现以除去罪}是什么意思？\Quote{除去}是什么意思？是使之受轻蔑。因为罪不再有任何胆量；因为它被废止了，在它应当要求惩罚时，却没有要求：就是说，它受了强暴：当它期待毁灭所有人时，它自己却被毁灭了。\par

\Quote{祂藉着献上自己而显现}（他说），也就是说，\Quote{祂显现}于天主面前，并亲近（祂）。因为不要（以为）大司祭惯常一年多次这样做……所以从此以后，这样做是徒然的，尽管仍在做；因为哪里没有伤口，何必需要药物？为此，祂规定了\Quote{不断}的祭献，因为它们的无力，且为使人记起罪。\par

6. 那么怎样呢？我们不是每日献祭吗？我们确实献祭，但这是纪念祂的死亡，而这份纪念是独一无二的，并非多次。为何是独一无二而非多次呢？因为那祭献曾一次而永远地被献上，并被带进至圣所。这是那祭献的预像，而这纪念是那祭献的纪念。因为我们总是献上同一祭物，并非今日一只羊，明日另一只，而总是同一事物：因此祭献是唯一的。然而，按此推论，既然献祭在多个地方举行，岂非有多个基督？但基督处处唯一，在此圆满，在彼也圆满，同一个身体。因此，正如祂在多处被献上，仍是同一个身体，而非多个身体；同样，也是同一个祭献。祂是我们的\OriginalTerm{大司祭}{High Priest}，祂献上了洁净我们的祭献。我们现在所献的，就是那时所献的，这祭献不会耗尽。这是为纪念那时所做的事而行。因为祂说：\Quote{你们要这样做，来纪念我。}\BibleRef{路 22:19} 这不是另一个祭献，如同大司祭所做的那样，而是我们总是献上同一祭献，或者更确切地说，我们举行一个祭献的纪念。\par

7. 既然我提到了这祭献，我愿对你们这些已受入门圣事的人说一点话；话虽不多，却具有极大的力量和益处，因为这不是我们的话，而是圣神的话。那么是什么呢？许多人每年领受这祭献一次，另一些人两次，还有一些人多次。我们的话是对所有人说的；不仅是对这里的你们，也是对那些居住在旷野中的人。因为他们一年一次领受，而且往往隔两年一次。\par

那么怎样呢？我们该赞同那些一年领受一次的人吗？还是那些多次领受的人？或是那些极少领受的人？不，不是一年一次的人，也不是经常领受的人，也不是稀少领受的人，而是那些以纯洁的良心、从纯洁的心、以无可指摘的生活前来领受的人。让这样的人不断亲近吧；但那些不是这样的人，连一次也不要亲近。你问：为什么？因为他们领受的是对自己的审判，是的，还有谴责、惩罚和报复。不要惊讶。因为正如滋养性的食物，若被一个没有胃口的人吃下，会败坏和腐蚀整个系统，并成为疾病的起因；同样，对于可畏的奥迹也是如此。你在属灵的筵席、君王的筵席上赴宴，却又用淤泥玷污你的口？你以馨香的油膏抹自己，却又使自己充满臭气？\par

告诉我，我恳求你，当你一年之后领受圣体时，你认为那四十天足够你洁净那整个时期的罪过吗？而一周之后，你又重蹈覆辙？请告诉我，如果你在久病之后康复了四十天，然后又放任自己吃那致病的食物，你岂不是也白费了先前的努力吗？因为如果自然的事物会改变，那么依赖于选择的事物更是如此。举例来说，我们靠本性观看，我们天生有健康的眼睛；但往往因不良的体质，我们的视力受损。如果自然的事物会改变，那么选择的事物更是如此。你为灵魂的健康指定四十天，或许连四十天也不到，却期望平息天主的忿怒？告诉我，你在开玩笑吗？\par

我说这些话，并非禁止你们一年一次前来，而是希望你们不断亲近。\par

8. 这些事是赐给圣洁的人的。当执事呼召圣洁的人时，他也这样宣告；甚至藉着这呼召，他查究所有人的过失。因为如同在羊群中，许多羊固然健康，但许多羊满身疥癣，就必须将这些病羊与健康的羊分开；同样，在教会里，既然有些羊健康，有些羊患病，他就藉着这呼声将二者分开，司铎（我说）藉着这最可畏的喊声四处巡视，呼召并吸引圣洁的人。因为人不可能知道邻人的事，（因为他说：\Quote{有谁知道人的事，除了人里面的那人的神？}\BibleRef{格前 2:11}）他在整个祭献完成后发出这呼声，使无人草率或偶然地来到这属灵的水泉。因为在羊群中（没有什么阻止我们再次使用同一例子），我们把病弱的羊关在里面，使它们处于黑暗中，给它们不同的食物，不允许它们享受纯净的空气、单纯的青草或圈外的泉水。同样，在这里，这呼声就如锁链一般。\par

你不能说：\Quote{我不知道，我不晓得这事有危险。}不，保禄尤其为此作了见证。但你或许会说：\Quote{我从未读过。}这不是借口，反而是指控。你每天进堂，却对此一无所知？\par

然而，为使你们甚至没有这个借口，为此，他高声呼喊，用可畏的喊声，像传令官一样高举着手，高高站立，众目所瞩，在那可畏的静默之后大声呼喊，邀请一些人，又禁止一些人；他不是用手做这事，而是用舌头，比手更清楚。因为那声音落在我们耳中，就像一只手，推开并驱逐一些人，又引入并呈献另一些人。\par

请你告诉我，我恳求你，在奥林匹克运动会上，难道传令官不是站着，用高昂的声音呼喊说：\Quote{有人控告这人吗？他是奴隶吗？他是贼吗？他是品行恶劣的人吗？}然而，这些争夺奖品的竞赛并非关乎灵魂，亦非关乎良好品德，而是关乎力量和身体。那么，如果在那里有身体的锻炼，尚且要对品性进行许多审查，那么在这里，灵魂独自是战斗者，更是何等需要呢？我们的传令官现今也站着，不是抓住每个人的头把他拉向前，而是从内部一齐抓住所有人的头；他不把别的控告者摆在他们面前，而是把他们自己摆在自己面前。因为他没有说：\Quote{有人控告这人吗？}而是说什么？\Quote{若有人控告自己。}因为当他说：\Quote{神圣之物归于神圣之人}，他意思是：\Quote{若有人不圣洁，就不要近前。}\par

他不是简单地说：\Quote{无罪，}而是：\Quote{圣洁。}因为使人圣洁的不仅仅是脱离罪，还有圣神的临在以及善工的富足。他说：我不仅希望你们脱离泥沼，还希望你们光明美丽。因为如果巴比伦王从俘虏中挑选青年人时，选中那些形貌俊美、面容秀丽的，那么当我们站在王桌前时，更需要形貌美丽——我指的是灵魂的形貌——有金饰、衣服纯洁、鞋子属于王室、灵魂的面容端正、金饰环绕，就是真理的腰带。让这样的人近前，触摸王杯。\par

但如果有人衣衫褴褛、污秽肮脏，想要进入王桌，要想想他将遭受多大的痛苦，四十天不足以洗去所有时间所犯的过犯。因为如果地狱尚且不够——尽管它是永恒的（正因此它才\Emph{是}永恒的）——更何况这短暂的时间呢？因为我们没有表现出强烈的悔改，而是软弱的。\par

9. 尤其太监应当侍立在王旁：我所说的太监，是指那些心地清明、没有皱纹斑点、思想崇高、灵魂的眼睛温柔而敏锐、活跃而犀利、不昏睡也不懈怠；充满自由，又远避无礼和放肆；警醒、健康，既不十分阴郁消沉，也不放荡软弱。\par

这个眼睛，我们有权自己创造，并使其敏锐美丽。因为当我们引导它，不朝向烟尘（因为所有属人之事都是如此），而是朝向和风、轻气、属天和高处之事，充满宁静纯洁和极大喜乐，我们将迅速恢复它，并使之强壮，如同它在这样的默观中畅快。你见过贪婪和巨富吗？不要举目仰望它们。那东西是泥沼，是烟，是恶气，黑暗，极大的痛苦和令人窒息的重担。你见过一个培养正义、自足、有充足休闲空间、有忧虑却不专注于世物的人吗？把你的眼睛放在那里，高举它；你将使它远为美丽辉煌，不喂养它以地上之花，而是以美德之花、节制、适度及所有其余。因为没有什么像邪恶的良心那样困扰眼睛（经上说：\Quote{我的眼睛，因愤怒而困扰}（见：\BibleRef{咏 6:7}））；没有什么像它那样使之昏暗。使眼睛摆脱这种伤害，你将使它充满活力而强壮，常以美好希望为养料。\par

愿我们所有人都使这眼睛以及灵魂的其他能力，如同基督所愿的那样，使我们堪当被置于我们之上的元首，前往祂所愿之处。因为祂说：\Quote{我愿意我在哪里，他们也同我在一起，好使他们看见我的荣耀。}（\BibleRef{若 17:24}）愿我们众人都在我们的主耶稣基督内享此恩典，愿光荣、权能、尊荣归于祂，及圣父和圣神，从今世直到永远，世世无尽。阿们。\par

\SourceRecord{Translated by Frederic Gardiner.；Nicene and Post-Nicene Fathers, First Series；Vol. 14.；Edited by Philip Schaff.；Buffalo, NY: Christian Literature Publishing Co.,；1889.；<http://www.newadvent.org/fathers/240217.htm>.}

% structure: p0001 source=h1 target=section reason=page-title

\section{\WorkTitle{希伯来书讲道集 第十八篇}}

% structure: p0002 source=h2 target=subsection reason=relative-heading-rank; base=h2

\subsection{\BibleBook{希伯来}{书} 10:8-13}

\BibleQuote{\Emph{上面当他说：\Quote{祭献和供物，全燔祭和赎罪祭，你既不喜欢，也不要}，这些都是按照}\Emph{法律所奉献的。然后又说：\Quote{看，我来是要承行你的旨意，天主}。他废除了先前的，为要建立以后的。藉着这旨意，我们}\Emph{得以圣化，就是藉着身体}耶稣\Emph{基督，一次而永远的。}\Emph{每位司铎天天侍立，屡次奉献同样的祭献，这些祭献决不能除去罪恶。但这个人奉献了一次赎罪的祭献，永远地，然后坐在天主右边，从此等待，直到他的仇敌被放在他的脚下。}}\par

1. 在前面，他已经表明那些祭献对于成全的洁净是无效的，是预像，且大有缺陷。既然对此论点有这样的反对：若是预像，为何在真理来临后，它们没有停止，也没有让位，反而仍旧被举行？他在这里就此着力，展示它们不再被举行，即使作为预像，因为天主不接受它们。而这一点他又不是从新约，而是从先知来展示，从古时带来最有力的证言，即它[旧制度]要结束并终止，他们所做的一切都是徒然，他们\Quote{时常反抗圣神}。\BibleRef{宗 7:51}\par

而且他还进一步指出，它们停止不但是现在，而是在默西亚来临之时，更确切地说，甚至在他来临之前；并且基督是如何没有在最后废除它们，而是它们先被废除，然后他来了；它们先被终止，然后他显现。免得他们说：即使没有这祭献，藉着那些祭献，我们也能中悦天主，他就等待这些祭献被定罪[为软弱]，然后显现；因为（他说）：\Quote{祭献和供物，你既不喜欢。}他由此除去了一切；说了普遍性的之后，他也特别地说：\Quote{在全燔祭和赎罪祭中，你既不喜欢。}但\Quote{供物}是指除祭献以外的一切。\Quote{于是我说：看，我来。}这是论谁说的？无非是基督。\par

此处他并不责难那些奉献的人，表明天主不接受他们，不是因他们的邪恶，正如他在别处所说，而是因为这事本身已被定罪，将来被证明没有力量，也不合时宜。那么这与\Quote{祭献}被\Quote{屡次}奉献有何关系？不仅从它们被\Quote{屡次}奉献（他的意思是）显出它们是软弱的，且一无所成；而且从天主不接受它们，因为它们是无益无用的。在另一处说：\Quote{你若想要祭献，我必会献上。}\BibleRef{咏 50:16}因此，由此他也表明他不想要它。所以祭献不是天主的旨意，而是废除祭献。因此他们奉献祭献是违反他的旨意。\par

\Quote{承行你的旨意}是什么？就是舍弃我自己，他意味：这是天主的旨意。\Quote{藉这旨意，我们得以圣化。}或者他甚至意味更深一层：祭献并不能使人洁净，而是天主的旨意。因此奉献祭献不是天主的旨意。\par

2. 为什么你们现在奇怪这不是天主的旨意，既然从起初就不是他的旨意？因为他说：\Quote{谁向你们要求了这事？}\BibleRef{依 1:12}\par

那么他亲自如何吩咐这事？是出于迁就。因为保禄说：\Quote{我愿众人像我一样}\BibleRef{格前 7:7}，关于节欲的事；又说：\Quote{我愿年轻的妇人结婚，生养儿女}\BibleRef{弟前 5:14}；并且制定了两个意愿，然而这两个都不是他自己的，尽管他命令；但一个确实是他自己的，因此他制定时不加理由；而另一个不是他自己的，虽然他希望这事，因此加上了理由。因为先指责了他们，因为他们\Quote{背弃了基督而放纵}\BibleRef{弟前 5:11}，然后他说：\Quote{我愿年轻的妇人结婚，生养儿女。}\BibleRef{弟前 5:14}同样在此处，奉献祭献不是他首要的意愿。因为，正如他所说：\Quote{我不愿恶人丧亡，但愿他归向我而得生}\BibleRef{则 33:11}；在另一处他说他不仅愿意，甚至是渴望这事：然而这些是彼此矛盾的：因为强烈的意愿就是渴望。那么你如何\Quote{不愿意}？你如何在另一处\Quote{渴望}，这是强烈意愿的标志？在此事上也是如此。\par

\Quote{藉这旨意，我们得以圣化}，他说。如何圣化？\Quote{藉着耶稣基督的身体，一次而永远的奉献。}\par

3. \Quote{每位司铎天天侍立，屡次奉献同样的祭献。}（因此站立是服务者的标记；同样，坐是被服务者的标记。）\par

% structure: p0012 source=h2 target=subsection reason=relative-heading-rank; base=h2

\subsection{\BibleBook{希伯来}{书} 10:14-15}

\Quote{但这个人奉献了一次赎罪的祭献，永远地，然后坐在天主右边，从此等待，直到他的仇敌被放在他的脚下。}\BibleRef{希 10:14}\Quote{因为藉着一次奉献，他成全了永远被圣化的人。圣神也为我们作证。}他说那些[祭献]没有被奉献；他从所写的和未写的推理；此外，他也引用了先知的话说：\Quote{祭献和供物，你既不喜欢。}他曾说天主赦免了他们的罪。\par

% structure: p0014 source=h2 target=subsection reason=relative-heading-rank; base=h2

\subsection{\BibleBook{希伯来}{书} 10:16-18}

祂又从所记载的证言证明这一点，因为（他说）\Quote{圣神}\InnerQuote{给我们作证：因为祂说过之后}（\BibleRef{希 10:16}）\Quote{这是自那时以后我要与他们订立的盟约——主说：我要将我的律法放在他们心里，写在他们脑中；他们的罪过和邪恶，我不再记忆。那里这些罪过得了赦免，便再没有赎罪的祭献了。}因此，当祂订立盟约时，就赦免了他们的罪过；而祂是通过祭献订立盟约的。所以，如果祂藉着一次祭献赦免了罪过，就不需要第二次了。\par

\Quote{祂坐在天主右边，从此等待。}为何等待？\Quote{直到祂的仇敌被放在祂脚下。因为祂藉着一次祭献，永远成全了那些被圣化的人。}但或许有人会说：为何不立即将他们放在脚下？为了后来将要产生并出生的信徒。那么，他们将被放在脚下是从哪里显明的呢？藉着\Quote{祂坐下了}这句话。祂再次引用那个证言说：\Quote{直到我把你的仇敌放在你的脚下。}（\BibleRef{希 1:13}）但祂的仇敌是犹太人。既然祂说了\Quote{直到祂的仇敌被放在祂脚下}，而这些（仇敌）却猛烈逼迫使祂，所以此后祂才引入所有关于信德的论述。但仇敌是谁？所有不信者：魔鬼。为表明他们屈服的程度，祂没有说\Quote{被屈服}，而是\Quote{被放在祂脚下}。\par

4. 因此，让我们不要成为祂仇敌中的一员。因为不仅是那些不信者和犹太人是仇敌，那些充满不洁生活的人也是。\Quote{因为属肉体的思念，是与天主为仇，因为它不服天主的法律，也是不能服。}（\BibleRef{罗 8:7}）那么，你说什么？这不是可责备的理由。不，这正是极大的可责备的理由。因为恶人，只要他是恶的，就不能服从天主的法律；但他却可以改变，成为善人。\par

那么，让我们抛弃属肉体的思念。但什么是属肉体的？凡使身体昌盛顺遂，却损害灵魂的：例如，财富、奢华、荣耀（这些都是属肉体的）、肉体的爱。那么，让我们不要爱获利，而要永远追求贫穷：因为这是极大的善。\par

但是（你说）这使人谦卑并微不足道。（的确：）因为我们需要这一点，它对我们大有裨益。（经上）说：\Quote{贫穷使人谦卑。}（\BibleRef{箴 10:4}）基督又说：\Quote{神贫的人是有福的。}（\BibleRef{玛 5:3}）那么，你因行在通往德行的道路上而忧伤吗？难道你不知道这给予我们极大的自信吗？\par

但是，有人说：\Quote{穷人的智慧被藐视。}（\BibleRef{训 9:16}）另一个人又说：\Quote{求你别使我贫穷，也别使我富裕。}（\BibleRef{箴 30:8}）以及\Quote{救我脱离贫穷的火炉。}（\BibleRef{依 48:10}）再者，如果财富和贫穷来自主，那么贫穷或财富怎能是恶呢？为什么说这些话？它们是在旧盟约下说的，那时十分看重财富，极度轻视贫穷，那时（前者）是咒骂，（后者）是祝福。但现在不再是如此了。\par

但你们愿意听贫穷的赞美吗？基督追求贫穷，并说：\Quote{人子却没有枕头的地方。}（\BibleRef{玛 8:20}）祂又对门徒说：\Quote{不要带金、银，也不要带两件内衣。}（\BibleRef{玛 10:9}）保禄写信说：\Quote{似一无所有，却无所不有。}（\BibleRef{格后 6:10}）伯多禄对那位生来跛脚的说：\Quote{银子和金子，我没有。}（\BibleRef{宗 3:6}）是的，即使在旧盟约本身中，当财富被赞赏时，谁是受赞赏的？难道不是只有一件羊皮袄的厄里亚吗？难道是厄里叟吗？难道是若翰吗？\par

那么，没有人该因他的贫穷而自卑：不是贫穷使人卑微，而是那迫使我们需求许多人、并使我们欠许多人情的财富。\par

还有什么比雅各伯更贫穷的呢（告诉我），他说：\Quote{如果主赐给我食物吃、衣服穿}（\BibleRef{创 28:20}）？厄里亚和若翰难道缺乏勇气吗？难道前者没有责备阿哈布，后者没有责备黑落德吗？后者说：\Quote{你占有你兄弟斐理伯的妻子是不合法的。}（\BibleRef{谷 6:18}）厄里亚勇敢地对阿哈布说：\Quote{扰乱以色列的不是我，而是你和你的父家。}（\BibleRef{列上 18:18}）你看到这特别产生勇气：贫穷（我说）吗？因为富人是个奴隶，易受损失，并且处在每个想伤害他之人的权力下；而那什么都不拥有的人，不怕没收，也不怕罚款。如果贫穷使人缺乏勇气，基督就不会派遣祂的门徒带着贫穷去从事需要极大勇气的工作。因为穷人非常坚强，没有可被错待或虐待的地方。但富人在各方面易受攻击：正如一个人很容易抓住一个拖着许多长绳子的人，却不能轻易抓住一个赤身裸体的人。这里，就富人的情况也是如此：奴隶、金子、土地、无数事务、无数忧虑、困难处境、需要，使他易为所有人攻击。\par

5. 那么，从今以后，没有人该把贫穷当作耻辱的原因。因为如果有德行在，世上的所有财富与之相比，连泥土、甚至尘埃都不如。因此，如果我们想要进入天国，就让我们追求这一点。因为祂说：\Quote{变卖你所有的，施舍给穷人，你必有财宝在天上。}（\BibleRef{玛 19:21}）又说：\Quote{富人进入天国，是难的。}（\BibleRef{玛 19:23}）你看到吗？即使我们没有拥有它，我们也该把它吸引到自己身上。贫穷是这么大的善！因为它如同引导我们走上通往天堂之路，它是战斗的涂膏，是伟大而可赞美的操练，是宁静的港湾。\par

但你却说：我需要许多东西，又不愿受人恩惠。然而，即使在这方面，富人也比你逊色；因为你或许为生计而求人，而他却因贪婪而厚颜无耻地求取无数事物。由此可见，富人才是需要许多人的帮助，且常常是那些配不上他们的人。例如，他们常常需要士兵或奴隶阶层的人；但穷人甚至不需要皇帝本人，即使他需要皇帝，也会因本可富有却自甘卑微而受人钦佩。\par

因此，不要让任何人指责贫穷是无数罪恶的根源，也不要让人违背基督的话，祂宣称贫穷是德行的圆满，说：\BibleQuote{你若愿意是成全的。}\BibleRef{玛 19:21} 因为祂不仅在言语上如此说，也在行动上彰显，并通过祂的门徒教导我们。因此，让我们追求贫穷，对清醒的人来说，这是最大的善。\par

也许我听众中有些人视贫穷为不祥之事而避之。我不怀疑这一点。因为这种疾病在多数人中十分严重，财富的暴政如此之大，以至于他们连在言语上都无法忍受放弃财富，反而视之为不祥。愿这种思想远离基督徒的灵魂：因为没有人比那自愿并欣然选择贫穷的人更富有了。\par

6. 如何证明？我将告诉你，若你愿意，我将证明那自愿选择贫穷的人甚至比国王本人更富有。因为国王确实需要许多东西，且焦虑不安，担心军需供应不足；而贫穷的人却一无所缺，无所忧虑，即便忧虑，也不为如此重大的事。那么，请告诉我，谁是富人？是那每日祈求、竭力聚敛、唯恐匮乏的人，还是那毫无聚敛却丰足有余、无需依赖任何人的人？因为给人信心的是德行和对天主的敬畏，而非财产。财产反而使人受束缚。经上记载：\BibleQuote{礼物与馈赠，使明智者的眼目失明，如同口罩掩口，阻止责斥。}\BibleRef{德 20:29}\par

请看，贫穷的\Person{伯多禄}如何惩治富有的\Person{阿纳尼雅}。难道不是一人富、一人穷吗？但请看，\Person{伯多禄}以权威说话：\BibleQuote{你告诉我，你们卖田地的价钱就是这些吗？}\BibleRef{宗 5:8} 而后者顺从地回答：\BibleQuote{是的，就是这些。} 你或许会说：谁能使我也像\Person{伯多禄}一样？你若愿意，你也能像\Person{伯多禄}；舍弃你所有的一切。\BibleQuote{他施舍钱财，赒济贫穷。}\BibleRef{咏 111:9} 跟随基督，你便也能像他。如何？你或许会说：他行了奇迹。那么，我问你，是奇迹使\Person{伯多禄}受人钦佩，还是他生活方式的勇气？难道你没有听见基督说：\BibleQuote{不要因为魔鬼屈服于你们而欢喜；你若愿意是成全的……}\BibleRef{路 10:20} 听\Person{伯多禄}说：\BibleQuote{银子和金子，我没有；但把我所有的给你。}\BibleRef{宗 3:6} 如果一个人拥有金银，他就没有那些其他的恩赐。\par

你或许会说：既然如此，为什么许多人两者都没有？因为他们并非自愿贫穷；因为自愿贫穷的人拥有一切美善。虽然他们不能复活死人或瘸子，但比这一切更大的，是他们拥有对天主的信赖。他们将在那一日听到那有福的声音：\BibleQuote{你们蒙我父降福的，来吧！}（还有什么比这更好呢？）\BibleQuote{承受自创世以来为你们预备的国度：因为我饿了，你们给了我吃的；我渴了，你们给了我喝的；我作客，你们收留了我；我赤身露体，你们给了我穿的；我患病，你们看顾了我；我在监里，你们来探望了我。承受自创世以来为你们预备的国度吧。}\BibleRef{玛 25:34} 因此，让我们逃避贪婪，好能获得天国。让我们喂养穷人，好能喂养基督；使我们能藉着基督耶稣我们的主，与祂同为承继者。愿光荣、权能、尊崇归于祂及圣父与圣神，从今世直到永远，万世无疆。阿们。\par

\SourceRecord{Translated by Frederic Gardiner.；Nicene and Post-Nicene Fathers, First Series；Vol. 14.；Edited by Philip Schaff.；Buffalo, NY: Christian Literature Publishing Co.,；1889.；<http://www.newadvent.org/fathers/240218.htm>.}

% structure: p0001 source=h1 target=section reason=page-title

\section{希伯来书第十九篇讲道}

% structure: p0002 source=h2 target=subsection reason=relative-heading-rank; base=h2

\subsection{希伯来书 10:19-23}

\BibleQuote{所以，弟兄们，我们既然靠着耶稣的血，得以坦然进入至圣所，是藉着祂给我们开启的一条又新又活的路，从幔子经过，这幔子就是祂的血肉，又既然有一位大司铎掌管天主的家，我们就当以真诚的心，在信德的确证下走近，同时我们的心已洒净，脱离了邪恶的良心，身体也用清水洗净了。我们要坚定不移地持守我们所宣认的望德。}\par

1. \BibleQuote{所以，弟兄们，我们靠着耶稣的血，得以坦然进入至圣所，是藉着祂给我们开启的一条又新又活的路。} 他既已表明了大司铎、祭品、会幕、盟约和应许的差异，且这差异是巨大的，因为那些是暂时的，这些是永恒的；那些\BibleQuote{快要消失}，这些是永久的；那些没有能力，这些是完美的；那些是预像，这些是实体；因为（他说）\BibleQuote{不是按照属血肉的诫命，而是按照无尽生命的能力。}\BibleRef{希 7:16} 以及\BibleQuote{祢是永远的司铎。}\BibleRef{希 5:6} 论到盟约，他说那一者是旧的（因为\BibleQuote{那逐渐衰败、变旧的事物，就快要消失}\BibleRef{希 8:13}），但这一者是新的，且有罪过的赦免，而那一者完全没有这类事物：因为（他说）\BibleQuote{法律从未使任何事物达到完美。}\BibleRef{希 7:19} 又说，\BibleQuote{祭品和供物祢不想要。}\BibleRef{希 10:5} 那一者是手造的，而这一者是\BibleQuote{非人手所造的}\BibleRef{希 9:11}；那一者\BibleQuote{有山羊的血}\BibleRef{希 9:12}，这一者是有主的血；那一者有的司铎\BibleQuote{站着}，这一者\BibleQuote{坐着}。既然那些都是低劣的，这些更为伟大，因此他说：\BibleQuote{所以，弟兄们，我们得着放胆。}\par

2. \OriginalTerm{放胆}{Boldness}：从何而来？正如罪（他意谓）产生羞耻，同样，我们的一切都被赦免，成为同继承者，并享受如此伟大的爱，便产生放胆。\par

\Quote{因为进入至圣所。} 他这里所说的\Quote{进入}是什么意思？是指天堂，以及接近属灵的事物。\par

\Quote{祂所启用的，} 即祂所预备的、祂所开始的；因为从那时起，使用的开始就称为启用；祂所预备的（他意谓），并且祂亲自经过的那条路。\par

\Quote{一条又新又活的路。} 这里他表达了\Quote{望德的确证}。他说\Quote{新的}。他急于表明我们拥有更伟大的一切；因为如今天堂的门已经打开，这是连亚巴郎也未尝得到的。他说\Quote{又新又活的路}，因为先前的路是死亡之路，通向阴府，但这一条是生命之路。然而他没有说\Quote{生命}，却称之为\Quote{活}（即指那常存的定规）。\par

\Quote{从幔子经过}（他说）\Quote{就是祂的血肉。} 因为这血肉首先开辟了那条路，藉着它祂开启了祂所行走的路。他称[血肉]为\Quote{幔子}，是很有道理的。因为当它被高举时，天上的事物就显现了。\par

\Quote{让我们走近}（他说）\Quote{带着真诚的心}。 我们应该\Quote{走近}什么？走近圣物、信仰、属灵的敬拜。\Quote{带着真诚的心，在信德的确证下}，因为什么也看不见：从此既没有司铎，也没有祭品，也没有祭台。然而，那时的司铎也不可见，他站在里面，而所有人都站在外面，全体百姓都在外面。但在这里，不仅发生了司铎进入至圣所，我们也进入。因此他说：\Quote{在信德的确证下}。 因为疑惑的人可能以某种方式相信，正如甚至现在也有许多人说，有些人复活，有些人却不复活。但这并不是信德。他说：\Quote{在信德的确证下}；因为我们应当如同看见事物那样相信，不，甚至更当如此；因为\Quote{在这里}，在可见的事物上可能受骗，但在那里却不会；\Quote{在这里}，我们信赖感官，但在那里，我们信赖圣神。\par

\Quote{我们的心已洒净，脱离了邪恶的良心。} 他表明不仅需要信德，也需要有德行的生活，以及自觉没有邪恶。因为至圣所不会\Quote{在确证下}接纳那些没有这样预备的人。因为它们是圣的，是至圣所；但这里没有凡俗之人进入。他们是在身体上被洒净，我们是在良心上，以至于我们现在甚至可以被德行本身洒净。\Quote{身体也用清水洗净了。} 这里他指的是圣洗，这不再洁净身体，而是洁净灵魂。\par

\Quote{因为那应许的是信实的。} \Quote{那应许的}什么？我们即将离开此地，进入天国。所以不要过分好奇，也不要要求推理。我们的[信仰]需要信德。\par

% structure: p0013 source=h2 target=subsection reason=relative-heading-rank; base=h2

\subsection{希伯来书 10:24-25}

3. \Quote{又要}（他说）\Quote{彼此相顾，激发爱心，勉励行善。不可放弃我们自己的聚会，像有些人惯常这样，反当彼此劝勉，既看见那日子临近，就更当如此。} 又在别处说：\Quote{主已经近了；应当一无挂虑。}\BibleRef{斐 4:5} \Quote{因为现今我们的救恩比初信时更近了；从此时候不多了。}\BibleRef{罗 13:11}\par

什么是\BibleQuote{不要放弃我们自己的聚会}？\BibleRef{格前 7:29} 他知道聚集在一起能产生很大的力量。\BibleQuote{因为无论在哪里，有两三个人}（经上说）\BibleQuote{因我的名聚在一起，我就在他们中间}\BibleRef{玛 18:20}；又说：\BibleQuote{愿他们合而为一，如同我们也}\BibleRef{若 17:11}；以及：\BibleQuote{他们都是一心一［一］魂}\BibleRef{宗 4:32}。不仅如此，而且通过聚会，爱德得以增长；爱德增长了，属天主的事也必然随之而来。\BibleQuote{热切的祈祷}（经上说）\BibleQuote{由人民}\BibleRef{宗 12:5}。\Quote{如同有些人的习惯那样}。在这里，他不仅劝勉，也责备［他们］。\par

\BibleQuote{我们也当彼此顾念，}他说，\BibleQuote{激发爱德和善行。}他知道这也出自\BibleQuote{聚会。}因为如\BibleQuote{铁磨铁}\BibleRef{箴 17:17}，同样，交往也增长爱德。如果石头摩擦石头能产生火焰，何况灵魂与灵魂相融呢！但不是为了竞争，他说，而是\BibleQuote{激发爱德。}什么是\BibleQuote{激发爱德}？就是更多地爱与被爱。\BibleQuote{以及善行，}好使他们获得热忱。因为若实行比言谈更有教导力，你们中间也有许多教师，他们以行动实现了这一点。\par

什么是\BibleQuote{以真诚的心走近}？那就是没有虚伪；因为\BibleQuote{畏惧的心和发软的手有祸了}\BibleRef{德 2:12}：意思是，我们中间不可有虚伪；不可口是心非，因为这就是虚伪；也不可胆怯，因为这不是\BibleQuote{真诚的心}的标志。胆怯出自不信。但这如何成就呢？若我们藉着信德完全确信。\par

\BibleQuote{心已洒过}：他为什么不说\Quote{已被洁净}？［因为］他愿指出两种洒的区别：一种说是天主的，另一种是我们自己的。因为洗净并洒净良心是天主的事；但\BibleQuote{以}真理和\BibleQuote{在信德的完全确信中走近}是我们的事。然后他也从那应许者的真实坚固他们的信德。\par

什么是\BibleQuote{并用清水洗净身体}？是用那使之洁净的水；或是那没有血的水。\par

然后他加入那完美的事，就是爱。\BibleQuote{不要放弃我们自己的聚会，}有些人（他说）这样做，分裂聚会。因为\BibleQuote{兄弟帮助兄弟，有如坚城}\BibleRef{箴 18:19}。\par

\BibleQuote{但我们当彼此顾念，激发爱德。}什么是\BibleQuote{我们当彼此顾念}？例如若有人有德行，我们当效法他，看他以便去爱和被爱。因为善行出自爱德。聚会是一件大善事：因为它使爱更炽热；从爱生出一切美善。因为凡不是出于爱所做的，没有一件是善的。\par

4. 那么，让我们彼此\Quote{坚固}。\BibleQuote{因为爱是法律的满全}\BibleRef{罗 13:10}。如果我们彼此相爱，就不需要劳苦和流汗。这是一条自动导向德行的道路。就像在大道上，若有人找到了起点，就被它引导，不需要人拉着他的手；同样，对于爱，只要抓住起点，就立刻被它引导和指引。\BibleQuote{爱不加害于人}\BibleRef{罗 13:10}；\BibleQuote{不图谋恶事}\BibleRef{格前 13:5}。愿每个人自己反省，他是如何对待自己的。他不嫉妒自己；他愿为自己谋求一切美好；他把自己放在一切之前；他愿意为自己做一切事。如果我们对别人也如此，一切痛苦就结束了；没有仇怨；没有贪婪：因为谁会占自己的便宜呢？没有人；相反，我们会拥有一切共有的，并不停止聚会。如果我们这样做，怨恨记忆就没有位置：因为谁会让自己记恨自己呢？谁会对自己生气呢？我们岂不最宽容自己吗？如果我们对邻人也如此，就永不会有怨恨的记忆。\par

你（说）人怎么可能爱邻人如同自己？若别人没有做到，你固然可以认为不可能；但若有人做到了，显然是由于懈怠，才未被我们自己做到。\par

此外，基督命定的没有一件是不可能的，因为许多人甚至超过了祂的命令。谁做到了这一点？\Person{保禄}、\Person{伯多禄}和全体诸圣。其实，若我说他们爱了邻人，这不算什么：他们如此爱仇敌，无人会像爱自己那样爱同道。因为谁会为了同道的情面，在即将进入天国时，选择下地狱呢？没有人。但\Person{保禄}为了他的仇敌，为了那些用石头打他、鞭打他的人，选择了这个。那么，如果我们不向朋友表现出\Person{保禄}对敌人所表现之爱的极小部分，我们还能有何宽恕，有何借口呢？\par

在祂以前，真福梅瑟也甘愿为那些用石头砸他的仇敌而被涂抹于天主的书上。达味看见那些起来反对他的人被杀时，也说：\Quote{我是牧人，犯了罪，但这些羊做了什么？}\BibleRef{撒下 24:17}他手中握有撒乌耳时，也不肯杀他，反而救了他；那时他自己正处于危险之中。如果这些事是在旧盟约之下成就的，我们这些活在新盟约之下、连与他们同等的标准也达不到的人，还有什么借口呢？因为如果\Quote{除非你们的义德超过经师和法利塞人的义德，你们决不能进入天国}\BibleRef{玛 5:20}，那么当我们还不如他们时，又怎能进入呢？\par

5. \Quote{爱你们的仇人，}\BibleRef{玛 5:44}祂说。所以你应爱你的仇人：因为你行善不是为他，而是为你自己。何以见得？你正变得像天主。他若被你爱，并没有多大益处，因为他被一个同为奴仆的人所爱；但你若爱你的同为奴仆者，就获益良多，因为你正变得像天主。你看见你行善不是为他，而是为你自己吗？因为祂立定奖赏不是为他，而是为你。\par

那么，倘若他是恶的，你该怎么办（你说）？奖赏就更大。你甚至该为他的邪恶而感激他：即使他在受了万千恩惠之后仍是恶的。因为如果他不是极其邪恶，你的奖赏就不会极其增加；所以你用来不爱他的理由，即说他邪恶，正是爱他的理由。除去对手，你就除去了得冠冕的机会。你没看见那些运动员吗？他们装了沙袋来练习？但你不需要这样做。生活充满了操练你、使你强壮的事。你没看见树木，越是受风摇动，就变得越发坚固？我们若忍耐，也将变得强壮。因为经上说：\Quote{忍耐的人富于智慧，但心量狭小的人极其愚妄。}\BibleRef{箴 14:29}你看他对前者是多么称赞，对后者是多么谴责？\Quote{极其愚妄，}即非常愚蠢。所以我们不要彼此胆怯：因为这不是出于敌意，而是出于心量狭小。如果灵魂强健，它就能轻而易举地忍受一切，没有什么能使之沉没，反而会引领它进入宁静的海港。愿我们众人因我们的主耶稣基督的恩宠和慈爱，偕同祂，与圣父及圣神，享光荣、权能、尊威，自今至永远，永无穷尽。阿们。\par

\SourceRecord{Translated by Frederic Gardiner.；Nicene and Post-Nicene Fathers, First Series；Vol. 14.；Edited by Philip Schaff.；Buffalo, NY: Christian Literature Publishing Co.,；1889.；<http://www.newadvent.org/fathers/240219.htm>.}

% structure: p0001 source=h1 target=section reason=page-title

\section{\WorkTitle{致希伯来人书讲道集 第二十篇}}

% structure: p0002 source=h2 target=subsection reason=relative-heading-rank; base=h2

\subsection{\BibleRef{希 10:26-27}}

\BibleQuote{\Emph{因为如果我们明知真相之后，仍然故意犯罪，就不再有任何赎罪的祭品了，} \Emph{只有一种可怕的审判的期待，和那将要吞灭敌人的烈火。}}\par

1. 树木被栽种，得到一切其余的照料，以及园丁的双手和劳作，却不为这些劳作结出果实，就会被连根拔起，丢进火里。同样，在我们的照耀中也发生类似的事。因为当基督栽种了我们，我们享受了圣神的浇灌，然后却不结出果实；那么，火，甚至地狱的火，在等待着我们，以及不灭的火焰。\par

保禄于是先劝勉他们要以爱德和善工结出果实，并且从更慈祥的考量来敦促他们（这些考量是什么？就是我们得以进入至圣所，\BibleQuote{祂为我们开创的那条新路} \BibleRef{希 10:20}），然后又从更阴郁的方面同样做了，如此说话。因为他说了\BibleQuote{不要放弃我们自己的聚会，像有些人习惯的那样，而要彼此劝勉，并且越发这样，因为你们看见那日子临近} \BibleRef{希 10:25}，这足以安慰人，他又加上说，\BibleQuote{因为如果我们明知真相之后，仍然故意犯罪。} 他的意思是，需要善工，是的，非常大的需要。\BibleQuote{因为如果我们明知真相之后，仍然故意犯罪，就不再有任何赎罪的祭品了。} 你们曾被洁净；你们已从指控中被释放；你们成了儿子。如果你们再回到从前的呕吐物中，那么在另一边等待着你们的是绝罚、火，以及一切这类的东西。因为没有第二次祭献。\par

2. 此刻我们再次受到那些否认悔改的人，以及那些推迟领受圣洗之人的攻击。一些人说，他们领受圣洗并不安全，因为没有第二次赦免；另一些人断言，如果施予那些犯了罪的人奥迹是不安全的，如果没有第二次赦免。\par

那么，我们要对两者说什么呢？他没有否认悔改，也没有否认经由悔改的赎罪，更没有把跌倒的人推开、丢入绝望。他并非如此敌视我们的救恩；而是怎样呢？他否认了第二次圣洗。因为他没有说\Quote{不再有悔改}，也没有说\Quote{不再有赦免}，而是说\Quote{不再有}\Quote{祭献}，也就是说，不再有第二次十字架。因为他所说的祭献就是指这个。他说：\BibleQuote{因为祂藉一次的祭献，永远使被圣化的人得以成全} \BibleRef{希 10:14}，不像犹太人的仪式。因此他全篇如此多地论述祭献，说它是一次，且仅是一次；不但要表明在这点上它与犹太人的仪式不同，也要使人们更加坚定，使他们不再期待依照犹太法律再有别的祭献。\par

\BibleQuote{因为}，他说，\BibleQuote{如果我们故意犯罪。} 他多么倾向于宽恕。他说，\BibleQuote{如果我们故意犯罪，} 这样，那些并非故意犯罪的人就能得到宽恕。\BibleQuote{在得知真理之后}：他可能指关于基督，或关于一切教义。\BibleQuote{就不再有任何赎罪的祭品了，} 但那是什么呢？\BibleQuote{只有一种可怕的审判的期待，和那将要吞灭敌人的烈火。} 藉着\Quote{敌人}，他不仅指不信者，也指那些行事违背德性的人；或者，他也指同样的火要接纳那些家里的人，就像接纳\Quote{敌人}一样。然后，表达其吞噬的本性，他说，仿佛给它生命，\BibleQuote{那将要吞灭敌人的烈火。} 因为如同一只被激怒、极其凶猛残酷的野兽，不等到抓住某人吃掉就不会安息；同样，那火，如同被义愤所驱策，无论抓住什么都不会放手，而是吞灭、撕碎。\par

3. 接着他也加上威胁的理由，就是它有充分的根据，它是公正的；因为这有助于信心，当我们表明它是公正的时候。\par

% structure: p0010 source=h2 target=subsection reason=relative-heading-rank; base=h2

\subsection{\BibleRef{希 10:28}}

因为，他说，\BibleRef{希 10:28} \BibleQuote{凡违背梅瑟法律的，凭两三个见证人，就不得怜悯而死。} \BibleQuote{不得怜悯，} 他说；因此那里没有赦免，没有怜悯，尽管那是梅瑟的法律；因为他订立了其中的大部分。\par

\Quote{凭两三个}是什么意思？如果两三个人作证，他说，他们立即就受惩罚。\par

% structure: p0013 source=h2 target=subsection reason=relative-heading-rank; base=h2

\subsection{\BibleRef{希 10:29}}

如果这样，在旧盟约之下，当梅瑟的法律被藐视时，就有如此大的惩罚，\BibleRef{希 10:29} \BibleQuote{你们想，那践踏了天主子，将那盟约的血视为普通之物，又侮慢了恩宠之圣神的人，应受怎样更严厉的惩罚呢？}\par

一个人如何\Quote{践踏天主子}？当他在奥迹中分享祂时（他可能会说）却犯了罪，难道不是践踏了祂吗？难道不是藐视了祂吗？因为正如我们不把那些被践踏的人当回事，同样，那些犯罪的人也不把基督当回事；所以他们犯了罪。你成了基督的身体，却把自己交给了魔鬼，以致它践踏你。\par

\BibleQuote{且将血视为普通之物，} 他说。什么是\Quote{普通}？就是\Quote{不洁}，或者指它并不比其他事物优越。\par

\BibleQuote{又侮慢了恩宠之圣神。} 因为那不接受恩惠的人，就是对施惠者侮慢。祂使你成为儿子：你却想成为奴隶。祂来与你同住：你却向祂引入邪念。基督愿意与你同住：你却藉着暴食、醉酒践踏祂。\par

凡不相称地领受圣事的人，请听；凡不相称地接近那祭台的人，请听。\BibleQuote{不可把圣物给狗}（祂说）\BibleQuote{免得日后它们践踏在脚下}\BibleRef{玛 7:6}，也就是免得他们轻视，免得他们唾弃。然而祂说的不是这个，而是比这更可怕的。因为祂用可怕的事来约束他们的灵魂。因为这事也适于使人皈依，并不亚于安慰。同时，祂既显示了区别，也显示了惩罚，并摆明了他们的审判，好像是一件明显的事。\Quote{你们想，那该受怎样更重的惩罚呢？}这里，在我看来，祂也暗示了圣事。\par

% structure: p0019 source=h2 target=subsection reason=relative-heading-rank; base=h2

\subsection{希伯来书 10:30-31}

4. 然后祂引用证据说：\BibleRef{希 10:30}\BibleQuote{落在永生天主手中，真是可怕！}\Quote{因为}经上记载：\BibleQuote{报复归我，我要还报，主说。又说：主要审判自己的百姓。}\Quote{我们落在主手中，总比落在人手中好。}\BibleRef{德 2:18} 但你们若不悔改，必要\BibleQuote{落在}天主\BibleQuote{手中}：那是可怕的；落在人手中，算不得什么。当他说，我们看见有人在此受罚时，不要为眼前的事惊慌，而要为将来的事战栗。\BibleQuote{因为按祂的仁慈，祂的义怒也是如此。} 且\BibleQuote{祂的愤怒必落在罪人身上。}\BibleRef{德 5:6}\par

同时，他也暗示了别的事。因为他说：\BibleQuote{报复归我，我要还报。}这是针对他们的仇敌说的，这些仇敌正在作恶，而不是针对那些受苦的人。这里他也在安慰他们，几乎是在说：天主永远存在且活着，所以即使他们现在得不到（赏报），将来必会得到。他们理应哀叹，而不是我们；因为我们的确要落在他们手中，但他们却要落在天主手中。因为受苦的人并未受害，而是那作恶的人受害；也不是那受惠的人受益，而是施惠的人受益。\par

5. 既然知道这些事，让我们在受苦上忍耐，在行善上积极。如果我们轻视财富和尊荣，就能做到这一点。那摆脱了这些情感的人，是所有人中最慷慨的，甚至比那穿紫袍的人更富有。你没有看见金钱带来多少祸患吗？我不是说贪婪带来多少，而只是说我们依恋这些东西。例如，如果一个人失去了钱财，他过的生活比任何死亡都更痛苦。你为什么悲伤，人啊？你为什么哭泣？因为天主使你从过度的看守中解脱出来？因为你不必坐立不安、战战兢兢？再者，若有人把你锁在宝藏旁，命令你永远坐在那里，为别人的财物看守，你会悲伤，会厌恶；而你，在给自己加上最沉重的锁链之后，竟然在从奴役中解脱时悲伤？实在，忧愁与欢乐都是幻想的。因为我们守护这些财物，好像是守护别人的东西一样。\par

现在我的讲论是针对妇女的。一个女人常有一件金线织成的衣裳，她把它抖开，包在细麻布中，小心保管，为它战栗，却无法享受它。因为她要么死去，要么成为寡妇。或者，即使这些事都没有发生，但因害怕经常使用而穿坏，从而剥夺自己享用它的机会，她便以另一种方式剥夺了自己，即节省不用。但是（你说）她把它传给另一个人。但这也不一定：即使她传给了别人，那另一个人也会同样使用它。若有人搜查她们的房屋，他会发现，最昂贵的衣裳和其他精选之物，被特别尊敬地对待，好像它们是活的主人一样。因为她不习惯使用它们，而是恐惧战栗，驱赶蛾子和其它会咬坏它们的东西，并把大部分放在香水和香料中，也不让所有人都配得见到它们，反而常常和丈夫一起亲自仔细整理它们。\par

请告诉我：保禄不是有理由称贪婪为\BibleQuote{偶像崇拜}吗？\BibleRef{哥 3:5} 因为这些女子对待她们的衣裳、她们的黄金，正如同对待她们的偶像一样尊敬。\par

6. 我们要搅动淤泥到几时呢？我们要粘在泥土和造砖上到几时呢？因为正如他们为埃及王劳苦，我们也为魔鬼劳苦，且遭受更严厉的鞭打。因为灵魂超越身体多少，忧虑也超越鞭笞的伤痕多少。我们每日被鞭打，充满恐惧、焦虑和战栗。但是，如果我们叹息，如果我们仰望天主，祂就派遣给我们，不是梅瑟，也不是亚郎，而是祂自己的圣言和痛悔。当这圣言来到，抓住我们的灵魂时，祂将我们从痛苦的奴役中释放出来，祂将我们领出埃及，脱离无益的虚荣，脱离无利的奴役。因为他们的确得到了至少金饰作为建造的工价后才出来，但我们却一无所获；但愿真的是一无所获！因为实际上我们也得到，不是金饰，而是埃及的罪恶，即罪恶和惩罚与刑罚。\par

那么，让我们学会被人利用，学会被人恶意对待；这是基督徒的本分。让我们轻视金衣，轻视金钱，免得我们轻视自己的救恩。让我们轻视金钱，而不轻视灵魂。因为灵魂受责罚，灵魂受惩罚；那些东西留在这里，而灵魂却前往那边。请告诉我，你为什么将自己撕成碎片，却不察觉呢？\par

7. 这些话我是对那些贪得无厌的人说的。对受欺压的人也该说几句。要慷慨地忍受他们的欺压；他们是在毁灭自己，而不是你。他们确实骗取了你的钱财，却剥夺了自己天主的恩宠和援助。一个人若失去这些，即使披上全世界的财富，也是众人中最贫穷的；而最贫穷的人若拥有这些，就是最富有的。因为\BibleQuote{主（经上说）是我的牧者，我必一无所缺}\BibleRef{咏 22:1}。\par

告诉我，如果你有一个丈夫，伟大而可敬，深切地爱你、照顾你，而且你知道他会永远活着，不会死在你前头，并让你安全地享用一切，如同己有：你还会想要拥有任何东西吗？即使你被剥夺了一切，难道你不认为自己因此更富有吗？\par

那么你为什么悲伤？因为你没有财产？但你要想，你已失去了犯罪的机会。或者因为你曾经拥有而被剥夺了？但你已获得了天主的恩宠。你怎么获得的呢（你说）？祂说过：\BibleQuote{你们为什么不宁愿受亏负呢？}\BibleRef{格前 6:7} 祂说过：\Quote{那些以感恩之心忍受一切的人是有福的。}所以你要想想，如果你以行动表明这些，你将享有何等大的恩宠。因为只要求我们一件事：\Quote{事事感谢}天主，我们就拥有一切丰盛。例如：你失去了一万磅金子？立刻感谢天主，你就通过那句话和感恩获得了十倍万金。\par

8. 告诉我，你什么时候认为约伯是有福的？是他拥有那么多骆驼、羊群和牛群的时候，还是他说出那句话的时候：\BibleQuote{主赐予，主收回。}\BibleRef{约 1:21} 所以魔鬼使我们遭受损失，不仅是为了夺走我们的财物（他知道这算不了什么），更是为了通过这些迫使我们说出亵渎的话。就像在约伯身上，他所求的不仅是使他变穷，更是使他成为亵渎者。总之，当魔鬼剥夺了他一切之后，看看他是怎样通过妻子对他说的：\BibleQuote{你说一句反对主的话，然后死去吧。}\BibleRef{约 2:9} 可是，你这被诅咒的，你已经剥夺了他一切。但（他说）这并不是我所追求的，因为我还没有达到我所做一切的目的。我追求的是剥夺天主的助佑；正是为此我也剥夺了他的财物。这才是我想要的，其他的算不了什么。如果这一点没有达到，他不仅没有受到任何损害，反而得到了益处。你看到吗？就连那邪恶的魔鬼也知道在这件事上损失有多大。\par

请注意他怎样通过妻子设下奸计。你们作丈夫的，要听这话，尤其是那些有贪财的妻子、迫使你们亵渎天主的人。要想想约伯。但是，如果你愿意，让我们看看他多么有节制，怎样使她闭口无言。他说：\BibleQuote{你为什么说话像愚妇一样？}\BibleRef{约 2:10} 诚然，\BibleQuote{邪曲的交谈会败坏良善的风俗}\BibleRef{格前 15:33}，无论在什么时候，尤其是在患难中；那时，进谗言的人很有力量。如果灵魂本身就容易急躁，更何况还有劝诫者呢？岂不是跌入深坑？妻子是很大的福，也是很大的祸。因为妻子是很大的（福），看看他从（撒殚）那里怎样试图攻破坚固的墙。他说：'剥夺他的财产没有把他击倒；损失没有产生多大效果。'因此他说：\BibleQuote{看他是否当面诅咒你。}\BibleRef{约 2:5} 你看他向往的是什么。\par

如果我们感恩地承受（损失），我们甚至会恢复这些东西；即使不恢复，我们的赏报会更大。因为当约伯英勇奋斗后，天主也将这些东西归还给了他。当祂向魔鬼表明，约伯不是为了这些东西而事奉祂时，祂也将这些东西归还给了他。\par

9. 因为天主就是这样。当祂看到我们没有执着于今生的事物时，祂就把它们赐给我们。当祂看到我们更重视属灵的事物时，祂也会赐给我们属肉体的事物。但不是先赐予，免得我们脱离属灵的事物；祂为了保全我们，不赐予肉体的事物，好让我们远离它们，即使我们不情愿。\par

并非如此（你说），如果我得到了（它们），我就满足了，并且更加感恩。这是假的，人啊，因为到那时你尤其会变得轻率。\par

那么为什么（你说）祂把（它们）赐给许多人呢？从哪里看得出是祂赐予的呢？但除了祂还有谁呢（你说）？他们的贪婪，他们的掠夺。那么天主为什么允许这些呢？正如祂（允许）谋杀、盗窃和暴力一样。\par

那么（你会说），那些通过继承从父亲那里获得遗产的人，他们自己却充满了无数的罪恶，这又怎么说呢？这又如何呢？天主怎会容忍他们（你说）享受这些呢？当然，就像祂容忍盗贼、杀人犯和其他作恶的人一样。因为现在不是审判的时候，而是选择最佳生活方式的时候。\par

我刚才所说的话，我再重复一遍：那些享受了一切福乐却仍不变得更好的人，将受更重的惩罚。因为不是所有的人都受同样的惩罚；而那些在天主的恩惠之后仍然作恶的人，要受更重的惩罚；而那些在贫穷之后[这样做的人]则不然。知道这是真的，请听祂对\Person{达味}所说的话：\BibleQuote{我岂没有将你主人的一切财产都赐给你吗？} \BibleRef{撒下 12:8} 所以，每当你看到一个年轻人不劳而获地继承了父家的遗产，却仍然作恶，你就该确信，他的惩罚增加了，报复更加严厉。那么，我们不要效法这些人；但若有人获得了美德，若有人得到了属灵的财富，[就让我们效法他吧]。因为（经上记载）：\BibleQuote{那些信赖自己财富的人有祸了！} \BibleRef{咏 48:6}；\BibleQuote{敬畏主的人是有福的。} \BibleRef{咏 127:1} 告诉我，你要属于哪一类呢？无疑是那些被称为有福的人。因此，要效法这些人，而不是那些人，好使你也获得为他们预备的福乐。愿我们大家都能获得，因我们的主耶稣基督，祂与圣父及圣神，同享光荣，从现在直到永远，世世无穷。阿们。\par

\SourceRecord{Translated by Frederic Gardiner.；Nicene and Post-Nicene Fathers, First Series；Vol. 14.；Edited by Philip Schaff.；Buffalo, NY: Christian Literature Publishing Co.,；1889.；<http://www.newadvent.org/fathers/240220.htm>.}

% structure: p0001 source=h1 target=section reason=page-title

\section{致希伯来人讲道集 第21篇}

% structure: p0002 source=h2 target=subsection reason=relative-heading-rank; base=h2

\subsection{希伯来书 10:32-34}

\BibleQuote{ \Emph{但你们要追忆往昔的日子，那时你们蒙受了光照以后，忍受了极大苦难的搏斗；} \Emph{一方面，你们由于讥讽和磨难而成了众人注目的对象；} \Emph{另一方面，你们成了那些同样受苦之人的同伴。} \Emph{因为你们同情了那些被捆绑的人，} \Emph{并且欣然接受了你们财物的被抢掠，知道你们在天上} \Emph{有更美好且永存的产业。}}\par

1. 最高明的医生在深切口、并以伤口加剧痛苦之后，在安抚患处、给受扰的灵魂带来休息和振奋时，不会立即进行第二次切口，而是以温和的、适合消除剧痛的药膏来安抚已切开的伤口。保禄在震撼了他们的灵魂、以对地狱的回忆刺痛了他们、并使他们确信轻视天主恩宠的人必会灭亡之后，也从梅瑟的法律中证明了他们也必灭亡，且更可怕，又用其他见证加以确认，并说过：\BibleQuote{落在永生天主手中真是可怕的事。}\BibleRef{希 10:31} 然后，惟恐灵魂因过度恐惧而沮丧，被忧伤所吞噬，他便通过称赞和劝勉来安抚他们，并以他们自己的行为给予他们热忱。因为他说道：\BibleQuote{你们要追忆往昔的日子，那时你们蒙受了光照以后，忍受了极大苦难的搏斗。} 从已经成就的事迹而来的劝勉是强有力的：因为开始一项工作的人应当继续前进并增加它。仿佛他说：\Quote{当你们被引入（教会）、处于学习者的行列时，你们显示了如此大的热忱、如此大的高尚；但现在却不再如此了。} 而激励者特别以自己的榜样来激励他们。\par

他并没有简单地说\Quote{你们忍受了搏斗}，而是说\Quote{极大的}搏斗。此外，他没有说\Quote{试探}，而是说\Quote{搏斗}，这是称赞和极大赞美的表达。\par

然后他还特别列举出来，扩大他的话语，增加称赞。怎样列举呢？\BibleQuote{一方面，你们由于讥讽和磨难而成了众人注目的对象}（他说）；因为讥讽是件大事，足以败坏灵魂，蒙蔽判断。请听先知所说的：\BibleQuote{他们终日对我说：\InnerQuote{你的天主在哪里？}}\BibleRef{咏 41:10} 又说：\BibleQuote{若是仇敌辱骂我，我还能忍受。}\BibleRef{咏 54:12} 因为人类极其虚荣，所以很容易被此胜过。\par

他并没有简单地说\BibleQuote{由于讥讽}，而是说甚至带着极大的强度，\BibleQuote{成了众人注目的对象}。因为当一个人在私下被讥讽时，固然痛苦，但在众人面前则更甚。请告诉我，当那些离开了犹太教的卑贱、仿佛转向了最佳生活方式、轻视了祖传习俗的人，被自己的人民虐待、且无人帮助时，那祸患是何等之大。\par

2. 我无法说（他说）你们确实受了这些苦并悲伤，而你们甚至大大喜乐。这一点他以\BibleQuote{你们成了那些同样受苦之人的同伴}来表达，并且他举出了宗徒们本人。他的意思是：你们不仅不以自己的受苦为耻，甚至与那些遭受同样事的人分享了。这也是激励者的话语。他没有说：\Quote{承担我的苦难，与我分享}，而是尊重你们自己的。\par

\BibleQuote{你们同情了那些被捆绑的人。} 你看，他是在谈论他自己和其他在监狱中的人。这样，你们并不把\Quote{捆绑}视为捆绑；而是像高贵的角力者那样站立：因为你们不仅在自己的困苦中不需要安慰，甚至成了别人的安慰。\par

并且\BibleQuote{你们欣然接受了你们财物的被抢掠。} 哦！何等\BibleQuote{信德的确证}\BibleRef{希 10:22}！然后他也提出了动机：不仅是为他们的搏斗而安慰他们，也是使他们不致被动摇而离开信德。当你们看到财产被抢掠时（他的意思是），你们忍受了；因为你们已经看见了那看不见的，如同看见的：这是真正信德的效果，并且你们以自己的行为显示了它。\par

那么，抢掠或许是由于抢掠者的暴力，无人能阻止；因此至今尚不清楚你们是为了信德而忍受抢掠。（尽管这也是清楚的。因为如果你们选择不信，你们有能力不被抢掠。）但你们所做的远比这更伟大：甚至\Quote{欣然}忍受了这些事；这完全是出于高尚的灵魂，他们因被鞭打而喜乐。因为经上记载：\BibleQuote{他们离开公议会，心里欢喜，因被算是配为这名受辱。}\BibleRef{宗 5:41} 但\Quote{欣然}忍受的人表明他有某种赏报，并且这事不是损失，而是收益。\par

此外，\Quote{你们取了}一词显示了他们自愿的忍受，因为他的意思是：你们选择了并接受了。\par

\BibleQuote{知道你们在天上有更美好且永存的产业}（他说）；即，坚实的，不像此世的会朽坏。\par

% structure: p0014 source=h2 target=subsection reason=relative-heading-rank; base=h2

\subsection{希伯来书 10:35}

3. 随后，在赞美他们之后，他说：\BibleRef{希 10:35} \BibleQuote{所以，你们不要丧失那使你们得到大赏报的信赖。} 他是什么意思？他没有说，\Quote{你们已经丢弃了它，又恢复它}；而是，更增强他们的话，\Quote{你们拥有它}。因为再次恢复那已丢弃的，需要更多劳苦；但不失去那紧握的，则不然。但致迦拉达人时，他说了完全相反的话：\BibleRef{迦 4:19} \BibleQuote{我的子女！我为你们再受产痛，直到基督在你们内形成}；所以有理由；因为他们更懈怠，故需要更严厉的话；但这些人心灰意冷，所以更需要安慰的话。\par

\Quote{所以不要丧失}（他说）\Quote{你们的信赖}，这样他们对天主有极大的信赖。\Quote{它拥有}（他说）\Quote{极大的赏报。}\Quote{我们何时接受到它们（有人可能说）？看！我们方面的一切都已完成。}因此，他按照他们自己的假设预见了他们，意思是：如果你们知道在天上拥有更好的产业，就不要在这里寻求任何东西。\par

\Quote{因为你们需要忍耐}，不是任何[劳苦的]添加，好使你们继续在同样状态，免得丢弃那已交付你们手中的。你们不需要别的，只需要像你们已站立的这样站立，好使你们到达终点时，能领受恩许。\par

% structure: p0018 source=h2 target=subsection reason=relative-heading-rank; base=h2

\subsection{希 10:36}

\BibleRef{希 10:36} \BibleQuote{因为（他说）你们需要忍耐，以便在遵行天主旨意后，能领受恩许。} 你们只需要一件事，就是忍受延迟；不是要再争战。你们已接近冠冕（他意指)；你们已忍受了一切搏斗：锁链、迫害；你们的财产被抢夺。然后呢？从此你们站着等待加冕：只忍耐这事，冠冕的延迟。啊，安慰何其伟大！就像对一位已击败所有对手、再无对手的运动员说话，然后他将要戴冠冕，却等不及比赛裁判到来并为他戴上花冠的片刻；他不耐烦，想要走出去、逃离，仿佛不能忍受干渴和炎热。\par

% structure: p0020 source=h2 target=subsection reason=relative-heading-rank; base=h2

\subsection{希 10:37}

然后，他也暗示这一点，他怎么说？\BibleRef{希 10:37} \BibleQuote{还有片刻，那将要来的必来，决不迟延。} 免得他们说：\Quote{他何时来？}他用圣经安慰他们。因为同样，当他在另一处说 \BibleRef{罗 13:11} \BibleQuote{现在我们的救恩更近了} 时，他安慰他们，因为剩余的时间是短暂的。他说这话不是出于自己，而是出于圣经。但如果从那时就说了\Quote{还有片刻，那将要来的必来，决不迟延}，那么现在显然他更近了。因此，等待并非小赏报。\par

% structure: p0022 source=h2 target=subsection reason=relative-heading-rank; base=h2

\subsection{希 10:38}

\BibleRef{希 10:38} \BibleQuote{现在那义人（他说）将因信德而生活；但若有人退缩，我的灵魂不喜悦他。} 这是一个极大的鼓励，当人指出他们已在全局中成功，却因一点懈怠而失去。\BibleRef{希 10:39} \BibleQuote{但我们不是退缩以致丧亡的人，而是有信德以致灵魂的救恩的人。}\par

% structure: p0024 source=h2 target=subsection reason=relative-heading-rank; base=h2

\subsection{希 11:1-2}

4. \Quote{现在信德是所希望之事的实体，是未见之事的确证。因这信德，先人们都得了称赞。}哦，他用了多么强而有力的表达，说\Quote{未见之事的确证}。因为[我们说]在非常明显的事物中有\Quote{确证}。那么信德就是看见不明显的事（他意指），并使未见之事与可见之事具有同样的确据。因此，既不可能不信可见之事，另一方面，也不能有信德，除非人对不可见之事比最清楚可见之事更加确信。因为希望的对象似乎无实体，信德给它们实体，或者更确切地说，并非给予实体，而是它本身就是它们的实体。例如，复活尚未到来，也不实质性地存在，但希望在灵魂中使其成为实体。这就是\Quote{事物的实体}的意思。\par

因此，如果它是\Quote{未见之事的确证}，为何你们竟想看见它们，以至于从信德和成为义人上堕落？既然\Quote{义人将因信德而生活}，而你们，如果你们想看见这些事，就不再是信者。你们已经劳苦（他说），你们已经奋斗：我也承认这点，尽管如此，等待吧；因为这就是信德：不要寻求\Quote{这里}的全部。\par

5. 这些话确实是对希伯来人说的，但也是对许多在此聚集之人的一般劝勉。如何并以什么方式？对心灰意冷的人；对气馁的人。因为当他们看到恶人亨通，自己却遭厄运时，他们感到困扰，不耐烦地忍受；他们渴望惩罚和报复他人；同时等待自己受苦的赏报。\Quote{因为还有片刻，那将要来的必来。}\par

那么让我们对懈怠者这样说：无疑将有惩罚；无疑他将到来，此后复活的事件甚至就在门口。\par

从何处可见此事（你说）？我并不是说，从先知；因为我现在也不是只对基督徒说话；但即使有外邦人在此，我完全有信心，并拿出证据，指示他。如何（你说）？\par

基督预言了许多事。如果那些先前的事没有应验，那么不要相信它们；但如果它们都已应验，为何怀疑那些剩余的呢？事实上，完全没有应验就相信其一，或当一切都应验了却不信其余，都是极不合理的。\par

但我将用一个例子使这事更清楚。\Person{基督}曾说，\Place{耶路撒冷}将被攻陷，且要被如此攻陷，如同以前没有一座城那样，并且它永远不会再被重建：这预言确实应验了。祂说，将有\Quote{大灾难}\BibleRef{玛 24:21}，也成就了。祂说，一粒芥菜子被撒下，\Quote{福音的宣讲}也要如此扩展：我们每天都看到这传到全世界。祂说，凡离开父亲或母亲、兄弟或姊妹的，必得着许多父亲和母亲：这事我们看见藉事实应验了。祂说：\Quote{在世上你们有苦难，但你们放心，我已经胜了世界}\BibleRef{若 16:33}，也就是说，没有人能胜过你们。我们看见这事藉事件成就了。祂说：\Quote{阴间的门也不能胜过教会}\BibleRef{玛 16:18}，即使受迫害，也没有人能熄灭\Quote{福音的宣讲}：事件的经历也证实了这个预言：然而当祂说这些话时，很难相信祂。为什么？因为这些都是言语，祂尚未对所讲的事提出证据。所以如今它们变得远更可信。祂说：\Quote{当福音传遍万民时，然后结局才来到}\BibleRef{玛 24:14}；看！如今你们已到了终点：因为世界大部分已听过宣讲，所以结局就在眼前。让我们战栗，亲爱的。\par

6. 但是，请告诉我，你担心结局吗？结局本身确是临近，但各人的生命与死亡更近。因为经上说：\Quote{我们一生的年日是七十岁，若是强壮可到八十岁。}\BibleRef{咏 89:10}审判的日子近了。让我们敬畏。\Quote{弟兄不能赎回，人岂能赎回？}\BibleRef{咏 48:7}在那里我们将大为忏悔，\Quote{但在死亡中无人赞美祂。}\BibleRef{咏 6:5}因此他说：\Quote{让我们以感谢来到祂面前}\BibleRef{咏 93:2}，即祂的来临。因为在此世，无论我们做什么都有效验；但那里，不再如此。请告诉我，如果有人将我们暂时放在火窑里，我们岂不忍受一切以逃脱，即使要舍弃钱财，甚至卖身为奴？有多少人陷入重病，宁愿舍弃一切以得解脱，若给他们选择？若今世短短的病痛就这样折磨我们，到那时，我们又将如何，那时忏悔也无济于事？\par

7. 我们现在充满多少邪恶，而不自知？我们彼此相咬、相吞，在冤枉、控告、诽谤、因邻人的声望而烦恼中。\BibleRef{迦 5:15}\par

且看这困难。当人想破坏邻人的名誉时，他说：\Quote{某人说了关于他的这话；\Person{天主}啊，原谅我，不要严审我，我必须为所听见的交账。}那你为何要讲它，若你不相信？你为何讲它？你为何因多次传播而令它可信？你为何传那不实的故事？你不相信，却恳求\Person{天主}不严审你？那就别说，保持沉默，你便摆脱了一切恐惧。\par

但我不知这病从哪里降到人身上。我们成了饶舌者，心中毫无保留。听智者的话说：\Quote{你听到一句话，让它死在你内心，放胆吧，它不会涨破你。}\BibleRef{德 19:10}又说：\Quote{愚人听到一句话，就痛苦如临产的妇女。}\BibleRef{德 19:11}我们准备好指控，准备好定罪。即使我们没有犯别的恶行，这已足以毁灭我们，将我们带入地狱，这使我们陷入千万邪恶。为使你确切知道这点，听先知说：\Quote{你坐着毁谤你的兄弟。}\BibleRef{咏 49:20}\par

但你说：不是我，而是那人\Quote{[告诉我的]}。不，是你自己；因为\Emph{你}若不说，别人就不会听见；即使他听见了，你也不会为那罪负责。我们本应遮掩并隐藏邻人的过失，但你却假借热心的外衣来宣扬它们。你成了，不是控诉者，而是长舌者、轻浮者、愚人。哦，何等聪明！你未意识到，你既羞辱了那人，也羞辱了你自己。\par

且看由此产生了多大的恶。你惹动了天主的义怒。你岂没听见\Person{保禄}关于寡妇说：\Quote{她们不但}（这是他原话）\Quote{学习懒惰，而且成了长舌者、好管闲事者，挨家挨户闲游，说些不该说的话。}\BibleRef{弟前 5:13}所以即使你相信别人对你兄弟所说的事，你也不该说它；更何况当你不相信时呢。\par

但你还顾全自己的利益？你怕受\Person{天主}审问？那就怕因你的饶舌而被审问吧。因为在这里，你不能说：\Quote{\Person{天主}，不要因轻浮话语审问我}：因为整件事是轻浮话语。你为何传布它？你为何增加那恶？这足以毁灭我们。为此\Person{基督}说：\Quote{你们不要判断人，免得你们被判断。}\BibleRef{玛 7:1}\par

但我们不留意这事，也没从\Person{法利塞人}的事吸取教训。他说了真话：\Quote{我不像这个税吏}\BibleRef{路 18:11}，而且还是在无人听见时说的；但他却被定罪。若他在说真话、无人听见时仍被定罪，那些像长舌妇一样，到处传播连自己也不相信的谎话的人，将受何等可怕的惩罚？他们岂能逃脱？\par

8. 从此以后，让我们在口前设置\Quote{门和门闩}。（\BibleRef{德 28:25}）因为闲谈产生了无数的恶事；家庭被毁，友谊破裂，无数其他的灾祸发生。人啊，不要忙碌于你邻人的事务。\par

但你多言多语，且有此软弱。向天主谈论你自己的[过犯]吧：如此这软弱不再是软弱，反而成为益处。向你的朋友，向那些彻底的知己、义人、你所信赖的人谈论你自己的[过犯]，好使他们为你的罪祈祷。若你谈论他人的[罪]，你毫无获益，一无所获，反而毁了自己。若你向主承认自己的[罪]，你将获得巨大的赏报：因为经上说：\Quote{我说，我要向主承认我的过犯，你赦免了我心中的不敬。}（\BibleRef{咏 31:5}）\par

你愿意审判吗？审判你自己的[罪]吧。没有人会控告你，若你定自己的罪；但若你不定罪，他就要控告；他必控告你，除非你定罪自己；他将控告你麻木不仁。你看见某人发怒、恼怒、做出其他不合时宜的事？立刻想到你自己的[过犯]：这样你就不会过于谴责他，也使自己从过去过犯的重担中解脱。若我们如此规范自己的行为，如此管理自己的生活，若我们定自己的罪，就很可能不会犯许多罪，并且会做许多善事，成为公正温和的人；并享受一切应许给爱天主之人的恩惠：愿我们众人因着我们的主耶稣基督的恩宠和慈爱而获得这一切，愿光荣、权能和尊威归于祂，偕同圣父及圣神，从今直到永远，世世无尽。阿们。\par

\SourceRecord{Translated by Frederic Gardiner.；Nicene and Post-Nicene Fathers, First Series；Vol. 14.；Edited by Philip Schaff.；Buffalo, NY: Christian Literature Publishing Co.,；1889.；<http://www.newadvent.org/fathers/240221.htm>.}

% structure: p0001 source=h1 target=section reason=page-title

\section{\BibleBook{}{希伯来书}第二十二篇讲道}

% structure: p0002 source=h2 target=subsection reason=relative-heading-rank; base=h2

\subsection{\BibleRef{希 11:3-4}}

\BibleQuote{\Emph{因着信德，我们就明白宇宙是由天主的话造成的，这样，看得见的是由看不见的成的。} \Emph{因着信德，亚伯尔向天主奉献了比加音更美善的祭献，藉这祭献他得了见证，称他为义人，因为天主为他的礼物作了证；他虽然死了，却因这信德仍旧说话。}}\par

1. 信德需要一个慷慨而刚毅的灵魂，一个超越一切感官事物、并胜过人类理性软弱之灵魂。因为除非人提升自己超过世俗的常规，否则不可能成为信徒。\par

既然希伯来人的灵魂已经完全软弱了，他们虽然始于信德，但因环境——我指的是苦难、灾难——后来变得灰心丧气、精神不振，并从他们的立场摇动了，他便首先从这些事本身鼓励他们，说：\BibleQuote{你们要追忆往昔的日子}\BibleRef{希 10:32}；然后从圣经说：\BibleQuote{义人将因信德而生活}\BibleRef{希 10:38}；之后从论证说：\BibleQuote{信德是所希望之事的担保，是未见之事的确证}\BibleRef{希 11:1}。现在又从他们的祖先——那些伟大而可钦佩的人——出发，仿佛在说：如果当美善近在眼前时，所有人都因信德而得救，那么我们更当如此了。\par

因为当灵魂找到一个与自己同受苦难的人时，它便得安慰、恢复气息。这我们既可在信德上看到，也可在苦难中看到：\BibleQuote{为叫你们藉着我们彼此共有的信德获得安慰}\BibleRef{罗 1:12}。因为人类非常不自信，不能信赖自己，对所拥有的一切都战战兢兢，并高度重视多数人的意见。\par

2. 那么保禄做了什么？他用祖先来鼓励他们；此前则用普遍的常识。因为请告诉我，由于信德被毁谤为没有证明、更像是一种欺骗，所以他指出最伟大的事物是通过信德而非理性的推论获得的。那么，他如何指出这一点呢？请告诉我。他说得很明显：天主从无中造出了有，从不显现的造出了显现的，从不存在的造出了存在的。但祂\Quote{藉着圣言}这样做了，这从哪里显明呢？因为理性不会提示这类事；恰恰相反，理性提示：显现的事物是从显现的事物产生的。\par

因此，哲学家们明确地说：\Quote{无中不能生有}，他们是\Quote{属血气的}\BibleRef{犹 1:19}，丝毫不信任信德。然而，当这些人碰巧说出任何伟大而高尚的话时，就被发现是把这话交给信德。例如：\Quote{天主无始、无生}；因为理性并不提示这一点，反而提示相反。请想一想，我求你们，他们的大愚昧。他们说天主无始；然而这比从无中造出万物更加惊人。因为说祂无始、无生，既非由自己生，也非由祂而生，比说天主从无中造出了现存之物更加充满困难。因为在这方面有许多不确定：例如，有人造了它，被造的有起始，总之，它是被造的。但在另一方面呢？祂自存、无生，祂既无起始也无时间；请告诉我，这些事岂不需要信德吗？但他没有断言这一点，那远为重大，而是断言了那较小的。\par

那么，从哪里显明天主造了这些呢？他会说：理性并未提示；没有人在场看着这事发生。从哪里显明呢？这显然是信德的结果。\BibleQuote{因着信德，我们明白诸世界是造成的。}为何\BibleQuote{因着信德}？因为\BibleQuote{看得见的是由看不见的成的}。这就是信德。\par

3. 陈述了这普遍原则之后，他随后用个人来检验它。因为一个知名的人相当于整个世界。他后来至少暗示了这一点。因为当他把世界与一二百人对等，然后看到数量之小时，他随后说：\BibleQuote{世界是不配他们的}\BibleRef{希 11:38}。\par

且看他把谁放在第一位：是那受了虐待的人，而且是被兄弟虐待。这正是他们自己的苦难，\BibleQuote{因为你们也受了本地人的苦害}\BibleRef{得前 2:14}。并且是被一个没有受任何委屈、却因天主的事而嫉妒他的兄弟所害；表明他们也是被人以恶眼相看、遭嫉妒的。他尊敬了天主，并因尊敬天主而死；尚未获得复活。但他的预备是明显的，他的部分已完成，但天主的部分尚未向他实施。\par

此处所称的\Quote{更美善的祭献}是指更尊贵、更辉煌、更必要的祭献。\par

我们不能说（他说）那祭献没有被接受。天主确实接受了它，并对加音说：\BibleQuote{你若是做得对，岂不抬起头？但若做得不对，罪就伏在门前。}\BibleRef{创 4:7} 所以亚伯尔既做得对，也分得对。然而为此他得了什么报酬？他被他兄弟的手杀死了；他父亲因犯罪所受的判决，现在这正直的人首先承受了。他因这来自兄弟的伤害而受苦更甚，并且他是第一个受苦的人。\par

他做这些事做得对，不仰望任何人。因为当他如此尊敬天主时，他还能仰望谁呢？他的父亲和母亲吗？但他们曾以恶报善，辜负了祂的恩惠。那么仰望他的兄弟吗？但兄弟也羞辱了天主。所以他是独自寻求了美善。\par

及那堪受如此大荣的人，究竟遭受了什么？他被处死。此外，他又如何被另样\Quote{见证他为义人}？据说，有火降下烧尽了祭品。因为取代了[\Quote{上主看顾了\Person{亚伯尔}和他的祭品}]（见：\BibleRef{创 4:4}），叙利亚文译本说：\Quote{祂用火焚烧了它们。}因此，那位用言行见证义人、并看见他为祂被杀的人，并没有为他报仇，而是任他受苦。\par

但你的情况并非如此：又怎能呢？你既有先知、榜样、无数的鼓励，又有完成的征兆和奇迹？由此，那才是真正的信德。因为他见到了什么奇迹，才相信他会得到善事的报酬呢？难道他不是单单由于信德而选择了德行吗？\par

\Quote{并且因着它，他虽然死了，却仍在说话}是什么意思？为了不使他们陷入极大的沮丧，他指出他在某种程度上已得到了报酬。怎样呢？\Quote{从他所发出的影响是巨大的，}他说，\Quote{他仍在说话}；也就是说，[\Person{加音}]杀了他，但并没有同时消灭他的光荣和记忆。他没有死；因此你们也不会死。因为一个人的苦难越严重，他的光荣就越大。\par

他如何\Quote{仍在说话}？这既是他活着的标志，也是他被所有人称颂、赞美、视为有福的标志。因为鼓励他人成为义人的人，就是在说话。因为没有什么言词比那个人的受苦更有效。正如诸天仅凭其面貌就在说话，同样他也因被人纪念而说话。即使他自我宣扬，即使他有一万个舌头，并且活着，也不会像现在这样受人钦佩。也就是说，这些事情不会毫无报应地发生，也不会轻易消散。\par

% structure: p0019 source=h2 target=subsection reason=relative-heading-rank; base=h2

\subsection{\BibleRef{希 11:5-6}}

4. \BibleQuote{因着信德，哈诺客被接去，使他不见死亡，人也找不着他，因为天主将他接去了。} 这人显示了比亚伯尔更大的信德。怎么呢（你问）？因为他虽然在后，但[亚伯尔]所遭遇的事足以引导他回转。怎样呢？天主预知[亚伯尔]将被杀。因为他对加音说：\Quote{你犯了罪：不要再增添。} 天主虽受他尊敬，却没有保护他。然而，这并没有使[哈诺客]变得漠不关心。他没有对自己说：\Quote{何必劳苦和危险？亚伯尔尊敬了天主，天主却没有保护他。那已逝的人，从他兄弟的惩罚中得了什么益处？他又能从中得到什么好处？就算他受重罚，那与被杀的人有何相干？} 他既没有说也没有想任何这类事，反而超越这一切，知道若有一位天主，就必定也有一位赏报者：虽然那时他们还不知道复活的事。但如果那些还不知道复活、并在此世看到矛盾之事的人，尚且这样讨[天主]喜悦，何况我们呢？因为他们既不知道复活，也没有任何榜样可看。正是这件事使[哈诺客]蒙悦纳，即他什么都未得到。因为他知道[天主]\Quote{是赏报者}。从哪里他知道的呢？你说，\Quote{因为他赏报了亚伯尔}？如此，理性暗示了其他事，但信德却与所见相反。即使那时（他可能会说）如果你在此世一无所获，也不要烦恼。\par

如何是\Quote{因着信德}而使\Quote{哈诺客被接去}呢？因为他蒙[天主]悦纳是他被接去的原因，而信德又是他蒙悦纳的原因。因为他若不知道将要获得赏报，又怎能蒙悦纳呢？\Quote{但没有信德，是不可能蒙悦纳的}。怎样呢？如果一个人相信有一位天主，并且有赏罚，他就会获得赏报。那么，这悦纳从何而来呢？\par

5. 有必要\Quote{相信祂是}，而不是\Quote{祂是什么}。如果\Quote{祂是}需要信德，而非推论，那么就无法通过推论来理解\Quote{祂是什么}。如果\Quote{祂是赏报者}需要信德，而非推论，那么又怎可能通过推论来把握祂的本质？因为什么推论能达到这一点？因为有人说，存在的事物是自因的。你不见，除非我们在一切事上，不仅对赏罚，而且对天主本身的存在都有信德，否则一切都将丧失吗？\par

但许多人问，哈诺客被接去了哪里，为什么他被接去，为什么他没有死，他和厄里亚都没有死，以及他们是否还活着，如何活着，以什么形态。但问这些是多余的。因为一个被接去，另一个被提去，圣经已经说了；但他们在哪里，如何存在，圣经并未添加；因为他们只说必要之事。因为这件事，我是指他的被接去，确实在起初就发生了，人的灵魂因而得到希望，即死亡将被毁灭，魔鬼的暴政将被推翻，死亡将被消除；因为他被接去，不是死了，而是\Quote{使他不见死亡。}\par

因此，祂又补充说：他被提升去，因为他曾中悦天主。正如一位父亲，在恐吓了他的儿子之后，确实愿意立即收回他的威胁，但他忍耐并坚持决心，以便暂时管教和纠正他，让威胁保持坚定；同样，天主——仿佛照人的方式来说——并没有继续坚持，而是立即表明死亡已被废除。祂先允许死亡发生，愿意藉着儿子来恐吓父亲：因为为了表明判决确实已经确定，祂不是立刻将恶人置于这一惩罚之下，而是将那位中悦天主的人，我说的是有福的亚伯尔，置于此惩罚下；几乎紧接着他之后，祂就提升了哈诺克。此外，祂并没有使前者复活，以免他们立刻变得胆大；但祂却提升了仍活着的后者：藉着亚伯尔激起恐惧，又藉着后者给予殷勤中悦祂的热忱。因此，那些说一切事物都是自行管理、自行支配，并不期待赏报的人，是不会中悦天主的；外邦人也是如此。因为\Quote{祂成了那些殷勤寻求祂者的赏报者}（藉着行为和知识）。\par

6. 既然我们有了\Quote{赏报者}，就让我们做一切事，以免被剥夺德行的赏报。因为忽视这样的赏报、轻视这样的赏报，实在值得许多眼泪。因为正如祂是那些\Quote{殷勤寻求祂者}的赏报者，同样对那些不寻求祂的人，则是相反的。\par

祂说：\BibleQuote{你们求，必要给你们}\BibleRef{玛 7:7}；但我们怎能找到上主呢？请想想金子是如何找到的：须付出大量辛劳。[\BibleQuote{我夜间用手寻求了主，在祂面前，我并没有受骗}\BibleRef{咏 76:2}]，也就是说，正如我们寻找失去的东西，让我们也这样寻找天主。我们岂不集中精神吗？岂不询问每一个人吗？岂不离开家吗？岂不应许钱财吗？\par

例如，假设我们中有人失去了他的儿子，我们什么不做？我们走遍了哪块土地、哪片海洋？我们不把金钱、房屋以及一切其他事物都视为次要于寻找他吗？如果我们找到了他，我们就紧抓住他，牢牢抱住他，不让他离开。当我们打算寻找任何东西时，我们千方百计地忙碌，以找到所寻找的。对于天主，我们岂不更应当这样做，因为所寻找的是必不可少的？更确切地说，不是以同样的方式，而是更多！但既然我们软弱，至少像你寻找金钱或儿子那样寻找天主吧。你岂不为了祂而离开你的家？你从未为了钱而离开家吗？你岂不千方百计地忙碌吗？当你找到它时，你岂不充满信心吗？\par

7. 祂说：\Quote{你们求，必要给你们。}因为寻求的事物需要很多关注，尤其是在天主的事上。因为有许多障碍，许多遮蔽，许多妨碍我们感知的事物。因为正如太阳是光明的，公开显现在众人面前，我们无需寻找它；但是，如果我们把自己埋葬起来，颠倒一切，就需要很多劳动才能看见太阳；同样，在这里，如果我们把自己埋葬在邪恶欲望的深处、情欲和现世事务的黑暗中，我们就很难向上看，很难抬起头，很难看得清楚。被埋在地下的人，无论他向上看的程度如何，他就那样接近太阳。因此，让我们抖掉尘土，让我们打破笼罩在我们身上的迷雾。这迷雾又厚又密，不允许我们看清楚。\par

那么，你怎么说，这云如何被冲破？如果我们把\Quote{正义的太阳}的光线吸引到自己身上。\BibleQuote{愿我双手的举起，如同晚间的祭献}\BibleRef{咏 140:2}。让我们也用手举起我们的心思：你们这些已受入门奥迹的人，知道我指的是什么，也许你们也认出了这句话，一看就猜到了我暗示的事。让我们高举我们的思念于高处。\par

我自己认识许多人，几乎悬空于地面，过度地伸展双手，因不能升到空中而灰心，并这样热切地祈祷。我希望你们常这样，若不常这样，至少时常这样；若不十分时常，至少偶尔，至少在早晨，至少在晚祷中。因为，请告诉我，你不能伸展双手吗？伸展意志，伸展到你愿意的程度，甚至直达天上。即使你愿意触摸最高的顶点，即使你愿意升得更高并走在上面，这对你也是开放的。因为我们的心思比任何有翅膀的生物更轻、更高。当它从圣神获得恩宠时，啊！它何等敏捷！何等迅速！它周游一切！它从不沉没或跌落地面！让我们为自己预备这些翅膀：藉着它们，我们将能飞越今世汹涌的海洋。最快的鸟儿在片刻之间飞越山岭、森林、海洋和岩石而毫发无损。心思也是如此；当它有了翅膀，当它脱离了现世的事物，就没有什么能抓住它，它高于一切，甚至高于魔鬼的火箭。\par

魔鬼并非如此出色的射手，能射达如此的高度；他固然射出他的箭矢，因他毫无羞耻，却未能射中目标；箭矢无效地折回，不仅无效，反而落在他自己头上。因为凡由他射出的，必然击中某物。正如人所射出的箭，或射中他所瞄准之人，或刺中飞鸟、篱笆、衣服、木头，或仅仅射中空气，魔鬼的箭矢也是如此。它必然击中；若它击不中所射之人，便必然击中发射者本身。我们从许多事例中可以学到：当我们未被击中时，无疑他自己被击中了。例如，他谋害约伯：他没有击中约伯，反被击中。他谋害保禄，他没有击中保禄，反被击中。若我们留意，到处可见此事。因为即使他击中了，自己也被击中；何况当他未击中时，就更是如此了。\par

8. 那么，让我们将他的武器转向他自己，并以信德盾牌武装和坚固自己，以坚定之心守卫，成为不可攻破的。魔鬼的箭矢就是邪情。愤怒尤其是一种火、一种火焰；它引人、毁灭、吞噬；让我们以坚忍和宽容将它熄灭。因为正如烧红的铁浸入水中，失去其火，愤怒之人遇到忍耐之人，对忍耐之人无害，反而有益于他，并更完全地制伏自己。\par

因为没有什么能与坚忍相比。这样的人永不受侮辱；如同金刚石的身体不受创伤，这样的灵魂也是如此。因为他们超越箭矢的射程。坚忍之人是崇高的，崇高到不受射击的伤害。当一个人暴怒时，你当嘲笑；但不要公开嘲笑，以免激怒他；要在心里为他而嘲笑。因为对孩童，当他们激怒地打击我们，仿佛在报复时，我们嘲笑。如果你嘲笑，那么你和他的差别，就如同成人与孩童一样；但若你发怒，你便使自己成了孩童。因为愤怒的人比孩童更无知。若看一个发怒的孩童，岂不嘲笑他？经上说：\BibleQuote{心窄的人，实是愚昧。}\BibleRef{箴 14:29}愚昧的人就是孩童：又经上说：\BibleQuote{坚忍的人，富有智慧。}\BibleRef{箴 14:29}让我们追随这\BibleQuote{富有智慧}，好能获得在基督耶稣我们的主内所应许的福乐；愿光荣、权能、尊崇归于他，连同圣父及圣神，从今世直到永远，世世无穷。阿们。\par

\SourceRecord{Translated by Frederic Gardiner.；Nicene and Post-Nicene Fathers, First Series；Vol. 14.；Edited by Philip Schaff.；Buffalo, NY: Christian Literature Publishing Co.,；1889.；<http://www.newadvent.org/fathers/240222.htm>.}

% structure: p0001 source=h1 target=section reason=page-title

\section{希伯来书讲道集 第二十三篇}

% structure: p0002 source=h2 target=subsection reason=relative-heading-rank; base=h2

\subsection{希伯来书 11:7}

\BibleQuote{\Emph{因着信德，诺厄蒙天主警告}\Emph{关于尚未看见的事，怀着敬畏，预备了方舟，为救他一家的性命；借此他定了世界的罪，并成为因信德而得的义的承继者。}}\par

1. \BibleQuote{因着信德}（他说）\BibleQuote{诺厄蒙天主警告。}正如天主圣子谈及自己的来临所说：\BibleQuote{在诺厄的日子，人照常嫁娶} \BibleRef{路 17:26}，因此宗徒也使他们想起一个合适的图像。因为哈诺客的例子只是信德的例子；而诺厄的例子则也是不信的例子。这不但是一种完全的安慰和劝勉，当时不但信者被证明为蒙悦纳，而且不信者也遭受相反的事。\par

因为他怎么说？\BibleQuote{因着信德，蒙天主警告。}什么是\BibleQuote{蒙天主警告}？就是\BibleQuote{预先告知了他。}但为什么要使用\OriginalTerm{神谕}{divine communication}这个说法？因为在另一处经上也说：\BibleQuote{他蒙圣神告知，}以及\BibleQuote{神谕怎么说？} \BibleRef{路 2:26} \BibleRef{罗 11:4} 你看见圣神的同等尊荣吗？因为正如天主启示，圣神也如此启示。但他为什么这样说？预言被称为\Quote{神谕}。\par

\BibleQuote{关于尚未看见的事，}他说，即关于雨水。\par

理性并未提示任何这类事；因为\BibleQuote{他们照常嫁娶}；天空晴朗，没有变化的迹象：然而他仍然敬畏：\BibleQuote{因着信德（他说）诺厄蒙天主警告关于尚未看见的事，怀着敬畏，预备了方舟，为救他一家的性命。}\par

\BibleQuote{借此他定了世界的罪}是什么意思？他表明他们应受惩罚，因为他们甚至没有因这预备而醒悟。\par

\BibleQuote{他成为}（他说）\BibleQuote{因信德而得的义的承继者}：即因相信天主，他被证明为义。因为这是一颗真诚归向祂的灵魂的部分，不认为任何事比祂的话更可靠，正如\OriginalTerm{不信}{Unbelief}恰恰相反。信德显然产生义。因为我们曾蒙天主警告关于地狱的事，他也同样：但那时他被嗤笑；他受辱骂和讥讽；但他毫不介意。\par

% structure: p0010 source=h2 target=subsection reason=relative-heading-rank; base=h2

\subsection{希伯来书 11:8-9}

2. \BibleQuote{因着信德，亚巴郎蒙召往将要去承受为产业的地方去，就听命前往，并且不知道往哪里去。因着信德，他在预许之地侨居，如在异乡，住在帐幕内，与同沐同一预许的承继人依撒格和雅各伯。}[\BibleQuote{因着信德}]：因为（请告诉我）他效法了谁？他有一个外邦人和偶像崇拜者为父；他没有听过先知；他不知往哪里去。因为希伯来信者仰望这些祖先，以为他们已享有无数的福乐，但他显示他们中无人获得任何东西；所有人都未得赏报；无人已接受他的酬报。他离开了他的故乡和家园，\BibleQuote{并且不知道往哪里去。}\par

如果他本人如此，他的后裔也同样居住，有何可奇？因为看见预许似乎被否定（因祂曾说：\BibleQuote{我要将这地赐给你和你的后裔} \BibleRef{创 12:7}），他看见他的儿子住在那里；又看见他的孙子住在非己之地；然而他毫不动摇。因为亚巴郎的事如我们所料的发生，因为预许日后要在他的家族中成就（虽然对本人也说：\BibleQuote{给你和你的后裔，}不是\BibleQuote{通过你的后裔给你，}而是\BibleQuote{给你和你的后裔}）：但无论他，还是依撒格，还是雅各伯，都没有享受预许。因为一个为雇工服务，另一个被赶出；他本人甚至因恐惧而失败：他在战争中夺取了一些物品，若无天主的帮助，另一些已被毁灭。因此[宗徒]说：\BibleQuote{与同沐同一预许的承继人，}不是仅他本人，他也指承继人。\par

% structure: p0013 source=h2 target=subsection reason=relative-heading-rank; base=h2

\subsection{希伯来书 11:13}

3. \BibleQuote{这些人全都因信德死了，}他说，\BibleQuote{没有获得所预许的。}在此值得提出两个疑问：既然他说天主\BibleQuote{接去了哈诺客，人找不着他，以致他没有见死，}他又如何说\BibleQuote{这些人全都\Emph{死了}于信德中。}再者，他说\BibleQuote{他们未获得所预许的，}却宣称诺厄已得了赏报：\BibleQuote{为救他一家的性命，}哈诺客被\BibleQuote{接去，}亚伯尔\BibleQuote{仍然说话，}亚巴郎已握有那地，然而他说：\BibleQuote{这些人全都因信德死了，没有获得所预许的。}那么这是什么意思？\par

必须先解决第一个[疑难]，然后第二个。\BibleQuote{这些人全都}（他说）\BibleQuote{\Emph{死了}于信德中。}这里\Quote{全都}一词的用法不是因为他们都死了，而是除了那一个例外，\Quote{这些人全都死了}，我们知道他们是死了的。\par

而那一句\Quote{未获得所应许的}，是真的：因为给诺厄的应许当然不是这里所说的。然而，他所说的\Quote{应许}是哪一种？依撒格和雅各伯获得了关于土地的应许；但对于诺厄、亚伯尔和哈诺客，他们获得了哪种应许？那么，他要么是谈到这三个人；要么，如果也谈到另外那些人，那应许并不是指亚伯尔应受赞美，也不是指哈诺客被提升，也不是指诺厄被保全；这些事是为了他们的德行而临到他们的，仿佛是未来之事的预尝。因为天主从起初就知道人类需要极大的迁就，祂不仅赐给我们来世的事物，也赐给现世的；例如，基督甚至对门徒说：\Quote{凡舍弃了房屋、或兄弟、或姊妹、或父亲、或母亲的，必要获得百倍，并承受永生。}\BibleRef{玛 19:29}又说：\Quote{你们先寻求天主的国，这一切自会加给你们。}\BibleRef{玛 6:33}你看见这些是作为附加物赐下的，好叫我们不致灰心吗？因为正如运动员在比赛时虽享有精心照料，但并不在此刻享受完全安逸，而是生活在纪律之下，之后才完全享受；同样，天主也不让我们在此世享受\Quote{完全}的安逸。因为即使在此世，祂也确实赐予了一些。\par

4. \Quote{他们从远处看见这些应许，}他说，\Quote{且拥抱了它们。}在此他暗示了一件奥秘的事：他们预先领受了关于未来之事的一切所说的话；关于复活，关于天国，关于基督来时宣讲的其他事，因为这就是他所指的\Quote{应许}。那么，他要么是指这个意义，要么是指他们确实没有领受这些应许，但在对它们的信赖中死去，并且他们只是借着信德而如此信赖。\par

\Quote{他们从远处看见}：在四代之前；因为过了那么多代，他们才从埃及上去。\par

\Quote{并且拥抱了它们，}他说，\Quote{且欢喜。}他们如此确信这些应许，以至于甚至拥抱〔或\Quote{致意}〕它们，这是借用了船上的人从远处看见渴望已久的城市的比喻：他们在进入城市之前，就以问候的话语掌握和占据它们。\par

% structure: p0020 source=h2 target=subsection reason=relative-heading-rank; base=h2

\subsection{希伯来书 11:10}

\Quote{因为他们期待着}（他说）\Quote{那座有根基的城，其设计者和建造者是天主。}你看见他们是在这个意义上领受的，即他们已经接纳它们并对它们满怀信心。那么，如果充满信心就是领受，你也可以领受。因为这些人，虽然未曾享受那些福分，却因着他们的强烈渴望而看见了它们。现在这些事为何发生？为了使我们羞愧，因为事实上，当地上的事物被应许给他们时，他们并不在意，而是寻求将来的\Quote{城}；而天主却再三向我们谈论天上的城，我们却仍在寻求现世的那座。祂对他们说：我要把现世的事物赐给你们。但当祂看见——更确切地说，当他们自己显明为配得上更大的事物时，祂就不再让他们领受这些较小的，而是那些更大的；这是要向我们显示，他们配得上更大的事物，因为他们不愿受这些束缚。好像一个人向一个懂事的孩子承诺玩具，不是要他接受，而是要借此展现他的智慧，当他要求更重要的东西时。因为这正表明，他们以极大的热忱远离了土地，以至于甚至不接受所赐的。因此，他们的后裔为此领受了土地，因为祖先本身就配得上那土地。\par

\Quote{那座有根基的城}是什么意思？难道这些〔可见的城〕没有\Quote{根基}吗？与那座相比，它们没有。\par

\Quote{其设计者和建造者是天主。}哦！对那座城是何等的赞美！\par

% structure: p0024 source=h2 target=subsection reason=relative-heading-rank; base=h2

\subsection{希伯来书 11:11}

5. \Quote{因着信德，连撒辣自己，}他说。这里他开始用一种使他们羞愧的方式说话，以防他们表现得比一个女人更怯懦。但或许有人会说：她怎么是\Quote{因着信德}，既然她笑了？不，她的笑虽来自不信，她的恐惧却来自信德，因为说\Quote{我没有笑}\BibleRef{创 18:15}，是出于信德。由此可见，当不信被清除后，信德就取而代之。\par

\Quote{因着信德，连撒辣也获得了能力，得以怀孕生子，尽管她已过了生育年龄。}\Quote{怀孕生子}是什么意思？她已成了死人，是不生育的，却获得了保留种子、即受孕的能力。因为她的缺陷是双重的：第一是因年龄，因为她确实老了；第二是因天性，因为她本来不育。\par

% structure: p0027 source=h2 target=subsection reason=relative-heading-rank; base=h2

\subsection{希伯来书 11:12}

\Quote{因此，从一个人}他们全\Quote{都生出来，如同天上的繁星，又如同海边的沙粒。}\Quote{因此}（他说）\Quote{从一个人}他们全\Quote{都生出来。}此处他不仅说撒辣生了孩子，而且说她还成为如此多人的母亲，其数目之多，连多产的子宫也无法生育。\Quote{如同繁星，}祂说。那么，为何祂常数算他们，尽管祂说：\Quote{你的后裔要像天上的星那样不可数算}\BibleRef{创 15:5}？祂要么是指其数量极多，要么是指那些不断出生的人。因为，请告诉我，你能数算某一家族中的人，如某人是谁的儿子，某人又是谁的儿子吗？但这就是天主的应许，祂的工程如此巧妙地安排。\par

6. 但如果祂附加应许的事物尚且如此令人赞叹、如此超乎期望、如此宏伟，那么那些以此为基础、在此之上所增添的事物，又将是如何呢？那么，还有什么比那些获得这些事物的人更有福呢？还有什么比那些失落它们的人更不幸呢？因为，一个人若被驱逐出祖国，人人都怜惜他；若丧失产业，人人都视他为可怜；那么，对于丧失天堂及其中的宝藏的人，岂不更应当用泪如雨下的哀哭来为他悲悼吗？或者，他甚至不值得被哭泣：因为一个人若因并非自身的原因而受苦，才值得被哭泣；但若因自己的选择而陷入邪恶，他便不值得流泪，只值得哀号；或者更确切地说，只值得哀悼。因为连我们的主\Person{耶稣基督}也曾为耶路撒冷哀悼哭泣，虽然它是那么不虔敬。实在，我们配得上无数的流泪、无数的哀号。如果整个世界都获得声音，无论是石头、木头、树木、野兽、飞鸟、鱼类，总之，整个世界若获得声音而为我们这些失落了那些福乐的人哀哭，它仍不足以哀哭和悲叹。因为，有什么言语、什么理智，能够描述那种福乐与美德、那种愉悦、那种光荣、那种幸福、那种光辉呢？\Quote{眼所未见，耳所未闻，人心所未曾想到的}\BibleRef{格前 2:9}（他没说它们仅仅超越我们所想象的，而是说无人曾想到过）\Quote{天主为爱祂的人所预备的。}因为，那些福乐该是怎样的呢？既然是那位天主作预备和建立者？因为，若在祂造了我们之后，我们尚未做什么事，祂就白白赐予如此大的恩惠——乐园、与祂本人的亲密交往、应许我们不死、幸福而无忧虑的生命——那么，对那些为祂如此劳苦奋斗、忍受一切的人，祂岂不更要赐予什么呢？祂为了我们，连自己的独生子也没有怜惜；为了我们这些曾是仇敌的人，祂将自己的儿子交付死亡；如今我们已成为祂的朋友，祂岂不视我们为配得一切呢？祂既已使我们与祂和好，岂不将一切分施给我们呢？\par

7. 祂既极其丰盛又无限富有；祂渴望并竭力获得我们的友谊；我们却并不如此竭力追求。我说什么呢？\Quote{并不竭力追求}？我们并不像祂那样愿意获得那些福乐。而祂所做的事表明，祂比我们更愿意。因为，我们为了自己，尚且难以舍弃一点点金子；而祂为了我们，竟舍弃了祂自己的儿子。让我们妥善利用天主的爱，让我们收获祂友谊的果实。因为\Quote{你们是我的朋友}（他说）\Quote{如果你们遵行我所吩咐你们的。}\BibleRef{若 15:14}多么奇妙！祂的仇敌，那些与祂相距无限遥远、在一切方面都被祂无可比拟地超越的人——祂使他们成为祂的朋友，并称他们为朋友。那么，为了这友谊，有什么不可忍受的呢？为了人的友谊，我们常常冒险；但为了天主的友谊，我们连钱财也不肯舍弃。我们的状况确实需要哀悼，需要哀哭、流泪、哀号、捶胸顿足。我们已从我们的希望中坠落，从崇高地位上被贬抑，我们显示出自己不配天主的尊荣；即使在祂的恩惠之后，我们仍变得麻木不仁、忘恩负义。魔鬼已剥夺了我们一切福乐。我们曾被视为配得成为儿子；我们是祂的弟兄和同继承人，如今却与辱骂祂的仇敌毫无区别。\par

从此以后，我们还有什么安慰呢？祂召叫我们上天，我们却将自己推入地狱。\Quote{起誓、说谎、偷窃、奸淫，都充斥在大地上。}\BibleRef{欧 4:2}有些人\Quote{流血又流血}；另一些人则做出比流血更恶劣的事。许多受屈受压的人，宁死也不愿遭受这些；若非敬畏天主，他们甚至可能自杀，对自己如此凶狠。这些事难道不比流血更糟糕吗？\par

8. \Quote{祸哉，我的灵魂！因为虔诚人从地上灭绝了，人间没有正直人。}\BibleRef{米 7:1}让我们现在也如此哭喊，首先为我们自己；但请协助我的哀悼。\par

也许有人甚至厌恶并嗤笑。正因如此，我们更应加深哀叹，因为我们如此疯狂、迷失自己，竟然不知道自己在疯狂，反而对那些本该哀叹的事嗤笑。人啊！\Quote{天主的义怒从天上显露在一切不虔敬和不义的人身上}\BibleRef{罗 1:18}；\Quote{天主必要显现，烈火在祂面前燃烧，周遭有狂风大作。}\BibleRef{咏 49:3}\Quote{烈火在祂面前燃烧，烧灭祂四围的敌人。}\BibleRef{咏 96:3}\Quote{上主的日子有如火炉。}\BibleRef{拉 4:1}但没有人将这些记在心里；这些可畏可惧的道理反而比神话更被轻视，被人践踏。听者无人；而嗤笑作乐者却比比皆是。我们将有何等办法？从哪里找到安全？\Quote{我们灭亡了，我们完全被吞灭了}\BibleRef{户 17:12}，我们成了仇敌的笑柄，外邦人和魔鬼的讥讽。如今魔鬼大得志，他得意洋洋，欢喜快乐。那受托照管我们的天使们都蒙羞悲伤；没有人能转变你们；我们一切努力都归于徒然，在你们看来，我们好像是空谈者。现在甚至适宜呼唤上天，因为没有人听；呼求元素作证：\Quote{天哪，请听！地啊，请侧耳！因为上主发言了。}\BibleRef{依 1:2}\par

伸出援手吧，你们这些尚未被压倒的人，给那些因酗酒而堕落的人。你们健康的人，给那些患病的人；你们头脑清醒的人，给那些疯狂、头晕目眩的人。\par

我恳求你们，不要让任何人将朋友的喜爱置于自己的救恩之上；让暴力和责备只看一件事——他的益处。当一个人被热病抓住时，甚至奴隶也会抓住他们的主人。因为当那病症紧迫地折磨他，使他的心智混乱，一群奴隶站在旁边时，他们在主人的灾难中，不再认识主仆的律法。\par

我劝告你们，让我们冷静下来：每日都有战争，城邑被淹没，四周有无数的毁灭，天主的义怒像网一样从四面围困我们。而我们，好像中悦了他一样，安然而居。我们都准备好手去行不义之财，没有一人去帮助他人；全都为了掠夺，没有一人去保护；每个人都热切地想着如何增加自己的财产，没有一人想着如何帮助穷人；每个人都焦虑如何增加财富，没有一人想着如何拯救自己的灵魂。一种恐惧占据了所有人，唯恐（你说）我们变穷；没有人在痛苦和战栗中唯恐我们堕入地狱。这些事需要哀哭，这些事需要谴责，这些事需要责备。\par

9. 但我不愿谈论这些事，但我被悲伤所迫。原谅我：我被忧伤迫使说出许多甚至我不愿说的事。我看到我们的伤口是沉重的，我们的灾祸无法安慰，苦难超过了安慰。我们完了。\BibleQuote{但愿我的头是水，我的眼是泪水的泉源} \BibleRef{耶 9:1}，使我得以哀哭。让我们哭泣吧，亲爱的，让我们哭泣，让我们呻吟。\par

或许这里有人说，他除了哀哭什么也不对我们说，除了眼泪什么也不说。相信我，这并不是我的意愿，并非我的意愿，我宁愿进行一连串的称赞和赞美，但现在不是做这些事的时候。亲爱的，令人忧伤的不是哀哭，而是做那些招致哀哭的事。痛苦不是应当回避的，而是犯下那些招致痛苦的事。你不要受罚，我就不会哀悼。你不要死，我就不会哭泣。如果身体确实躺卧死了，你叫所有人与你一同悲伤，并认为那些不哀悼的人没有同情；而当灵魂灭亡时，你告诉我们不要哀悼吗？\par

但如果我不哭泣，我就不能成为父亲。我是一个充满慈爱的父亲。请听保禄如何呼喊：\BibleQuote{我的孩子们，我为你们再受产痛} \BibleRef{迦 4:19}。有哪位临盆的母亲发出像他这样痛苦的呼喊？但愿你们能看见我心中的火，你们就会知道，我燃起的悲痛比任何妇女或遭受过早寡妇之苦的人更炽烈。她不会如此哀悼她的丈夫，任何父亲也不会如此哀悼他的儿子，就像我为我们这里的这群人哀悼一样。\par

我看不到任何进步。一切都变成了毁谤和指责。没有人以悦乐天主为己任；反而（他说）\Quote{让我们说这个人或那个人的坏话。}\Quote{这个人不适合作圣职人员。}\Quote{这个人过着不体面的生活。}当我们应当为自己的恶行悲伤时，我们却判断他人，而我们不应该这样做，即使我们清白无罪。\BibleQuote{因为谁使你与众不同呢？}（他说）\BibleQuote{你有什么不是领受的呢？既然是领受的，为什么你夸耀，好像不是领受的呢？} \BibleRef{格前 4:7} \BibleQuote{而你，为什么判断你的弟兄呢？} \BibleRef{罗 14:10}，你自己充满无数的恶行吗？当你说，\Quote{这个人是个坏人，挥霍无度，行为不端}时，想想你自己，严格省察你自己的状况，你就会后悔你所说的话。因为没有什么，绝对没有什么，比回忆我们自己的罪过更能有力地激发德行。\par

如果我们心中反复思量这两件事，我们就能获得所应许的福分，我们就能洁净自己，擦去过失。只要我们有时认真思考；让我们忧虑这件事，亲爱的。让我们在这里因反省而悲伤，以免以后在惩罚中悲伤，而能享受永恒的福分，那里\BibleQuote{痛苦、忧伤和叹息都已逃遁} \BibleRef{依 35:10}。使我们能获得超越人理解的恩赐，在我们主基督耶稣内，愿荣耀和权能归于他，直到永远。阿们。\par

\SourceRecord{Translated by Frederic Gardiner.；Nicene and Post-Nicene Fathers, First Series；Vol. 14.；Edited by Philip Schaff.；Buffalo, NY: Christian Literature Publishing Co.,；1889.；<http://www.newadvent.org/fathers/240223.htm>.}

% structure: p0001 source=h1 target=section reason=page-title

\section{\WorkTitle{希伯来书讲道集 第二十四篇}}

% structure: p0002 source=h2 target=subsection reason=relative-heading-rank; base=h2

\subsection{\BibleRef{希 11:13-16}}

\BibleQuote{ \Emph{这些人都怀着信德死去}，\Emph{没有获得所恩许的，却从远处看见了}，\Emph{且拥抱了它们，并承认自己在世上只是旅客和寄居者。因为说这样话的人，显明自己是在寻求一个家乡。他们如果想念自己离开的那个家乡，他们还有回去的机会；但是如今他们渴慕一个更美好的家乡，即天上的；因此天主不以} \Emph{被称为他们的天主为耻，因为祂已为他们预备了一座城。} }\par

1. 第一德行，甚至整个德行，就是成为这个世界的旅客和外人，与这里的事物毫无共通之处，而是悬空脱离它们，如同与我们陌生的事物一样；正如那些有福的门徒，关于他们，他说：\BibleQuote{他们披着绵羊皮、山羊皮，各处奔跑，受穷乏、患难、苦害：世界不配拥有他们。}\EditorialNote{见：\BibleRef{希 11:37-38}}\par

他们称自己是\Quote{旅客}；但保禄说得更有过之：因为他不仅称自己为旅客，而且说自己对世界是死的，世界对他是死的。\BibleQuote{因为世界（他说）已被钉死于我，我也被钉死于世界。}\BibleRef{迦 6:14} 但我们既是公民，又非常活跃，忙于这里的一切事物如同公民。义人对世界是何等，\Quote{旅客}和\Quote{死人}，我们对天堂也是如此。他们对于天堂是何等，活着的和作为公民的，我们对于世界也是如此。因此我们是死的，因为我们拒绝了那真正的生活，而选择了这暂时的生活。因此我们惹动了天主的忿怒，因为当天上的享受摆在我们面前时，我们却不愿意与地上的事物分离，反而像虫子一样，从地上到地上转来转去，再从这块到那块；总之，我们甚至不愿仰望一小会儿，也不愿从人间事务中抽身，而是如同沉浸在麻木、睡眠和醉态中，被幻想所麻痹。\par

2. 正如那些被甜蜜睡眠控制的人，不仅夜间躺在床上，甚至早晨追上他们、白昼亮光来临，他们仍不害羞地沉浸在享乐中，把工作和活动的时节变成酣睡和怠惰的时刻；同样，我们也真是如此，当白昼临近，黑夜已深，或者更确切地说，白昼已到时，因为\BibleQuote{趁着白日做工}（经上说）\BibleRef{若 9:4}；当是白昼时，我们却实行一切属于黑夜的事：睡觉、做梦、沉迷于奢侈的幻想；我们悟性的眼睛和身体的眼睛都闭上；我们说话不当，言语荒谬；即使有人给我们重重一击，夺去我们所有财物，甚至放火烧屋，我们也不会察觉。\par

或者不如说，我们甚至不等别人这样做，而是自己每天刺伤、伤害自己，不成体统地躺着，被剥夺了一切信用、一切荣誉，既不能隐藏自己的羞耻行为，也不允许别人这样做，反而暴露于众目睽睽之下，成为旁观者和路人无穷的嘲笑和讥讽的对象。\par

3. 你们是否认为，恶人自己也嘲笑那些与他们性情相似的人，并且谴责他们？因为天主在我们内安置了一个不受贿赂、永远不会完全毁灭的审判庭，即使我们堕入极深的恶习；因此连恶人自己也审判自己，如果有人称他们为所是的，他们就羞愧、生气，说这是侮辱。这样，他们谴责自己的行为，即使不是通过行为，也是通过言语、通过良心，甚至通过行为。因为当他们避开目光、秘密地进行他们的勾当时，他们最强有力地证明了他们自己对那事本身所持的看法。因为邪恶如此明显，所有人都指责它，甚至那些追随邪恶的人也是如此；而德行的品质是如此，即使那些不效法它的人也赞美它。因为连私通者也会赞美贞洁，贪婪者也会谴责不义，易怒者也会赞美忍耐，指责争吵，放荡者也会谴责放荡。\par

那么，你说，他为什么还追求这些东西？由于过度的懒惰，并非因为他认为那是好的；否则他不会因事情本身羞愧，也不会在别人指责时否认。实际上，很多人被抓住后，因忍受不了羞耻而自缢。在我们内里美好和正当的事物的见证是如此强烈。因此，善比太阳更明亮，而恶比任何事物更丑陋。\par

4. 圣人们是\Quote{旅客和寄居者}。如何以何种方式？亚巴郎在哪里承认自己是\Quote{旅客和寄居者}？可能他自己也承认了；但达味承认说：\BibleQuote{我是旅客}，以及什么？\BibleQuote{如同我所有的祖先。}\BibleRef{咏 38:12} 因为那些住在帐篷里的人，那些甚至用钱购买墓地的人，显然在某种意义上他们是旅客，因为他们甚至没有埋葬死者的地方。\par

那么，他们岂是说他们在巴勒斯坦地是\Quote{旅客}？绝不是：而是就整个世界而言；这很有道理，因为他们在那里没有看到任何他们所希望的事物，而是一切都是陌生的。他们希望实践德行，但这里有许多邪恶，与他们格格不入。他们没有朋友，没有熟人，只有很少几个。\par

但他们是怎样成为\Quote{旅客}的？他们对这里的事物毫不关心。这不仅是言语，更是行为。以什么方式？\par

祂对\Person{亚巴郎}说：\Quote{离开那似乎是你的故乡的地方，来到一个外方之地}。他并没有依恋他的亲属，反而毫不在意地放弃了他们，好像要离开一个外方之地一样。祂对他说：\Quote{献上你的儿子}，他就献上了他，好像没有儿子一样；好像摆脱了本性一样，他就献上了他。他所获得的财富对所有过路人是共有的，他视之为无物。他把首位让给别人；他投身于危险；他忍受了数不清的患难。他没有建造华丽的房屋，没有享受奢靡，不关心衣着——这些都是世俗之事；而是在各方面都如同属于那上面的城邑一样生活；他展现了款待、手足之爱、怜悯、宽容、对财富和现世荣耀以及其它一切的轻视。\par

他的儿子也像他一样：当他被赶走时，当战争临到他时，他退让让步，如同在外方之地。因为外方人无论遭受什么，都忍受着，好像不在自己的家乡一样。甚至当他的妻子被夺走时，他也忍受了，如同在异乡；他在各方面都像一个家乡在天上的人那样生活，表现出节制和有秩序的生活。因为他生了一个儿子后，就不再与妻子同房，而且是在青春之花已过时才娶了她，表明他这样做不是出于情欲，而是顺服天主的应许。\par

至于\Person{雅各伯}呢？他不是只求食粮和衣着吗？这正是那些真正是陌生人、那些极度贫穷的人所求的。当他被赶出时，他不是像陌生人一样让位吗？他不是为工价而服侍吗？他不是到处都像陌生人一样忍受了数不清的苦难吗？\par

5. 而这些事（他说）他们说了，\Quote{寻求}他们的\Quote{本国}。啊！这差别多大！他们每天在产痛中，渴望脱离这个世界，返回他们的家乡。但相反，如果我们患上热症，我们就忽略一切，像小孩子一样哭泣，害怕死亡。\par

我们如此受影响并非没有理由。因为我们在此世生活不像陌生人，也不像赶回家乡的人，而是像要去受罚的人，所以我们悲伤，我们应该喜乐时却悲伤；所以我们战栗，像杀人犯或强盗头目当被带到审判台前时，会想起自己所做的一切，因而恐惧发抖。\par

但他们并非如此，而是奋勇前进。连\Person{保禄}也叹息；\Quote{我们}（他说）\Quote{在这帐棚里的人，叹息沉重。}\BibleRef{格后 5:4} 那些与\Person{亚巴郎}在一起的人也是如此；\Quote{他们是客旅}，他说，就全世界而言，而且\Quote{他们寻求一个家乡。}\par

这是什么样的\Quote{家乡}？是他们所离开的那个吗？绝不是。因为有什么阻止他们，如果他们愿意，可以回去再次成为公民呢？但他们寻求的是天上的家乡。因此他们渴望离开此地，并这样中悦了天主；因为\Quote{天主不以被称为他们的天主为耻。}\par

6. 啊！这是何等尊荣！祂屈尊\Quote{被称为他们的天主}。你说什么？祂被称为大地的主，诸天的主，而你把\Quote{祂不以被称为他们的天主为耻}当作一件大事吗？这是伟大而真正伟大的事，是极有福的证明。何以见得？因为祂被称为大地和诸天的主，也如外邦人的主：因祂创造并形成了它们；但对那些圣者而言，［被称为］天主不是这个意思，而是如同一位真正的朋友。\par

我举一个例子向你们说明：如同在大户人家中，当那些管家非常受尊敬，自己管理一切，并对主人有极大自由时，主人就以他们命名，人们可能发现许多这样称呼的。但我说什么呢？正如我们可以说天主，不是外邦人的天主，而是世界的天主，所以我们可以说\Quote{亚巴郎的天主}。但你们不知道这是何等尊荣，因为我们尚未达到。因为现在祂被称为所有基督徒的主，然而这名号超过了我们的功绩：想一想，如果祂被称为一个人的天主，那是何等伟大！祂——被称为整个世界的主——\Quote{不以被称为三个人的天主为耻}。这合情合理：因为圣者们能扭转天平，我不说胜过世界，而是胜过一万个这样的世界。\Quote{因为一个遵行主旨的人，胜过一万个违命之徒。}\BibleRef{德 16:3}\par

现在，他们称自己为\Quote{客旅}就是这个意思，这是明显的。但假设他们说自己是\Quote{客旅}是因为异乡，那么为什么\Person{达味}也称自己为客旅呢？他不是君王吗？不是先知吗？他不是在自己的家乡生活吗？那么他为什么说：\Quote{我是客旅，是寄居的}\BibleRef{咏 38:12}？你怎么是客旅？\Quote{如同}（他说）\Quote{我所有的祖先。}你看见他们也是客旅吗？我们有一个家乡，他意思是，但不是真正的家乡。但你自己怎么是客旅呢？相对于大地而言。所以他们也［是客旅］，相对于大地：因为\Quote{正如他们一样}，他说，我也是；正如他一样，他们也是如此。\par

7. 让我们现在就作成外方人，好使天主\Quote{不以称我们为祂的天主为耻}。因为当祂被称作恶人的天主时，对祂而言是一种羞辱，祂也以他们为耻；正如当祂被称作善良、仁慈和修养德行者的天主时，祂受光荣一样。因为如果我们拒绝被称作恶奴的主人，并丢弃他们；若有人来对我们说：\Quote{某人做了无数坏事，他是你的奴隶，不是吗？}我们立刻说：\Quote{决不是}，以摆脱羞辱：因为奴隶与主人关系密切，羞辱从一人传到另一人。——但他们如此显赫，如此充满信心，以致祂不但不以被称为他们的天主为耻，甚至亲自说：\BibleQuote{我是亚巴郎的天主，依撒格的天主，雅各伯的天主。}\BibleRef{出 3:6}\par

让我们也，我亲爱的，作成\Quote{外方人}；好使天主\Quote{不以我们为耻}；使祂不致以我们为耻，并将我们交付给地狱。那些人曾这样说：\BibleQuote{主啊，我们不是因祢的名预言，因祢的名行了许多奇事吗？}\BibleRef{玛 7:22}但请看基督对他们说什么：\BibleQuote{我不认识你们}：这正如主人对待恶奴的做法，他们跑到主人面前，希望摆脱羞辱。\BibleQuote{我不认识你们}，祂说。那么，祢如何惩罚那些祢所不认识的人呢？我说\Quote{我不认识}，是另一种意思：就是说，\Quote{我否认你们，弃绝你们。}但愿我们不会听到这句致命而可怕的话。因为如果他们驱魔和预言的人都被否认，因为他们的生活不相称；何况我们呢？\par

8. 也许你会问：那些显过预言能力、行过奇迹、驱过魔的人，怎么可能被否认呢？很可能他们日后改变了，变成了邪恶的人；因此他们从先前的德行中一无所获。因为不仅我们的开端要光辉，结局更要光辉。\par

请问，演说家不也竭尽全力使演讲的结尾精彩，以便在掌声中退场吗？公职人员在任期结束时不也作出最精彩的展示吗？摔跤手若不能在最后也更精彩地战胜对手，若他战胜了所有人却在最后被击败，岂不全都徒劳无功？舵手若已横渡整个海洋，却在港口失事，岂不丧失之前所有的辛劳？医生呢？若他使病人摆脱疾病，在即将治愈出院时却毁灭了他，岂不摧毁了一切？同样，在德行方面，凡没有以与开端相称、和谐一致的结局来增补的人，都将毁坏、沦丧。那些从起点光彩夺目、欢欣鼓舞地跃出，后来却变得软弱无力的人，正是如此。因此他们既被剥夺了奖赏，也不被主人承认。\par

让我们听这些事，我们当中贪爱财富的人：因为这是最大的不义。\BibleQuote{贪爱钱财是万恶的根源。}\BibleRef{弟前 6:10}让我们听，我们当中想增加现有财产的人，让我们听，并有时停止贪婪，以免我们听到他们将要听到的同样的话。现在让我们怀着敬畏听，以免那时受到惩罚地听：\BibleQuote{离开我}（祂说），\BibleQuote{我从来不认识你们}\BibleRef{玛 7:23}，就连那时（祂的意思是）当你们显示预言、驱魔的时候也不认识。\par

祂在此也暗示另一件事：即使那时他们也是邪恶的；从一开始，恩宠甚至通过不配的人工作。因为若它通过巴郎工作，那么通过不配的人，为了那些将要受益的人，就更是如此了。\par

但如果连征兆和奇迹都不能使人免受惩罚；那么，若一个人恰好处在司铎的尊位上，就更不用说了：即使他达到最高的尊荣，即使恩宠在他内为圣职而工作，即使为了那些需要他领导的人而做其他一切事，他也要听到：\BibleQuote{我从来不认识你们}，不，连那时当恩宠在你内工作时也不认识。\par

9. 哦！在那里对生活的纯洁将有多严格的审查！单是生活本身如何足以领我们进入天国！而缺乏生活则将人交付给毁灭，即使他有千万个奇迹和征兆。因为没有什么比卓越的生活更令天主喜悦。\BibleQuote{如果你们爱我}\BibleRef{若 14:15}，祂宣布；祂没有说\Quote{行奇迹}，而是什么？\BibleQuote{遵守我的诫命。}又说：\BibleQuote{我称你们为朋友}\BibleRef{若 15:14}，不是当你们驱魔时，而是\BibleQuote{如果你们遵守我的话。}因为那些事来自天主的恩赐；但这事在恩赐之后，也来自我们自己的勤勉。让我们努力成为天主的朋友，而不是继续作祂的仇敌。\par

我们总是说这些事，总是发出这些劝勉，既对自己也对你们：但再无更多收获。因此我也害怕。我本希望保持沉默，以免增加你们的危险。因为当一个人经常听到，却仍不行动，这就会激起主的愤怒。但我也惧怕另一种危险，即沉默的危险，如果我被任命为宣讲圣言的人却保持沉默。\par

那么，我们该怎么做才能得救呢？让我们在有机会时开始实践德行；让我们为自己分派德行，如同农夫分派农活一样；在这个月，让我们克服诽谤、伤害、不义的愤怒；让我们为自己立下规矩，说：今天让我们改正这一点。再次，在这个月，让我们操练忍耐，在另一月，操练另一种德行：当我们习惯了这一德行后，再去学习另一种，如同在学校里学习功课一样，守护已经获得的，并获取其他的。\par

此后，让我们进而轻视财富。首先约束我们的手，不攫取，然后施舍。不要简单地将一切混淆，用同一双手既杀人又施恩。此后，让我们转向另一种德行，再从那一德行转向另一种。\BibleQuote{淫秽、愚妄的谈吐和戏谑，连提都不可在你们中间。}（\BibleRef{弗 5:4}）让我们就这样走在正道上。\par

无需花费金钱，无需劳力，不必流汗，只要有愿意的心，一切便成就。无需长途跋涉，也无需跨越无垠的海洋，只需认真、思想预备，给舌头套上笼头。让我们抛弃不合时宜的责备、愤怒、无节制的欲望、奢侈、挥霍，并从灵魂中除去对财富的渴望、伪誓和惯常的誓言。\par

如果我们这样修德，拔除先前的荆棘，撒下天国的种子，我们就能获得所应许的美善事物。因为农夫会来，将我们收进祂的谷仓，我们必获得一切美善事物，愿我们都能获得，借着我们的主\Person{耶稣基督}的恩宠和慈爱，愿光荣、权能、尊崇归于祂，偕同\Person{圣父}及\Person{圣神}，从今世直到永远。阿们。\par

\SourceRecord{Translated by Frederic Gardiner.；Nicene and Post-Nicene Fathers, First Series；Vol. 14.；Edited by Philip Schaff.；Buffalo, NY: Christian Literature Publishing Co.,；1889.；<http://www.newadvent.org/fathers/240224.htm>.}

% structure: p0001 source=h1 target=section reason=page-title

\section{《希伯来人书》讲道辞 25}

% structure: p0002 source=h2 target=subsection reason=relative-heading-rank; base=h2

\subsection{希 11:17-19}

\Quote{ \Emph{因着信德，} \Emph{［亚巴郎］在受试探的时候，献上了依撒格；那领受了恩许的人，献上了自己的独生子，关于这儿子曾说过：\InnerQuote{藉依撒格你的后裔将得名号。}他以为天主甚至能使人从死者中复活，因此他也从死者中得回了他的儿子，这只是一个预像。}}\par

1. 亚巴郎的信德确实伟大。因为在亚伯尔、诺厄和哈诺客的情况下，只有理智的相反，必须超越人的理智；而在亚巴郎的情形中，不仅必须超越人的理智，还要显明更多的事。因为天主的作为似乎与天主的作为相矛盾；信德与信德相矛盾，诫命与恩许相矛盾。\par

我是这样说的：他曾说：\Quote{离开你的故乡、你的家族，我必将这地方赐给你。} \BibleRef{创 12:1} \Quote{他在这地方没有赐给他产业，连立足之地也没有。} \BibleRef{宗 7:5} 你看见所做的事与恩许相矛盾吗？他又说：\Quote{藉依撒格你的后裔将得名号。} \BibleRef{创 21:12} 亚巴郎相信了；可是天主又说：你将这将要使全世界由他的后裔充满的人献祭给我。你看见诫命与恩许之间的矛盾吗？他命令的事似乎与恩许相抵触，但义人却没有动摇，也没有说他自己受了欺骗。\par

因为，他的意思是，你们不能说：他应许了安逸，却赐予了困苦。因为在我们的情形中，他所应许的，他也实现了。怎么实现的呢？\Quote{在世界上}（他说）\Quote{你们将受苦难。} \BibleRef{若 16:33} \Quote{不背起自己的十字架跟随我的，不配是我的。} \BibleRef{玛 10:38} \Quote{憎恨自己性命的，必不得保全。} \BibleRef{若 12:25} 以及\Quote{凡不舍弃一切所有的跟随我的，不配是我的。} \BibleRef{路 14:27} 又说：\Quote{你们要为了我的缘故被带到总督和君王面前。} \BibleRef{玛 10:18} 又说：\Quote{人的仇敌就是自己家里的人。} \BibleRef{玛 10:36} 但属于安息的事是在来世。\par

但就亚巴郎而言，却不同。他奉命行事，与恩许相抵触；然而即使如此，他没有不安，没有动摇，也不以为受了欺骗。而你们除了所应许的以外，什么也没有忍受，却不安了。\par

2. 他亲耳听到那应许者所说的话与恩许相反，却没有被打扰，反而当作一致的事去做。因为它们确是和谐的；按人的计算虽相矛盾，但藉着信德却是和谐的。宗徒自己教导了我们这是如何做到的，他说：\Quote{他以为天主甚至能使人从死者中复活。} 他的意思是：凭着他相信天主赐予那本不存在的东西并使死者复活的同一信德，他也深信天主会使他在牺牲被杀后复活。因为（我说的是按人的计算）从已死、衰老、已不能生育的子宫中赐予一个孩子，与复活一个被杀的人，都是不可能的。但他以前的信德为将来的事预备了道路。\par

请留意：好事在先，艰难在后，且是在他年老之时。但相反地，对你们而言（他说），悲痛的事在先，好事在后。这是对那些胆敢说\Quote{他应许了我们死后福乐；或许他欺骗了我们}的人说的。他证明\Quote{天主甚至能使人从死者中复活}，如果天主能使人从死者中复活，毫无疑问，他必偿还所应许的一切。\par

但如果亚巴郎这么多年以前，就相信\Quote{天主能使人从死者中复活}，我们更该相信。你看见（如我起初所说）死亡尚未进入，他却立即引导他们盼望复活，并使他们如此确信，以致他们受命时，竟亲手杀死自己的儿子，并甘心献出那些他们期望借以充满世界的人。\par

他也藉着说\Quote{天主试探了亚巴郎} \BibleRef{创 22:1} 显明了另一件事。那么怎样？天主不知道那人高贵且蒙悦纳吗？他为什么试探他？不是为他自己要知道，而是为向别人显明，并使他的刚毅向众人显扬。这里也说明试炼的原因，使他们不致以为这些苦难是因为被天主遗弃。因为对宗徒们而言，必须受试炼，因为有许多人迫害或反对他们；但在亚巴郎的情形，有什么必要为他设计不存在的试炼呢？显然，这试探是藉着他的命令。其他的事是藉着他的许可发生，但这件甚至藉着他的命令。如果试探这样使人蒙悦纳，甚至在没有机会时，天主也锻炼他的运动员，那么我们更该勇敢地忍受一切。\par

在此他着重地说：\Quote{因着信德，当受试探的时候，他献上了依撒格，} 因为除了信德，没有别的理由使他奉献。\par

3. 此后，他继续同一个思想。没有人（他说）能声称，他还有另一个儿子，并且期望恩许由他成就，因此他放心地献上了这个。\Quote{并且}（他的话是）\Quote{他献上了自己的独生子，就是领受恩许的那一位。} 你为什么说\Quote{独生子}？那么，依市玛耳是从谁生的？我的意思是，就恩许之言而言，他是\Quote{独生子}。因此，在说\Quote{独生子}之后，为了表明他是为此而说，他加上了：\Quote{关于这儿子曾说过：\InnerQuote{藉依撒格你的后裔将得名号。}} 就是说，\Quote{从他。} 你看见他是多么钦佩圣祖所作的事？\Quote{藉依撒格你的后裔将得名号，} 他却带了那儿子来献祭。\par

之后，为使没有人以为他这样做是出于绝望，并因这一命令而抛弃了那信德，但能明白这也是真正出于信德，他说他也持守了那信德，尽管它似乎与此相左。但其实并不相左。因为他没有以人的推理来衡量天主的能力，而是将一切交托于信德。因此他敢于说，天主\Quote{能使人从死者中复活}。\par

\BibleQuote{因此他也以预像领回了他。}即借着一只公绵羊，他说。怎样？公绵羊被宰杀后，他得以保全；这样，藉着公绵羊他再次领回了他，以公绵羊代替他献祭。但这些事是预像：因为此处是天主子被杀。\par

我恳求你们，要留意祂的慈爱何其伟大。因为既然要赐予人类一项大恩惠，祂愿意如此行，不是出于恩惠，而是如同欠债者，安排一个人先因天主的命令献出自己的儿子，这样祂自己在献出祂的独生子时看上去就不是做了大事，因为一个人已先于祂行了这事；这样祂就可以被认为不是出于恩宠，而是出于亏欠。因为我们也愿意向所爱的人施以这样的恩惠，希望先显得从他们那里得到了一些小东西，然后再赐予他们一切；我们更夸耀接受而非施与；我们不说\Quote{我给了他这个}，而是说\Quote{我从他那里得到了这个}。\par

\BibleQuote{因此他以预像领回了祂。}即如同一个谜语（因为公绵羊可说是依撒格的预像），或如同一个预表。因为祭献既已完成，依撒格在用意上已被杀，所以天主将他给了圣祖。\par

4. 你们看，我常常说的，在此也显明了吗？当我们证明我们的心意已臻于完美，并显示出我们轻视属世之事时，那时属世之事也赐给了我们；但不会在此之前；以免我们既已受其束缚，接受它们后仍被束缚。你首先摆脱你的奴役（祂说），然后接受，这样你接受时就如同主人而非奴隶。轻视财富，你便会富有。轻视光荣，你便会荣耀。轻视向敌人报仇，你便会达到。轻视安息，你便会得到它，这样你在接受时就如同自由人，而非囚犯或奴隶。\par

因为如同对待小孩，当孩子急切地渴望幼稚的玩物时，我们小心翼翼地藏起它们，比如一个球和诸如此类的东西，免得他因这些事物而妨碍必需的事；但当他不看重它们，不再渴望时，我们就毫无顾虑地给他，因为知道从此它们不会对他造成危害，欲望已没有足够力量使他离开必需之事；同样，当天主看到我们不再急切地渴望今世的事物时，祂就允许我们使用它们。因为我们拥有它们就像自由人和成年人，而非像孩子。\par

因为（作为证据）如果你轻视向敌人报仇，你便会得到它，请听他说：\BibleQuote{如果你的仇人饿了，给他吃；如果他渴了，给他喝。因为你这样做，就是把炭火堆在他头上。}\BibleRef{罗 12:20} 再者，如果你轻视财富，你便会得到它们，请听基督说：\BibleQuote{凡为我的名，舍弃了房屋、或兄弟、或姊妹、或父亲、或母亲、或妻子、或儿女、或田地的，必要领取百倍的赏报，并承受永生。}\BibleRef{玛 19:29} 如果你轻视光荣，你便会得到它，再听基督亲自说：\BibleQuote{谁若愿意在你们中成为大的，就当作你们的仆役。}\BibleRef{玛 20:26} 又说：\BibleQuote{凡高举自己的，必被贬抑；凡贬抑自己的，必被高举。}\BibleRef{玛 23:12}\par

你说什么？我若给我的仇人喝，就惩罚了他吗？我若放弃我的财物，就拥有了它们吗？我若降卑自己，就会被高举吗？是的，祂说，因为我的权能正是如此，以相反者给予相反者。我富有资源和计谋：不要害怕。万物的本性随从我的旨意，而非我依从本性。我做一切事，我不受它们控制；因此我也能改变它们的形态和秩序。\par

5. 你为何对这些事例感到惊奇呢？因为你在所有其他事例中也会发现同样的情况。你若伤害人，便受伤害；你若受伤害，便未受伤害；你若惩罚人，便不是惩罚了别人，而是惩罚了自己。因为经上说：\BibleQuote{喜爱不义的人，也憎恨自己的灵魂。}\BibleRef{咏 10:5} 你看见你并未伤害人，反而受伤害了吗？因此保禄也说：\BibleQuote{你们为什么不情愿受欺呢？}\BibleRef{格前 6:7} 你看见这不算是受欺吗？\par

当你侮辱人时，你便受侮辱。大部分人也部分知道这一点：就如他们彼此说：\Quote{我们走吧，不要自取其辱。}为什么？因为你和他的差别很大：无论你怎样侮辱他，他都认为那是荣耀。让我们在所有情况下都思考这一点，并且超脱侮辱。我会告诉你们如何做。\par

我们若与那穿紫袍的人争辩，应当考虑：侮辱他，就是侮辱我们自己，因为我们变得该受羞辱。告诉我，你是什么意思？当你是一位天上公民，拥有那高超的哲学时，你竟与那\BibleQuote{以地上的事为念}\BibleRef{斐 3:19}的人一起自取其辱吗？因为虽然他拥有无数财富，虽然他有权力，他却还不知道其中的益处。不要因侮辱他而侮辱你自己。宽恕你自己，不是宽恕他。尊重你自己，不是尊重他。难道没有这样的谚语：尊重人的，尊重自己？这是合理的：因为他尊重的不是别人，而是自己。请听某位智者的话：\BibleQuote{你要按你的灵魂的尊荣对待它。}\BibleRef{德 10:28} \BibleQuote{按它的尊荣}指的是什么？如果那人欺骗（意思是），你不要欺骗；如果他侮辱，你不要侮辱。\par

6. 请告诉我，如果有穷人拿了你院子里扔出来的泥土，你会为此而告上法院吗？当然不会。为什么？免得你羞辱自己；免得众人谴责你。同样的事也发生在这里。因为富人是贫穷的，他越富有，在真正贫穷的事上就越贫穷。金子是泥土，被扔在院子里，不在你家里，因为你的家是天堂。那么，你会为此而告上法院吗？天上的公民岂不谴责你？他们岂不将你逐出他们的国度，你这如此卑鄙、如此猥琐的人，竟然选择为一小块泥土争斗？因为即使世界属于你，有人拿走了它，你难道还应该在意吗？\par

你们难道不知道，即使你们拿走世界十倍、百倍、万倍，甚至两倍于此，也无法与天堂中最小的美善相比？那么，羡慕现今之事的，就轻视了那边的事，因为他判断这些事值得努力，尽管远不及另一处。不，他其实无法羡慕那些事。因为当他狂热于这些世俗之事时，又怎能呢？让我们斩断绳索和纠缠：因为世俗之事就是如此。\par

我们要弯腰多久？我们要互相谋划多久，如同野兽，如同鱼类？不，野兽并不互相谋划，而是针对不同族类的动物。比如，熊轻易不杀熊，蛇不杀蛇，它们尊重同族的相同。但是你，与同族的人，有无数的联系：共同的起源、理性能力、认识天主、上万其他事物、自然的力量，他对你是亲属，分享相同本性——你却杀害他，使他陷入无数的邪恶。因为，即使你没有拔剑，没有将右手插入他的脖子，当你让他陷入无数邪恶时，你做的比这更严重。因为如果你做了其他事，你使他免于焦虑，但如今你使他陷入饥饿、奴役、沮丧、许多罪恶。我说这些，且不会停止说，不是准备让你们去谋杀，也不是催促你们犯下比那轻的罪行；而是为了让你们不要自信，好像你们不必交账。\Quote{因为}（它说）\Quote{夺取他人生计的}\BibleRef{德 34:22}，又说求食的，它说。\par

7. 我们终于要管住自己的手，或者不如说，我们不要管住它们，而是荣耀地伸出它们，不是为了抓取，而是为了施舍。不要让我们的手不结果子，也不枯干；因为不行施舍的手是枯干的；而抓取的手是污染和不洁的。\par

不要让任何人用这样的手吃饭；因为这是对受邀者的侮辱。因为，请告诉我，如果有人让我们躺在挂毯、软榻和织金细麻布上，在一座宏大华丽的宅邸中，旁边有一大群侍从，准备了金银托盘，装满了各种昂贵、多样的佳肴，然后催我们吃，只要我们忍受他把手沾上泥浆或人的粪便，然后这样坐下来与我们一同用餐——谁会忍受这种折磨？他难道不会觉得这是侮辱吗？我想他肯定会，并立刻离开。但事实上，你看到的不是手上沾满了真正的污秽，而是连食物本身也沾满了，然而你却不离开、不逃避、不指责。不，如果他是一个有权势的人，你甚至把它当作大事，在吃这些东西时毁灭自己的灵魂。因为贪得无厌比任何泥浆更糟糕；它污染的不是身体，而是灵魂，使之难以洗净。因此，尽管你看见那个坐着吃饭的人，手上脸上都沾满了这种污秽，他的房子充满了污秽，连他的桌子也充满了污秽（因为粪土，即使有比这更不洁的东西，也不如那些食物不洁和污染），你难道还觉得自己受了大尊荣，好像要享受一番吗？\par

难道你不怕保禄吗？他允许我们随意去外邦人的宴席，如果愿意，但甚至不愿意我们去贪得无厌者的宴席。因为他说：\Quote{若有称为弟兄的}\BibleRef{格前 5:11}，这里他意指每一位信徒，不仅仅是过独身生活的人。因为是什么构成弟兄关系？重生的洗；能够呼求天主为父。所以，一位隐修士，如果他是慕道者，就不是弟兄；但信徒即使在世间，也是弟兄。他说：\Quote{若有称为弟兄的}\BibleRef{格前 5:11}。因为那时根本没有过隐修生活的痕迹，但这蒙福的宗徒将他的所有论述指向世间的人。他说：\Quote{若有称为弟兄的，是淫乱的、或贪得无厌的、或醉酒的，这样的人，连吃饭也不可同他一起吃。}但对外邦人则不同：而是\Quote{若有不信的人}，意指外邦人，\Quote{请你们，你们若愿意去，凡摆在你们面前的，只管吃}\BibleRef{格前 10:27}。\par

8. 哦！何等严格！然而我们不仅不避开醉酒者，反而去他们的家，分享他们摆在我们面前的东西。\par

因此一切都颠倒了，一切皆混乱、倾覆、毁灭。因为，请告诉我，如果这样一个人邀请你赴宴，你被视为贫穷和卑微，然后他听见你说：\Quote{既然摆在我面前的是欺诈的果实，我不愿玷污自己的灵魂。}他岂不感到羞辱？岂不窘迫？岂不羞愧？单是这一点就足以纠正他，使他因自己的财富而自认可怜，因你的贫穷而钦佩你，如果他看见自己被你如此郑重地藐视。\par

但我们\Quote{成了}（我不知为何）\Quote{\OriginalTerm{人的奴仆}{servants of men}}，\BibleRef{格前 7:23} 尽管保禄大声呼喊说：\Quote{你们不要成为人的奴仆。} 那么我们是如何成为\Quote{人的奴仆}的呢？因为我们首先成了肚腹、金钱、光荣和所有其余事物的奴仆；我们放弃了基督赐予我们的自由。\par

那么，那成为奴仆的人（请告诉我）将有何等待遇？请听基督说：\Quote{奴仆不永远住在家里。} \BibleRef{若 8:35} 你有一个自足的宣言，即他永不进入天国；因为\Quote{家}正是指此。因为他说：\Quote{在我父的家里有许多住处。} \BibleRef{若 14:2} 那么\Quote{奴仆}就\Quote{不永远住在家里}。祂所说的奴仆是指那\Quote{罪的奴仆}。但那\Quote{不永远住在家里}的，将永远住在地狱里，得不到任何安慰。\par

不，事态已恶化到如此地步，他们甚至用这些不义之财施舍，而许多人接受。因此我们的勇气已崩溃，我们无法责备任何人。但无论如何，至少从今以后，让我们逃避由此产生的祸患；你们这些在泥沼中打滚的人，停止这种玷污，节制你们对这些宴会的狂热，但愿我们如今还能设法使天主对我们仁慈，并得享所应许的福乐。愿我们众人都在我们的主耶稣基督内获得这一切，愿光荣、权能、尊崇归于祂，偕同圣父及圣神，从今世直到永远，万世无疆。阿们。\par

\SourceRecord{Translated by Frederic Gardiner.；Nicene and Post-Nicene Fathers, First Series；Vol. 14.；Edited by Philip Schaff.；Buffalo, NY: Christian Literature Publishing Co.,；1889.；<http://www.newadvent.org/fathers/240225.htm>.}

% structure: p0001 source=h1 target=section reason=page-title

\section{《希伯来人书》第廿六篇讲道}

% structure: p0002 source=h2 target=subsection reason=relative-heading-rank; base=h2

\subsection{\BibleRef{希 11:20-22}}

\BibleQuote{因着信德，依撒格祝福了雅各伯和厄撒乌论及将来的事。因着信德，雅各伯临死时祝福了若瑟的两个儿子，并倚靠着杖头朝拜。因着信德，若瑟临终时提到以色列子民出离的事，并嘱咐关于他的骨骸。}\par

1. \BibleQuote{许多先知和义人}（经上说）\BibleQuote{切愿看见你们所看见的，却没有看见；愿意听见你们所听见的，却没有听见。}\BibleRef{玛 13:17}那么，那些义人岂不知道一切将来的事吗？是的，极其确定。因为若因那些不能接受他的人软弱，圣子未被启示——他以正当的理由向那些在德行上杰出的人启示了。保禄也就此说，他们知道\Quote{将来的事}，即基督的复活。\par

或他不是这个意思：而是说\Quote{因着信德，关于将来的事}[意思是]不是[关于]未来世界，而是关于今世中\Quote{将来的事}。因为一个人寄居异地，怎能赐予如此祝福，若非藉着信德？\par

但另一方面，他获得了祝福，却没有领受。你看，我关于亚巴郎所说的，也可用于雅各伯：他们并未享用祝福，但祝福归于他的后裔，他自己却获得了\Quote{将来的事}。因为我们发现他的兄弟反而享用了祝福。因为[雅各伯]一生都在奴役和雇佣劳作中度过，[处于]危险、阴谋、欺骗和恐惧中；当法郎问他时，他说：\BibleQuote{我寄居的日子又少又苦}\BibleRef{创 47:9}；而另一个却独立自主、非常安全地生活，后来还让[雅各伯]畏惧。那么祝福在哪里成就呢？岂不只在[未来]世界吗？\par

你看，从一开始，恶人享用此世之物，义人却相反吗？然而并非全部。因为看，亚巴郎是义人，他也在此世享用，但伴随着痛苦和试探。因为财富是他所有的一切，其他关于他的事充满痛苦。因为义人不可能不受苦，即使他富有：因为当他愿意被欺骗、受冤屈、忍受一切其他事时，他必定受苦。所以，即使他享用财富，[却]不是没有悲伤。为何？你问。因为他处于痛苦和困境中。但如果那时义人受苦，何况现在呢？\par

\Quote{因着信德，}他说，\Quote{依撒格祝福了雅各伯和厄撒乌论及将来的事}（然而厄撒乌是长子；但他因雅各伯的卓越而先提他）。你看他的信德何其大！他那时从哪里应许给儿子们如此大的祝福？完全出自他对天主的信德。\par

2. \Quote{因着信德，雅各伯临死时祝福了若瑟的两个儿子。}在此我们应当列出完整的祝福，以彰显他的信德和预言。\Quote{并倚靠着杖头，}他说，\Quote{朝拜。}这里，他意思是，他不仅说了话，甚至对未来之事如此确信，以至于也藉行动表明。因为另一位君王将要出自厄弗辣因，因此经上说\BibleQuote{他在杖头俯伏。}也就是说，即使他现在是老人，\Quote{他俯伏}向若瑟，显示全体百姓将要向他行的敬礼。这事已经发生，当他的兄弟\Quote{俯伏}向他时；但以后将通过十支派应验。你看他如何预言日后的事？你看他们的信德何其大？他们多么相信\Quote{关于将来的事}？\par

因为此处有些事，即现在的事，只是忍耐和忍受恶待、领受无益之物的例子，例如亚巴郎和亚伯尔的事例。但另一些是信德的[例子]，如诺厄的事例：有天主，有报应。（因为此处的信德是多方面的：既有报应，也有等待它，不在相同条件下，以及在奖品之前角力。）还有关于若瑟的事，也仅属信德。若瑟听说天主曾应许亚巴郎，祂承诺\BibleQuote{将这地赐给你和你的后裔；}尽管在异地，尚未看见承诺实现，但他仍不犹豫，反而如此相信，以至于\BibleQuote{提到出离的事，并嘱咐关于他的骨骸。}他不仅自己相信，也引导其余人走向信德：使他们常记念出离的事（因为他若非完全确信[此事]，就不会\Quote{嘱咐关于他的骨骸}），盼望他们回归[客纳罕]。\par

因此，当有人说：\Quote{看，连义人也关心自己的坟墓！}让我们回答他们：是为了这个理由：因为他知道\BibleQuote{大地属于上主，及其中所充满的一切。}\BibleRef{咏 24:1}他岂能不知道这事？他生活在如此大的哲学中，终生在埃及度过。然而，如果他愿意，他本可回去，不哀伤或烦恼。但当他将父亲运到那里后，为何也命令他们把自己的骨骸运上去？显然是为了这个理由。\par

但什么？请告诉我，梅瑟自己的骨头不是安放在异地吗？亚郎的、达尼尔和耶肋米亚的，不也是如此吗？至于宗徒们的，我们不知道他们大多数安放在哪里。因为伯多禄、保禄、若望和多默的坟墓是众所周知的；但其余那么多宗徒的坟墓却无人知晓。所以，我们完全不要为此哀叹，也不要如此心胸狭隘。因为无论我们葬在哪里，\Quote{大地属于主，且其中一切所有。} \BibleRef{咏 23:1} 的确，该发生的事必然发生；然而，哀悼、痛哭、悲叹亡者，乃是出于心胸狭隘。\par

% structure: p0013 source=h2 target=subsection reason=relative-heading-rank; base=h2

\subsection{希伯来书 11:23}

3. \BibleRef{希 11:23} \Quote{因着信德，梅瑟出生后，被他的父母隐藏了三个月。} 你们看见吗？在此情况下，他们寄望于死后地上的事。许多事在他们死后才实现。这是针对那些说：\InnerQuote{死后那些在他们活着时未能得到的事为他们成就了；他们也不相信死后会实现。}的人。\par

此外，若瑟并没有说：\InnerQuote{他没有在我生前将土地赐给我，也没有赐给我父亲或祖父，他们的贤德也应受尊敬；难道他会将这些连他们都没有领受的恩惠赐给这些可怜人吗？}他完全没有这样说，而是以信德克服并超越了这一切。\par

他列举了亚伯尔、诺厄、亚巴郎、依撒格、雅各伯、若瑟，全都是卓越可敬的人。他又将话题降到普通人身上，使劝勉更加强烈。因为卓越之人有此类感受并不稀奇，显得不如他们也并不可怕；但若显得连无名之辈都不如，这才是可怕的。他先从梅瑟的父母开始，他们是默默无闻的人，没有像他们的儿子那样伟大。因此，他继续列举甚至妓女和寡妇来增加他所说之事的奇特。因为（他说）\Quote{因着信德}，\Quote{妓女辣哈布没有与那些不信的人一同灭亡，因为她曾和和平平地接待了侦探。}他不仅提及信德的赏报，也提及不信的赏报，如诺厄的例子。\par

但现在我们必须谈梅瑟的父母。法郎下令所有男婴都要被消灭，无人能逃脱危险。他们从何处指望能救自己的孩子？从信德。什么样的信德？（他说）\Quote{他们看见}，\Quote{这孩子俊美。}那景象吸引他们走向信德：如此，从一开始，甚至从襁褓之时，极大的恩宠就倾注在这位义人身上，这不是自然的作用。因为请看，这孩子一出生就显俊美，悦人眼目。这是谁的作为？不是自然的，而是天主的恩宠，这恩宠也激发并坚固了那个外邦女子埃及人，并牵引她前进。\par

然而事实上，在他们身上，信德并没有足够的基础。因为从眼见而相信算什么？但你们（他会说）是从事实相信，并且有许多信德的保证。因为\Quote{甘心忍受你们的财物被抢掠} \BibleRef{希 10:34} 和其他类似的事，正是信德和忍耐的明证。但既然这些希伯来人也曾相信，后来却变得灰心，他便表明那些古圣的信德也是持久不变的，例如亚巴郎的信德，尽管环境似乎与之相抗。\par

（他说）\Quote{而且}，\Quote{他们不怕王的命令，} 尽管那道命令正在执行，但这（他们对孩子的希望）只是一种单纯的期望。这固然是他父母的行动；但梅瑟自己又贡献了什么？\par

% structure: p0020 source=h2 target=subsection reason=relative-heading-rank; base=h2

\subsection{希伯来书 11:24-26}

4. 接下来又是一个适合他们的例子，或者说更伟大的例子。因为他说：\BibleRef{希 11:24} \Quote{因着信德，梅瑟长大以后，拒绝被称为法郎女儿之子，宁愿与天主的子民同受苦害，也不愿暂时享受罪中之乐；他以基督的凌辱比为埃及的宝藏更宝贵的财富，因为他注视着赏报。} 就好像他对他们说：\InnerQuote{你们中没有人离开过王宫，甚至辉煌的王宫，也没有离开过这样的宝藏；也没有人本可成为王子却像梅瑟一样轻看这些。}他不仅简单离开，他用\Quote{拒绝}一词来表达，即憎恶、躲避。因为当天堂摆在眼前时，羡慕埃及的王宫便是多余的了。\par

请看保禄多么巧妙地表达。他没有说：\InnerQuote{以天堂和天堂中的事物}，\InnerQuote{比埃及的宝藏更宝贵}，而是说什么？\Quote{基督的凌辱}。因为为了基督而受凌辱，他认为比这样安逸更好；而这本身即是赏报。\par

（他说）\Quote{宁愿选择}，\Quote{与天主的子民同受苦害。}因为你们确实是为了自己受苦，而他却\Quote{选择}为别人受苦；当他能够虔诚生活并享受福乐时，他自愿投身于如此多的危险之中。\par

（他说）\Quote{而不愿}，\Quote{暂时享受罪中之乐。}他把不愿意\Quote{与其他人同受苦害}称为\Quote{罪}：他说，这就是梅瑟所视为的\Quote{罪}。如果他把不情愿\Quote{与其他人同受苦害}视作\Quote{罪}，那么受苦害必然是一大善事，因为他从王宫中投身于此。\par

但他这样做，是因为他看到眼前一些伟大的事物。\Quote{他以基督的凌辱比为埃及的宝藏更宝贵的财富。} 什么是\Quote{基督的凌辱}？就是像你们这样受凌辱，是基督所受的凌辱；或者说是他为基督的缘故所受的凌辱：因为\Quote{那磐石就是基督} \BibleRef{格前 10:4}；是像你们这样受凌辱。\par

但\Quote{基督的凌辱}是什么呢？那就是，我们弃绝祖先的道路而受凌辱；当我们奔向天主时受虐待。很可能他也受了凌辱，当人对他说：\BibleQuote{你昨天杀那埃及人，今天也要杀我吗？}\BibleRef{出 2:14} 这就是\Quote{基督的凌辱}，至终，至最后一息受虐待：正如祂自己受凌辱，听见那些祂为之被钉十字架的人，那些同族的人说：\BibleQuote{祢如果是天主子，就……}\BibleRef{玛 27:40} 这就是\Quote{基督的凌辱}，当一个人受自己家族的人，或受他正在施恩的人凌辱时。因为\Person{梅瑟}也从那曾受他恩惠的人那里遭受了这些事。\par

他用这些话鼓励他们，指明连\Person{基督}和\Person{梅瑟}这两位卓越人物都受了这些苦。所以这与其说是\Person{梅瑟}的凌辱，不如说是\Quote{基督的凌辱}，因为祂是从\BibleQuote{自己的人}\BibleRef{若 1:11}那里受这些苦的。但两者都没有发闪电，也没有感到愤怒，反而祂受辱骂并忍受一切，而他们却\BibleQuote{摇头}\BibleRef{玛 27:39}。既然读者也很可能听见这样的事，并且盼望赏报，所以他说连\Person{基督}和\Person{梅瑟}都受了类似的苦。如此，安逸是罪的份，而受凌辱是基督的份。那么你愿意要什么呢？\Quote{基督的凌辱}，还是安逸？\par

% structure: p0028 source=h2 target=subsection reason=relative-heading-rank; base=h2

\subsection{希伯来书 11:27}

5. \BibleRef{希 11:27} \BibleQuote{因着信德，他离开了\Place{埃及}，不怕王的愤怒；因为他恒心忍耐，如同看见了那不可见者。}你怎么说？他不害怕吗？然而经上说，当他听见时，他\BibleQuote{害怕了}\BibleRef{出 2:14}，并因此通过逃跑寻求安全，悄悄溜走，秘密撤退；之后他极其害怕。仔细观察措辞：他说\BibleQuote{不怕王的愤怒}，是指他再次挺身而出。因为害怕之人不会再次承担护卫之责，也不会插手此事。但他确实再次承担了，那是将一切交托给天主之人的作为：因为他没有说：\Quote{他在找我，忙于搜索，我不能再参与此事了。}\par

所以连逃跑也是信德之举。那么为什么他没有留下（你说）？为不让自己陷入预先看见的危险。因为最终这会是试探天主：跳进危险之中，说：\Quote{让我们看天主是否会拯救我。}这正是魔鬼对\Person{基督}说的话：\BibleQuote{祢跳下去吧。}\BibleRef{玛 4:6} 你看见这是魔鬼的事吗？无缘无故且毫无目的地将自己投入危险，以试探天主是否会拯救我们？因为他（\Person{梅瑟}）不能再作他们的保护者，当时那些受恩惠的人如此忘恩负义。因此留在那里会是愚蠢无知之举。但这一切之所以发生，是因为\Quote{他恒心忍耐，如同看见了那不可见者。}\par

6. 如果我们也用我们的心思一直看见天主，如果我们常在记念中思念祂，万事在我们眼中都将可忍受，万事可承担；我们将轻易担起一切，超乎一切之上。因为如果一个人看见他所爱的人，更确切地说，记念他而心灵振奋，思想提升，轻易承担一切，并在记念中喜悦；那么一个心中常思念那位曾赐予我们切实之爱的主，并记念祂的人，何时会感到痛苦，或畏惧任何可怕或危险的事？何时他会胆怯？永不。\par

因为万事在我们眼中显得困难，是因我们没有按应有的方式记念天主；因为我们没有常常将祂怀在心中。因为祂很可能公义地对我说：\Quote{你忘记了我，我也要忘记你。}于是祸患成为双重的：既是我们忘记祂，也是祂忘记我们。因为这两件事彼此纠缠，然而又是两件。因为天主的记念效果伟大，祂被我们记念的效果也伟大。前者的结果是我们选择善事；后者的结果是我们成就善事，并使其完成。因此先知说：\BibleQuote{我要从\Place{约旦}地，从\Place{赫尔孟}小山记念你。}\BibleRef{咏 41:6} 在\Place{巴比伦}的人民这样说：在那里，我要记念你。\par

7. 因此我们也要这样，如同在\Place{巴比伦}一样。因为我们虽然没有坐在好战的敌人中间，但仍是在仇敌中间。因为有些人确实坐着作俘虏，但另一些人甚至感觉不到被掳，如\Person{达内尔}，如三个青年\BibleRef{咏 136:1}；他们虽然在充军中，在本地反而成为比掳走他们的君王更光荣的人。那掳走他们的人竟向俘虏致敬。\par

你看见美德多么伟大吗？当他们实际被掳时，他服侍他们如同主人。因此他才是俘虏，而不是他们。如果当他们在本国时，他前来在他们自己的土地向他们致敬，或者如果他们在那里作统治者，那就不如此神奇了。但神奇的是，在他捆绑他们、掳掠他们、将他们安置在自己国内之后，他不羞于在众人面前向他们致敬，并\BibleQuote{献上供物}\BibleRef{达 2:46}。\par

你看见真正辉煌的是那些与天主有关的事，而人间之事如同影子吗？他似乎不知道，他正为自己领来了主人，也不知道他将那些他即将朝拜的人投进了火窑。但在他们看来，这些事如同梦一般。\par

我们当敬畏天主，亲爱的，我们当敬畏祂：即使我们被掳，也比众人更荣耀。愿对天主的敬畏常与我们同在，就无任何事会令我们愁苦，纵然你说到贫穷、疾病、被掳、为奴，或任何其他愁苦之事；甚至这些事本身反会为我们效力，转成益处。那些人曾是俘虏，君王却朝拜了他们；\Person{保禄}是制帐棚的，众人却向他献祭，如同对神一般。\par

8. 这里产生一个问题：你问，为何宗徒们阻止献祭，撕裂衣裳，劝阻他们的企图，并沉痛哀叹地说：\BibleQuote{你们做什么？我们也是和你们有同样性情的人。} \BibleRef{宗 14:15}；而\Person{达尼尔}却没有做任何这类事？\par

因为他也是谦卑的，并将荣耀归于天主，不亚于他们，这从许多地方显而易见。尤其从他是天主所爱这一事实，更显明可见。因为若他将属于天主的尊荣归于自己，天主必不会容他存活，更不会使他得尊荣。其次，因为他甚至极其坦率地说：\BibleQuote{至于我，大王，这奥秘并非因我内有智慧而启示给我。} \BibleRef{达 2:30} 再者，他在狮洞中是为了天主的缘故，当先知给他送食物时，他说：\BibleQuote{因为天主记念了我。} \BibleRef{达 14:38} 他是如此谦卑而痛悔。\par

他在狮洞中是为了天主的缘故，却仍自认不配得祂的记念，不配蒙垂听。然而我们，虽胆敢犯下无数玷污，且是众人中最污秽的，若第一次祈祷未蒙垂听，便退缩了。真的，他们与我们之间的距离是何等巨大，如同天地之间的距离，甚至更大。\par

你说什么呢？在如此多的成就之后，在狮洞中所行的奇迹之后，你竟自认如此谦卑吗？是的，他说；因为无论我们做了什么，\BibleQuote{我们是无用的仆人。} \BibleRef{路 17:10} 如此，他预先实践了福音的训诫，并自视无有。因为他说：\BibleQuote{天主记念了我。} 他的祈祷又充满了何等伟大的谦卑之心。再者，那三个青年也如此说：\BibleQuote{我们犯了罪，行了不义。} \BibleRef{达 4:6} 他们处处显示自己的谦卑。\par

然而\Person{达尼尔}有无数可令他自高自大的机会；但他知道这些机会也是因他不\Emph{自高自大}才临到他，因此他并未毁掉自己的宝藏。因为在所有人中，并在全世界，他都受赞颂，不仅因为君王俯伏在地，向他献祭，视他为神——而这位君王本人在世界各地被尊为神，因为他统治全地（这事从\Person{耶肋米亚}可见：他说：\BibleQuote{谁给大地穿上衣裳，} \BibleRef{耶 43:12}；又说：\BibleQuote{我将它赐给了我的仆人拿步高} \BibleRef{耶 27:6}；此外也从君王在他的信中所说的话可见）。又因为他不只在所在之处，而是在各处受人钦佩，甚至比其余万民都在场亲眼见他更伟大；因为连君王也在信中承认他的顺服与那奇迹。然而，又因他的智慧，他也受人钦佩，因为经上说：\BibleQuote{你比\Person{达尼尔}更明智吗？} \BibleRef{厄 28:3} 然而在这一切之后，他仍是如此谦卑，为主而死万死不辞。\par

那么，你问，既然他如此谦卑，为何不拒绝君王对他的敬拜或祭献呢？\par

9. 这一点我不说，因为我只提这问题就够了，其余的我留给你们，好使至少借此激发你们的思想。（然而，我恳求你们，具备这样的榜样，为了敬畏天主而选择一切事；因为事实上，如果我们真诚地持守将来的事，我们在此世也将得着。）因为他并非出于傲慢而这样做，这从他所说的话显然可见：\BibleQuote{你的赠礼归你自己吧。} \BibleRef{达 5:17}\par

因为除此之外，还有一个问题：他如何在言语上拒绝，却在行动上接受了尊荣，并戴上了金链。\BibleRef{达 5:29}\par

此外，\Person{黑落德}听到呼喊\BibleQuote{这是神的声音，不是人的声音}，\BibleQuote{他没有归光荣于天主，便身体破裂，肠子流了出来} \BibleRef{宗 12:22}，而这个人却接受了甚至属于天主的尊荣，不单是言语。\par

然而有必要说明这是怎么回事。在那个情况（在\Place{吕斯特辣}），人们陷入了更大的偶像崇拜，但在\Person{达尼尔}的情况并不如此。怎么呢？因为他的受尊崇本身，正是天主的荣耀。因此他预先说：\BibleQuote{至于我，不是因我内有任何智慧。} \BibleRef{达 2:30} 此外，他甚至似乎并没有接受祭献。因为君王（如经上所记）说他们应该献祭，但似乎并未付诸行动。但在那里（\Place{吕斯特辣}），他们甚至献上了公牛，并称呼一个\BibleQuote{\Person{则乌斯}}，另一个\BibleQuote{\Person{赫尔默斯}}。\BibleRef{宗 14:12}\par

那时他接受了金链，为使自己被认识；但为何不见他拒绝祭献呢？因为在另一个情况中，他们也没有实际献祭，只是试图这样做，而宗徒们阻止了他们；因此在这里，他也应当立刻拒绝（敬拜）。在那里是全体民众，在这里是君王。为何他没有阻止（\Person{达尼尔}）？他预先表达了：即（\Person{拿步高}）并非以敬拜神的方式向他献祭以推翻宗教崇拜，而是为更伟大的奇迹。如何呢？那（\Person{拿步高}）颁布谕令乃是为天主的缘故；因此（\Person{达尼尔}）没有削减所给的尊荣。但那些人（在\Place{吕斯特辣}）并非如此行事，而是以为他们真是神。因此他们被拒绝了。\par

而在这里，在向他行了敬礼之后，他做了这些事：因为他并非把他当作神敬拜，而是当作一个智者。\par

但他是否献了祭，并不清楚；即便他真的献了，也未必是\Person{达尼尔}所接纳的。\par

他称他为\Quote{\Person{贝尔特沙匝}，的名字}他自己的\Quote{神}，这又有什么呢？看来他们并不认为自己的神有什么了不起，因为他竟这样称呼一个俘虏；他命令所有人都朝拜那形状各异、色彩斑斓的像，又崇拜那龙。\par

此外，巴比伦人比\Place{吕斯特拉}的那些人愚昧得多。因此，不能立刻引导他们走向这事。还可以说许多事，但到此为止就够了。\par

因此，如果我们希望获得一切美善，就让我们寻求天主的事物。因为正如寻求世界事物的人，既失去世界，也失去天主的事物；同样，偏爱天主事物的人，两者都得着。那么，让我们不要寻求这些，而要寻求那些，好使我们也能获得在我们主基督耶稣内所应许的美善；愿光荣、权能、尊崇归于祂，与圣父及圣神，自今至永远，世世无穷。阿们。\par

\SourceRecord{Translated by Frederic Gardiner.；Nicene and Post-Nicene Fathers, First Series；Vol. 14.；Edited by Philip Schaff.；Buffalo, NY: Christian Literature Publishing Co.,；1889.；<http://www.newadvent.org/fathers/240226.htm>.}

% structure: p0001 source=h1 target=section reason=page-title

\section{\WorkTitle{论希伯来书 第27篇}}

% structure: p0002 source=h2 target=subsection reason=relative-heading-rank; base=h2

\subsection{\BibleRef{希 11:28-31}}

\Quote{\Emph{因着信德，他守了逾越节和洒血礼，免得那毁灭长子者触摸他们。} \Emph{因着信德，他们过了红海，如行干地；埃及人试着过去，就被淹死了。} \Emph{因着信德，耶里哥城墙被绕行七天后，倒塌了。} \Emph{因着信德，妓女辣哈布没有与那些不信的人一同灭亡，因为她曾平安地接待了探子。}}\par

1. 保禄惯于在论述中顺便确立许多道理，并且思想极其丰富。因为圣神的恩宠就是如此。他并非用多言囊括少许思想，而是以简练的表达包含伟大而多样的思想。请看，至少在他劝勉并论及信德时，他是如何提醒我们关于何等类型和奥秘的事，而我们是拥有其实体的。\Quote{\Emph{因着信德}，他说，\Emph{他守了逾越节和洒血礼，免得那毁灭长子者触摸他们。}}\par

但\Quote{洒血礼}是什么？每家宰杀一只羔羊，把血涂在门框上，这是抵御埃及毁灭者的方法。如果一只羔羊的血在埃及人中间并在如此大的毁灭中保护了犹太人安然无恙，那么基督的血岂不更能拯救我们吗？我们不是把血涂在门框上，而是涂在灵魂上。因为即使在现今，这深夜中毁灭者仍四处游荡；但让我们以那祭献装备自己。（他称\Quote{洒血}为傅油。）因为天主已把我们领出埃及，领出黑暗，领出偶像崇拜。\par

虽然所做的事本身算不得什么，但所成就的却是伟大的。因为所做的事是血，但所成就的是救恩，是阻止并防止毁灭。天使惧怕那血；因为他知道那是何等的预像；他想到了主的死亡，因而战栗；所以他不敢触摸门框。\par

梅瑟说：涂抹，他们就涂抹了，并且充满信心。而你，拥有羔羊本身的血，难道不该充满信心吗？\par

2. \Quote{因着信德，他们过了红海，如行干地。} 他再次将一个整个民族与另一个民族相比，免得他们说：我们不能像圣人们一样。\par

\Quote{因着信德}（他说）\Quote{他们过了红海，如行干地；埃及人试着过去，就被淹死了。} 这里他也引导他们回忆起在埃及所受的苦。\par

怎么\Quote{因着信德}？因为他们盼望能过海，因此他们祈祷；或者更确切地说，是梅瑟祈祷。你可看见，信德处处超越人的理性、软弱和卑微？你可看见，他们同时既相信又惧怕惩罚，无论是门上的血还是在红海？\par

他清楚表明那水是真实的，通过那些跌入其中并窒息的人；那并非幻象：但如狮子坑中，被吞噬者证明了事实的真实性；又如火窑中，被焚烧的人证明了事实；同样在这里，你看到同样的事物对一些人成为救恩和荣耀的原因，对另一些人则成为毁灭的原因。\par

信德是如此大的善。当我们陷入困惑时，我们就被拯救，即使我们走到死亡本身，即使我们的处境绝望。因为那时还有什么出路？他们没有武装，被埃及人和大海包围；逃跑必被淹死，否则落入埃及人之手。但祂从不可能中拯救了他们。那在一种情况下如干地铺展的，在另一种情况下如大海吞没了敌人。前者它忘了自己的本性；后者它甚至武装起来对抗它们。\BibleRef{智 19:20}\par

3. \Quote{因着信德，耶里哥城墙被绕行七天后，倒塌了。} 因为确实，号角声不能推倒石头，即使吹上一万年；但信德能成就万事。\par

你可看见，在所有情形下，都不是出于自然的顺序，也不是由于任何自然法则而改变，而是事事超乎意料？因此，在这件事上也是如此，事事超乎意料。因为他再三说过，我们应当信赖未来的希望，因此他合理地引入所有这些论证，表明不仅现在，而且从一开始，所有的奇迹都是借着它而成就和完成的。\par

\Quote{因着信德，妓女辣哈布没有与那些不信的人一同灭亡，因为她曾平安地接待了探子。} 那么，如果你们比一个妓女更不信，那就可耻了。她仅仅听了那些人所说的话，就立即相信了。结果也随即而来；当众人都灭亡时，只有她得救。她没有对自己说：我要和许多朋友在一起。她没有说：难道我可能比这些不信的明智之人更明智——而我竟要相信吗？她没有说这样的话，而是相信了所发生的事，而那些事正是他们很可能要遭受的。\par

% structure: p0016 source=h2 target=subsection reason=relative-heading-rank; base=h2

\subsection{\BibleRef{希 11:32}}

4. \BibleRef{希 11:32} \Quote{我还要说什么呢？时间不够我叙述了。} 此后他不再列出名字；而是以一名妓女结束，并以那人的身份羞辱他们之后，他不再详述那些历史，免得被认为冗长。然而他并没有弃置不顾，而是匆匆掠过，两者都处理得非常明智：避免令人厌烦，又不破坏论述的紧凑；他既没有完全沉默，也没有说得使人烦恼；因为他实现了两个目的。因为当一个人在辩论中激烈争辩时，如果他继续争辩，他会使听者疲倦，在听者已被说服时还烦扰他，并博得好胜虚荣的名声。因为他应当适应有益的事。\par

\Quote{我还要说什么呢}（他说）？\Quote{时间不够我叙述基德红、巴辣克、三松、依弗大、达味和撒慕尔以及众先知的事。}\par

有些人指责保禄，因为他把巴辣克、三松和依弗特放在这里。你怎么说？在引进了那个妓女之后，难道他不该引进这些人吗？请不必告诉我他们其余的生活，只问他们是否不曾相信并在信德中闪耀。\par

% structure: p0020 source=h2 target=subsection reason=relative-heading-rank; base=h2

\subsection{希伯来人书 11:33}

他说：\Quote{\BibleQuote{众先知}}，\BibleRef{希 11:33}\Quote{\BibleQuote{因着信德，征服了列国}}。你看，他在这里并没有见证他们的生活是辉煌的；因为这并不是问题的焦点：而至今所探讨的是关于他们的信德。请告诉我，他们岂不是因着信德成就了一切吗？\par

他说：\Quote{\BibleQuote{因着信德}}，\Quote{\BibleQuote{他们征服了列国}}；那些跟随基德红的人。\Quote{\BibleQuote{行了正义}}；谁？同样是那些人。显然他在这里指的是仁慈。\par

我认为他说\Quote{\BibleQuote{获得了恩许}}是指达味说的。但这些恩许是哪种呢？就是那些话，其中天主说他的\Quote{\BibleQuote{后裔要坐在}}他的\Quote{\BibleQuote{宝座上}}。\BibleRef{咏 131:12}\par

% structure: p0024 source=h2 target=subsection reason=relative-heading-rank; base=h2

\subsection{希伯来人书 11:34}

\Quote{\BibleQuote{堵住了狮子的口}}，\BibleRef{希 11:34}\Quote{\BibleQuote{熄灭了烈火的力量，逃脱了剑刃的锋刃}}。看，他们如何在死亡本身之中：达尼尔被狮子包围，三个青年在火窑中，以色列人、亚巴郎、依撒格、雅各伯在各种试探中；然而即使如此，他们也没有绝望。因为这就是信德：当事情变得不利时，我们应当相信并没有发生什么不利的事，而是一切都井然有序。\par

\Quote{\BibleQuote{逃脱了剑刃的锋刃}}。我认为他又在说那三个青年。\par

\Quote{\BibleQuote{从软弱中变为刚强}}。这里他暗示了从巴比伦回归时发生的事。因为\Quote{\BibleQuote{从软弱中}}，就是从被掳中。当犹太人的境况已变得绝望，当他们不比枯骨好多少时，谁能预料他们会从巴比伦归来，而且不仅是归来，还要\Quote{\BibleQuote{成为英勇}}并\Quote{\BibleQuote{使外邦的军队逃跑}}？但有人说：\Quote{我们没有遇到这样的事}。然而这些是\Quote{\BibleQuote{将来之事}}的预象。\par

% structure: p0028 source=h2 target=subsection reason=relative-heading-rank; base=h2

\subsection{希伯来人书 11:35}

\Quote{\BibleQuote{妇人得了她们已死的亲人复活}}。\BibleRef{希 11:35}他这里说的是关于先知们的事，厄里叟、厄里亚；因为他们曾使死人复活。\par

5. \BibleRef{希 11:35}\Quote{\BibleQuote{又有人受了酷刑，不肯接受释放，为获得更好的复活}}。但我们没有获得复活。然而，他的意思是，我能证明他们也被斩首，且\Quote{\BibleQuote{不肯接受[释放]，为获得更好的复活}}。因为，请告诉我，当他们可以活着时，为什么他们不选择活着？他们岂不是显然在寻求更好的生命吗？那些使他人复活的人，自己却选择死亡，为\Quote{\BibleQuote{获得更好的复活}}，不像那些妇人的孩子所得的复活。\par

这里我认为他暗示了若翰和雅各伯。因为斩首被称为\Quote{\BibleQuote{酷刑}}。他们仍能看见太阳，仍能停止指责[罪人]，然而他们却选择了死亡；那些使他人复活的人，自己选择了死亡，\Quote{\BibleQuote{为获得更好的复活}}。\par

% structure: p0032 source=h2 target=subsection reason=relative-heading-rank; base=h2

\subsection{希伯来人书 11:36}

\BibleRef{希 11:36}\Quote{\BibleQuote{又有人受了戏弄和鞭打的试探，更有甚者，还有锁链和监禁}}。他以这些结束；这些是更切近的事。因为这些[例子]特别带来安慰，当苦难来自同一原因时，因为即使你提到更极端的事，但若不是来自同一原因，你就毫无效果。因此他以这个结束了他的讲论，提到\Quote{\BibleQuote{锁链、监禁、鞭打、石击}}，暗示斯德望的例子，也暗示匝加利亚的例子。\par

因此他补充说：\Quote{\BibleQuote{他们被刀杀死}}。你说什么？有些\Quote{\BibleQuote{逃脱了剑刃的锋刃}}，而有些\Quote{\BibleQuote{被刀杀死}}。\BibleRef{希 11:34}这是什么？你称赞哪一个？你钦佩哪一个？后者还是前者？不，他说：前者固然适合你，而后者，因为信德甚至在死亡本身上也是坚强的，并且这是将来之事的预象。因为信德有两种奇妙特性：它既能成就大事，又能承受大事，并把自身看作没有承受什么。\par

而且你不能说（他说）这些人是罪人和无价值的。因为即使你把整个世界放在他们对面，我发现他们压低了天平，更有价值。那么他们在这生命中要接受什么呢？这里他提升他们的思想，教导他们不要固着于眼前的事物，而要思念比这今生一切更大的事，因为\Quote{\BibleQuote{世界配不上}}他们。那么你希望在这里接受什么呢？因为接受你在这里的赏报是对你的侮辱。\par

6. 因此，让我们不要思念世俗的事物，也不要在此寻求我们的赏报，也不要如此卑鄙。因为如果\Quote{\BibleQuote{整个}} \Quote{\BibleQuote{世界配不上}}他们，你为什么寻求它的一部分呢？这是有理由的；因为他们是天主的朋友。\par

现在这里\Quote{\BibleQuote{世界}}他是指人，还是指受造物本身？两者都是：因为圣经习惯用这个词指两者。如果整个受造物，他会说，连同属于它的人类，放在天平里，它们仍不能与这些人的价值相等；这是有理的。因为一万斗糠和草不能与十颗珍珠相等，所以他们也不能；因为\Quote{\BibleQuote{一个遵行主旨的人，胜过一万个违命者}} \BibleRef{德 16:3}；这里\Quote{\BibleQuote{一万个}}指的是[不仅]许多，而是无限的众多。\par

试思义人有多大的价值。农的儿子若苏厄说：\BibleQuote{太阳停在基贝红，月亮停在厄隆山谷}\BibleRef{苏 10:12}，事情就这样成了。让整个世界来，或者两三个、四个、十个、二十个世界，让他们说并做这事，他们却无能为力。但天主的友人命令他友人的受造物，或者更确切地说，他恳求他的友人，仆人们便顺从了，他在地上向天上的发令。你看见这些事是为服役而完成它们被指定的进程吗？\par

这比梅瑟的奇迹更大。为何？因为命令海与命令天体不是相同的事。前者固然伟大，非常伟大，但仍不能与后者相比。\par

这是为什么？若苏厄的名字是一个预像。因此，正因为这名字本身，受造物敬畏他。那么怎样？难道没有别的人叫耶稣吗？有；但这个人之所以如此称呼，是作为预像；因为他原被称为曷舍亚。因此名字被改了：原来是一个预告和预言。他带领人民进入应许之地，正如耶稣带领人进入天堂；法律却没有，因为连梅瑟也没有带他们进去，而是留在了外面。法律没有能力带入，恩宠才有。你看见从起初就预先勾勒的预像吗？他向受造物发令，或者更确切地说，向受造物的主要部分，向头部本身发令，而他站在下面；好叫当你看见取了人形的耶稣说同样的话时，不致惊慌，也不致以为奇怪。他在梅瑟还活着的时候，就扭转了战争。这样，法律还活着的时候，祂管理万有，但并非公开地。\par

7. 但我们该思想圣人们的美德有多大。如果\Emph{在这里}他们行这样的事，如果\Emph{在这里}他们做这样的事，像天使所做的一样，那么在天上呢？他们的光辉该有多大？\par

也许你们每人都会希望自己能命令太阳和月亮。（此时，那些主张天是球形的人会说什么？因为他为何不只说\Quote{太阳停住}，而加上了\BibleQuote{太阳停在厄隆山谷}，就是他要使白日更长吗？这事在希则克雅时代也发生过。太阳后退了。这比前者更奇妙，因为它走向相反的方向，尚未完成它的行程。）\par

如果我们愿意，我们将获得比这些更大的事。因为基督向我们应许了什么？不是使我们叫太阳或月亮停住，也不是使太阳后退，而是什么？\BibleQuote{我和圣父要到祂那里去}，祂说，\BibleQuote{并要在他那里作我们的住所。}\BibleRef{若 14:23} 当万有的主亲自下来与我同在时，我还要太阳、月亮和这些奇迹做什么？我不需要这些。因为其中哪一样是我需要的呢？祂自己将成为我的太阳和光明。因为请告诉我，如果你进入一座宫殿，你会选择什么：能够重新安排其中一些固定的事物呢，还是使君王成为亲密的朋友，以致说服他与你同住？后者远胜于前者。\par

8. 但有的人说，人命令什么，基督也命令什么，这有什么稀奇？但你说，基督不需要圣父，祂凭自己的权能行事。好。那么你首先承认并说祂不需要圣父，凭自己的权能行事；然后我要问你，祂的祈祷岂不是迁就和安排（因为基督当然不会比农的儿子若苏厄低劣），并且是为教导我们？因为当你听见一位教师咿呀学语，念着字母，你并不会说他是无知的；当他问某个字母在哪里时，你知道他并非出于无知而问，而是希望引导学生；同样，基督作祈祷并非因为需要祈祷，而是渴望引导你，使你不断地致力于祈祷，恒心不懈、清醒而极其警醒地祈祷。\par

我所说的警醒，不仅指夜间起来，也指白天祈祷时要清醒。因为这样的人被称为警醒。因为在夜间祈祷也可能睡着，在白天祈祷也可能清醒，当灵魂向天主伸展，当它思念与谁交谈、向谁说话，当它想到天使们战战兢兢地侍立，而他却张着嘴、搔着痒走近时。\par

9. 祈祷是一强大的武器，如果以适当的心意献上。为叫你知道它的力量，持续的恳求胜过了无耻、不义、野蛮的残忍和傲慢的鲁莽。因为祂说：\BibleQuote{听听那不义的法官说了什么。}\BibleRef{路 18:6} 再者，它又胜过了懒惰，友谊未能成就的，持续的恳求却成就了：\BibleQuote{他虽然不因他是朋友而给他}（祂说），\BibleQuote{但因他切求，就起来给他。}\BibleRef{路 11:8} 持续的坚持使那不配的妇人成了配得的。\BibleQuote{拿儿女的饼扔给小狗，是不对的。}她说：\BibleQuote{是啊！主，可是小狗也吃主人桌子上掉下来的碎屑。}\BibleRef{玛 15:26} 让我们致力于祈祷。如果带着诚恳、不图虚荣、以纯朴的心献上，它是一强大的武器。它曾扭转战争，造福了整个本不配得的民族。\BibleQuote{我听到了他们的哀号，}祂说，\BibleQuote{就下来拯救他们。}\BibleRef{宗 7:34} 它本身就是一救命的良药，有能力预防罪恶并医治过犯。孤独的寡妇曾勤于此道。\BibleRef{弟前 5:5}\par

如果我们谦卑地祈祷，捶着胸膛像税吏一样，说出他所说的，如果我们说：\BibleQuote{可怜我这个罪人吧！}\BibleRef{路 18:13}，我们必将获得一切。因为即使我们不是税吏，我们也有不亚于他的其他罪恶。\par

因为你不要对我说，你只是在某件小事上犯了错，因为事情具有相同的性质。正如一个人无论杀死的是孩童还是成人，都同样被称为杀人犯；同样，无论他在多还是少上欺诈，也都被称为欺诈。是的，记恨也不算小事，甚至是重大的罪。因为经上说：\\BibleQuote{记恨者的道路趋向死亡。}\\BibleRef{箴 12:28}并且\\BibleQuote{无缘无故向弟兄发怒的，应受地狱的裁判；并且称弟兄为傻瓜}\\BibleRef{玛 5:22}、愚昧和无数的此类事情。\par

但我们甚至不配地分享了那伟大的奥迹，我们嫉妒，我们辱骂。我们中有些人甚至常常醉酒。然而这些事中的每一件，即使单独来看，也足以将我们逐出天国；当它们都聚集在一起时，我们还有什么安慰呢？我们需要大量的补赎，亲爱的，大量的祈祷，大量的忍耐，大量的恒心，才能使我们获得那已应许给我们的美善事物。\par

那么，让我们也说：\\Quote{天主，可怜我这个罪人吧！}不，我们不仅要这样说，还要这样想；如果别人这样称呼我们，我们也不要生气。他听了\\BibleQuote{我不像这个税吏}\\BibleRef{路 18:11}这句话，并没有被激怒，反而痛悔。他接受了责难，也就消除了责难。那人说出了伤口，他就寻求药物。那么，让我们说：\\BibleQuote{天主，可怜我这个罪人吧！}\\BibleRef{路 18:13}但即使别人这样称呼我们，我们也不要愤慨。\par

但如果我们说了自己千万句坏话，而听到别人说时却感到烦恼，那就不是谦卑，也不是告解，而是炫耀和虚荣。你说，自称罪人是炫耀吗？是的，因为我们获得了谦卑的名声，我们被钦佩，被称赞；而如果我们说自己相反的话，我们就被鄙视。所以，我们这样做也是为了名声。但什么是谦卑呢？就是当别人责骂我们时，我们忍受，承认自己的过错，忍受恶言。然而，即使这样也不是谦卑的标志，而是坦率。但现在我们称自己为罪人、不配的人，以及成千上万其他的名字，但如果别人用其中一个称呼我们，我们就烦恼，就变得野蛮。你看，这不是告解，甚至不是坦率吗？你曾说自己是那样的人：如果听到别人也这样说，并且受到责备，不要愤慨。\par

这样，当别人责备你时，你的罪就变得轻省了：因为他们确实给自己加上了重担，却引导你走向智慧。听一听有福的达味在史米诅咒他时所说的话：\\BibleQuote{由他吧，因为上主吩咐他，好使他垂顾我的谦卑，并且上主今日必以善报还我这咒骂。}\\BibleRef{撒下 16:11}\par

但你在过分地数说自己坏话的同时，如果听不到别人对最正义者应有的称赞，就怒气冲冲。你看，你是在戏弄那些不可戏弄的事吗？因为我们甚至拒绝赞美，是为了得到更多的赞美，以便获得更高的颂扬，从而更加被钦佩。因此，当我们拒绝接受称赞时，我们其实是在增加它们。我们做一切事都是为了名声，而不是为了真理。因此，一切都是空洞的，一切都是不切实际的。所以我恳求你们，现在至少从这邪恶之母——虚荣中退出，按照天主所悦纳的生活而生活，这样你们就能获得将来的美善，在耶稣基督我们的主内，愿光荣归于祂和圣父，连同祂的圣神，从今世直到永远。阿们。\par

\SourceRecord{Translated by Frederic Gardiner.；Nicene and Post-Nicene Fathers, First Series；Vol. 14.；Edited by Philip Schaff.；Buffalo, NY: Christian Literature Publishing Co.,；1889.；<http://www.newadvent.org/fathers/240227.htm>.}

% structure: p0001 source=h1 target=section reason=page-title

\section{《希伯来书》第二十八篇讲道}

% structure: p0002 source=h2 target=subsection reason=relative-heading-rank; base=h2

\subsection{希 11:37-40}

\BibleQuote{\Emph{他们披着绵羊皮、山羊皮，四处漂流，受穷乏、患难、苦害（这世界配不上他们）；在旷野、山岭、山洞、地穴漂流无定。}}\par

1. 实在说来，每当我一思想圣人们的成就，便常为自己的处境感到沮丧，因为那些圣人终生所经历的事，我们连梦中都不曾经验过——他们并非因罪受罚，而是始终行义却常受苦难。\par

请想一想，我恳求你们，厄里亚——我们今天的话题又转到他身上；保禄在此处提到他，并以他作为例证的结尾。这个例子正适合他们的情形。在论及宗徒们的事迹，说\BibleQuote{他们被刀杀、被石头砸死}之后，他又回到厄里亚，因为厄里亚也遭受了同样的事。\BibleRef{列下 1:8}既然他们可能尚未十分重视宗徒们，他便从那位被提升天、受特别敬仰的厄里亚身上汲取劝勉和安慰。\par

因为他说：\BibleQuote{他们披着绵羊皮、山羊皮，四处漂流，受穷乏、患难、苦害，这世界配不上他们。}\par

他们连衣服都没有，他说，由于极度的苦难；没有城市、没有房屋、没有住处。正如基督所说的：\BibleQuote{人子却没有枕头的地方。}\BibleRef{玛 8:20}何止没有住处？连立足之地都没有：即使到了旷野，也得不到安息。因为经上不是说\Quote{他们在旷野坐下}，而是说即使在旷野，他们也逃亡、被迫离开，不仅从有人烟之地，而且从无人居住之地被赶出。他又提醒他们当时所处的境况以及在那里发生的事。\par

接着他说：他们为了基督的缘故而控告你们。他们对厄里亚又有何指控，竟驱赶他、迫害他、迫使他与饥荒搏斗？这些正是希伯来人当时所受的。至少，经上说弟兄们决定送救济给受难的门徒。\BibleQuote{每人都按自己的能力，决定送救济住在犹太的弟兄们}\BibleRef{宗 11:29}，这也是这些人的情形。\par

\Quote{受苦害}（或作\Quote{受虐待}），他说，即遭受痛苦，在旅程中、在危险中。\par

但\Quote{他们四处漂流}是什么意思？他说：\BibleQuote{在旷野、山岭、山洞、地穴漂流无定}，如同流亡者、被遗弃者，如同犯了最卑劣罪行的人、被认为不配见太阳的人；他们在旷野中也找不到避难所，必须不断逃亡、寻找藏身之处、将自己活埋在地下、始终活在恐惧中。\par

2. 那么，如此巨大转变的酬报是什么？补偿是什么？\par

他们还没有得到，仍在等待；在这样巨大的苦难中死去之后，他们还没有得到酬报。他们在许多世代之前就赢得了胜利，却还没有得到酬报。而你们仍在奋斗中，就感到烦恼吗？\par

你们也想想，亚巴郎和宗徒保禄坐着等待，直等到你们得以成全，然后他们才能得到酬报，这是何等重大的一件事。因为救主早已预先告诉他们：除非我们也到场，他不会给他们酬报。正如一位慈爱的父亲对表现良好、完成工作儿子们说：除非你们的兄弟来到，我不会给你们饭吃。你们还没有得到酬报，就烦恼吗？那么亚伯尔怎么办？他是最早的胜利者，却坐着未受冠冕。诺厄怎么办？那些早期的人怎么办？他们正在等待你们和你们以后的人。\par

你们看见吗？我们比他们更有优势。因为他说：\BibleQuote{天主为我们预备了更好的事。}\BibleRef{希 11:40}为了使他们在我们之前受冠冕而显得比我们更有优势，他指定了同一个加冕的时刻给众人：那在许多年前就取得胜利的人，将与你一同接受冠冕。你们看见他的慈爱眷顾吗？\par

他没有说：\BibleQuote{他们不经过我们就不能受冠冕}，而是说：\BibleQuote{他们不经过我们就不能得以成全。}意思是到那时他们才显得完全。他们在战斗上比我们早，但在冠冕上不比我们早。他没有亏待他们，而是尊荣了我们。因为他们也等待弟兄们。如果我们\BibleQuote{同属一个身体}，那么这身体全体一起受冠冕、而非一部分一部分受冠冕，快乐就更大。义人在这方面也值得钦佩：他们以弟兄的福祉为喜乐，如同自己的福祉。所以，对他们自己来说，与自己的肢体一同受冠冕是符合他们意愿的。全体一起受荣耀，是极大的喜悦。\par

% structure: p0016 source=h2 target=subsection reason=relative-heading-rank; base=h2

\subsection{希 12:1}

3. \BibleRef{希 12:1} \BibleQuote{所以，我们既有如此众多如云的证人围绕着我们。}（他说）在许多地方，圣经从对应的事物中汲取安慰。正如先知所说：\BibleQuote{避暑所和躲避风暴与雨水的藏身处}\BibleRef{依 4:6}。这里他也这样说：对圣人们的记忆，如同云彩对受烈日暴晒的人一样，使那被苦难压制的灵魂重新稳定和恢复。\par

他没有说：\Quote{高高悬在我们上方}，而是说：\Quote{围绕着我们}，这更甚于前者；这样我们更安全。\par

什么样的\Quote{云}？\BibleQuote{如云的证人。}他不仅称新约中的人，也称旧约中的人为\Quote{证人}（或\Quote{殉道者}）。因为他们也见证了天主的伟大，例如三个青年、厄里亚的同代人、所有先知。\par

\Quote{放下一切。} \Quote{一切}：什么？就是酣睡、冷漠、低劣的推理，一切属人的事物。\par

\Quote{且容易缠累我们的罪}；\OriginalTerm{容易缠累的}{\GreekText{εὐπερίστατον}}，就是说，要么是\Quote{容易围绕我们的}，要么是\Quote{容易被围绕的}，但更可能是后者。因为若我们愿意，战胜罪是容易的。\par

\Quote{让我们以忍耐}（他说）\Quote{跑那摆在我们前面的赛程。}祂没有说\Quote{让我们像拳击手那样搏斗}，也没有说\Quote{让我们摔跤}，也没有说\Quote{让我们作战}；而是把其中最轻的——赛跑——提了出来。祂也没有说\Quote{让我们加长路程}；而是说\Quote{让我们在此忍耐坚持，不要气馁。}\Quote{让我们跑}（他说）\Quote{那摆在我们前面的赛程。}\par

4. 接下来，作为他劝勉的总结，他既放在开始也放在结尾的，就是基督本身。\BibleRef{希 12:2} \Quote{仰望}（他说）\Quote{为我们的信德创始成终的耶稣}；这正是基督自己不断地对门徒说的话：\Quote{若有人称家主为贝耳则步，何况他的家人呢？}\BibleRef{玛 10:25} 又说：\Quote{门徒不能高过老师，仆人不能高过主人。}\BibleRef{玛 10:24}\par

% structure: p0024 source=h2 target=subsection reason=relative-heading-rank; base=h2

\subsection{希伯来书 12:2}

\Quote{仰望}（他说），就是说，让我们学习奔跑。因为在一切技艺和竞赛中，我们通过注视我们的师傅，通过视觉接受某些规则，从而将技艺印在心上；所以在这里，若我们愿意奔跑并学会跑得好，让我们仰望基督，就是耶稣，\Quote{为我们信德创始成终的那位。}这是什么意思？祂将信德放在我们内。因为祂对门徒说：\Quote{不是你们拣选了我，而是我拣选了你们}\BibleRef{若 15:16}；保禄也说：\Quote{到那时我就要认清，如同我被认清一样。}\BibleRef{格前 13:12} 祂将开端放在我们内，祂也将完成。\par

\Quote{祂}，他说，\Quote{因那摆在面前的喜乐，忍受了十字架，轻看羞辱。}这就是说，如果祂愿意，祂本可以完全不受苦。因为\Quote{祂没有犯过罪，口中也找不到诡诈}\BibleRef{伯前 2:22}；正如祂在福音中所说的：\Quote{世界的首领来到，在我内一无所获。}\BibleRef{若 14:30} 那么，如果祂本没有被迫被钉十字架，却为我们被钉，我们岂不更应当高贵地忍受一切吗！\Quote{我有权柄}，祂说，\Quote{舍掉我的生命，也有权柄再取回它。}\BibleRef{若 10:18}\par

\Quote{祂因那摆在面前的喜乐}（他说）\Quote{忍受了十字架，轻看羞辱。} 但什么是\Quote{轻看羞辱}？祂选择了，他意指，那耻辱的死亡。因为假设祂死了。为何还要耻辱地死呢？没有别的原因，只是为了教导我们不看重人的光荣。因此，虽然祂没有义务，祂却选择了它，教导我们勇敢面对它，并把它看为无物。为什么他说的是\Quote{羞辱}而不是\Quote{痛苦}？因为祂忍受这些事时，不是带着痛苦。\par

\Quote{祂坐在天主圣座的右边。}你看到保禄在书信中所说的奖赏吗？\Quote{为此，天主极其高举祂，赐给祂一个名字，超越一切名字，使因耶稣基督之名，一切膝都要屈膝。}\BibleRef{斐 2:9} 祂是在谈论肉体。那么，即使没有奖赏，这个典范也足以说服我们接受一切。但现在奖赏也摆在我们面前，且不是普通的，而是伟大而不可言喻的。\par

5. 所以让我们无论何时遭遇这类事，在宗徒之前先考虑基督。为什么？祂的整个一生充满了侮辱。因为祂不断听到人们称祂为疯子、骗子、术士；有一次犹太人说：\Quote{不}，(经上说)\Quote{祂迷惑众人。}\BibleRef{若 7:12} 又说：\Quote{那个骗子活着的时候说过：三天后我要复活。}\BibleRef{玛 27:63} 他们也以巫术毁谤祂，说：\Quote{祂藉贝耳则步驱魔。}\BibleRef{玛 12:24} 以及\Quote{祂疯了，有魔鬼附身。}\BibleRef{若 10:20} \Quote{我们说得岂不对吗？}（经上说）\Quote{祂有魔鬼附身且疯了？}\BibleRef{若 8:48}\par

祂从他们那里听到这些，当祂行善、行奇迹、显示天主的作为时。因为如果祂什么也没做时就被这样说，那便不那么奇妙；但奇妙的是，当祂教导真理时，被称为\Quote{骗子}；当祂驱魔时，被说成\Quote{有魔鬼}；当祂推翻一切反对天主的事物时，被称为术士。因为他们不断地用这些事来控告祂。\par

若你愿意知道他们针对祂的嘲笑和讽刺（那尤其伤害我们灵魂的事），首先听祂亲属的：\Quote{这不是}（经上说）\Quote{木匠的儿子吗？祂的父母我们不是认识吗？祂的兄弟姐妹不都在我们这里吗？}\BibleRef{玛 13:55} 又因祂的家乡而嘲笑祂，说祂是\Quote{纳匝肋人}。又说：\Quote{查考一下，}经上说，\Quote{你便知道，从加里肋亚没有出过先知。}\BibleRef{若 7:52} 忍受了如此大的毁谤。又说：\Quote{经上难道不是说，基督来自白冷城吗？}\BibleRef{若 7:42}\par

你们也要看见他们所发出的讥讽的嘲弄吗？经上说，他们来到十字架跟前，竟朝拜了他；他们打他，掴他的脸，说：\BibleQuote{告诉我们，打你的是谁？}\BibleRef{玛 26:68}；他们又拿醋给他，说：\BibleQuote{你若是天主子，就从十字架上下来吧。}\BibleRef{玛 27:40}。大司祭的仆人也用手掌打他，他说：\BibleQuote{我若说得不对，你指证不对之处；若说得对，你为什么打我？}\BibleRef{若 18:23}。他们又戏弄地给他披上袍子，吐唾沫在他脸上，不断地试探他。\par

你们也要看见那些控告，有的暗中的，有的公开的，有的来自门徒的吗？他说：\BibleQuote{你们也要离去吗？}\BibleRef{若 6:67}。而那句\BibleQuote{你附了魔鬼}\BibleRef{若 8:48}，是出自已经信了的人之口。他不是不断地逃亡，有时在加里肋亚，有时在犹太吗？他的考验不是从襁褓时就很大吗？当他还是幼童时，他的母亲不是带他下到埃及去了吗？为此一切，他说：\BibleQuote{仰望耶稣，他是我们信德的创始者和完成者，他为了那摆在他面前的喜乐，忍受了十字架，轻视了羞辱，如今坐在天主宝座的右边。}\par

那么，让我们仰望他，也仰望他门徒的[磨难]，阅读保禄的[著作]，听他所说：\BibleQuote{以极大的坚忍，在患难、困苦、迫害、艰难、鞭打、监禁中}\BibleRef{格后 6:4}。又说：\BibleQuote{直到此时此刻，我们仍是又饥又渴，赤身露体，受人拳打，居无定所，亲手劳苦作工。被人辱骂，我们就祝福；被人迫害，我们就忍受；被人毁谤，我们就劝慰。}\BibleRef{格前 4:11}。我们中有谁受过这些事最小的一部分呢？因为他说：[我们] \BibleQuote{像是迷惑人的，像是受轻视的，像是一无所有的。}\BibleRef{格后 6:8}。又说：\BibleQuote{我被犹太人鞭打了五次，每次四十下减去一下；被棍打了三次，被石头砸了一次，在深海中度过了一日一夜；多次行路，在患难、困苦、饥渴中。}\BibleRef{格后 11:24}。这些事似乎中悦天主，请听他所说：\BibleQuote{为此我三次求主，他对我说：我的恩宠为你够了，因为我的能力在软弱中才全显出来。}\BibleRef{格后 12:8}。\BibleQuote{所以，}他说，\BibleQuote{我甘心情愿在软弱、患难、困苦、迫害、鞭打、监禁中，以使基督的能力常在我身上。}此外，请听基督自己所说：\BibleQuote{在世界上你们要受迫害。}\BibleRef{若 16:33}\par

% structure: p0035 source=h2 target=subsection reason=relative-heading-rank; base=h2

\subsection{希 12:3}

6. \BibleQuote{因为，}他说，\BibleQuote{你们要仔细想想那忍受罪人这样反对他的，免得你们灰心丧胆。}\par

他略过一切，用\Quote{反对}一词概括了全部，并加上了\Quote{如此}。因为他用\Quote{反对}指明了掌掴、嘲笑、侮辱、辱骂、戏弄等等。还不仅这些，也包括他一生中教导时所遭遇的事。\par

因为基督的苦难和宗徒们的苦难确是极大的安慰。他深知这是更好的德行之路，以致他亲自走了这条路，虽然并不需要；他深知磨难对我们有益，它反而成为安息的基础。请听他所说：\BibleQuote{谁若不背起自己的十字架跟随我，就不配是我的。}\BibleRef{玛 10:38} 他的意思是，如果你是门徒，就要效法师傅；因为这才是作门徒。但如果他走了患难之路，你却走了安逸之路，你就不再走他所走的路，而是另一条。那么，你怎么跟随呢？既然不跟随，又怎能作门徒，不随从师傅呢？保禄也说：\BibleQuote{我们软弱，你们却强壮；我们受轻视，你们却受尊敬。}\BibleRef{格前 4:10} 他的意思是，我们追求相反的事，你们却作门徒，我们作师傅，这合理吗？\par

7. 所以磨难是一件大事，亲爱的，因为它成就两件大事：它消除罪过，并使人坚强。\par

那么，你说，如果它倾覆并毁灭呢？磨难不会如此，而是我们自己的懈怠。怎么（你说）？如果我们清醒警醒，如果我们恳求天主不\BibleQuote{许我们受试探超过我们能承受的}\BibleRef{格前 10:13}，如果我们始终持守他，我们必能屹立不动，对抗我们的仇敌。只要我们以他为助佑，即使试探比一切风更猛烈地吹来，它们也如同糠秕和轻飘的叶子。请听保禄所说：\BibleQuote{在这一切事上，我们已大获全胜。}\BibleRef{罗 8:37} 又说：\BibleQuote{我确信，现时的苦难不能与将来显示给我们的光荣相比。}\BibleRef{罗 8:18} 又说：\BibleQuote{我们这短暂轻微的苦难，正在为我们成就永远无比重大的光荣。}\BibleRef{格后 4:17}\par

试想，他将何等大的危险、海难、接踵而来的苦难，以及诸如此类的事称为\Quote{轻省}；并要效法这位坚忍不拔的人，他随随便便、毫不在意地穿戴这身体。你贫穷吗？但不如保禄那般贫穷，他曾受饥饿、口渴和赤身的考验。他遭受这些不是一日，而是持续忍受。何处可见呢？请听他亲口说：\BibleQuote{直到此时此刻，我们还是又饥又渴，又赤身露体。}\BibleRef{格前 4:11} 哦！当他忍受如此巨大的苦难时，在宣讲中已拥有何等大的荣耀！此时（他写此信时）已进入第二十年。因为他说：\BibleQuote{我认得一个在主内十四年的人，或在身内，或在身外，我不知道。}\BibleRef{格后 12:2} 又说：\BibleQuote{过了三年，}（他说）\BibleQuote{我上耶路撒冷去。}\BibleRef{迦 1:18} 再听他说的：\BibleQuote{我宁愿死，也不让人使我这夸耀落空。}\BibleRef{格前 9:15} 不仅如此，他在书信中还写道：\BibleQuote{我们成了世界上的污秽。}\BibleRef{格前 4:13} 还有什么比饥饿更难忍受？比严寒更难忍受？比被那些他后来称为\BibleQuote{假弟兄}的弟兄们所设的阴谋更难忍受？\BibleRef{格后 11:26} 难道他不被称为世上的祸害？骗子？颠覆者？难道他没有被鞭打切割？\par

8. 亲爱的，让我们把这些事放在心上，深思默想，常常记念，这样，即使我们受冤屈、被掠夺、遭遇无数祸患，也永不会气馁。只要我们得以在天堂蒙悦纳，一切就都能忍受。只要我们在那里面前顺利，此世的事就微不足道。这些事不过是影子和梦；无论它们是什么，在本性或持续上都不算什么，而那些是所盼望、所期待的。\par

因为你们愿将什么与那些可怖的事相比呢？与不灭的火？与不死的虫？你们能指出此世的哪一件事可与\Quote{咬牙切齿}、\Quote{锁链}、\Quote{外面的黑暗}、\Quote{忿怒}、\Quote{患难}、\Quote{困苦}相比？至于时间呢？一万年与无尽无终的岁月相比，又算什么？不过像一滴水掉进无垠的大海罢了。\par

至于福乐呢？那里的超越性更大。经上说：\BibleQuote{眼所未见，耳所未闻，人心所未想到的}\BibleRef{格前 2:9}，这些也要持续无尽的岁月。为了这些，就算被鞭切割无数次，被杀，被烧，经历万死，忍受言语和行为上一切可怕之事，岂不应当吗？因为即使有可能让人在火中活着活下去，为了得到那些应许的福乐，难道不应当忍受一切吗？\par

9. 但我为何向这些连轻视财富都不肯、反而紧抓财富如同不死之人说这些闲话呢？他们若从大量中拿出一点点，便以为做了一切？这不算是施舍。因为施舍是那位寡妇倾囊而出的施舍，她将\BibleQuote{所有的生活费}都投上了。\BibleRef{谷 12:44} 但倘若你未能做到像那寡妇那样捐献，至少也该捐献你所有的盈余：保留所需的，而非多余的。\par

然而，没有人捐献他的盈余。因为只要你拥有许多仆人和丝绸衣服，这些东西都是盈余。没有什么是必不可少或必需的，没有它们我们也能生活；这些东西是多余的，纯粹是额外添上的。那么，如果你愿意，我们来看看什么是我们不可或缺的。如果我们只有两个仆人，我们就能生活。既然有些人没有仆人也能生活，我们若不满足于两个仆人，还有什么借口呢？我们也可以有一栋砖造的三间房；这已足够我们使用。因为难道不是有些有子女有妻子的人只有一间房吗？如果你愿意，还可以有两个僮仆。\par

10. 也许你会说：一位贵妇人只带两个仆人外出，岂不是可耻？不，一位贵妇带两个仆人外出并不可耻，但带许多仆人外出才是可耻的。也许你听了会笑。相信我，这确实\Emph{是}可耻的。你以为带着许多仆人外出像贩羊的或贩奴隶的一样很了不起吗？这是骄傲和虚荣，而另一种做法则是哲学和体面。因为贵妇不应该从侍从的数目上被认出来。拥有许多奴隶算什么德行？这不属于灵魂，凡不属于灵魂的就不显示高贵。当她满足于少量事物时，她才是真正的贵妇；但当需要许多事物时，她就成了奴隶，甚至不如奴隶。告诉我，天使不也独自往来于全世界，不需要任何人跟随吗？他们因此就比我们低劣吗？他们不需要随从，而我们却需要。如果不需任何随从是天使般的，那么谁会更加接近天使的生活呢？是需要许多随从的人，还是需要少的人？这不是可耻吗？因为行不当之事就是可耻。\par

请告诉我，在公共场所，是谁吸引人们的目光？是那个带着许多随从的人，还是那个带着少数随从的人？那个独自一人的人，不是比那个带着少数随从的人更不引人注目吗？你看到这种（前者）行为是一种耻辱吗？在公共场所，是谁吸引人们的目光？是那个穿着华丽衣服的人，还是那个穿着朴素、不加修饰的人？再问，在公共场所，是谁吸引人们的目光？是那个骑着骡子、配有金饰鞍辔的人，还是那个简朴地、合宜地步行的人？我们甚至不看后者，即使我们看到了她；但众人不仅挤着去看前者，还问：她是谁？从哪里来？我不说这会引起多大的嫉妒。那么（告诉我），被人注视是可耻的，还是不被注视是可耻的？什么时候羞耻更大——当所有人都盯着她看时，还是当没有人看她时？当他们打听她的事时，还是当他们甚至不关心时？你看到我们做一切事，不是为了谦逊，而是为了虚荣吗？\par

然而，既然不可能把你们从那种（虚荣）中拉出来，我现在满足于让你们知道：这种（行为）并不是耻辱。惟独罪是耻辱，却没有人认为它是耻辱。惟独罪是耻辱，却没有人认为它是耻辱，反而认为其他一切才是耻辱。\par

第十一节。让你的衣服是必需的，而不是多余的。然而，为了不让你们过于受拘束，我向你们保证：我们不需要金饰或花边。这不是我说的。因为这些话不是我的，请听真福的保禄所说，他郑重地嘱咐妇女，\BibleQuote{要装饰自己，不要以编发、黄金、珍珠，或昂贵的衣服。}\BibleRef{弟前 2:9} 但是，保禄啊，你愿意告诉我们是哪一种（衣服）吗？因为她们或许会说，只有金饰才是昂贵的，丝绸并不昂贵。请告诉我们，你愿意哪一种。\BibleQuote{只要有衣有食，我们就当知足。}\BibleRef{弟前 6:8} 让我们的衣服仅仅足以遮盖我们。因为天主赐给我们衣服，正是为了这个原因：使我们遮盖赤身，而这任何衣服都能做到，即使花费很小。或许你们会笑，你们这些穿丝绸衣服的人；因为真的，人们可以好好地笑一笑，想想保禄吩咐的是什么，而我们实践的是什么！\par

但我的讲话不仅是针对妇女，也是针对男人。因为我们所有其他的东西都是多余的；只有穷人没有多余之物；或许他们也是出于不得已：因为如果他们有力量，他们也不会节制。然而，\BibleQuote{无论是假意，或是真心}\BibleRef{斐 1:18}，就目前而言，他们没有多余之物。\par

第十二节。那么，让我们穿那些足够我们需要的衣服。因为黄金有什么意义呢？这些东西适合于舞台上的人，这种服饰属于她们，属于娼妓，属于那些做一切事以吸引目光的人。让那些在舞台或舞池上的人美化自己吧。因为她想吸引所有人到她那里。但是，一个自称虔敬的妇女，不要这样美化自己，而要换一种方式。你有一种远比那更好的美化自己的方法。你也有一个舞台：为了那个舞台，你要使自己美丽：穿上那些装饰。你的舞台是什么？是天堂，是天主的众天使的会集。我不仅对童贞女说，也对世俗中的妇女说。所有相信基督的人都有那个舞台。让我们说那些话，好使那些观众喜悦。穿上那样的衣服，好使你们取悦他们。\par

因为请告诉我，如果一名娼妓放下她的金饰、长袍、笑声、机智和不洁的言谈，穿上廉价衣服，简朴地打扮自己（走上舞台），说出虔诚的话语，谈论贞洁，不说任何失礼的话，难道众人不会站起来吗？这个舞台不会散开吗？他们不会把她赶出去，如同一个不知道如何迎合众人、说出那个撒殚舞台所陌生的话语的人吗？同样，如果你穿着她的衣服进入天堂的舞台，观众也会把你赶出去。因为在那里，不需要这些金饰，而是需要不同的。是怎样的？正如先知所说的：\BibleQuote{她身穿金绣衣饰}\BibleRef{咏 44:13}，但不是为了使身体白皙发光，而是为了美化灵魂。因为在那舞台上竞技和角力的是灵魂。\BibleQuote{王女的一切光荣是在内里}\BibleRef{咏 44:13}，它说。你用这些来装扮自己；因为这样，你既使自己摆脱了其他无数的邪恶，也使你的丈夫摆脱了焦虑，使你自己摆脱了操心。\par

因为这样，当你不需许多东西时，你的丈夫就会尊重你。每一个男人通常对那些向他请求的人感到羞怯；但是当他看到他们不需要他时，他就放下骄傲，以平等的方式与他们交谈。当你的丈夫看到你在任何事情上都不需要他，你看不起从他那里来的礼物时，那么，即使他很傲慢，他也会比你穿金戴银时更尊重你；你就不再是他的奴隶。因为那些我们需要的人，我们不得已要屈从。但是，如果我们节制自己，我们就不再被视为罪犯，而他知道我们服从他是出于敬畏天主，而不是为了从他那里得到的东西。因为现在，当他给我们许多恩惠时，无论他得到多少尊敬，他都认为他还没有得到全部（应得的）；但那时，虽然他得到很少，他也会认为这是一种恩惠，他不会责备，他也不会因为你的缘故被迫欺诈。\par

13. 还有什么比预备金饰，在浴场和广场上佩戴，更不合理呢？然而，在浴场和广场上也许不足为奇，但一个女人竟这样盛装进入教会，实在可笑。因为，她究竟为了什么理由佩戴金饰来到这里呢？她本该来聆听\Quote{她们不以金饰、珍珠或华服装饰自己}的训诫。\BibleRef{弟前 2:9} 那么，女人啊，你来的目的是什么？是要与保禄争辩，表明即使他把这些话重复一万次，你也不在乎？还是想令我们这些教师因白费口舌而羞愧？请告诉我：若有一位外教不信徒，听了这段经文，知道蒙福的保禄说了这些，而他有一位信主的妻子，却见她十分注重打扮，戴上金饰来教会听保禄告诫妇女\Quote{不要以金饰}\BibleRef{弟前 2:9}、\Quote{珍珠}或\Quote{华服}装饰自己，当他在小房间里看见她穿戴这些、精心安排时，难道不会自言自语说：\Quote{我妻子为什么还待在屋里？她为何如此缓慢？为何戴上金饰？她要去哪里？进教会？去听道？\InnerQuote{不要华服}？}难道他不会讪笑，不会轰然大笑？难道他不会认为我们的信仰是儿戏和欺骗？因此，我恳求你们，让我们将金饰留给游行、剧场和店铺的招牌。不要让天主的肖像用这些东西来打扮：让淑女以文雅为装饰，而文雅就是没有骄傲和炫耀。\par

即使你想从人那里获得荣耀，你也会以此方式得到。因为我们不会那么惊讶于一个富人的妻子身穿金饰和丝绸（因为这是她们共同的做法），而是当她穿一件仅以羊毛制成的朴素简单的衣服时，人人都会钦佩，人人都会称赞。因为在那金饰和华服的装饰中，她有许多人与她共享。即使她超过一人，也会被另一人所超越。是的，即使她超过所有人，她也必须向皇后本人认输。但在另一种情况下，她甚至超过了皇后本人。因为她在财富中独自选择了穷人的衣服。所以，即使我们追求荣耀，这里的荣耀也更大。\par

14. 我这样说不仅是对寡妇和富人；因为这里寡居的需要似乎导致了这一点；也是对那些有丈夫的人说的。\par

但你会说：\Quote{如果我衣着朴素，我的丈夫就不喜悦我。}你并非想取悦你的丈夫，而是想取悦穷人的众多妇女；或者更确切地说，不是取悦她们，而是让她们痛苦嫉妒，使她们的贫困更加深重。因为你而发出了多少亵渎之语！\Quote{愿没有贫穷}他们会说，\Quote{天主憎恨穷人；天主不爱那些贫穷的人。}因为，你并非想取悦你的丈夫，并为此打扮自己，这一点你通过自己的行为向所有人表明了。因为你一跨过卧室的门槛，就立刻脱下所有袍子、金饰和珍珠；而在家里，在所有地方，你却不穿戴它们。\par

但如果你真的想取悦你的丈夫，有些方式可以取悦他：温柔、和善、端庄。因为请相信我，女人啊，即使你的丈夫极其堕落，这些——温柔、端庄、不自傲、不奢华、不浪费——才是更能赢得他的东西。因为即使你设计一万种这类东西，你也不会约束那放荡的人。这一点，拥有这类丈夫的人都知道。因为无论你如何美化自己，他若是放荡之人，仍会去找妓女；而那贞洁规矩的丈夫，你不会通过这些手段赢得他，而是通过相反的方式：是的，通过这些你甚至使他痛苦，使自己披上爱世界的名声。因为如果你的丈夫出于尊重，且因他是个清醒的人而不说话，但内心他会谴责你，不会掩饰敌意和嫉妒。难道你不会因引起人对你的敌意而驱散未来的一切喜悦吗？\par

15. 也许你为听到的话而烦恼，并愤慨地说：\Quote{他更加激怒丈夫们反对他们的妻子。}我这样说，并非要激怒你们的丈夫，而是希望你们自愿这样做，为了你们自己，而不是为了他们；不是为了使你们免于他们的嫉妒，而是为了使你们免于此生的炫耀。\par

你想要显得美丽吗？我也希望如此，但要拥有天主所寻求的美丽，即\Quote{君王所贪慕的}。\BibleRef{咏 44:11} 你愿意谁作你的爱人？天主还是人？如果你拥有那种美丽，天主会\Quote{贪慕你的美貌}；但如果你拥有另一种美丽而没有这种美丽，祂会憎恶你，而你的爱人将是放荡之人。因为没有一个爱有夫之妇的人是好人。即使就外在的装饰而言，也要考虑这一点。因为另一种装饰，即灵魂的装饰，吸引天主；但这一种吸引放荡之人。你看到我关心你，我为你焦虑，希望你美丽，真正美丽，灿烂，真正灿烂，使得你不再是放荡之人的爱人，而是万有的主天主的爱人？而那位以祂为爱人的人，将像谁呢？她将在天使的歌团中占有一席之地。因为如果一位被君王所爱的人被视为超乎所有人的幸福，那么被天主以深厚的爱所爱的人，其尊荣又会如何呢？即使你把整个世界放在天平上，也抵不过那种美丽。\par

那么，让我们培养这种美；用这些装饰来装扮我们自己，好使我们能进入天堂，进入属灵的厅室，进入那无玷的洞房。因为这美容易被任何东西摧毁；即使它保持得好，且疾病和焦虑都不能损害它（这是不可能的），它也不能持续二十年。但那另一种美却永远绽放，永远茂盛。\Emph{那里}，没有变化可惧；没有衰老带来皱纹，没有潜在的疾病使之枯萎；没有沮丧的焦虑毁损它；它远远超越这一切。但这（地上的美）尚未出现就飞逝了，即便出现，也没有多少仰慕者。因为心灵有序的人不仰慕它；而那些仰慕它的人，是以纵欲来仰慕。\par

16. 因此，我们不要培养这种（美），而要培养那另一种：让我们拥有那一种，好使我们能拿着明亮的火炬进入洞房。因为这不只是向童贞女应许的，也是向贞洁的灵魂应许的。因为如果这仅仅属于童贞女，那五个（童贞女）就不会被关在门外了。所以这属于所有灵魂贞洁的人，即那些摆脱世俗幻想的人：因为这些幻想腐蚀我们的灵魂。因此，如果我们保持不被玷污，我们就会离开那里，并被接纳。\Quote{我曾把你们许配给一个丈夫，}他说，\Quote{要把你们如同贞洁的童贞女献给基督。}\BibleRef{格后 11:2}他说这些话，不是指童贞女，而是指整个教会的全体。因为不被玷污的灵魂就是贞洁的，即使她有丈夫：她在那真正的贞洁方面是贞洁的，值得钦佩。因为身体的贞洁只是那另一种贞洁的伴随和影子，而那另一种才是真正的贞洁。让我们培养这个，这样我们就能以喜乐的面容看见新郎，拿着明亮的火炬进入，如果灯油不缺，如果我们熔化我们的金饰，获得使我们的灯明亮的油。这油就是仁爱。\par

如果我们把所拥有的分施给别人，如果由此制作灯油，那么它就会保护我们，那时我们不会说，\Quote{请给我们油，因为我们的灯快要灭了}\BibleRef{玛 25:8}，也不会向别人乞讨，也不会当我们去卖油的人那里时被关在门外，也不会在我们敲门时听到那可怕而恐怖的声音，\Quote{我不认识你们。}\BibleRef{玛 25:12}但祂会承认我们，我们将与新郎一同进去，进入属灵的新房，享受数不尽的福乐。\par

因为如果这里的洞房如此明亮，房间如此辉煌，没有人会厌倦观看它们，那么那里更是如此。天堂是厅室，而洞房比天堂更美好；那时我们将进入。但如果洞房如此美丽，新郎又会是怎样呢？\par

而为什么我说，\Quote{让我们放下我们的金饰，施舍给穷人}？因为如果你们甚至应该出卖自己，如果你们应该成为奴隶而非自由人，以便能与那位新郎在一起，享受那美，甚至仅仅是为了看那面容，你们难道不该以乐意的心欢迎一切吗？我们在地上观看并仰慕一位君王，但当我们看到一位君王和一位新郎同时出现，我们更该以乐意之心欢迎祂。的确，这些是影子，而那些是实体。在天上有了一位君王和一位新郎！也被认为值得拿着火炬走在祂前面，靠近祂，永远与祂同在，我们有什么不该做的？我们有什么不该完成的？我们有什么不该忍受的？我恳求你们，让我们对那些福乐产生一些渴望，让我们渴慕那位新郎，让我们在真正的贞洁上成为贞洁的。因为主寻求灵魂的贞洁。怀着这个，让我们进入天堂，\Quote{没有斑点、皱纹或任何这类东西}\BibleRef{弗 5:27}；好使我们也能获得所应许的美好事物，愿我们众人借着我们的主耶稣基督的恩宠和仁慈都分沾，愿光荣、权能、尊崇归于祂，与圣父及圣神，从现在直到永远，永无穷尽。阿们。\par

\SourceRecord{Translated by Frederic Gardiner.；Nicene and Post-Nicene Fathers, First Series；Vol. 14.；Edited by Philip Schaff.；Buffalo, NY: Christian Literature Publishing Co.,；1889.；<http://www.newadvent.org/fathers/240228.htm>.}

% structure: p0001 source=h1 target=section reason=page-title

\section{\WorkTitle{希伯来人书}第二十九篇讲道}

% structure: p0002 source=h2 target=subsection reason=relative-heading-rank; base=h2

\subsection{希伯来人书 12:4-6}

\BibleQuote{你们与罪恶争斗，还没有抵抗到流血的地步。你们竟全然忘了那劝勉你们如同劝勉儿子的話，说：\InnerQuote{我儿，不可轻看主的惩戒，被他责备的时候，也不要灰心；因为主所爱的，他必惩戒；他鞭打他所收纳的每一个儿子。}}\par

1. 有两种表面相反、却彼此大大增益的安慰，他都在此提出。一种是说某人已受了许多苦；因为灵魂有自己受苦的许多见证时，便得振奋；他在前面也引用了这一点，说：\BibleQuote{你们要回想往日，蒙光照以后，忍受了大爭战的各种苦难。}\BibleRef{希 10:32} 另一种是说：\Quote{你们没有受什么大事。}前者在灵魂疲惫时使之振作、恢复气息；后者在灵魂怠惰松懈时使之回转并压下骄傲。因此，为使那见证（对他们的受苦）不致在他们心中生出骄傲，请看他的做法。他说：\BibleQuote{你们还没有抵抗到流血的地步。}他没有立刻接着说下去，而是在向他们显示所有那些\Quote{抵抗到流血}的人，并引出基督的荣耀、他的苦难之后，才轻易地继续他的论述。他在写给格林多人的信中也这样说：\BibleQuote{你们所遇见的试探，无非是人所能受的}\BibleRef{格前 10:13}，即小的。因为当灵魂思想到自己还没有达到完全的（考验），并从已遭遇的事鼓励自己时，这就足以激发和纠正灵魂。\par

他所说的是：你们还没有遭受死亡；你们的损失涉及钱财、名誉、被赶逐离开地方。然而基督为你们流了他的血，而你们还没有为自己流血。他为了真理争战，甚至至于死，为你们而战；而你们还没有进入危及死亡的险境。\par

\BibleQuote{你们竟全然忘了那劝勉。}即你们放松了双手，你们变得灰心。他说：\BibleQuote{你们与罪恶争斗，还没有抵抗到流血的地步。}这里他指明罪恶既是极其强壮的，也是武装起来的。因为\Quote{你们抵抗}（坚定地站立）一词用于那些坚定站立的人。\par

2. 他说：\Quote{这劝勉}是\BibleQuote{对你们如同对儿子说的话：\InnerQuote{我儿，不可轻看主的惩戒，被他责备的时候，也不要灰心；因为主所爱的，他必惩戒；他鞭打他所收纳的每一个儿子。}}他从事实本身得了鼓励；此外还把从论证、从这经文中得来的鼓励加了上去。\par

\BibleQuote{被他责备的时候，也不要灰心。}由此可知，这些事是出于天主的。因为知道这些事之所以有力是天主的作为，他允许（它们），这也是不小的安慰，正如保禄所说：\BibleQuote{他对我说：\InnerQuote{我的恩寵够你用的，因为我的能力是在软弱上显得完全。}}\BibleRef{格后 12:9} 是他允许（它们）。\par

\BibleQuote{因为主所爱的，他必惩戒；他鞭打他所收纳的每一个儿子。}你不能说任何一个义人是没有患难的：即使他表面上似乎没有，但我们\Emph{我们}不知道他其他的患难。所以每个义人必须经过患难。因为基督声明：\BibleQuote{那宽而大的路通向丧亡，但窄而小的路通向生命。}\BibleRef{玛 7:13} 如果藉着那条路可以进入生命，并且没有别的路，那么所有进入生命的都藉着窄路（进入），许多已进入生命的人也是如此。\par

% structure: p0010 source=h2 target=subsection reason=relative-heading-rank; base=h2

\subsection{希伯来人书 12:7}

他说：\BibleQuote{你们为受惩戒而忍受}，不是为了刑罚，不是为了报复，也不是为了受苦。看，从他们以为自己被（天主）离弃的那一点，他却说他们可以确信自己没有被离弃。他的意思好比说：因为你们遭受了这么多祸患，你们以为天主离开了你们、恨你们吗？如果你们不曾受苦，那么这样认为是对的。因为如果\BibleQuote{他鞭打他所收纳的每一个儿子}，那么不受鞭打的，也许就不是儿子。那么，你会说，恶人不受痛苦吗？他们确实受苦；那么如何？他没有说：每一个受鞭打的都是儿子，而是说：每一个儿子都受鞭打。因为在所有情况下，他鞭打他的儿子；因此需要证明的是，是否有儿子不受鞭打。但你不能说出一个：有许多恶人也受鞭打，如杀人犯、强盗、巫术师、盗墓者。但这些人是为他们的恶行付上代价，不是作为儿子受鞭打，而是作为恶人受惩罚；但你们是作为儿子。\par

% structure: p0012 source=h2 target=subsection reason=relative-heading-rank; base=h2

\subsection{希伯来人书 12:8}

3. 然后他又从一般习俗（论证）。你看他是如何从各方面提出论证的：从圣经的事实、从圣经的话语、从我们自己的观念、从日常生活的榜样。\BibleRef{希 12:8} \BibleQuote{如果你们不受惩戒……}。你看见我说了我刚才所说的吗：不受惩戒就不能成为儿子？因为在家庭中，父亲不关心私生子，即使他们什么也不学，即使他们不杰出；但他们为自己的合法儿子担心，怕他们懒惰。那么，如果不受惩戒是私生子的（标记），我们应当为惩戒而喜乐，如果这是合法性的（标记）。天主待你们如同待儿子正是为此。\par

% structure: p0014 source=h2 target=subsection reason=relative-heading-rank; base=h2

\subsection{希伯来人书 12:9}

\BibleQuote{此外，我们曾有生身的父亲管教我们，我们尚且敬重他们。}再次，他从他们自己的经验、从他们自己所受的（论证）。正如上面他说：\BibleQuote{你们要回想往日}\BibleRef{希 10:32}，同样这里他说：\BibleQuote{天主待你们如同待儿子}，你们不能说：我们受不了；是的，\BibleQuote{如同待儿子}是慈爱地。因为如果他们敬重\BibleQuote{生身的父亲}，你们岂不更要敬重你们天上的父亲吗？\par

然而，这差异不仅来自于此，也不仅来自人，更来自原因本身和事实。因为祂和他们施加惩罚的根据并不相同：他们是为了\Quote{他们以为好的事}，即常常是满足自己的喜好，并非总是着眼于利益。但在此，不能这样说。因为祂这样做并非为自己利益，而是为了你，且只为你的益处。他们这样做是为了你也能对他们有用，常常毫无理由；但在此，没有这类事。你看到这也带来安慰吗？因为我们最深地依恋那些现世父母，当我们看到他们命令或劝诫我们并非为自己利益，而是完全且唯独为我们着想。这才是真爱，真实的爱，当我们被爱却对爱我们的人毫无用处时——不是为收取，而是为施与。祂责罚，祂做一切，祂尽心竭力，为使我们能堪当承受祂的恩惠。\par

% structure: p0017 source=h2 target=subsection reason=relative-heading-rank; base=h2

\subsection{希伯来书 12:10}

\BibleQuote{他们固然}（他说）\BibleQuote{在数日之间，随己意管教我们，但祂管教我们是为我们的益处，使我们能分享祂的圣洁。}\BibleRef{希 12:10}\par

什么是\Quote{祂的圣洁}？就是祂的纯洁，使我们按自己的力量堪当与祂相配。祂切愿你领受，祂做一切为能赐予你：你们岂不应切切努力以领受吗？\BibleQuote{我向主说}（某人说）\BibleQuote{祢是我的主，因为祢不需要我的美物。}\BibleRef{咏 15:2}\par

\BibleQuote{此外，}他说，\BibleQuote{我们曾有过肉身的父亲管教我们，我们尚且敬重他们；难道我们不该更顺服众灵之父，从而得生吗？}（众灵之父，无论是指神恩的父，还是祈祷的父，或无形体能力的父。）如果我们这样而死，那么\Quote{我们将活着。}\BibleQuote{因为他们确实在数日之间随己意管教我们}，因为看似[如此]的并非总是有益，但\BibleQuote{祂是为我们的益处。}\par

4. 因此，管教是\Quote{有益的}；因此，管教是\Quote{分享圣洁}。诚然，这益处甚大：因为当它驱除懒惰、邪欲和对今生的爱慕，当它助佑灵魂，当它使人轻视一切在此世的事物（因为苦难确实如此），它岂不圣洁？它岂不吸引圣神的恩宠？\par

让我们思想义人，他们因何缘故全部闪耀发光？岂非来自苦难？并且，如果你愿意，让我们从最初开始历数他们：\Person{亚伯尔}、\Person{诺厄}本人；因为他不可能是那众多恶人中唯一的，却未受过苦难；因为经上说：\BibleQuote{诺厄独自在他的世代中是完美的人，中悦了天主。}\BibleRef{创 6:9} 试想，我恳求你，如果现在，当我们有无数人可以效法其美德，有父亲、儿女和教师，我们尚且如此困苦，那么我们可以设想他独自在这么多人中承受了什么？但我何必谈论那奇异奇妙之雨的境况？或者我该谈谈\Person{亚巴郎}，他不断的流浪、妻子被掳、危险、战争、饥荒？我该谈谈\Person{依撒格}，他经历了何等可怕的事，被逐出各地，徒劳劳作，为他人劳累？或者\Person{雅各伯}？因为确实没有必要列举他所有的苦难，但可以提出他自己与法郎谈话时的见证：\BibleQuote{我寄居的日子又少又苦，不及我祖先寄居的日子。}\BibleRef{创 47:9} 或者我该谈谈\Person{若瑟}本人？或者\Person{梅瑟}？或者\Person{若苏厄}？或者\Person{达味}？或者\Person{厄里亚}？或者\Person{撒慕尔}？或者你愿意我谈论所有先知吗？你岂不会发现所有这些人都因苦难而光彩夺目？那么告诉我，\Emph{你}渴望因安逸和奢侈而光彩夺目吗？但你不能。\par

或者我该谈谈宗徒们？不，他们超越了所有人。基督曾这样说：\BibleQuote{在世界上你们将有苦难。}\BibleRef{若 16:33} 又说：\BibleQuote{你们将要痛哭哀号，世界却要喜乐。}\BibleRef{若 16:20} 还有，\BibleQuote{通往生命的门是窄而狭的。}\BibleRef{玛 7:14} 那道路的主说，那是\BibleQuote{窄而狭的}；而你却寻求\Quote{宽广的}路？这岂非不合理？结果，你走另一条路，不会到达生命，而是毁灭，因为你选择了通向那里的路。\par

你要我把那些奢华生活的人带到你面前吗？让我们从末后到最初追溯。那在火窑中焚烧的富人；为肚腹而活的犹太人，\BibleQuote{他们的神是自己的肚腹}\BibleRef{斐 3:19}，他们在旷野中始终寻求安逸，最终被毁灭；同样，\Place{索多玛}人因贪食而毁灭；\Person{诺厄}时代的人，岂不也是因为他们选择了这种安逸放荡的生活？因为\BibleQuote{他们奢靡，}经上说，\BibleQuote{在饱食无度中。}\BibleRef{厄 16:49} 这是指索多玛人。但如果\Quote{饱食无度}造成了如此大的恶，我们该如何谈论其他美味？\Person{厄撒乌}，他岂不在安逸中？而那些\BibleQuote{天主的儿子们}\BibleRef{创 6:2}，看到女子，便坠入深渊？那些被无节制的欲望逼疯的人又如何？所有异邦的君王，\Place{巴比伦}人、\Place{埃及}人，他们岂非悲惨地灭亡？他们岂不在痛苦中？\par

5. 至于现在的事，告诉我，岂不也是如此？请听基督说：\BibleQuote{穿软衣服的人是在王宫里}\BibleRef{玛 11:8}，但不穿这类衣服的人在天上。因为软衣服甚至使严肃的灵魂松懈，摧毁它，使之乏力：是的，即使遇到粗糙坚硬的身体，经过这般轻柔处理，也很快变得柔软软弱。\par

因为，请告诉我，你认为女人为何如此软弱？仅仅因为她们的性别吗？绝不是：而是由于她们的生活方式及其教养。因为她们避免日晒、不活动、沐浴、涂油、大量香水、床榻的柔软细嫩，最终使她们变得如此。\par

为了使你明白，请听我说。告诉我：从园中取一棵树，它生长在未开垦之处，饱受风吹，把它种在湿润阴暗的地方，你会发现它远不如当初取来的那棵。事实如此，［从以下事实可见］在乡间长大的女人比城镇里的更强壮：她们在摔跤中能胜过许多城里女人。因为当身体变得愈发娇弱，灵魂也必然受到损害，因为其活动大多随［身体］而变化。因为生病时，因软弱我们变了个人；康复时，又变了回来。正如琴弦，当音调微弱、松弛、调音不当，艺术之美也就破坏了，不得不屈从于琴弦的恶劣状态；同样，身体也是如此，灵魂从中受到许多伤害，承受许多痛苦。因为当身体需要大量照料时，灵魂便忍受着痛苦的奴役。\par

6. 因此，我恳求你们，让我们通过劳动使身体强壮，不要把它当作病人来照料。我的话不仅是对男人，也是对女人。因为，女人啊，你为何不断以奢华削弱［身体］、耗尽它？你为何用脂肪毁坏体力？这脂肪是松弛，而非力量。然而，如果你戒除这些事，以不同方式管理自己，那么，当力量和良好的身体习惯存在时，你的容貌也会如你所愿地改善。如果你用千万种疾病困扰它，那么既不会有红润面色，也不会有健康；因为你将总是情绪低落。你知道，正如天空晴朗时使美丽的房屋显得辉煌，同样，心灵的快乐加在美貌上，使之更佳；但如果［一个女子］情绪低落、痛苦，她就变得更难看。疾病和痛苦产生低落情绪；而疾病源于身体因极度奢华而过于娇弱。所以，即使为此，你若听从我的忠告，也将逃避奢华。\par

\Quote{但是，你会说，奢华带来快乐。} 是的，但这快乐不如烦恼之大。而且，快乐仅限于上颚和舌头。因为当餐桌撤去，食物咽下，你就像从未吃过一样，或者更糟，因为你从此忍受压迫、胀气、头痛、如同死亡的睡眠，并且常常因饱食而失眠、呼吸阻塞、打嗝。你会痛苦地咒骂自己的肚子，而你本应咒骂无节制的饮食。\par

7. 那么，我们不要养肥身体，而要听从保禄的话：\BibleQuote{不要为肉体安排，去满足私欲，}\BibleRef{罗 13:14} 就如一个人取食物倒入阴沟，那将食物倒入肚腹的人也是如此：或者不如说，并非如此，而是更糟。因为在前者，他使用阴沟对自己无害；但在后者，他产生无数疾病。因为滋养身体的是适量且能消化的食物；但超过我们所需的，不但不滋养，反而败坏其他。但无人看见这些事，由于某种偏见和不合时宜的快乐。\par

你希望滋养身体吗？除去多余之物；给予足够且能消化之物。不要装载过多，免得压垮它。适量既是营养，也是快乐。因为没有什么比消化良好的食物更能带来快乐；没有什么比它更能带来健康；没有什么比它更能带来官能的敏锐，没有什么比它更能远离疾病。因为适量既是营养，又是快乐，又是健康；但过度是伤害、不快和疾病。因为饥荒所做的，饱食同样做到；或者更严重的祸害。因为前者在几天内使人离世、得释放；但后者侵蚀、腐烂身体，使身体陷入长期疾病，然后最痛苦的死亡。但\Emph{我们}，尽管认为饥荒是极可怕的事，却追求饱食，而那比饥荒更令人痛苦。\par

这疾病从何而来？这疯狂从何而来？我并不是说我们应该消瘦自己，而是说我们应该吃多少食物才能给我们带来快乐，那是真正的快乐，并能滋养身体，使身体井然有序，适合灵魂的活力，紧密结合在一起。但是，当身体被奢华浸透时，在洪流中就无法如人们所说的，保持住连接骨架的螺栓和关节。因为洪水涌入，整体便分崩离析。\par

\BibleQuote{不要为肉体安排}（他说）\BibleQuote{去满足私欲。}\BibleRef{罗 13:14} 他说得好。因为奢华是不合理私欲的燃料；即使奢华的人本是最有哲学的，他也必然受到酒和饮食的影响，必然放松，必然忍受更大的火焰。因此［产生］奸淫，因此产生通奸。因为饥饿的肚子不能产生情欲，或者说，只用适量食物的肚子不能。但产生不雅情欲的是因奢华而松弛的肚子。正如非常潮湿的土地和充满水分、保留大量湿气的粪堆会生虫，而摆脱了这种湿气、没有过度之处的土地则结出丰硕的果实：即使不耕种，也长出草，若耕种，则结出果实：［我们也是如此。］\par

我们不要让我们的肉体变得无用、无益或有害，而要在其中栽种有用的果实和结果子的树木；不要用奢侈来软化它们，因为当它们腐烂时，也会生出虫子来代替果实。同样，被植入的欲望，如果你过分滋润它，就会产生无理性的快乐，甚至是最为无理性的。因此，让我们除去这有害的邪恶，好能获得那应许给我们的美善，就在我们的主耶稣基督内，愿光荣偕同祂，并与圣神一同，归于圣父，自今世以至无穷世。阿们。\par

\SourceRecord{Translated by Frederic Gardiner.；Nicene and Post-Nicene Fathers, First Series；Vol. 14.；Edited by Philip Schaff.；Buffalo, NY: Christian Literature Publishing Co.,；1889.；<http://www.newadvent.org/fathers/240229.htm>.}

% structure: p0001 source=h1 target=section reason=page-title

\section{希伯来书第三十篇讲道}

% structure: p0002 source=h2 target=subsection reason=relative-heading-rank; base=h2

\subsection{希伯来书 12:11-13}

\BibleQuote{\Emph{一切的管教，当时看来不是喜乐，而是忧愁}，\Emph{后来却为那些由此受过操练的人，结出平安的义果}。\Emph{所以，你们要把下垂的手、发软的膝挺起来；也要为你们的脚修直道路，使那瘸腿的不至脱臼，反而得到痊愈。}}\par

1. 服用苦药的人，先是感到一些不适，随后才体会到益处。德性如此，恶行也如此。后者先有快乐，后有沮丧；前者先有沮丧，后有快乐。但二者并不相等：先苦后甜与先甜后苦，岂是一回事？为何呢？因为在后一种情形，对将来沮丧的预期使当下的快乐减少；而在前一种情形，对将来快乐的预期削减了当下沮丧的强度；结果是，在一种情形中我们从未有过快乐，在另一种情形中我们从未有过忧伤。差异不仅在此，还在于其他方面。例如，二者的持续时间并不相等，而是后者远为长久、更为丰富。在此处，在属灵之事上更是如此。\par

由此，保禄便着手安慰他们；并且再次采纳众人的共同判断——当一个人说出众所公认的事理时，无人能够反对，也无人能与这共同的决定相争。\par

他说：你们正在受苦。因为管教就是这样，这是它的开端。因为\BibleQuote{一切的管教，当时看来不是喜乐，而是忧愁}。他说得好：\Quote{看来不是}。他的意思是，管教不是忧愁的，而是\Quote{看来}如此。\Quote{一切的管教}：并非这种或那种，而是\Quote{一切}，包括人间的和属灵的。你看见他是从我们共同的观念来论证的吗？\Quote{看来}（他说）\Quote{是忧愁的}，所以它并非[真正]如此。因为什么样的忧愁能生出喜乐？同样，快乐也不会生出沮丧。\par

\BibleQuote{后来却为那些由此受过操练的人，结出平安的义果。}不是\Quote{果}，而是\Quote{果子}，表示极大的丰盛。\par

他说：\BibleQuote{为那些由此受过操练的人}。\Quote{由此受过操练的人}是什么意思？就是那些长久忍受、恒心忍耐的人。他使用了一个吉利的说法。那么，管教就是操练，使运动员强壮，在战斗中不可战胜，在战争中势不可挡。\par

如果\BibleQuote{一切的管教}都是如此，那么这一次的管教也是如此；所以我们应当期待美好的事物，期待甜蜜而平安的结局。不要因为管教本身艰难，就诧异它有甜美的果实；因为树木的树皮几乎毫无滋味且粗糙，但果实却是甜美的。他是从共同的观念引出这个论证的。因此，如果我们应当期待这样的事物，你们为何还烦恼呢？你们已经忍受了痛苦，为何对将来的好处灰心呢？你们已经忍受了令人不快之事，就不要对赏报灰心。\par

他像是对赛跑者、拳击手和战士说话。你们看见他如何装备他们，如何鼓励他们吗？他说：\Quote{走正直的路}。这里他是针对他们的思想说的；就是说，不要怀疑。因为如果管教是出于爱，出于关怀，并以美好的结局告终（他用事实、言语和一切理由证明了这一点），你们为何灰心呢？因为灰心丧气的人，就是那些不被将来的希望所坚固的人。他说：\Quote{走正直的路}，好使你们的瘸腿不致加重，反而恢复原状。因为一个瘸腿的人奔跑，会加重伤处。你们看见吗？我们完全有可能得到彻底的痊愈。\par

% structure: p0011 source=h2 target=subsection reason=relative-heading-rank; base=h2

\subsection{希伯来书 12:14}

2. \BibleQuote{你们要追求与众人和平，并追求圣德，没有圣德，没有人能看见主。}他在上面所说的\Quote{不可放弃聚会} \BibleRef{希 10:25}，这里也暗示了。因为没有什么比独处更容易使人受到诱惑的打击和征服。因为，请告诉我：在战争中分散阵形，敌人便毫不费力，只需逐个捕捉他们，他们就更加无助。\par

他说：\BibleQuote{你们要追求与众人和平，并追求圣德}。那么也要与作恶的人和平吗？他说：\BibleQuote{若是可能，就要尽力与众人和睦} \BibleRef{罗 12:18}。他的意思是：你们要尽力\Quote{和睦生活}，不要损害信仰；但无论你们受到什么恶待，都要高尚地忍受。因为忍耐恶事是试炼中的一大武器。基督也使祂的门徒强壮，说：\BibleQuote{看，我派遣你们好像羊进入狼群；所以你们要机警如同蛇，纯朴如同鸽子} \BibleRef{玛 10:16}。祢说什么？我们在\Quote{狼群}中，祢却叫我们像\Quote{羊}和\Quote{鸽子}？是的，祂说。因为没有什么比高尚地承受加于我们身上的恶事更能使作恶者羞愧的了；并且不以言语或行为报复。这既使我们自己更有哲学精神，也为我们赢得更大的赏报，并且有益于他们。有人傲慢无礼吗？你当祝福他。看，你将由此得到多少益处：你熄灭了邪恶，为自己获得了赏报，使他惭愧，而你自己并未受到任何严重损害。\par

\Quote{與眾人和睦，並追求聖潔。} 他所說的\Quote{聖潔}是什麼意思？是指在婚姻中過貞潔而有規律的生活。他說：若有人未婚，讓他保持純潔，讓他結婚；或者若他已婚，讓他不要行淫，而要與自己的妻子同住：因為這也是\Quote{聖潔}。如何呢？婚姻不是\Quote{聖潔}，但婚姻保存了從信德而來的聖潔，不允許與娼妓結合。因為\Quote{婚姻是尊貴的}\BibleRef{希 13:4}，但不聖潔。婚姻是純潔的：但它本身並不賜予聖潔，除非它禁止褻瀆由我們的信德所賜予的那聖潔。\par

\Quote{沒有它，（他說）沒有人能見主。} 他也對格林多人說：\Quote{你們不要自欺：無論是淫亂的、拜偶像的、奸淫的、作娈童的、亲男色的、偷竊的、貪婪的、醉酒的、辱罵的、勒索的，都不能承繼天主的國。}\BibleRef{格前 6:9} 因為那成了娼妓身體的人，怎能成為基督的身體呢？\par

% structure: p0016 source=h2 target=subsection reason=relative-heading-rank; base=h2

\subsection{希伯來書 12:15}

\Quote{你們要謹慎，免得有人得不到天主的恩寵；免得有苦根生出來攪擾你們，因而使許多人被玷污；免得有淫亂的或褻慢的人。} 你看他如何處處將共同的救恩放在每個人手中？\Quote{天天彼此勸勉，趁著還稱為\InnerQuote{今天}時。}\BibleRef{希 3:13} 那麼，不要將所有重擔都丟給你們的教師；不要全丟給管理你們的人：你們也能彼此建立（他說）。這點他在給得撒洛尼人的信中也寫道：\Quote{你們要彼此勸勉，正如你們所行的。}\BibleRef{得前 5:11} 又說：\Quote{你們要用這些話彼此安慰。}\BibleRef{得前 4:18} 這也是我們現在勸勉你們的。\par

如果你們願意，你們彼此間的成就會比我們更大。因為你們彼此相處的時間更長，你們比我們更了解彼此的事務，你們也知道彼此的缺失，你們有更大的言論自由、愛和親密；這些對於教導來說並非小事，而是偉大而適宜的入門：你們比我們更能責備和勸勉。不僅如此，因為我只是一個人，而你們是許多人；你們無論有多少，都能成為教師。因此，我懇求你們，不要\Quote{忽略了這恩賜}\BibleRef{弟前 4:14}。你們每一個人都有妻子、有朋友、有僕人、有鄰居；讓他責備他，讓他勸勉他。\par

因為，在肉體的營養方面，人們結伴共食、共飲，並定下日子彼此湊份子（如他們所說），透過結伴來補足每人獨自一人的不足——例如，若必須去參加葬禮、宴席，或幫助鄰居處理任何事情——而不為德行的教導這樣做，豈不是荒謬嗎？是的，我懇求你們，不要讓任何人忽略它。因為他從天主那裡領受的賞報是大的。為了讓你們明白，那受托五個塔冷通的是教師；那受托一個塔冷通的是學生。如果學生說：我是學生，我不冒風險，並將他所領受的天主的才能、那共同而簡單的才能隱藏起來，不給建議，不明說，不責備，不勸勉（若他能夠），卻把它埋在地裡（因為那隱藏天主恩賜的心腸確實是塵土和灰燼）：那麼，如果他因懶惰或邪惡而將其隱藏，那麼說\Quote{我只有一個塔冷通}對他毫無辯護。你有一個塔冷通。你本該再帶來一個，使塔冷通加倍。如果你帶來了一個額外的，你就不會受責備。因為祂沒有對那帶來兩個的說：\Quote{你為什麼沒有帶來五個？}反而看他和那帶來五個的價值相同。為什麼？因為他賺得了和他所有的同樣多。而且，因為他比那受托五個的領受得更少，他並沒有因此疏忽，也沒有因他的微小（受托）而藉口懶惰。你不該看那有兩個的；或者更確切地說，你應該看他，正如那有兩個的效法了有五個的，你也該效法那有兩個的。因為，如果對於有財物而不施捨的人有刑罰，那麼對於能夠以任何方式勸勉而不這樣做的人，豈不是有最大的刑罰嗎？在前者，身體得到營養；在後者，靈魂得到營養；那裡你阻止了暫時的死亡，這裡你阻止了永久的死亡。\par

但你說：我沒有口才。但不需要口才或雄辯。如果你看見一個朋友去行淫，對他說：你在追求一件惡事；你不羞愧嗎？你不臉紅嗎？這是錯的。你難道不知道這是不對的嗎？但他被情慾拖著走。患病的人也知道喝冷水不好，但他們需要有人阻止他們。因為受苦的人不容易在病中幫助自己。所以需要你這健康的人來醫治他。如果他不被你的話說服，就在他離開時監視他，拉住他；也許他會感到羞愧。\par

\Quote{那有什麼益處呢？}（你說）\Quote{當他為了我這樣做，因為被我攔住了？}不要過於細算。暫時，無論用什麼方法，使他離開惡習；讓他習慣於不再掉入那個坑中，無論是藉著你，還是藉著任何方法。當你使他習慣了不去，那麼在他稍微喘口氣之後，你就能教導他，他應該為天主而做，而不是為人。不要想一次就全部糾正，因為你做不到；而要溫和地、逐步地做。\par

如果你看见他去喝酒，或去除了酗酒别无他物的宴会，那时你也照样行；另一方面，如果他发现你有任何缺点，也要恳求他帮助你、纠正你。因为这样，他甚至会自己承受规劝，当他看见你同样需要规劝，并且你不是作为凡事都做对了的人，也不是作为老师，而是作为朋友和兄弟帮助他时。要对他说：我为你做了服务，提醒你合宜的事；你也一样，无论你看见我有什么缺点，请阻止我，纠正我。若你看见我易怒、贪婪，请约束我，用劝诫来束缚我。\par

这就是友谊；这样，\Quote{兄弟帮助兄弟，成了一座坚城。} \BibleRef{箴 18:19} 因为不是吃喝构成友谊；这样的友谊连强盗和杀人犯也有。但如果我们真是朋友，如果真的彼此关心，就让我们在这些方面互相帮助。这引导我们走向有益的友谊：让我们阻止那些引人走向地狱的事物。\par

7. 因此，被规劝的人不要愤慨：因为我们是人，都有缺点；规劝的人也不要像夸胜于人、炫耀自己那样去做，而要私下地、温和地做。规劝别人的人需要更大的温和，这样才能说服人忍受切割。你们没有看见外科医生，当他们烧灼、切割时，是何等温和地施用治疗吗？规劝别人的人更应当如此行。因为规劝比火和刀更锋利，使人惊跳。为此，外科医生费尽心力使他们安静地忍受切割，并尽可能轻柔地施行，甚至稍微让步，然后给予喘息的时间。\par

我们也应当如此提出规劝，使被规劝的人不致惊走。因此，即使必须受辱，甚至受击打，我们也不应推辞。因为那些被切割的人也向切割他们的人发出无数的喊叫；但他们不顾这些，只关心病人的健康。同样，在此情形中，我们也应当做一切事使我们的规劝有效，忍受一切，仰望那将临的赏报。\par

\Quote{你们要彼此担待重担，} 他说，\Quote{这样就满全了基督的律法。} \BibleRef{迦 6:2} 如此，我们彼此规劝、互相担待，就能完成建立的工作。这样，你们将使我们的劳苦变得轻松，凡事与我们同担，伸出援手，成为分享者与参与者，既在彼此的救恩上，也在各人自己的救恩上。让我们耐心忍受，彼此担待重担，互相规劝，以求达到在主耶稣基督内所应许的美善。愿光荣、权能、尊崇归于他与圣父及圣神，从今世直到永远。阿们。\par

\SourceRecord{Translated by Frederic Gardiner.；Nicene and Post-Nicene Fathers, First Series；Vol. 14.；Edited by Philip Schaff.；Buffalo, NY: Christian Literature Publishing Co.,；1889.；<http://www.newadvent.org/fathers/240230.htm>.}

% structure: p0001 source=h1 target=section reason=page-title

\section{希伯来书第三十一篇讲道词}

% structure: p0002 source=h2 target=subsection reason=relative-heading-rank; base=h2

\subsection{希伯来书 12:14-15}

\BibleQuote{追求与众人和平，并追求圣德，若无圣德，无人能见主。}\par

有许多事物是基督信仰的特色，但超越一切且更美好的，就是彼此相爱与和平。因此，基督也说：\Quote{我把我的平安赐给你们。}\BibleRef{若 14:27} 又说：\Quote{如果你们彼此相爱，众人因此就可认出你们是我的门徒。}\BibleRef{若 13:35} 因此，保禄也说：\Quote{追求与众人和平，并追求圣德}，即纯洁，\Quote{若无此，无人能见主。}\par

\Quote{仔细察看，免得有人失落天主的恩宠。} 他这样说，仿佛他们正结伴同行一段漫长的旅程，在一大群人中。你们要注意，不可让任何人落后：我不仅寻求你们自己抵达，也寻求你们仔细照料其他人。\par

他说：\Quote{免得有人}失落天主的恩宠。（他指的是未来的美善事物、福音的信德、最佳的生活方式：因为它们都是出于\Quote{天主的恩宠}。）不要对我说，那只是一个人丧亡。即使为一个人，基督也死了。难道你不关心\Quote{基督为他而死的}那人吗？\BibleRef{格前 8:11}\par

他说：\Quote{仔细察看}，即仔细寻找、思考、彻底查证，如同对待病人一样，并以各种方式检查、彻底查证。\Quote{免得有苦根子长起来困扰你们。}\BibleRef{申 29:18} 这是出自《申命记》；他借用植物的比喻。他说：\Quote{免得有苦根子}；他在另一处写道：\Quote{一点面酵能使整个面团发酵。}\BibleRef{格前 5:6} 他的意思是，我不仅为他本人希望这样，也因由此产生的危害。这就是说，即使存在这种根，也不要让任何嫩芽长出来，而要把它砍掉，免得它结出果实，免得它玷污和污染别人。因为他说：\Quote{免得有苦根子长起来困扰你们；并因此使许多人受玷污。}\par

他称罪为\Quote{苦}是很有道理的：因为确实没有比罪更苦的事，那些犯罪后良心受谴责、忍受许多苦楚的人知道这一点。因为罪极其苦，它甚至歪曲了理性本身。这就是苦的性质：它毫无益处。\par

他说\Quote{苦根子}也很好。他没有说\Quote{苦的}，而是说\Quote{苦的根源}。因为一颗苦的根有可能结出甜果实；但一个苦的根源、泉源和基础，却绝不可能结出甜果实；因为一切都是苦的，没有任何甜味，都是苦的，都是令人不快的，都充满了憎恶和厌恶。\par

他说：\Quote{并因此使许多人受玷污。} 也就是说，要断绝那些好色的人。\par

% structure: p0011 source=h2 target=subsection reason=relative-heading-rank; base=h2

\subsection{希伯来书 12:16}

2. \Quote{恐怕有淫乱者，或亵圣的人，如厄撒乌，他为了一餐食物，出卖了长子的名分。}\par

厄撒乌在哪方面是\Quote{淫乱者}呢？他并未说厄撒乌是淫乱者。他说：\Quote{恐怕有淫乱者}，然后又说：\Quote{追求圣德：恐怕有人像厄撒乌一样，是亵圣的}：即贪食、无节制、属世、出卖属灵之事。\par

\Quote{他为了一餐食物出卖了长子的名分}，他因自己的懒惰出卖了从天主而来的这个尊荣，为了一点快乐，失去了最大的尊荣和荣耀。这正适合他们。这是可憎恶的、不洁之人的行为。所以，不仅淫乱者是不洁的，贪食者、肚腹的奴隶也是不洁的。因为他也是一种快感的奴隶。他被迫欺诈、被强迫贪婪、以千万种方式行为不端，成了那情欲的奴隶，且时常亵渎。所以他认为\Quote{长子的名分}毫无价值。也就是说，为了暂时的满足，他甚至牺牲了\Quote{长子的名分}。所以从此以后，\Quote{长子的名分}属于我们，不属于犹太人。同时，这也加重了他们的灾难：首先的成了最后，最后的成了首先：前者因勇敢的忍耐，后者因怠惰而成为最后。\par

% structure: p0015 source=h2 target=subsection reason=relative-heading-rank; base=h2

\subsection{希伯来书 12:17}

3. 他说：\Quote{你们知道，后来他想要承受祝福，竟被拒绝了。因为他没有找到悔改的机会，虽然他流泪切求。} 这是怎么回事？难道他真的排除了悔改吗？绝不是。但是，你们可能会问，他怎么\Quote{没有找到悔改的机会}？因为他若谴责自己，若大声哀号，为何\Quote{没有找到悔改的机会}？因为那并不是真正的悔改。正如加音的悲伤不是悔改，谋杀证明了这一点；同样，在此案中，他的话也不是悔改，后来的谋杀证明了这一点。因为他也曾意图杀害雅各伯。他曾说：\Quote{为父亲守丧的日子近了，那时我要杀死我的弟弟雅各伯。}\BibleRef{创 27:41} \Quote{眼泪}不能给他\Quote{悔改}的能力。并且（宗徒）没有简单地说\Quote{藉悔改}，而是说\Quote{虽流泪切求，仍没有找到悔改的机会}。为什么呢？因为他没有按应当的方式悔改，这才是悔改；他没有按应有的方式悔改。\par

因为他是（宗徒）如何说出这话的呢？他们已变得\BibleQuote{怠惰}\BibleRef{希 6:12}，他是如何再次劝诫他们的呢？当他们已变得\BibleQuote{瘸腿}时\BibleRef{希 12:13}，他是如何劝诫的呢？当他们\BibleQuote{瘫痪}时\BibleRef{希 12:12}，他是如何劝诫的呢？当他们\BibleQuote{松弛}\BibleRef{希 12:12}时，他是如何劝诫的呢？因为这是跌倒的开端。在我看来，他似乎暗示他们中间有些行淫的人，但当时不愿纠正他们；他佯装不知，好让他们自己改正。因为起初确实应该佯装不知；但后来，当他们继续犯罪时，再加以责备，这样他们就不至于变得不知羞耻。梅瑟在齐默黎和科斯比的女儿的事上也这样做了。\par

\Quote{因为他没有找着}——他说——\Quote{悔改的机会}，他没有找着悔改；或者说，他犯罪到了无可悔改的地步。那么，确有罪过到了无可悔改的地步。他的意思是：让我们不要陷入无法挽回的跌倒。只要还是瘸腿的问题，就容易重新站直；但如果我们误入歧途，那还剩下什么呢？因为他这样对那些尚未跌倒的人说话，以恐惧震慑他们，说跌倒的人不可能得到安慰；但对那些已经跌倒的人，为了不使他们陷入绝望，他却说相反的话，如此说：\BibleQuote{我的孩子们！我愿再为你们受产痛，直到基督在你们内形成。}\BibleRef{迦 4:19} 又说：\BibleQuote{凡是靠法律成义的，你们就与基督隔绝，从恩宠上跌了下来。}\BibleRef{迦 5:4} 看！他证实他们已经跌倒了。因为站立的人听见跌倒后不可能获得宽恕，就会更加热心，更加谨慎地站稳；然而，如果你对已经跌倒的人也使用同样的严厉，他就永远不会再站起来。因为凭着什么希望他会表现出改变呢？\par

但他不仅哭泣（你们说），而且\Quote{恳切寻求}。那么，他并没有排除悔改；而是叫他们谨慎，不要跌倒。\par

4. 因此，凡不信地狱的人，让他们想起这些事；凡以为犯罪不会受罚的人，让他们考虑这些事。为什么厄撒乌没有获得宽恕？因为他没有照所应当的悔改。你愿看见完美的悔改吗？请听伯多禄背主后的悔改。因为圣史在叙述关于他的事时说：\BibleQuote{他出去，痛哭起来。}\BibleRef{玛 26:75} 因此，甚至这样的大罪也被赦免了，因为他照所应当的悔改了。虽然牺牲尚未奉献，祭献尚未作成，罪过尚未被除去，它仍然掌权作主。\par

并且为叫你们知道，这背主（的起因）与其说是由于懈怠，不如说是由于他被天主舍弃，天主教训他认识人的限度，不要反驳师傅的话，也不要比别人更高傲，而要晓得没有天主什么也不能作，并且\BibleQuote{若不是上主建造房屋，建造的人徒然劳苦}\BibleRef{咏 126:1}；因此基督也唯独向他说：\BibleQuote{撒殚想要筛你们像筛麦子一样，我不容许，为使你的信德不致丧失。}\BibleRef{路 22:31} 因为他很可能自高自大，自己意识到他比众人更爱基督，因此\Quote{他痛哭起来}；并且他在哭泣之后，也做了同类的事。因为他做了什么？此后他把自己置于无数危险中，并以许多方式显示了他的刚毅和勇气。\par

犹达斯也后悔了，但却是以邪恶的方式：因为他上吊自缢了。厄撒乌也后悔了，正如我所说；或者更确切地说，他连后悔也没有；因为他的眼泪不是悔改的眼泪，而是骄傲和愤怒的眼泪。接下来发生的事证明了这一点。有福的达味悔改了，如此说：\BibleQuote{我每夜要洗涤我的床铺，以眼泪浸湿我的卧榻。}\BibleRef{咏 6:6} 并且那许久以前犯的罪，过了这么多年，这么多世代，他仍然哀悼，仿佛刚刚发生一样。\par

5. 因为悔改的人不应当发怒，也不应当凶暴，而应当痛悔，如同被定罪的人，如同没有胆量的人，如同已被判决的人，如同只能靠怜悯得救的人，如同对自己的恩人忘恩负义的人，如同不知感恩的人，如同被弃绝的人，如同该受无数惩罚的人。如果他想到这些，就不会发怒，不会愤怒，而要悲伤、哭泣、呻吟，昼夜哀叹。\par

悔改的人决不应当忘记自己的罪过，一方面恳求天主不要记念它，另一方面他自己却永不忘记。如果我们记念它，天主就会忘记它。让我们对自己施行惩罚，让我们控告自己；这样我们就会取悦审判者。因为承认的罪过会变小，不承认的会变得更坏。如果罪过加上不知羞耻和忘恩负义，那么连自己从前犯过罪都不知道的人，怎能防备自己再次陷入同样的罪恶呢？\par

因此，我恳求你们，不要否认我们的罪过，也不要不知羞耻，免得我们被迫受罚。加音听见天主说：\BibleQuote{你的弟弟亚伯尔在那里？他回答：我不知道，我岂是看守我弟弟的人？}\BibleRef{创 4:9} 你们看见这如何使他的罪更重吗？但他的父亲却没有这样做。那么怎样呢？当他听见：\BibleQuote{亚当，你在那里？}\BibleRef{创 3:9} 他说：\BibleQuote{我听见了你的声音，就害怕起来，因为我赤身露体，就躲藏了。}\BibleRef{创 3:10} 承认我们的罪过，并不断记念它们，实是一大善事。没有什么比不断记念罪过更能有效地治好过犯。没有什么能使人如此不轻易作恶。\par

6. 我知道良心会退避，不愿因回忆恶行而受鞭挞；但要紧紧抓住你的灵魂，给它戴上笼头。因为如同一匹难驯的马，它不耐忍受加于它的事，不愿说服自己犯了罪：但这全是撒殚的作为。让我们说服它犯了罪；让我们说服它犯了罪，好使它也悔改，以便悔改后逃脱痛苦。告诉我，当你尚未告明你的罪时，你以为如何能获得赦免？的确，犯罪的人值得怜悯与仁慈。但你这尚未说服自己（犯了罪）的人，既然对某些事毫无羞愧，怎能期望被怜悯呢？\par

让我们说服自己犯了罪。不仅用口舌说，也要用心思说。不要只称自己是罪人，还要细数我们的罪，逐一过目。我不是要你炫耀自己，也不是要你在别人面前控告自己；而要听从先知的话，他说：\Quote{向主陈明你的道路。} \BibleRef{咏 36:5} 在天主面前告明这些事。在审判者面前以祈祷告明你的罪；若不能用口舌，也要用记忆，这样你才能堪受怜悯。\par

如果你不断记念自己的罪，就永不会记念邻人的过犯。我不是说，如果你说服自己是个罪人；这并不像单独检验罪过本身那样能谦卑灵魂。如果你不断记念这些事，就不会记念所受的委屈；你不会发怒，不会辱骂，不会有高傲的思想，不会再次陷入同样的罪，反而会更热切地趋向善事。\par

7. 你看见从记念我们的罪产生了多少美妙的效果吗？那么，让我们把这些罪写在我们的心思里。我知道灵魂不能忍受如此痛苦的回忆；但让我们约束并强迫它。现在被这回忆啮咬，总比那时被惩罚啮咬更好。\par

现在，如果你记念它们，并不断将它们呈现在天主面前\EditorialNote{参看第448页}，为它们祈祷，你必能迅速抹去它们；但如果你现在忘记它们，那时你将在不愿的情况下被记起，当它们在全世界面前，在众朋友、仇敌和天使面前被公开展示时。因为天主并非只对达味说：\Quote{你所行的事，我必在} \BibleRef{撒下 12:12} \Quote{众人面前显明}，而是对我们每个人说的。你怕人（他说），看重他们过于天主；而天主看见你，你却不关心，却因人在面前而羞愧。因为经上说：\Quote{人的眼目，这是他们的恐惧。} 因此你将在这点上受罚；因为我必责斥你，将你的罪摆在众人眼前。因为这事是真实的，在那日我们所有人的罪都将被公开展示，除非我们现在借着不断的记念将它们除去。请听残酷和无情是如何被公开揭露的：\Quote{我饿了}（祂说）\Quote{你们没有给我吃的。} \BibleRef{玛 25:42} 这些话是在何时说的？是在角落里？是在隐秘处？绝对不是。那时是什么时候？\Quote{当人子在祂的光荣中来临时} \BibleRef{玛 25:31}，\Quote{万民}都聚集在一起，祂将彼此分开，然后在众人面前说话，将\Quote{他们} \Quote{安置在祂的右边}和\Quote{左边} \BibleRef{玛 25:33}：\Quote{我饿了，你们没有给我吃的。}\par

再看那五个童女也在众人面前听到：\Quote{我不认识你们。} \BibleRef{玛 25:12} 因为五个和五个不仅指五个的数量，而是指那些邪恶、残忍、无情的童女和那些并非如此的童女。同样，那埋没一个塔冷通的仆人也在带了五个和两个塔冷通的人面前听见：\Quote{你这邪恶懒惰的仆人。} \BibleRef{玛 25:26} 但祂不仅用言语，也用行为定他们的罪；正如福音作者所说：\Quote{他们要瞻望他们所刺透的那一位。} \BibleRef{若 19:37} 因为复活将是同时发生，包括罪人和义人。祂也同时出现在审判中给所有的人。\par

8. 因此，想想那时谁将惊慌，谁将悲伤，谁被拖去受火刑，而另一些人则被加冕。\Quote{来吧}（祂说），\Quote{我父所祝福的，承受自创世以来为你们预备的国度吧。} \BibleRef{玛 25:34} 又说：\Quote{离开我，到那为魔鬼和他的使者预备的火里去。} \BibleRef{玛 25:41}\par

让我们不仅听这些话，还要把它们写在眼前，想象祂现在就在说这些话，而我们正被带去受那火刑。我们那时会有什么心肠？什么安慰？当我们被劈开时？当我们被指控贪婪时？我们能说出什么托辞？什么巧辩？毫无；只能被捆绑、弯腰、拖到火炉口、火河、黑暗、永远不止的刑罚，而且无人可求。因为经上说：不可能从这边到那边去；\Quote{在我们与你们之间有一道深渊} \BibleRef{路 16:26}，即使有人愿意，也无法过去伸出援手；我们必须不断焚烧，无人援助，即使是父亲或母亲，或任何谁，即使他对天主有很大的胆量。因为经上说：\Quote{兄弟不能赎回；人能赎回吗？} \BibleRef{咏 48:8}\par

既然人不能将得救的希望寄托于他人，而（必须）在自己身上，靠着天主的慈爱，我恳求你们，让我们做一切事，使我们的品行纯洁，生活最美善，甚至从一开始就不沾染任何污点。但若非如此，无论如何，我们不要在沾染污点后沉睡，而要不断藉着悔改、眼泪、祈祷和慈悲工作洗去污秽。\par

那么，你或许会说：若我不能行慈悲工作怎么办？然而，无论你多么贫穷，你总有\BibleQuote{一杯凉水}\BibleRef{玛 10:42}。无论你多么贫乏，你总有\BibleQuote{两个小钱}\BibleRef{谷 12:42}。你还有双脚，可以探访病人，可以进入监牢；你还有屋顶，可以接待旅客。因为凡不行慈悲工作的人，绝无宽恕，绝无宽恕。\par

我们不断地对你们说这些事，为的是藉着反复重申，哪怕能产生一点效果也好。我们这样说，并非那么关心那些接受益处的人，而是关心你们自己。因为你们确实给了他们今世的事物，却换得了天上的事物：愿我们都能获得这一切，在基督耶稣我们的主内，愿光荣归于圣父，偕同圣神，自今而始，至于永远，世世无穷。阿们。\par

\SourceRecord{Translated by Frederic Gardiner.；Nicene and Post-Nicene Fathers, First Series；Vol. 14.；Edited by Philip Schaff.；Buffalo, NY: Christian Literature Publishing Co.,；1889.；<http://www.newadvent.org/fathers/240231.htm>.}

% structure: p0001 source=h1 target=section reason=page-title

\section{希伯来书第三十二篇讲道}

% structure: p0002 source=h2 target=subsection reason=relative-heading-rank; base=h2

\subsection{希伯来书 12:18-24}

\BibleQuote{\Emph{因为你们并非走近了} \Emph{一个可触摸且燃烧的火}，以及黑暗、昏黑、暴风、号角的声音和说话的声音；听见这声音的人都恳求不要再对他们说话。\Emph{（因为他们不能承受所命令的：\InnerQuote{连走兽若触摸这山，也该用石头砸死。}} \Emph{并且这景象如此可怕，以致\Person{梅瑟}说：\InnerQuote{我非常恐惧战栗。}）但是你们却走近了\Place{熙雍山}和永生天主的城、天上的\Place{耶路撒冷}，以及千万天使、欢聚的集会}，\Emph{和登记在天上的首生者的教会，并审判众人的天主、已臻完美的义人的灵魂}，\Emph{以及新约的中保\Person{耶稣}和所洒的血，这血比\Person{亚伯尔}的血说话更美。}}\par

1. 圣殿内的事物，至圣所，诚然神奇；而在\Place{西乃}山上所行的事也同样是可畏的：\Quote{火，黑暗，昏黑，暴风。} \BibleRef{申 33:2} 因为经上记载：\Quote{天主在\Place{西乃}显现，}这些事在古代就被传颂。然而，新约并非借由这些事物赐下，而是由天主以简单的言词赐予。\par

请看，他如何也在这些方面进行比较。他随后提出这些是很有理由的。因为当他用无数的[论据]说服他们，并已显示出每个盟约之间的区别之后，然后，既然那一个已被定罪，他便容易地进入这些方面。\par

他说什么呢？\BibleQuote{因为你们并非走近了可触摸且燃烧的火，以及黑暗、昏黑、暴风、号角的声音和说话的声音；听见这声音的人都恳求不要再对他们说话。}\par

他意思是说，这些事是可畏的；如此可畏，以致他们甚至不能忍受听见，连\Quote{走兽}也不敢上山。（但后来之事并非如此。\Place{西乃}之于天何有？可触摸之火与不可触摸的天主何有？因为\BibleQuote{天主是吞噬的火。} \BibleRef{希 5:29}）因为经上记载：\Quote{不要让天主对我们说话，让\Person{梅瑟}对我们说话。所命令的事如此可畏：即使走兽触摸这山，也必用石头砸死；\Person{梅瑟}说：我非常恐惧战栗。} \BibleRef{出 20:19}至于百姓，有何可怪？他自己进入\BibleQuote{天主所在的黑暗中}，却说：\BibleQuote{我非常恐惧战栗。} \BibleRef{出 20:21}\par

2. \BibleQuote{但是你们却来到了\Place{熙雍山}和永生天主的城、天上的\Place{耶路撒冷}，以及千万天使、欢聚的集会、和登记在天上的首生者的教会，并审判众人的天主、已臻完美的义人的灵魂，以及新约的中保\Person{耶稣}和所洒的血，这血比\Person{亚伯尔}的血说话更美。}\par

取代\Quote{梅瑟}的是\Person{耶稣}。取代百姓的是\Quote{千万天使}。\par

他所说的\Quote{首生者}是指谁？是指信友们。\par

\Quote{并已臻完美的义人的灵魂。}他说，你们将与这些人同在。\par

\Quote{并新约的中保耶稣和所洒的血，这血比亚伯尔的血说话更美。}那么，\Person{亚伯尔}的[血]会说话吗？他说：\BibleQuote{是的，他虽死了，却藉着它仍说话。} \BibleRef{希 11:4} 天主又说：\BibleQuote{你弟弟的血由地里向我喊冤。} \BibleRef{创 4:10} 或者是这个[意义]或者是那个；因为至今它仍被传颂，但不如基督的那样。因为基督的血洁净了众人，并发出更清晰、更分明的声音，因为它有更大的见证，即事实的见证。\par

% structure: p0013 source=h2 target=subsection reason=relative-heading-rank; base=h2

\subsection{希伯来书 12:25-29}

\BibleQuote{你们要谨慎，不要拒绝那说话的那位。因为如果那些拒绝在地上说话的那位的人不能逃脱，何况我们如果背弃从天上说话的那位，就更不能逃脱了。当时他的声音震动了地，但现在他应许说：\InnerQuote{再一次我不单要震动地，而且也要震动天。}这\InnerQuote{再一次}的话，是指明那些被震动的、受造之物要被除去，好让那些不被震动的得以常存。因此，我们既然领受了一个不能震动的国度，就应当怀着恩宠，藉此以虔诚和敬畏的心事奉天主，为中悦他所用。因为我们的天主是吞噬的火。}\par

3. 那些事是可畏的，但这些事却更加可敬和荣耀。因为在这里没有\Quote{黑暗}，也没有\Quote{昏黑}，也没有\Quote{暴风}。在我看来，他藉这些言语暗示旧约的晦涩，以及律法的蒙蔽与遮掩的特性。此外，律法的颁布者以烈火显现，可畏且适于惩罚违命者。\par

但\Quote{号角的声音}是什么？大概是如同某位君王来临。这无疑也将发生在第二次来临时。经上说：\BibleQuote{在最后号角响起时} \BibleRef{格前 15:52}，众人都要复活。但成就此事的是他声音的号角。那时，一切事物都是感官的、可见的、可听的；如今一切事物都是理智的、不可见的。\par

并且经上记载：\BibleQuote{有浓烟。} \BibleRef{出 19:18}。因为天主被称为火，并在荆棘丛中如此显现，他藉着烟来表示火。而\Quote{昏黑与黑暗}是什么？他再次表达其可畏。\Person{依撒意亚}也说：\BibleQuote{圣殿充满了烟云。} \BibleRef{依 6:4}。而\Quote{暴风}的目的是什么？人类原本粗心大意。因此有必要藉这些事来唤醒他们。因为当这些事发生、律法颁布时，没有人迟钝到不提升自己的心。\par

\Quote{梅瑟说话，天主用声音回答了他} \BibleRef{出 19:19}：因为天主的声音必须被发出。既然祂要藉梅瑟颁布祂的法律，因此祂使他堪受信赖。他们看不见梅瑟，因为浓厚的黑暗；也听不见他，因为他声音微弱。那么，又怎样？\Quote{天主用声音回答}，向群众讲话：是的，他的名要被传扬。\par

\Quote{他们恳求}（他说）\Quote{不再向他们说话。}\par

起初，他们自己就是天主藉肉体显现的原因。让梅瑟对我们说话，\Quote{不要让天主对我们说话} \BibleRef{出 20:9}。作比较的人更加抬高一边，为显示另一边远为伟大。在这方面，我们的［特权］更为温和、更为可赞叹。因为它们在两方面是伟大的：不但光荣且更伟大，而且更容易接近。他在致格林多人书中也这样说：\Quote{带着揭开的脸帕} \BibleRef{格后 3:18}，以及\Quote{不像梅瑟将帕子蒙在脸上} \BibleRef{格后 3:13}。他们，他说，不配得到我们所［得到的］。他们配得到什么？他们看见了\Quote{黑暗、幽暗}，听见了\Quote{一个声音}。但你也听见了一个声音，不是透过黑暗，而是透过肉体。你没有被打扰，没有受困扰，而是站着与中保交谈。\par

在另一种方式下，藉着\Quote{黑暗}他显示了不可见性。\Quote{黑暗}（经上说）\Quote{在祂脚下} \BibleRef{咏 17:9}。\par

\Emph{那时}连梅瑟都害怕，但现在无人害怕。\par

如同那时百姓站在下面，我们也这样。他们并非在下面，而是在天之下。圣子亲近天主，但不似梅瑟。\par

\Emph{那里}是旷野，这里是城。\par

4. \Quote{并接近千万天使。}这里他显示喜乐、欢愉，以代替\Quote{幽暗}、\Quote{黑暗}和\Quote{风暴}。\par

\Quote{并接近总集会，和记载在天上的首生者的教会，以及审判万有的天主。}他们没有走近，而是远远站着，连梅瑟也如此；但\Quote{你们已经走近了。}\par

这里他使他们害怕，说：\Quote{并接近审判万有的天主}；不单审判犹太人、信徒，而是审判全世界。\par

\Quote{并接近被成全的义人的灵魂。}他指那些蒙悦纳者的灵魂。\par

\Quote{并接近新约的中保耶稣，以及洒的血}，即洁净的血，\Quote{这血比亚伯的血说得更好}。如果血说话，更何况那被杀后活着的那位呢？但他说什么？\Quote{圣神}（他说）\Quote{也以无可言喻的叹息代求} \BibleRef{罗 8:26}。祂如何代求？每当祂进入一颗诚心，就提升它，并使之说话。\par

5. \Quote{你们要谨慎，不可拒绝那发言的}；即不可拒绝［祂］。\Quote{因为若是他们拒绝了那在地上发言的，不能逃脱。}他指谁？我想是梅瑟。他所说的是：如果他们\Quote{拒绝了}那在地上颁法律的\Quote{祂}而\Quote{不能逃脱}，我们若拒绝那从天上发言的，怎能逃脱？他这里并非说祂是另一位；绝无此事。他没有陈述一位和另一位，而是祂发出\Quote{从天上}的声音时显得可畏。祂自己既是这一位，也是那一位：但这一位是可畏的。因为他表达的并非位格的差异，而是恩赐的差异。这从何显明？\Quote{因为若是他们}（他说）\Quote{拒绝了那在地上发言的，不能逃脱，我们若背弃那从天上发言的，就更不能逃脱。}那么，这一位与那一位不同吗？但他为何说：\Quote{当时祂的声音震动了地}？因为那\Quote{当时}颁布法律的祂的\Quote{声音}\Quote{震动了地。但如今祂应许说：再一次，我不单震动地，而且也要震动天。这\InnerQuote{再一次}这句话，指明被震动的，就是受造之物，都要被除去。}因此一切都要被挪去，并重新塑造得更好。这正是他这里所暗示的。那么，你何必在不能长存的世界受苦时悲伤？在这世界很快就要过去时受苦？如果我们的安息是在世界的末期，那么人应当注视终结而受苦。\par

\Quote{使那}（他说）\Quote{不能震动的得以长存。}但\Quote{那些不能震动的}是怎样的？是那将来的事物。\par

6. 那么，让我们为获得那［安息］、享受那些福乐而做一切。是的，我祈祷并恳求你们，让我们为此努力。没有人在一座将要倒塌的城里建造。请告诉我，如果有人说过一年后这座城要倒塌，而另一座城绝不倒塌，你会在那将要倒塌的城里建造吗？同样，我现在也说这番话：我们不要在这世界里建造；它不久就要倒塌，一切都要被毁。但我为何说：它要倒塌？在它倒塌之前，我们就要被毁灭，遭受可怕的事；我们将被从它们那里移去。\par

为何我们建造在沙土上？让我们建造在磐石上：因为无论发生什么，那建筑屹立不倒，什么都不能摧毁它。有理由。因为对于所有这样的攻击，那区域是难以接近的，正如这区域是容易接近的。因为地震、火灾、敌人入侵，甚至在我们活着时就将它夺去；并且常常连同我们一同毁灭。\par

即使产业得以保留，疾病也迅速将我们带走，或我们若留下，也不得公正地享用。因为哪里有疾病、诬告、嫉妒、阴谋，还有何快乐可言？即便没有这些，若我们无子女，也常会不安、急躁，因没有可继承房屋和一切事物的人；从此我们便因替他人劳作而憔悴。是的，往往遗产在我们去世后，甚至在我们活着时就落入敌人之手。那么，还有什么比为敌人劳碌，而自己却积蓄罪过好让他们安息更悲惨呢？在我们的城市里常见到许多这样的例子。我不再多言，免得使那些被掠夺的人悲伤。因为我可以指名道姓地提到其中一些人，讲述许多故事，向你们展示许多房屋，其主人竟成了为它们劳作者的敌人：不仅是房屋，连奴隶和全部遗产也常常落入敌人之手。人间之事，就是如此。\par

但在天堂，无需惧怕这些——怕人死后敌人来继承遗产。因为那里没有死亡，没有仇恨；圣人们的帐幕是永久的居所；在那些圣人中间有欢腾、喜乐、欢乐。因为经上记载：\BibleQuote{义人的帐幕中有欢呼的声音。}（\BibleRef{咏 117:15}）这些帐幕是永恒的，没有尽头。它们不会因年久而倒塌，不会更换主人，而是永久保持最佳状态。这理所当然。因为那里没有可朽坏、可毁灭之物，一切都是不朽而纯洁的。让我们倾尽所有财富来建造这房屋。我们无需木匠、无需工人。穷人的手建造这样的房屋：跛子、瞎子、残废者，他们建造那些房屋。毫不奇怪，因为他们甚至为我们赢得天国，并使我们敢于面对天主。\par

7. 因为慈悲就像一门最卓越的艺术，是致力于它的人的保护者。它亲近天主，随时站在祂身旁，愿意为任何人求情，只要不被我们亏待；当我们以勒索的方式行善时，就是亏待了它。因此，若它是纯洁的，它便给奉献它的人极大的信心。它甚至为得罪了天主的人、犯罪的人代祷。它的力量如此伟大：它打断锁链，驱散黑暗，熄灭烈火，杀死虫子，赶走咬牙切齿声。天堂之门因它而大开，极其安全：如同一位王后进入时，门卫无人敢问她是何人、来自何方，全都立刻迎接她一样；慈悲也是如此。因为她确实是一位王后，使人肖似天主。因为祂说：\BibleQuote{你们应当慈悲，如同你们的天父是慈悲的一样。}（\BibleRef{路 6:36}）\par

她生有翅膀，轻快自如，金色的羽翼，其飞翔令天使们大为喜悦。经上记载：\BibleQuote{鸽子的翅膀镀有白银，翎毛闪耀着黄金的光泽。}（\BibleRef{咏 67:13}）她像一只活的金色鸽子，神态温柔，目光柔和，翱翔着。没有什么比那目光更美好。孔雀虽美，与她相比不过是只穴鸟。这只鸟如此美丽，令人赞叹。她不断仰望高处；天主的荣耀丰沛地环绕着她；她是一位金翼的贞女，装扮得体，面容美丽而温柔。她有翅膀，轻快自如，站在王座之旁。当我们受审判时，她突然飞临，现身，用自己的翅膀庇护我们，救我们脱离刑罚。\par

天主喜爱她胜过祭献。祂多番谈论她：祂如此爱她。经上记载：\BibleQuote{祂必看顾寡妇和孤儿}（\BibleRef{咏 145:9}）以及穷人。天主愿意因她而被称为\Quote{有慈悲的}。\BibleQuote{上主有怜悯，有恩典，宽恕慢怒，且有丰盛的慈爱}（\BibleRef{咏 144:8}），并真实。天主的仁慈覆盖全地。她拯救了人类（\BibleRef{咏 144:9}）：因为若非祂怜悯我们，万物早已灭亡。\BibleQuote{当我们还是仇敌时}（\BibleRef{罗 5:10}），她使我们和好，成就了无数的福泽；\BibleQuote{她促使天主之子成为奴仆，空虚了自己。}（\BibleRef{斐 2:7}）\par

让我们切望效法那拯救我们的慈悲；让我们爱她，珍惜她，胜过财富，并且超越财富，拥有一颗慈悲的心灵。没什么比慈悲更体现基督徒的特质。外教人和所有人最敬佩的，莫过于我们富有慈悲。因为我们自己也常需要这种慈悲，向天主说：\BibleQuote{天主，求祢按照祢丰厚的慈爱怜悯我们。}（\BibleRef{咏 50:1}）让我们先开始实践：其实并非我们先开始，因祂已向我们显示了自己的慈悲；但至少我们随后跟上。因为若人们怜悯一个有慈悲心的人，即使他做了无数错事，更何况天主呢？\par

8. 请听先知的话：\BibleQuote{我（他的话是）如同天主殿中的茂盛橄榄树。}（\BibleRef{咏 51:8}）让我们成为这样：让我们成为\Quote{橄榄树}；让我们在各方面都结满诫命的果子。因为作橄榄树还不够，还要结果实。有些人施舍时给得很少，一年或一周只一次，或随意给一点。这些人是橄榄树，但不是结果的，甚至是枯干的。因为他们显示同情，所以是橄榄树；但因他们不慷慨施予，所以不是结果实的橄榄树。但让我们结果实吧。\par

我曾屡次说过，现在仍要说：爱德的伟大，不在于所施与的多少，而在于施与者的心意。你们知道那寡妇的事迹。我们应当不断提起这个榜样，使穷人不致因看见她投了两个小钱而灰心。有些人甚至在装饰圣殿时献上毛发，这些也没有被拒绝。\BibleRef{出 35:23} 但如果他们本有金子，却献上毛发，那便是可咒骂的；但若他们只有这个，献上它，便被悦纳。为此，加音也受责备，并非因为他献了无价值之物，而是因为他所献的是他所有中最无价值的。\BibleQuote{可咒骂的}（经上说）\BibleQuote{是那有公畜却向天主祭献有残疾之物的。}\BibleRef{拉 1:14} 他并非绝对地说，而是说\BibleQuote{那有}（他说）\BibleQuote{却吝惜不舍的}。因此，若一人一无所有，他就不受责备，甚或反而有赏报。有什么比两个铜钱更微小？比毛发更无价值的？比一升面粉更微薄的？然而，这些与牛犊和金子同样蒙悦纳。因为\BibleQuote{人若有愿做的心，必蒙悦纳，乃是照他所有的，不是照他所无的。}\BibleRef{格后 8:12} 经上又说：\BibleQuote{你手若有行善的力量，就不可推辞。}\BibleRef{箴 3:27}\par

因此，我恳求你们，让我们欣然为穷人倾囊相助。即使我们献的很少，也必与那些投献最多的人获得同样的赏报；甚或比那些投献一万塔冷通的人所得更多。我们若这样做，必将获得天主那不可言喻的宝藏；若我们不仅听道，而且行道，不仅赞美爱德，也以行动表现出来。愿我们众人皆能获此恩赐，因我们的主耶稣基督，祂与圣父及圣神，永受光荣、权能、尊崇，从现在直到永远，世世无穷。阿们。\par

\SourceRecord{Translated by Frederic Gardiner.；Nicene and Post-Nicene Fathers, First Series；Vol. 14.；Edited by Philip Schaff.；Buffalo, NY: Christian Literature Publishing Co.,；1889.；<http://www.newadvent.org/fathers/240232.htm>.}

% structure: p0001 source=h1 target=section reason=page-title

\section{希伯来书第三十三篇讲道}

% structure: p0002 source=h2 target=subsection reason=relative-heading-rank; base=h2

\subsection{希伯来书 12:28-29}

\BibleQuote{\Emph{为此，我们既然是蒙受一个不可动摇的国度，就应该怀有恩宠（即感激），}\Emph{借此我们事奉}\Emph{天主，以虔诚和敬畏之情。因为我们的天主是吞灭的烈火。}}\par

1. 他在另一处也说了同样的话：\BibleQuote{看得见的是暂时的，看不见的才是永恒的}（见：\BibleRef{格后 4:18}）；并由此就我们在现世所忍受的苦难作出劝勉；在此他也这样做，并说：让我们坚定不移；\Quote{让我们怀有感恩之心}，意即让我们感谢天主。因为我不但不应因现世之事而沮丧，反而应为将来的事向祂表达最大的感谢。\par

\BibleQuote{借此我们中悦地事奉天主}，就是说：\Quote{因为只有这样才有可能中悦地事奉天主}，就是在一切事上感谢祂。\BibleQuote{你们做一切事}（他说）\BibleQuote{总不可抱怨，也不可争论}（见：\BibleRef{斐 2:14}）。因为凡抱怨所做的工作，都会削减并失去奖赏；就如以色列人——他们因抱怨付出了多么大的代价。因此他说：\BibleQuote{你们也不要抱怨}（见：\BibleRef{格前 10:10}）。因此，若不因一切事——包括我们的考验和它们的减轻——而感谢祂，就不可能\Quote{事奉}祂\Quote{中悦地}。这就是说，我们不要说出任何草率、无礼的话，而要谦卑自下，以便保持敬畏。因为这正是\Quote{以虔诚和敬畏之情}。\par

% structure: p0006 source=h2 target=subsection reason=relative-heading-rank; base=h2

\subsection{希伯来书 13:1-2}

\BibleQuote{弟兄相爱要持久。不可忘记款待旅客，因为藉此，有些人竟在不知不觉中款待了天使。}看他如何叮嘱他们持守已有的：他没有添加别的。他没有说：\Quote{要彼此相爱如弟兄}，而是说\Quote{弟兄相爱要持久}。同样，他也没有说：\Quote{要款待旅客}，好像他们不款待似的，而是说\Quote{不可忘记款待旅客}，因为因着他们的苦难，这很容易发生。\par

因此（他说）\BibleQuote{有些人竟在不知不觉中款待了天使}。你看见这是多大的荣耀，多大的益处！\par

什么是\Quote{不知不觉中}？就是他们款待了他们，却不知道。因此奖赏也很大，因为他款待他们，却不知道他们是天使。因为若是知道，那就不足为奇了。有人说他在这里也暗指\Person{罗特}。\par

% structure: p0010 source=h2 target=subsection reason=relative-heading-rank; base=h2

\subsection{希伯来书 13:3-6}

2. \BibleQuote{你们要记念被囚禁的人，好像与他们一同被囚禁；要记念受虐待的人，好像你们自己也亲身受虐待。婚配在众人中应受尊重，床第应是无玷的，因为淫乱和犯奸的人，天主必要审判。你们的生活不要贪爱钱财，要以你们所有的一切为满足。}\par

你看他对贞洁的谈论有多大。他说：\BibleQuote{你们要追求和平与圣洁；免得有淫乱或亵圣的人}（见：\BibleRef{希 12:14}）；又说：\BibleQuote{天主必要审判淫乱和犯奸的人}（见：\BibleRef{希 12:16}）。每次禁止都伴随着惩罚。他说：\Quote{你们要同众人追求和平与圣洁，没有圣洁，谁也不能见到主；但天主必要审判淫乱和犯奸的人。}\par

而他先写下\Quote{婚配在众人中应受尊重，床第应是无玷的}，表明他正确添加了后面的话。因为如果婚姻已被准许，那么淫乱者受罚是公正的，犯奸者受报应是正当的。\par

此处他是在揭露异端者。他没有再说：不要有人淫乱；但一旦说过了，他就像做一般劝勉那样继续下去，而不是专门针对他们。\par

\Quote{你们的生活不要贪爱钱财}，他说。他没有说：不要拥有任何东西，而是说：\Quote{你们的生活不要贪爱钱财}，就是让生活显出你们心智的哲学品格。[它将会显出，如果我们不寻求多余之物，只保留必需之物。]因为他在上面也说：\BibleQuote{你们欣然接受了财产的掠夺}（见：\BibleRef{希 10:34}）。他给出这些劝勉，为叫他们不要贪财。\par

\Quote{要以你们所有的一切为满足}（他说）。然后此处也有安慰（见：\BibleRef{希 13:5}）：\Quote{因为祂曾说过：我决不离开你，也决不弃舍你}（见：\BibleRef{希 13:6}）；\Quote{所以我们可放心大胆地说：主是我的助手，我决不害怕，人能把我怎么样？}再次在考验中给予安慰。\par

% structure: p0017 source=h2 target=subsection reason=relative-heading-rank; base=h2

\subsection{希伯来书 13:7}

3. \Quote{你们要记念那些统治你们的人。}这就是他上面努力要说的：因此\BibleQuote{你们要同众人追求和平与圣洁}（见：\BibleRef{希 12:14}）。他也向得撒洛尼人给出这个劝勉：要\BibleQuote{极其尊重他们}（见：\BibleRef{得前 5:13}）。\par

\BibleQuote{你们要记念（他说）那些统治你们的人，他们曾给你们讲过天主的道理；要效法他们的信德，细心观察他们行为的结局。}这是怎样的效法呢？确实是最好的：因为他说，观察他们的生活，\Quote{效法他们的信德}。因为纯洁的生活[产生]信德。\par

或者，他所谓的\Quote{信德}是指坚定。如何呢？因为他们相信将来的事。因为他们若怀疑将来的事，若犹豫不决，就不会表现出纯洁的生活。所以此处他也是在治疗同样的[邪恶]。\par

% structure: p0021 source=h2 target=subsection reason=relative-heading-rank; base=h2

\subsection{希伯来书 13:8-9}

\BibleQuote{耶稣基督昨天、今天，直到永远，常是一样。你们不可被各种异端邪说引入歧途。因为心靠恩宠得以坚固，才是好的；靠食物，决不是好的，因为靠食物的人，从未因它们获得益处。}\par

在这些话中，\BibleQuote{耶稣基督昨日、今日、直到永远，常是一样}，\Quote{昨日}指过去所有时间：\Quote{今日}指现在：\Quote{直到永远}指无尽的将来。这就是说：你们听说了有一位大司祭，但不是会衰败的大司祭。他永是相同的。好像有些人说：\Quote{他不在了，另有一位要来}，他说这话，表明那位曾是\Quote{昨日和今日}的，也\Quote{直到永远也是相同}。因为甚至现在，犹太人还说另有一位要来；他们既使自己远离了那一位，就会落入假基督的手中。\par

\BibleQuote{不要被各种怪异的道理所摇动。}不仅\Quote{怪异的道理}，也不要\Quote{各种道理}。\par

\BibleQuote{因为靠恩宠坚固人心是好的，而不是靠食物，那在食物上用心的人从未得益。}这里他柔和地暗示那些引进遵守\Quote{食物}的人。因为靠信德万物都是洁净的。所以需要的是信德，不是\Quote{食物}。\par

% structure: p0026 source=h2 target=subsection reason=relative-heading-rank; base=h2

\subsection{希 13:10}

因为\BibleRef{希 13:10} \BibleQuote{我们有一座祭台，那在会幕内侍奉的人没有权利吃}。不像犹太人的[规条]，我们当中的这些事，连大司祭也不可分享。所以既然他说了\Quote{不要遵守}，这话似乎是拆毁自己建筑的人说的，他又转过来。怎么，难道我们没有遵守的事吗？我们也有，而且我们也非常认真地遵守，甚至不与司祭们分享。\par

% structure: p0028 source=h2 target=subsection reason=relative-heading-rank; base=h2

\subsection{希 13:11-13}

4. \BibleQuote{因为那些牲口的身体，其血由大司祭为赎罪带进圣所，身体却在营外焚烧。为此耶稣亦然，为以自己的血圣化百姓，在}（他说）\BibleQuote{城门之外受苦。}你看见预表大放光彩吗？他说：\BibleQuote{为赎罪，}和\BibleQuote{在城外受苦。}\BibleRef{希 13:13}\BibleQuote{因此，让我们到营外就祂去，忍受祂所受的凌辱。}这就是说，受同样的苦；与祂的苦难相通。祂像罪犯一样在城外被钉死：我们也不要羞于\Quote{出去}[世界]。\par

% structure: p0030 source=h2 target=subsection reason=relative-heading-rank; base=h2

\subsection{希 13:14-15}

\BibleQuote{因为我们在这里没有常存的城}（他说）\BibleQuote{而是寻求那将来的。因此，借着祂，让我们不断向天主奉献赞颂的祭献，就是颂扬祂名之嘴唇的果实。}\par

\Quote{借着祂，}按照血肉，如同借着大司祭一样。\BibleQuote{颂扬}（他说）\BibleQuote{祂的名。}（参看第514页）让我们不要说出任何亵渎、急躁、狂妄、傲慢或绝望的话。这就是\BibleQuote{怀着敬畏和虔敬的恐惧。}\BibleRef{希 12:28}因为灵魂在磨难中变得沮丧和鲁莽。但让我们不要[这样]。你看，他再次说了之前说过的话：\Quote{不要放弃我们自己的聚会，}因为这样我们就能以敬畏做一切事。因为很多时候，即使出于对人的尊重，我们也避免做许多恶事。\par

% structure: p0033 source=h2 target=subsection reason=relative-heading-rank; base=h2

\subsection{希 13:16}

\BibleQuote{但不要忘记行善和施舍。}我[不仅]指着在座的弟兄说，也指着那些不在的。但如果别人掠夺了你们的财产，就用你们所拥有的一切显示好客。那么，当我们看到他们即使在财产被掠夺之后仍受到这样的劝诫，我们还有什么借口呢？\par

5. 他并没有说\Quote{不要忘记}款待陌生人，而是说\Quote{不要忘记好客}：也就是说，不仅是要款待陌生人，而且要怀着对陌生人的爱。此外，他没有提到将来的、为我们存留的报酬，免得使他们更懈怠，而是提到已经给予的。因为\BibleQuote{因此有些人不知不觉款待了天使。}\par

但让我们看看在什么意义上说\BibleQuote{婚姻在众人中应受尊重，床是无污的。}因为（他意指）它使信徒保持贞洁。这里他也暗示犹太人，因为他们认为妇女生产后是不洁的：并且\Quote{凡从床上来的人}，经上说，\Quote{是不洁的。}那些源于本性的事并非不洁，忘恩负义、毫无知觉的犹太人，而是源于选择的事。因为如果\BibleQuote{婚姻应受尊重}而且纯洁，你怎么还能认为人会被它玷污呢？\par

\BibleQuote{让你们的品行}（他说）\BibleQuote{不要贪婪}：因为许多人耗尽财产后，后来又想在施舍的伪装下重新获取，所以他说：\BibleQuote{让你们的品行不要贪婪}；也就是说，我们应该只渴望必要和不可或缺的东西。那么（你说）如果我们连这些都没有供应呢？这是不可能的；确实不可能。\BibleQuote{因为祂说过}，祂不撒谎，\BibleQuote{我绝不离开你，也绝不弃舍你。所以我们可以放胆说：主是我的帮助者，我必不怕人能对我做什么。}你从祂自己得到了应许：从此不要怀疑。祂应许了；不要质疑。但这一点，\Quote{我绝不离开你}，他不仅指钱财，也指其他一切。\Quote{主是我的帮助者，我必不怕人能对我做什么}；很有道理。\par

因此我们也当在一切试探中这样说；只要天主眷顾我们，我们就嘲笑人间之事。因为当祂是我们的仇敌时，即使所有人都作我们的朋友，也毫无益处；同样，当祂是我们的朋友时，即使全体人类一同攻击我们，也没有损害。\BibleQuote{我必不怕人能对我做什么。}\par

6. \BibleQuote{要纪念那些统治你们的人，他们曾向你们宣讲了天主的话。}在此我认为他也是指帮助。因为这是[隐含在] \Quote{他们曾向你们宣讲了天主的话}这些[话]中。\par

\Quote{你们要效法他们的信德，留心他们品行的结局。} 什么是\Quote{留心}？就是不断思考，自己省察，推理，准确地查究，按你们的选择来考验。\Quote{他们品行的结局}，即他们的品行直到结局：因为\Quote{他们的品行}有了善终。\par

\Quote{耶稣基督昨天、今天、直到永远，常是一样。} 不要以为那时祂确实行了奇事，但现在却不行事。祂是同一的。这就是：\Quote{你们要纪念那些管辖你们的人。}\par

\Quote{不要被各种异样的教训所迷惑。} \Quote{异样}，就是不同于你们从我们所听的；\Quote{各种}，就是各式各样的：因为它们没有稳定性，而是彼此不同。尤其关于食物的教训是多样的。\par

\Quote{因为心靠恩宠坚固是好的，而不是靠食物。} 这些就是\Quote{各种}、\Quote{异样}的教训：尤其如基督所说：\BibleQuote{不是入口的污秽人，而是出口的污秽人。} \BibleRef{玛 15:11} 并且请注意，他不敢公开说这话，而是仿佛暗指。\Quote{因为心靠恩宠坚固是好的，而不是靠食物。}\par

信德就是一切。如果信德建立，心就安稳。因此信德建立，而理性动摇。因为信德与理性相反。\par

\Quote{这些}（他说）\Quote{没有使那些拘守于此的人获益。} 告诉我，遵守这些有什么益处？难道不是反而毁灭？难道不是使人处于罪下？如果必须遵守，我们就必须警惕。\par

\Quote{这些}（他说）\Quote{没有使那些拘守于此的人获益。} 这就是那些一直谨守的人。\par

只有一种遵守，就是远离罪恶。有什么益处呢？当有些人如此污秽，以致不能分食祭品时，这并没有拯救他们丝毫；虽然他们热心于规条，但因没有信德，即使如此也一无所获。\par

7. 接下来，他从预象中除去祭献，将论述转向原型，说：\Quote{那些大司祭将血带入圣所的牲畜，其身体在营外焚烧。} 那时那些是这些事的预象，同样，基督在\Quote{外}受苦，成就了一切。\par

在此他也清楚表明祂是自愿受苦，指出那些事并非偶然，甚至神圣的安排本身就是一种在\Quote{外}的苦难。祂在\Quote{外}受苦，但祂的血被带入了天堂。你看，我们分尝的血已被带入圣所，即真正的圣所；分尝的祭献是唯独司祭有特权的。因此我们分尝的是真理（实体）。如果分尝的不仅是\Quote{羞辱}，还有圣化，那么\Quote{羞辱}就是圣化的原因。因为祂如何被羞辱，我们也一样。所以，如果我们出到\Quote{外}，我们就与祂共融。\par

但\Quote{让我们到祂那里去}是什么意思？就是让我们在祂的苦难中与祂共融，让我们承担祂的羞辱。因为祂并非简单地命令我们住在\Quote{城门外}，而是如同祂被当作罪人受辱，我们也一样。\par

又\Quote{借着祂，让我们向天主献上祭献。} 他所指的是哪种祭献？\Quote{嘴唇的果实，就是称谢祂的圣名。} 他们（犹太人）带来绵羊、牛犊，交给司祭；\Quote{我们}却不必带这些东西，而是带感恩。这\Quote{果实}让\Quote{我们的嘴唇}结出。\par

\Quote{因为这样的祭献是天主所喜悦的。} 让我们向祂献上这样的祭献，好使祂能献给父。因为除了借着圣子，或也借着痛悔的心，别无他途。这一切都是为软弱的人说的。因为感恩属于圣子是明显的：否则，尊荣如何相等？\BibleQuote{众人}（祂说）\BibleQuote{应尊敬子如同尊敬父一样。} \BibleRef{若 5:23} 尊荣如何相等？\Quote{我们嘴唇的果实，就是称谢祂的圣名。}\par

8. 让我们凡事感恩，无论是贫穷、疾病，或任何其他事：因为只有祂知道什么为我们有益。\BibleQuote{因为我们不知道该如何祈求才对。} \BibleRef{罗 8:26} 我们连如何求合宜的事都不知道，除非领受了圣神，就要留心凡事献上感恩，并崇高承受一切。我们贫穷吗？要感恩。我们生病吗？要感恩。我们受诬告吗？要感恩。当我们受苦时，要感恩。\par

这使我们亲近天主：那时我们甚至使天主成为我们的债务人。但当我们顺遂时，我们却是债务人，并且要受审。因为顺遂时，我们欠天主的债；而这些事常常带来审判，而苦难却是赎罪。苦难吸引怜悯，吸引慈爱；而顺遂却使人升高到疯狂的骄傲，也导致懈怠，使人自视过高，使人膨胀。为此先知也说：\BibleQuote{主，祢苦难我有益，为叫我能学习祢的法度。} \BibleRef{咏 118:71} 希则克雅蒙恩脱离灾难后，心就高傲；当他患病时，就被贬抑，亲近天主。经上说：\BibleQuote{祂击杀他们时，他们就寻求祂，回转，清早亲近天主。} \BibleRef{咏 77:34} 又说：\BibleQuote{蒙爱者肥胖粗壮时，就踢跳。} \BibleRef{申 32:15} 因为\BibleQuote{主在施行审判时为人所知。} \BibleRef{咏 9:16}\par

9. \OriginalTerm{患难}{Affliction} 是一件大善事。\BibleQuote{窄狭是路}\BibleRef{玛 7:14}，因此患难将我们推入这窄路。不受患难压迫的人无法进入。因为那在窄路上刻苦己身的人，也是那享受安逸的人；但那把自己扩张开的人，不得进入，且如同被夹住般受苦。请看保禄如何进入这窄路。他\BibleQuote{制服}\BibleQuote{自己的身体}\BibleRef{格前 9:27}，以便能够进入。因此，在他一切患难中，他不断感谢天主。你失去了财产吗？这已减轻了你大部分的宽大。你失去了荣耀吗？这是另一种宽大。你受了诬告吗？别人相信了那些针对你、而你毫不知情的事吗？\BibleQuote{你们欢喜踊跃罢！}\BibleQuote{因为你们是有福的}\BibleRef{玛 5:11}，祂说，\BibleQuote{几时人辱骂你们，捏造一切坏话毁谤你们，为了我的缘故，你们欢喜踊跃罢！因为你们在天上的赏报是丰厚的。}\par

你为何惊奇，若你忧伤，并希望从诱惑中解脱？保禄也曾希望解脱，并屡次恳求天主，却未获得。因为\BibleQuote{关于这事，我三次求了主}\BibleRef{格后12:8}，这是多次；祂对我说：\BibleQuote{我的恩宠够你用，因为我的德能在软弱中才全显出来。}\BibleRef{格后12:9}。这里\Quote{软弱}指\Quote{患难}。那么，当他听到这话时，他感恩地接受，并说：\BibleQuote{所以我甘心情愿夸耀我的软弱}\BibleRef{格后12:10}；也就是说，我喜欢，我安息于我的患难中。为此，让我们为一切事感谢，既为安慰，也为患难。让我们不抱怨，不不知恩。\BibleQuote{我赤身出于母胎，也赤身归去。}\BibleRef{约1:21}你出来时并非荣耀，不要寻求荣耀。你被带到生命中是赤身的，不仅无钱财，也无关荣耀和尊荣。\par

请思量从财富中常生出多大的灾祸。因为\BibleQuote{骆驼穿过针眼，比富人进入天国还容易。}\BibleRef{玛 19:24}你看财富阻碍了多少善事，而你竟寻求富有？你岂不因这障碍已被推翻而欢喜吗？那通往天国的路是如此窄狭，财富却是如此宽阔，且充满肿胀和膨胀。因此祂说：\BibleQuote{变卖你所有的}\BibleRef{玛 19:21}，使那路能接纳你。你为什么渴求财富？正因如此，祂从你那里取走了它，为要释放你脱离奴役。因为真正的父亲，当儿子被某个情妇败坏，虽多方劝诫仍不能使他离开她时，就会将情妇放逐。财富的丰裕也是如此。因为主顾念我们，并解救我们脱离由此产生的害处，就从我们那里取走了财富。\par

因此，我们不要视贫穷为恶事：唯有罪是恶事。因为财富本身也不是善事：唯有中悦天主才是善事。那么，让我们寻求贫穷，追求贫穷；这样我们就能抓住天堂，这样我们就能获得其他善事。愿我们众人靠着我们的主耶稣基督的恩宠和慈爱，偕同圣神，归于圣父，获得荣耀、权能、尊崇，从现在直到永远，永世无尽。阿们。\par

\SourceRecord{Translated by Frederic Gardiner.；Nicene and Post-Nicene Fathers, First Series；Vol. 14.；Edited by Philip Schaff.；Buffalo, NY: Christian Literature Publishing Co.,；1889.；<http://www.newadvent.org/fathers/240233.htm>.}

% structure: p0001 source=h1 target=section reason=page-title

\section{希伯来人书讲道 第三十四}

% structure: p0002 source=h2 target=subsection reason=relative-heading-rank; base=h2

\subsection{\BibleRef{希 13:17}}

\Quote{\Emph{要服从你们的长上，并要顺从。因为他们为了你们的灵魂警醒，将来必须交代，使他们能欣然这样做，而不致悲伤}，\Emph{因为这对你们是无益的。}}\par

1. 无政府状态是一种邪恶，是许多灾难的缘由，是混乱和骚乱的根源。因为，如果你从合唱团中拿走领唱，合唱团就不会和谐有序；如果你从军队方阵中撤去指挥官，操练就无法准时有序进行；如果你从船上取走舵手，你就会使船沉没；同样，如果你从羊群中除去牧人，你就倾覆并毁灭了一切。\par

无政府状态是一种邪恶，是毁灭的原因。但不服从长上也同样是邪恶。因为结果又回到相同之处。因为一个不服从长上的民族，就像没有长上一样；或许甚至更糟。因为在前一种情况，他们至少还有混乱的借口，但在后一种情况，就不再有了，而是要受惩罚。\par

但也许有人会说，还有第三种邪恶，就是当长上是坏的时候。我自己也知道，这不是小的邪恶，甚至比无政府状态更坏。因为无人领导往往可得救，也往往陷于危险；但由邪恶的人领导，则完全陷于危险，被引向毁灭的坑中。\par

保禄怎会说：\Quote{要服从你们的长上，并要顺从}？他之前说了：\Quote{效法他们的信德，留意他们生活的结局}（见：\BibleRef{希 13:7}），然后他说：\Quote{要服从你们的长上，并要顺从。}\par

那么（你说），当他是邪恶的时候，我们也应该服从吗？\par

邪恶？在什么意义上？如果确实在信德方面，就逃离并躲避他；不仅他若是人，甚至若是从天上降下的天使，也如此；但如果是在行为方面，就不要过分好奇。我举这个例子不是出于自己的头脑，而是出于圣经。因为请听基督说：\BibleQuote{经师和法利塞人坐在梅瑟的座位上}。祂先前说了许多关于他们的可怕事情，然后说：\BibleQuote{他们在梅瑟的座位上：所以凡他们所吩咐你们观察的，你们就观察并实行；但不要效法他们的行为。}（见：\BibleRef{玛 23:2}）他们（祂的意思是）拥有职务的尊严，但生活不洁。然而你要注意，不要注意他们的生活，而要注意他们的话。因为关于他们的品行，没有人会因此受害。这如何成立？既因为他们的品行是显明给众人看的，也因为即使他万倍邪恶，他也绝不会教导邪恶的事。但在信德方面，邪恶并不显明给众人，而且邪恶的（长上）也不会克制不教导它。\par

此外，\BibleQuote{你们不要判断，免得你们受判断}（见：\BibleRef{玛 7:1}）是关于行为，而不是关于信德；无疑，下文也说明了这一点。因为（他说）\BibleQuote{为什么你看见你兄弟眼中的木屑，却不理会自己眼中的大梁？}（见：\BibleRef{玛 7:3}）\par

\BibleQuote{所以凡他们所吩咐你们观察的，你们就实行}（现在\Quote{实行}属于行为而不属于信德）\BibleQuote{但不要效法他们的行为。}你看见（这段话）不是关于教义，而是关于生活和行为吗？\par

保禄先前称赞了他们，然后说：\BibleQuote{要服从你们的长上，并要顺从，因为他们为了你们的灵魂警醒，将来必须交代。}\par

让那些长上也听吧，不仅是那些在他们管辖之下的人；这样，正如属下应当服从，长上也应当警醒和清醒。你说什么？他警醒；他冒着掉头的危险；他因你的罪而受惩罚，为了你的缘故要面对如此可怕的事，而你却懒惰、漠不关心、安逸吗？因此他说：\BibleQuote{使他们能欣然这样做，而不致悲伤；因为这对你们是无益的。}\par

你看见，被轻视的长上不应为自己复仇，他的大复仇就是哀哭和叹息。因为对于被病人轻视的医生，也不可能为自己复仇，而只能哀哭和叹息。但如果（长上）叹息（他意思是），天主就向你施行报复。因为如果我们为自己的罪哀哭时能吸引天主到我们这里，那么当我们为别人的傲慢和轻蔑哀哭时，岂不更能？你看见他不允许自己被引向辱骂吗？你看见他的哲学多么伟大？应当哀哭的是被轻视、被践踏、被唾弃的人。\par

不要因为他没有向你报复就自信，因为哀哭比任何报复更糟。因为当他哀哭对自己毫无益处时，他就呼求主；就像老师和保姆的情况，当孩子不听从他时，就有人被叫来更严厉地对待孩子，这里也是如此。\par

3. 哦！多么大的危险！对这些可怜的人，那些把自己投入如此大惩罚深渊的人，该说什么呢？你必须为你所管辖的所有人——妇女、儿童和男人——交账；你把头放进这么大的火里。如果任何长上能得救，我会感到惊奇，因为面对这样的威胁和现在的漠不关心，我看到有些人仍然甚至冲上去，把自己投入如此大的权柄重担下。\par

因为那些被勉强拉来的人，既没有避难所也没有辩护，如果他们玩忽职守、疏忽大意；既然连亚郎也是被勉强拉来的，却仍处于危险之中；梅瑟也曾多次推辞，却又陷于危险；撒乌耳被托付了另一种治理权，他在推辞之后，仍因管理不善而陷入危险；那么，那些费尽心力去争取职务并主动投身其中的人，岂不更加使自己失去一切借口？这样的人更完全剥夺了自己的辩解。因为人应当既因良心、也因职务的重担而恐惧战栗；既不应在被勉强拉来时就断然推辞，也不应在未被拉来时就主动投身，反而应当逃避，预见这尊位的伟大；而一旦被抓住，又应当再次表现他们的虔敬畏惧。凡事不可过度。你若事先察觉到，就退避；说服自己，你不配这职务。同样，你若被抓住，就当怀着敬畏，始终表现出正直的心志。\par

% structure: p0018 source=h2 target=subsection reason=relative-heading-rank; base=h2

\subsection{希伯来书 13:18}

4. \BibleQuote{请为我们祈祷}（他说）；\BibleQuote{因为我们自信在众人中有良好的良心，愿意诚实为人。}\par

你看见吗？他使用这些辩解，是写给那些对他心怀怨恨的人，写给那些离弃他、视他为犯法者、甚至不愿听他的名字的人？既然他向那些恨他的人请求所有别人从爱他们的人那里所求的（为他们祈祷），因此他在这里引入这句话，说：\BibleQuote{我们自信有良好的良心。}因为不要告诉我控告；他说，我们的良心没有伤害我们什么；我们也不自觉曾谋害过你们。\BibleQuote{因为我们自信}，他说，\BibleQuote{在众人中有良好的良心}，不仅在外邦人中，也在你们中间。我们没有行任何诡诈之事，没有行任何虚伪之事。因为很可能有人散布关于他的这些诽谤。\BibleQuote{因为他们曾被告知关于你的事}（据说）\BibleQuote{你教导背道。}\BibleRef{宗 21:21}他的意思是，我写这些事不是作为敌人或反对者，而是作为朋友。这一点他也通过下文表明。\par

% structure: p0021 source=h2 target=subsection reason=relative-heading-rank; base=h2

\subsection{希伯来书 13:19}

\BibleQuote{我更恳求你们这样做，使我得以更快地回到你们那里。}他这样祈祷是深爱他们的表现，并且不是简单如此，而是极其恳切，以便，他说，我能快快到你们那里去。渴望到他们那里去，是自觉无亏的表现，也是恳求他们为他祈祷的表现。\par

因此，他先求了他们的祈祷，然后自己也为他们祈求一切好事。\par

% structure: p0024 source=h2 target=subsection reason=relative-heading-rank; base=h2

\subsection{希伯来书 13:20-21}

\BibleQuote{愿赐平安的天主}（他说）（你们因此不要彼此不和）\BibleQuote{从地上领回羊群的大司牧}\BibleRef{希 13:20}（这是关于复活说的）\BibleQuote{大司牧}（又一个附加：在这里他再次向他们确认关于复活的话语，直到最后）\BibleQuote{藉着永久的盟约之血，我们的主耶稣基督，}\BibleRef{希 13:21} \BibleQuote{使你们在一切善工上完美，以承行祂的旨意，在你们身上成就祂眼中所喜悦的事。}\par

他又一次向他们作崇高的见证。因为那\BibleQuote{完美}的是指起初开始、后来完成的。他为他们祈祷，这是切慕他们的人的行为。在其它书信中，他在序言中祈祷，而在这里他在结尾祈祷。\BibleQuote{在你们身上成就}，他说，\BibleQuote{祂眼中所喜悦的事，藉着耶稣基督，愿光荣归于祂，直到永远。阿们。}\par

% structure: p0027 source=h2 target=subsection reason=relative-heading-rank; base=h2

\subsection{希伯来书 13:22-23}

5. \BibleQuote{弟兄们，我恳求你们，容受这劝勉的话，因为我只是简略地写信给你们。}你看见吗？他写给别人的情况（没有），却写给了你们？因为（他意思是）我甚至没有以长篇大论来烦扰你们。\par

我想他们对弟茂德并非没有好感，因此他也提到他。因为\BibleRef{希 13:23} \BibleQuote{你们要知道}，他说，\BibleQuote{我们的弟兄弟茂德已被释放；他若快来，我必同他一起见你们。} \BibleQuote{被释放}，他说；从哪里？我想他曾被投入监狱；或者若非如此，就是他被从雅典打发走了。这在\WorkTitle{宗徒大事录}中也有记载。\par

% structure: p0030 source=h2 target=subsection reason=relative-heading-rank; base=h2

\subsection{希伯来书 13:24-25}

6. 你看见他如何表明德行既不完全从天主而生，也不完全由我们自己？首先说\BibleQuote{使你们在一切善工上完美}；他意思是，你们确实有德行，但需要使之完全。\BibleQuote{善工善言}是什么？使生活和教义都正确。\BibleQuote{按照祂的旨意，在你们身上成就祂眼中所喜悦的事。}\par

\BibleQuote{在祂眼前}，他说。因为这是最高德行，做天主眼中所喜悦的事，正如先知也说：\BibleQuote{按我手中的清洁在祂眼前。}\BibleRef{咏 17:24}\par

写了这么多之后，他说这与他将要说相比仍是少的。正如他在另一处也说：\BibleQuote{正如我简略地写信给你们，你们念了，就能明白我对基督奥秘的知识。}\BibleRef{弗 3:3}\par

请注意他的智慧。他没有说：\BibleQuote{我恳求你们，容受这}训诫的话，而是\BibleQuote{劝勉的话}，即安慰、鼓励的话。他意思是，没有人会因所说话的长度而感到厌倦（难道这使他们疏远他吗？绝不是：他其实并不想明说）：即，即使你们是小气的人，因为这样的人的特点是不耐烦听长篇大论。\par

\BibleQuote{你们要知道，我们的弟兄弟茂德已被释放；他若快来，我必同他一起见你们。}这足以说服他们顺服，既然他准备与他的门徒一同前来。\par

\BibleQuote{请问候你们的领导者和所有圣人。}请看他是如何尊敬他们，因为他写信给他们而不是给那些（领导者）。\par

\Quote{意大利的弟兄姊妹问候你们。愿恩宠与你们众人同在。阿们。}这是为他们众人所共有的。\par

但是，\Quote{恩宠}如何\Quote{与}我们\Quote{同在}呢？如果我们不辜负这恩惠，如果我们不怠惰于这恩赐。那么，\Quote{这恩宠}是什么呢？罪过的赦免，洁净：这就是\Quote{与}我们\Quote{同在}。因为谁（他的意思是）能辜负这恩宠而不破坏它呢？例如：祂白白赦免了你的罪。那么，\Quote{恩宠}——无论是善意还是圣神的效果——如何\Quote{与}你\Quote{同在}呢？如果你藉着善行吸引它到你身边。因为一切美善的原因就是圣神的\Quote{恩宠}持续与我们同在。这引导我们走向一切[美善]，正如它离开我们时，就毁灭我们，使我们孤寂。\par

7. 那么，我们不要将它赶走。因为它存留与离开，全在乎我们自己。前者是由于我们思念天上的事，后者是由于[我们思念]今生的事。\Quote{世人是不能领受的，因为看不见祂，也不认识祂。}（\BibleRef{若 14:17}）你看见了吗？属世的灵魂不能拥有祂。我们需要极大的热忱，才能将祂紧握，好让祂引导我们的一切事务，并使我们安全、极其平安地行事。\par

因为如同一艘船顺风航行，只要享有顺利而稳定的风，就不会被阻碍或沉没，反而因其航程的进展令船员和乘客大为赞叹，使船员安歇，不必拼命划桨，让乘客免除一切恐惧，并欣赏其航程中最愉悦的景象；同样，被圣神坚固的灵魂远超今生的一切巨浪，比船更强劲地开辟通往天堂的航道，因为它不是被风推送，而是由\OriginalTerm{护慰者}{Paraclete}亲自充满所有纯净的风帆：祂将我们心中一切松弛懈怠的驱散。\par

因为风若吹在松弛的帆上，便毫无作用；同样，圣神也不愿停留在松弛的灵魂中；而是需要极大的紧张、极大的热忱，使我们的心灵燃烧，使我们的行为在任何情况下都紧绷而振奋。例如，当我们祈祷时，应当十分专注，将灵魂伸向天堂，不是用绳索，而是用极大的热忱。同样，当我们行慈善时，也需要专注，免得因家计、儿女、妻子、贫穷的恐惧而松弛我们的帆。因为如果我们藉着未来之事的希望从各方面绷紧它，它就能良好地领受圣神的动力；那些可朽坏的、可怜的事物都不会落在它上面，即使有落上的，也因紧绷而迅速被弹回、被抖落、坠落。\par

因此我们需要极大的专注。因为我们也是航行在辽阔的大海上，充满许多怪物、许多礁石，为我们兴起许多风暴，从晴空万里中掀起最猛烈的暴风雨。那么，如果我们想平安顺遂地航行，就必须扬起帆，也就是我们的决心：因为这对我们已足够。因为亚巴郎也是如此，当他将自己的情感伸向天主，将坚定的决心摆在祂面前时，他还需要别的吗？什么也不需要；而是\Quote{他相信了天主，天主就以此算为他的正义。}（\BibleRef{创 15:6}）但信德出于真诚的意愿。他献上了自己的儿子，虽然没有杀掉他，却如同杀掉了一样领受了赏报；虽然工作没有完成，赏报却给予了。\par

那么，让我们的帆妥善有序，不是陈旧的（因为一切\Quote{衰颓老旧的，就临近消失}（\BibleRef{希 8:13}）），不是破烂的，以便能承受圣神的动力。因为经上说：\Quote{属血肉的人，不能领受圣神的事。}（\BibleRef{格前 2:14}）因为正如蜘蛛网不能承受一阵风，同样，专心于今生的灵魂和\OriginalTerm{属血肉的人}{natural man}也绝不可能领受圣神的恩宠：因为我们的推理与它们无异，只维持表面上的连接，却毫无力量。\par

8. 然而，如果我们警醒，我们的状况并非如此：无论什么临到[基督徒]身上，他都能承受一切，超越一切，比任何漩涡都强。因为假设有一个属灵的人，无数灾祸临到他，他却不被任何灾祸击败。我说什么呢？让贫穷、疾病、侮辱、辱骂、嘲笑、鞭打、各种折磨、各种嘲弄、毁谤、攻击临到他吧；他仿佛置身世界之外，摆脱了身体的感觉，嘲笑一切。\par

我的话并非吹嘘，我想即使现在也有许多[这样的人]存在；例如，那些过隐修生活的人。然而，你们会说，这并不稀奇。但我认为，在城里生活的也有这样不为人知的人。但如果你们愿意，我也能展示一些古人中的例子。为使你们明白，请思考保禄。有什么可怕的事他没有经历、没有忍受？但他都高尚地承受了。让我们效法他，这样我们就能满载货物，驶入平静的港湾。\par

那么，让我们将心灵伸向天堂，被那渴望紧握，披上属灵的火焰，系上它的烈焰。一个带火的人不怕遇见的人：无论是野兽、是人、数不尽的陷阱，只要他以火武装，一切都会为他让路，一切退避。火焰无法忍受，火不可抗拒，它焚烧一切。\par

让我们披上这火，将光荣归于我们的主耶稣基督，愿光荣、权能、尊威归于祂，偕同圣父及圣神，从现在直到永远，永无穷尽。阿们。\par

感谢天主。\par

\SourceRecord{Translated by Frederic Gardiner.；Nicene and Post-Nicene Fathers, First Series；Vol. 14.；Edited by Philip Schaff.；Buffalo, NY: Christian Literature Publishing Co.,；1889.；<http://www.newadvent.org/fathers/240234.htm>.}
